\documentclass[11pt,oneside,openany]{book}
\usepackage[T1]{fontenc}\usepackage[utf8]{inputenc}
\usepackage{lmodern,amsmath,amssymb,amsthm,geometry}
\geometry{paperwidth=178mm,paperheight=254mm,inner=22mm,outer=18mm,top=21mm,bottom=23mm,headheight=14pt,headsep=7mm}
\usepackage{xcolor,tikz,pgfplots,longtable,array,enumitem}
\newcommand{\toprule}{\hline\noalign{\vskip4pt}}
\newcommand{\midrule}{\noalign{\vskip3pt}\hline\noalign{\vskip4pt}}
\newcommand{\bottomrule}{\noalign{\vskip3pt}\hline}
\pgfplotsset{compat=1.18}\usetikzlibrary{arrows.meta,calc}
\definecolor{navy}{HTML}{13283C}\definecolor{teal}{HTML}{116B70}
\definecolor{magenta}{HTML}{D74E72}\definecolor{gold}{HTML}{B17409}
\definecolor{blue}{HTML}{315FC2}\definecolor{purple}{HTML}{7952A4}
\usepackage[colorlinks=true,linkcolor=teal,citecolor=teal,urlcolor=teal,bookmarksnumbered=true]{hyperref}
\hypersetup{pdftitle={Triangle Centers Inspired by Physics: Lecture Notes},pdfsubject={An undergraduate companion to six lectures}}
\makeatletter
\def\ps@pitc{\def\@oddhead{\small\color{navy}Triangle Centers Inspired by Physics\hfil Lecture Notes}\let\@evenhead\@oddhead
\def\@oddfoot{\hfil\small\thepage\hfil}\let\@evenfoot\@oddfoot}
\makeatother\pagestyle{pitc}
\setlength{\parskip}{3pt}\setlength{\emergencystretch}{2em}
\setcounter{tocdepth}{1}\setcounter{secnumdepth}{2}\setlist[enumerate]{itemsep=5pt}
\newcommand{\TC}{\mathcal C}\newcommand{\mean}{\operatorname{mean}}
\newcommand{\argmax}{\operatorname*{arg\,max}}\newcommand{\argmin}{\operatorname*{arg\,min}}
\newcommand{\tr}{\operatorname{tr}}\newcommand{\dd}{\,\mathrm d}\newcommand{\R}{\mathbb R}
\newtheorem{proposition}{Proposition}[chapter]\newtheorem{definition}[proposition]{Definition}
\newtheorem{example}[proposition]{Example}\theoremstyle{remark}\newtheorem{remark}[proposition]{Remark}
\allowdisplaybreaks\begin{document}\frontmatter


\begin{titlepage}\centering\vspace*{14mm}
{\small\scshape Geometry through physical experiments\par}\vspace{12mm}
{\Huge\bfseries\color{navy} Triangle Centers\\[3mm] Inspired by Physics\par}
\vspace{8mm}{\LARGE Lecture Notes\par}\vspace{12mm}
\begin{tikzpicture}[x=1.7cm,y=1.7cm]
\fill[teal!6](0,0)--(4,0)--(0,3)--cycle;
\draw[teal,very thick](0,0)--(4,0)--(0,3)--cycle;
\draw[magenta,thick,rotate around={-29.87:(1.3333,1)}](1.3333,1)ellipse(1.470 and .786);
\fill[navy](1.3333,1)circle(2pt);
\draw[blue,thick](.05,1.72)--(2.65,.27);
\end{tikzpicture}
\vfill{\large A self-contained companion to six undergraduate lectures\par}
\vspace{5mm}{\small Balance and geometry $\cdot$ charges and currents\\
Membranes and fluids $\cdot$ quantum and random motion\\
Heat and reaction $\cdot$ optics\par}
\vspace{10mm}{\small First teaching edition\\1 October 2026\par}
\end{titlepage}
\chapter*{Preface}\addcontentsline{toc}{chapter}{Preface}
A triangular boundary is one of the simplest pieces of geometry, yet its
physical center depends on the question being asked. A sheet balances at
one point and its wire frame at another. A stretched membrane is deepest
at a point defined by a partial differential equation. A quantum ground
state has both a most probable position and a mean position. Cooling moves
a temperature maximum, while coherent light supplies an oscillatory field
whose dark regions can be measured in several ways.

These notes collect the complete conceptual, mathematical and numerical
material of a six-lecture miniseries. They assume vectors, elementary
calculus and the idea of a matrix eigenvector. The field equations and
numerical methods are introduced where they become useful. Each chapter
follows one lecture, but can be read without the slides, animations or
project atlas. The lectures last thirty minutes each; working through the
derivations and exercises naturally takes longer.

The common theme is a \emph{selection rule}: specify a distribution,
energy, boundary-value problem or measurement, and only then ask which
point it selects. Existence, uniqueness, interiority, a familiar geometric
identification and response to vertex motion are distinct questions.
An integral or variational definition can characterize a point exactly,
even when its barycentric function has no short explicit expression.

This teaching edition belongs to an ongoing research project. Identities
are proved or sourced, computations carry their resolution limitations,
and open questions are labeled. The attractivity status follows the
30 September 2026 addendum, including the certified E0 failure and two
eventual parameter-family failures. It supersedes earlier conjecture
labels. Absence from a finite ETC snapshot proves neither transcendence
nor exclusion from a future catalogue.

\chapter*{Reading guide}\addcontentsline{toc}{chapter}{Reading guide}
\begingroup\small
Chapter 1 introduces moments, covariance, barycentric coordinates and
invariant lines and conics. Chapter 2 contrasts prescribed charge with
mobile conductor charge and develops magnetic selectors. Chapter 3
derives the Poisson equation and its membrane--fluid--Brownian equivalence.
Chapter 4 distinguishes quantum means, modes and stochastic conditionings.
Chapter 5 develops heat flow, reaction and atlas methods. Chapter 6 treats
illumination and coherent propagation. The appendices supply proofs,
finite-element formulas, a glossary and answers to the exercises.

Demonstration paragraphs offer laboratory and coding investigations.
The companion teaching package contains editable slides, complete
lecturer notes and four offline browser demonstrations. Their controls
manipulate geometry, cooling, billiards and diffraction. This book does
not require them. All vector diagrams and plot coordinates are embedded
in the single LaTeX source; no image folder or external data is needed
to compile it.

\paragraph{Notation.}
Let $T=\operatorname{conv}\{A,B,C\}$, with area $\mathcal A$ and boundary
$\partial T$. The opposite side lengths are $a,b,c$, perimeter is
$P=a+b+c$, semiperimeter is $s=P/2$, and angles
$\theta_A,\theta_B,\theta_C$ sum to $\pi$.
The diameter is $D_T=\max(a,b,c)$ and
\[
L^2=(a^2+b^2+c^2)/3.
\]
The Laplacian is $\Delta$; dimensionless heat time is denoted $\Delta_h$
when a distinction is useful. Covariance is $\Sigma$ and the identity
matrix is $\mathbf1$. Our reference triangle is
\[
A=(0,0),\quad B=(4,0),\quad C=(0,3),\quad(a,b,c)=(5,3,4).
\]
Thus $\mathcal A=6$, $D_T=5$, $L^2=50/3$, $r=1$, and
\[
I=X_1=(1,1),\quad G=X_2=(4/3,1),\quad W=X_{10}=(3/2,1).
\]

\paragraph{Numerical figures and evidence.}
Colored cells display archived fields on a coarse 196-cell mesh; this is
not the solver-refinement mesh or a field enclosure. Each panel has its
own relative scale. Increasing teal saturation means increasing displayed
value. In a darkness plot it means greater darkness weight, rather than
greater intensity. Decimal coordinates are estimates unless an enclosure
is stated. Successive changes are diagnostics, not certified error bars.
The ETC neighbor rankings compare sampled locations, not exact identities.
\endgroup

\begingroup\small\tableofcontents\endgroup
\listoffigures\mainmatter

\chapter{Balance, shape and invariant geometry}\label{ch:1}

The same triangular boundary can support a sheet, a wire, a fluid, a quantum particle or a light aperture. Each experiment can select a different point. Our task is to say exactly what physical question selects it, write the governing mathematics, and decide which properties follow. We will use one triangle repeatedly so that changing the physics is visible. The course assumes vectors, elementary calculus and some familiarity with forces and waves. Eigenvectors and partial differential equations will be introduced where they become useful.


\section{Several centers of one triangle}\label{sec:slide2}

Mark the incenter I, area centroid G and wire centroid W. They coincide in an equilateral triangle, but that is a consequence of symmetry, not a universal identity. For a scalene triangle they usually separate. A balance point is a first moment of a distribution. A hottest point, deepest point or weakest-field point is an extremum of a scalar field. Later the same field will give both an extremum and a mean. Neither one automatically has priority. Always finish the sentence "center of what, under which assumptions?"


\begin{figure}[htbp]\centering
\begin{tikzpicture}[x=1.45cm,y=1.45cm]
\draw[teal,thick] (0,0) -- (4,0) -- (0,3) -- cycle;
\filldraw[fill=teal,draw=white,line width=.45pt] (1,1) circle(1.5pt) node[above left,font=\scriptsize,text=teal,fill=white,inner sep=.6pt] {$I=X_1$};
\filldraw[fill=navy,draw=white,line width=.45pt] (1.333333,1) circle(1.5pt) node[below left,font=\scriptsize,text=navy,fill=white,inner sep=.6pt] {$G=X_2$};
\filldraw[fill=gold,draw=white,line width=.45pt] (1.5,1) circle(1.5pt) node[above right,font=\scriptsize,text=gold,fill=white,inner sep=.6pt] {$W=X_{10}$};
\node[below left] at(0,0){$A$};\node[below right] at(4,0){$B$};\node[above] at(0,3){$C$};\end{tikzpicture}
\caption[Three physically different selection rules in the reference triangle]{Three physically different selection rules in the reference triangle.}\label{fig:geometry-2}
\end{figure}

\paragraph{Demonstration and investigation.} Reveal the three markers separately. Suspend the cardboard from two points, drawing the corresponding vertical lines. Explain that the laboratory suspension lines are not intrinsic central lines of the triangle.


\section{Triangle centers}\label{sec:slide3}

A triangle center is an intrinsic point chosen by the triangle geometry. Translating the apparatus translates the selected point. Rotating or reflecting it transforms the point the same way. Scaling it transforms the point when all physical lengths and parameters scale consistently. A lab-fixed magnetic direction, a fixed wavelength and screen distance, or an imposed anisotropic material can break this property. A parameter family can still be a family of triangle centers if its dimensionless parameter stays fixed. For example, heat time becomes diffusion coefficient times time divided by a squared triangle length.


\[\TC(sRT+t)=sR\TC(T)+t,\qquad s>0,\quad R^TR=\mathbf1.\]



\begin{definition}
A triangle center is a single-valued rule on nondegenerate triangles
that is equivariant under similarities and independent of vertex labels.
A parameterized physical rule is a family of centers when its parameters
are specified in dimensionless form.
\end{definition}
The orthogonal matrix $R$ includes reflections. Thus handedness introduced
by a current direction must disappear from the selected point: changing
orientation reverses a signed field but does not change its minimum
magnitude or a consistently normalized signed mean. If a model has two
equally valid extrema, it defines an invariant set until a justified
selection convention makes it a point. This matters for magnetic foci,
excited quantum states and buckled plates.




\section{Barycentric coordinates}\label{sec:slide4}

Imagine placing masses u, v and w at the vertices. After normalizing their total mass to one, their balance point is the displayed weighted average. These are barycentric coordinates. The projective notation (u:v:w) forgets the common scale, so (1:1:1) and (2:2:2) describe the same point. A zero coordinate places the point on the opposite side. Positive coordinates place it strictly inside. Negative coordinates are allowed and describe exterior points. The weights also equal signed subtriangle areas divided by the total area. The coordinate opposite A is the perpendicular distance to BC divided by the altitude from A.


\[p=\alpha A+\beta B+\gamma C,\quad\alpha+\beta+\gamma=1.\]


If $\mathcal A_A$ is the signed area of $pBC$, then
$\alpha=\mathcal A_A/\mathcal A=d_A/h_A$. An affine coordinate of a
mean equals the mean of that coordinate.


\begin{figure}[htbp]\centering
\begin{tikzpicture}[x=1.45cm,y=1.45cm]
\draw[teal,thick] (0,0) -- (4,0) -- (0,3) -- cycle;
\fill[magenta!15] (4,0) -- (0,3) -- (1.2,0.9) -- cycle;
\fill[gold!15] (0,3) -- (0,0) -- (1.2,0.9) -- cycle;
\fill[blue!15] (0,0) -- (4,0) -- (1.2,0.9) -- cycle;
\draw[teal,thick] (0,0) -- (4,0) -- (0,3) -- cycle;
\filldraw[fill=magenta,draw=white,line width=.45pt] (1.2,0.9) circle(1.5pt) node[above right,font=\scriptsize,text=magenta,fill=white,inner sep=.6pt] {$P$};
\node at (1.73,1.3) {$\alpha$};\node at (0.4,1.3) {$\beta$};\node at (1.73,0.3) {$\gamma$};\node[below left] at(0,0){$A$};\node[below right] at(4,0){$B$};\node[above] at(0,3){$C$};\end{tikzpicture}
\caption[Each vertex weight equals the opposite subtriangle area divided by the full area]{Each vertex weight equals the opposite subtriangle area divided by the full area.}\label{fig:geometry-4}
\end{figure}

\paragraph{Demonstration and investigation.} Move the marker from equal weights to a heavier vertex mass. Ask students to estimate the new location before it moves.


\section{Barycentric center functions}\label{sec:slide5}

The conventional center-function language packages a symmetric construction into one first-coordinate function. The same cyclic rule supplies the other coordinates. Scaling all side lengths multiplies all weights by a common factor when f is homogeneous. Swapping the two sides meeting at A should leave its weight unchanged. The incenter uses f=a, the centroid uses f=1, and the wire centroid uses f=b+c. A center function can be implicit or involve special functions. Elementary formulas are especially convenient, but their convenience is not a requirement imposed by physics.


\[\TC=\bigl(f(a,b,c):f(b,c,a):f(c,a,b)\bigr)_{\rm bary}.\]



Normalized barycentric weights sum to one; interior weights are positive.
Outside points are permitted, with some negative weights.
A homogeneous $f$ of any common degree gives scale-independent ratios.
Reflection and relabeling require $f(a,b,c)=f(a,c,b)$ and corresponding
permutations of the weights. For example $f=1$ gives $G=(1:1:1)$;
$f=a$ gives the incenter; $f=b+c$ gives the wire centroid.
Here $a$ is opposite $A$, not an edge issuing from $A$.
Trilinear coordinates are different: if $(x:y:z)$ are signed side
distances, the barycentrics are $(ax:by:cz)$.

An implicit definition still determines barycentric functions: solve
for the point and invert the affine coordinate map. The challenge is
a useful explicit expression over the whole shape space. The ETC is
a catalogue, not a test of whether an expression is algebraic \cite{etc}.




\section{Classical centers with physical meanings}\label{sec:slide6}

The Encyclopedia of Triangle Centers, or ETC, is Kimberling's catalogue. An X-number identifies a geometric rule, not an experimental substance. This short list already spans balance, equal distance, minimum spring energy and magnetic energy. The symmedian point X6 is especially useful: a point attached to three zero-rest-length springs whose anchors slide freely on the side lines minimizes the sum of squared perpendicular distances. Its barycentrics are a squared, b squared, c squared. For X3, equal distance to the vertices is exact, but the minimum enclosing-circle center equals X3 only for an acute triangle. The right and obtuse cases need separate treatment.



\paragraph{A spring construction for $X_6$.}
Three equal transverse springs penalize distances to supporting side
lines, with $U(p)=k\sum_i d_i(p)^2/2$. Signed interior distances obey
$ad_A+bd_B+cd_C=2\mathcal A$. Minimization gives
$d_A:d_B:d_C=a:b:c$; conversion to barycentrics gives
$(a^2:b^2:c^2)$, the symmedian point.
The Hessian $k\sum_i n_i n_i^T$ is positive definite because the
side normals span the plane. This is physical stability, which says
nothing yet about response to vertex motion.

$G$ also minimizes $\int_T|p-q|^2\dd A_q$, is the pivot of least polar
inertia, and is the first far-field location of a uniform lamina.
Three equal point masses have the same mean but covariance four
times that of the uniform area measure.



\begin{center}\small\begin{tabular}{p{0.19\linewidth}p{0.71\linewidth}}\toprule
ETC point & Physical rule\\\midrule
$X_1$ & Largest contained disk; equal side clearances\\
$X_2$ & Uniform lamina; equal vertex masses; mean-square pivot\\
$X_3$ & Equal vertex distances; acute-triangle circle center\\
$X_6$ & Minimum sum of squared side distances\\
$X_{10}$ & Uniform perimeter mean; medial incenter\\
$X_{360}$ & Equal-core vertex-current energy limit\\\bottomrule\end{tabular}\end{center}

\paragraph{Demonstration and investigation.} Use three sliding spring anchors as a thought experiment. Ask why a spring attached to a fixed vertex would define a different energy.


\section{Vertex attractivity}\label{sec:slide7}

Attractivity here is geometric response to changing the triangle, following the definition supplied for this project. Keep B and C fixed and move A. Compare the vertex displacement with the center displacement using their dot product. For G the center moves by exactly h/3, so the dot product is |h|\ensuremath{^2}/3. This proves strict attraction. The two arrows need not be parallel for another center. Only their projection matters. Physical stability instead concerns perturbing a field on a fixed domain. A stable conductor or membrane can select a geometrically nonattractive center.


\[h\cdot[\TC(A+h,B,C)-\TC(A,B,C)]\ge0.\]



Hold $B,C$ fixed and move $A(t)=A+th$. For a differentiable rule,
\[
h\cdot[\TC(A+h)-\TC(A)]
=\int_0^1 h^T S_A(A+th)h\dd t,\qquad
S_A=\tfrac12(D_A\TC+D_A\TC^T).
\]
Positive semidefiniteness of every vertex response matrix therefore
suffices for all admissible straight motions. It is locally necessary:
a negative eigenvector yields a negative small motion.
The path must remain nondegenerate; the two half-planes of $BC$ are
separate domains. Antisymmetric derivatives produce sideways response
and contribute nothing to this dot product.
Physical stability concerns an energy Hessian in a fixed triangle,
not these matrices. The project's finite-motion convention is
specified in \cite{attractive}.




\paragraph{Demonstration and investigation.} Run the centroid motion. Later we will see an E0 motion with a negative certified dot product. Tell students that a positive finite list of tests cannot prove attraction for every triangle.


\section{The 3--4--5 laboratory}\label{sec:slide8}

This right triangle has side lengths 3, 4 and 5, area 6, semiperimeter 6 and inradius 1. The incenter is (1,1), the area centroid is (4/3,1), and the wire centroid is (3/2,1). Its circumcenter is (2,3/2). These exact values give us checks on every numerical illustration. Axes and field scales will change from panel to panel, but the labels distinguish a physical field from its selected center. Numerical field colors are illustrations of archived computations. They are not additional interval certificates.


\[\mathcal A=6,\quad r=1,\quad G=(4/3,1),\quad W=(3/2,1).\]


The incenter weights are $(5:3:4)$ and the wire weights are
$(7:9:8)$, normalized by $24$. Substitution gives the quoted points.



\section{Uniform area and uniform wire}\label{sec:slide9}

For a uniform lamina, equal areas carry equal mass. For a uniform wire, equal arclengths carry equal mass. Integrating the area gives G. Integrating each straight edge gives its midpoint, weighted by edge length, and then gives W. Equal vertex masses also give G, despite having a different second moment. If charge is immobile and prescribed uniformly, its charge centroid follows exactly the same first-moment calculation. Conducting charge is mobile and will not generally stay uniform.


\[G=(A+B+C)/3,\quad W=\frac{a(B+C)+b(C+A)+c(A+B)}{2P}.\]



\paragraph{Performing the two integrals.}
Write $q=A+u(B-A)+v(C-A)$, where $u,v\ge0$, $u+v\le1$.
The Jacobian is $2\mathcal A$, and elementary integration gives
\[
\mathbb E u=\mathbb E v=\tfrac13,\quad
\mathbb E u^2=\mathbb E v^2=\tfrac16,\quad \mathbb E uv=\tfrac1{12}.
\]
Integrate $v$ from $0$ to $1-u$, then $u$ from $0$ to $1$.
The first moments give $G$. A segment's mean is its midpoint;
weighting the three midpoints by segment length gives $W$.
For the reference triangle $a=5,b=3,c=4$, so $W=(3/2,1)$.
Useful identities are
\[
W=\frac{3G-I}{2},\qquad
\Sigma_{\rm area}=\frac1{12}\sum_{V=A,B,C}(V-G)(V-G)^T.
\]
Thus $I,G,W$ always share a central line, independent of the Euler
line. In the reference triangle both lines are horizontal, a special
coincidence rather than a general identity.




\paragraph{Demonstration and investigation.} Allow thirty seconds for a prediction about a frame made of three rods with equal masses rather than equal linear density. Reveal that it gives G.


\section{Covariance and inertia}\label{sec:slide10}

The centroid removes the first moment. The covariance matrix records mean squared projections about it. In two dimensions it is a symmetric two-by-two matrix. Its eigenvectors give principal directions. The inertia about an in-plane axis e measures squared distance perpendicular to e, so covariance and inertia have the same principal axes but their eigenvalues are interchanged. This is why an object spread far along one direction has a smaller rotational moment about that direction. A polar inertia about the perpendicular z-axis is M times the covariance trace.


\[\Sigma=\int(q-G)(q-G)^T\dd\mu,\qquad\mathsf I=M(\tr\Sigma\,\mathbf1-\Sigma).\]



For a probability measure with mean $c$, translation of the pivot gives
$\int|q-p|^2\dd\mu=\tr\Sigma+|p-c|^2$.
For rotation about in-plane unit vector $e$, the perpendicular squared
lever arm is $|q-c|^2-[e\cdot(q-c)]^2$. Averaging proves the inertia
formula. The major covariance axis has the smaller in-plane inertia;
the normal-axis moment is $M\tr\Sigma$.

For the reference lamina,
\[
\Sigma_{\rm area}=\begin{pmatrix}8/9&-1/3\\-1/3&1/2\end{pmatrix},
\qquad \lambda_\pm=\frac{25\pm\sqrt{193}}{36}.
\]
The ellipse axis ratio is
$\sqrt{(25+\sqrt{193})/(25-\sqrt{193})}\simeq1.87$,
whereas the eigenvalue ratio is about $3.50$.

\paragraph{The far field.}
For $n>2$ take a potential kernel proportional to $-|x-q|^{2-n}$;
in two dimensions use $\log|x-q|$. Taylor expansion around $c$ gives
\[
\Phi(x)=M K_n(x-c)+
\frac M2\sum_{ij}\Sigma_{ij}\partial_i\partial_jK_n(x-c)
+\text{higher multipoles},
\]
with physical constants absorbed in $K_n$. The dipole vanishes exactly
at the mean. Thus the leading point source is at $G$ or $W$ according
to the distribution. A shifted point source cannot reproduce a general
quadrupole tensor. Differentiation gives the corresponding far field.




\section{Principal axes}\label{sec:slide11}

Rotate the measuring direction through the centroid. The variance of the projection is e\ensuremath{^T}\ensuremath{\Sigma}e. The two extremizing orthogonal directions are eigenvectors of \ensuremath{\Sigma}. They are intrinsic central lines selected by the mass distribution. They do not usually coincide with medians, angle bisectors, or the Euler line. A circular covariance at equilateral gives no preferred direction, so the individual axes are undefined there even though the center remains perfectly defined.


For $e=(\cos\theta,\sin\theta)$, stationarity of $e^T\Sigma e$
gives $\tan(2\theta)=2\Sigma_{xy}/(\Sigma_{xx}-\Sigma_{yy})$.
An \texttt{atan2} evaluation chooses the correct quadrant; diagonalization
avoids division when the denominator vanishes.



\paragraph{Demonstration and investigation.} Rotate the measuring line while the two principal axes remain visible. Have students identify the maximum and minimum spread directions.


\section{The Steiner inellipse}\label{sec:slide12}

The Steiner inellipse turns the second moment into a visible conic. Its center is G, its principal directions are those of the lamina, and its squared semiaxes are the eigenvalues of 2\ensuremath{\Sigma}. It is the unique ellipse of largest area contained in the triangle. Its contact points are side midpoints. Doubling the semiaxes gives the Steiner circumellipse, which passes through the vertices and is the minimum-area enclosing ellipse. The inellipse semiaxis ratio is the square root of the reversed principal-inertia ratio, not that ratio itself.


\[(p-G)^T(2\Sigma_{\rm area})^{-1}(p-G)=1.\]



For an equilateral triangle centered at the origin with circumradius
$R$, rotational symmetry and the moment formula give
$\Sigma=R^2\mathbf1/8$. Its incircle radius is $R/2$,
so its shape matrix is $2\Sigma$. Every nondegenerate triangle is an
invertible affine image of this one. Covariance transforms as
$B\Sigma B^T$, while affine maps preserve side midpoints and tangency.
This proves the Steiner formula.

If $Q=2\Sigma$, semiaxes are $\sqrt{\lambda_\pm(Q)}$ and foci are
$G\pm\sqrt{\lambda_+(Q)-\lambda_-(Q)}\,e_+$.
The conic, its axes and its unordered focal pair are different
invariant structures. The foci reappear as magnetic field zeros.
The appendix proves the maximum-area property.



\begin{figure}[htbp]\centering
\begin{tikzpicture}[x=1.45cm,y=1.45cm]
\draw[teal,thick] (0,0) -- (4,0) -- (0,3) -- cycle;
\draw[magenta,thick] (2.607971,0.2678861) -- (2.630831,0.314005) -- (2.648135,0.3630615) -- (2.659808,0.4148455) -- (2.665802,0.4691352) -- (2.666089,0.5256981) -- (2.66067,0.5842921) -- (2.649567,0.6446662) -- (2.632827,0.7065619) -- (2.610523,0.7697141) -- (2.58275,0.8338524) -- (2.549626,0.8987023) -- (2.511294,0.9639858) -- (2.467918,1.029424) -- (2.419684,1.094735) -- (2.366797,1.159642) -- (2.309485,1.223864) -- (2.247993,1.287128) -- (2.182585,1.349162) -- (2.113539,1.409702) -- (2.041153,1.468486) -- (1.965736,1.525265) -- (1.887611,1.579795) -- (1.807112,1.631841) -- (1.724584,1.681182) -- (1.640381,1.727606) -- (1.554863,1.770915) -- (1.468397,1.810922) -- (1.381352,1.847457) -- (1.294102,1.880362) -- (1.207019,1.909498) -- (1.120478,1.934739) -- (1.034848,1.955978) -- (0.9504959,1.973123) -- (0.8677834,1.986101) -- (0.7870645,1.994856) -- (0.7086847,1.999351) -- (0.6329798,1.999567) -- (0.5602739,1.995502) -- (0.4908784,1.987175) -- (0.4250904,1.97462) -- (0.3631916,1.957892) -- (0.3054471,1.937062) -- (0.2521042,1.91222) -- (0.2033913,1.883471) -- (0.159517,1.850939) -- (0.1206691,1.814763) -- (0.08701404,1.775098) -- (0.05869592,1.732114) -- (0.03583598,1.685995) -- (0.01853213,1.636938) -- (0.006858461,1.585154) -- (0.0008649599,1.530865) -- (0.0005772928,1.474302) -- (0.005996692,1.415708) -- (0.01709995,1.355334) -- (0.03383952,1.293438) -- (0.05614372,1.230286) -- (0.08391705,1.166148) -- (0.1170406,1.101298) -- (0.1553724,1.036014) -- (0.1987485,0.9705763) -- (0.2469831,0.9052645) -- (0.2998695,0.8403584) -- (0.3571814,0.7761359) -- (0.4186734,0.7128719) -- (0.4840821,0.6508376) -- (0.5531273,0.5902983) -- (0.6255136,0.5315135) -- (0.7009308,0.4747348) -- (0.7790561,0.4202054) -- (0.8595549,0.3681588) -- (0.9420825,0.3188178) -- (1.026285,0.2723937) -- (1.111803,0.2290853) -- (1.19827,0.1890782) -- (1.285314,0.1525435) -- (1.372565,0.1196377) -- (1.459647,0.09050183) -- (1.546189,0.06526053) -- (1.631819,0.04402194) -- (1.716171,0.02687699) -- (1.798883,0.0138991) -- (1.879602,0.005143846) -- (1.957982,0.0006487199) -- (2.033687,0.0004329696) -- (2.106393,0.004497519) -- (2.175788,0.01282496) -- (2.241576,0.02537964) -- (2.303475,0.04210779) -- (2.36122,0.06293779) -- (2.414562,0.08778043) -- (2.463275,0.1165293) -- (2.50715,0.1490614) -- (2.545998,0.1852373) -- (2.579653,0.2249021) -- (2.607971,0.2678861);
\draw[blue,thick] (0.1193343,1.697285) -- (2.547332,0.302715);
\draw[gold,thick] (0.6360483,-0.213999) -- (2.030618,2.213999);
\filldraw[fill=navy,draw=white,line width=.45pt] (1.333333,1) circle(1.5pt) node[below,font=\scriptsize,text=navy,fill=white,inner sep=.6pt] {$G$};
\filldraw[fill=teal,draw=white,line width=.45pt] (2,0) circle(1.5pt) node[above right,font=\scriptsize,text=teal,fill=white,inner sep=.6pt] {};
\filldraw[fill=teal,draw=white,line width=.45pt] (2,1.5) circle(1.5pt) node[above right,font=\scriptsize,text=teal,fill=white,inner sep=.6pt] {};
\filldraw[fill=teal,draw=white,line width=.45pt] (0,1.5) circle(1.5pt) node[above right,font=\scriptsize,text=teal,fill=white,inner sep=.6pt] {};
\node[below left] at(0,0){$A$};\node[below right] at(4,0){$B$};\node[above] at(0,3){$C$};\end{tikzpicture}
\caption[The midpoint-tangent Steiner inellipse and the principal inertia directions of the uniform lamina]{The midpoint-tangent Steiner inellipse and the principal inertia directions of the uniform lamina.}\label{fig:geometry-12}
\end{figure}

\paragraph{Demonstration and investigation.} Reveal the midpoint contacts, then the principal axes, then the larger circumellipse. Keep both conics large enough to see their shared center.


\section{Affine deformation of the conic}\label{sec:slide13}

Compare the equilateral triangle on the left with its affine image on the right. Lengths and angles change, but straight lines, ratios along a line, midpoint locations and tangency remain. A circle generally becomes an ellipse. The center maps to G. This is a concrete demonstration of an invariant structure that is more informative than a point alone. It also connects with an oblique cut through an equilateral triangular prism: the inscribed circular cylinder cuts into the Steiner inellipse of the section triangle.


\[\Sigma\mapsto B\Sigma B^T,\qquad p\mapsto Bp+t.\]



\paragraph{Demonstration and investigation.} Run the point motions from the circle construction to the deformed ellipse. Ask the audience to track one midpoint contact rather than the whole picture.


\section{The skeleton has a different inertia ellipse}\label{sec:slide14}

A wire frame has a covariance and principal axes, so it certainly has inertia ellipses. What fails is the coincidence enjoyed by the area distribution: its covariance ellipse is not generally the midpoint-tangent Steiner inellipse. Scaling a wire covariance ellipse until it first touches the boundary will usually give only one or two contacts. There can also be a different ellipse centered at W and tangent to all sides. Do not confuse that geometric ellipse with the one determined by wire moments. Equilateral symmetry makes all these examples circles, hiding the distinction.


\[\Sigma_{\rm wire}=\begin{pmatrix}7/4&-2/3\\-2/3&1\end{pmatrix}.\]



Uniform wire means uniform arclength. General affine maps stretch
different edge directions by different amounts, so the wire measure
does not transform like normalized area. A positive-definite covariance
ellipse still exists:
$(p-W)^T\Sigma_{\rm wire}^{-1}(p-W)=1$.
It generally lacks midpoint tangencies and is not a maximum-area
inellipse. For an edge with endpoints $U,V$ and midpoint $m$,
\[
\Sigma_{\rm wire}=\frac1P\sum_{\rm edges}\ell
\left[\frac{(V-U)(V-U)^T}{12}+(m-W)(m-W)^T\right].
\]
This follows by the segment variance and parallel-axis identity,
and yields the reference matrix. Absence of a Steiner identification
is not absence of an inertia ellipse.



\begin{figure}[htbp]\centering
\begin{tikzpicture}[x=1.45cm,y=1.45cm]
\draw[teal,thick] (0,0) -- (4,0) -- (0,3) -- cycle;
\draw[magenta,thick] (2.762737,0.2614917) -- (2.785823,0.3071706) -- (2.803404,0.3558163) -- (2.815403,0.4072204) -- (2.82177,0.461163) -- (2.822477,0.517413) -- (2.81752,0.5757294) -- (2.806922,0.6358627) -- (2.790727,0.6975552) -- (2.769005,0.7605429) -- (2.741849,0.8245559) -- (2.709376,0.8893202) -- (2.671723,0.9545585) -- (2.629053,1.019991) -- (2.581548,1.085339) -- (2.529412,1.15032) -- (2.472868,1.214659) -- (2.412158,1.278077) -- (2.347542,1.340306) -- (2.279296,1.401077) -- (2.207714,1.46013) -- (2.133101,1.517213) -- (2.055777,1.572081) -- (1.976073,1.6245) -- (1.89433,1.674244) -- (1.810899,1.721101) -- (1.726136,1.764871) -- (1.640405,1.805365) -- (1.554073,1.84241) -- (1.46751,1.875848) -- (1.381085,1.905535) -- (1.29517,1.931345) -- (1.210132,1.953167) -- (1.126335,1.970907) -- (1.044138,1.984489) -- (0.963893,1.993856) -- (0.8859439,1.998967) -- (0.8106243,1.9998) -- (0.7382566,1.996352) -- (0.6691509,1.988637) -- (0.6036031,1.976689) -- (0.5418937,1.960559) -- (0.4842871,1.940315) -- (0.4310299,1.916045) -- (0.3823502,1.887852) -- (0.3384565,1.855857) -- (0.2995367,1.820197) -- (0.2657574,1.781025) -- (0.2372634,1.738508) -- (0.2141766,1.692829) -- (0.1965959,1.644184) -- (0.1845965,1.59278) -- (0.1782299,1.538837) -- (0.1775234,1.482587) -- (0.1824799,1.424271) -- (0.1930782,1.364137) -- (0.209273,1.302445) -- (0.2309948,1.239457) -- (0.2581507,1.175444) -- (0.2906244,1.11068) -- (0.3282769,1.045442) -- (0.3709468,0.9800086) -- (0.4184515,0.9146614) -- (0.4705876,0.8496796) -- (0.5271318,0.7853414) -- (0.5878419,0.7219225) -- (0.6524581,0.6596944) -- (0.7207036,0.5989234) -- (0.7922861,0.53987) -- (0.8668991,0.4827869) -- (0.9442232,0.4279186) -- (1.023927,0.3755) -- (1.10567,0.3257557) -- (1.189101,0.2788985) -- (1.273864,0.2351292) -- (1.359595,0.1946352) -- (1.445927,0.1575899) -- (1.53249,0.124152) -- (1.618915,0.09446454) -- (1.70483,0.06865474) -- (1.789868,0.0468331) -- (1.873665,0.02909307) -- (1.955862,0.01551061) -- (2.036107,0.00614388) -- (2.114056,0.001032999) -- (2.189376,0.0001998472) -- (2.261743,0.003647994) -- (2.330849,0.01136267) -- (2.396397,0.02331085) -- (2.458106,0.03944136) -- (2.515713,0.05968512) -- (2.56897,0.08395547) -- (2.61765,0.1121484) -- (2.661543,0.1441433) -- (2.700463,0.1798032) -- (2.734243,0.2189752) -- (2.762737,0.2614917);
\draw[blue,thick] (0.2915067,1.706784) -- (2.708493,0.2932157);
\draw[gold,thick] (0.7932157,-0.2084933) -- (2.206784,2.208493);
\filldraw[fill=navy,draw=white,line width=.45pt] (1.5,1) circle(1.5pt) node[below,font=\scriptsize,text=navy,fill=white,inner sep=.6pt] {$W$};
\node[below left] at(0,0){$A$};\node[below right] at(4,0){$B$};\node[above] at(0,3){$C$};\end{tikzpicture}
\caption[The covariance ellipse of the wire has center $W$; its scale is $\Sigma_{\rm wire}$]{The covariance ellipse of the wire has center $W$; its scale is $\Sigma_{\rm wire}$. It generally lacks the Steiner tangencies.}\label{fig:geometry-14}
\end{figure}

\paragraph{Demonstration and investigation.} Overlay the wire covariance ellipse, scaled to its first contact, and the Steiner inellipse. Ask students to count contact points.


\section{Central lines and central conics}\label{sec:slide15}

The Euler line joins the circumcenter, centroid and orthocenter. In this right triangle the orthocenter is A, and G divides the O-to-H segment in the ratio one to two. The principal axes are another pair of central lines, usually different from the Euler line. The Steiner ellipse is a central conic. We will meet its two foci again as magnetic field zeros. A pair may be a natural result when no intrinsic rule distinguishes its two members. It is better to retain that pair than force an arbitrary single center.


\[H,G,O\ \hbox{are collinear};\qquad HG=2GO.\]



\paragraph{Demonstration and investigation.} End with a thirty-second retrieval exercise: name the area mean, wire mean, principal-axis object, and conic. Allow students to explain one of them to a neighbor.


\section*{Questions and exercises}\addcontentsline{toc}{section}{Questions and exercises}

\begin{enumerate}

\item\label{ex:2} What would equal point masses at the vertices select?

\item\label{ex:3} Does a fixed-height optical experiment automatically define a scale-invariant triangle center?

\item\label{ex:4} What point has weights (1:1:0)?

\item\label{ex:5} Why does multiplying all three weights by the same number change nothing?

\item\label{ex:6} Is an X-number evidence that two physical experiments are different centers?

\item\label{ex:7} Does the center have to move exactly toward the displaced vertex?

\item\label{ex:8} Which of I, G and W lies farthest along the base direction?

\item\label{ex:9} What does uniform mean for those three equal-mass rods?

\item\label{ex:10} If the largest covariance eigenvalue is along e, is the inertia about e the larger one?

\item\label{ex:11} Why can a center be unique while its principal axes are not?

\item\label{ex:12} Why do ordinary rotations alone not suffice for this proof?

\item\label{ex:13} Does uniform wire density survive this same affine deformation?

\item\label{ex:14} Is the correct conclusion that no ellipse exists for the skeleton?

\item\label{ex:15} Which construction will reappear as a magnetic zero pair?

\end{enumerate}

\chapter{Charges, currents and conformal geometry}\label{ch:2}

We now allow a field to redistribute itself. The geometry remains a triangle, but a conducting charge distribution and a magnetic field magnitude can weight different regions very differently. Our central questions are whether the selected point exists, whether it is unique, whether a numerical value is reliable, and whether it follows a vertex displacement attractively. These questions need separate arguments.


\section{Fixed charge and mobile charge}\label{sec:slide17}

With immobile static charge, a uniform lamina gives G and a uniform frame gives W. In a conductor, free charges repel and redistribute until the potential is constant on the conductor. The density is typically enhanced near edges and corners. A centroid of positive charge still lies in the triangle, but it need not be G or W. At very large observation distances, expansion of the potential singles out this charge centroid as the origin where the dipole term vanishes.


\[c_\mu=\frac{\int q\dd\mu(q)}{\int\dd\mu(q)}.\]


Expand the far-field kernel about the first moment. Translating
the origin to that moment cancels the dipole correction, just as for
the gravitational distributions in Chapter 1.



\section{Electrostatic equilibrium energy}\label{sec:slide18}

In planar electrostatics the Green kernel is logarithmic, conventionally \ensuremath{-}log r. In three dimensions it is proportional to 1/r. The constant scale factors do not change an equilibrium distribution or its centroid. Varying charge while keeping total charge fixed shows that the equilibrium potential is constant on the charged support, with the appropriate inequality elsewhere. The zero-thickness idealization matters: some sets have zero capacity in higher dimensions, so an equilibrium measure cannot simply be assumed. The project uses specified uniform-thickness or uniform-radius limits for those cases.


\[\mathcal E(\mu)=\iint K_n(p,q)\dd\mu(p)\dd\mu(q),\quad\mu\ge0,\quad\mu(T)=1.\]



Potential is $V_\mu(p)=\int K_n(p,q)\dd\mu(q)$; physical energy has
an additional constant $Q^2/2$. In two dimensions
$K_2=-\log(|p-q|/L)$; for $n>2$, $K_n=|p-q|^{2-n}$.
A small transfer from $q$ to $p$ changes energy by
$2[V_\mu(p)-V_\mu(q)]$ to first order. Equilibrium therefore has
constant potential on its charged support, apart from negligible
exceptional sets, and no smaller potential at an allowed empty point.
This is an integral equation, not uniform charge density.

Positive panel charges minimize $q^TKq$ subject to $\sum q_i=1$.
Integrate singular self-panel and near-panel entries; refine edges
and compare moments as well as energy.
Exactly zero-radius wires have zero Newtonian capacity in dimensions
three and above. A lamina has critical zero capacity in four dimensions.
The adopted uniform-radius and uniform-thickness limits select
$X_{10}$ for wires and $X_2$ for the critical plate. They are limiting
conventions, not finite-energy equilibria on the ideal sets.
Thin three-dimensional wires approach uniform density slowly
\cite{wire}; corners and precise thickening need separate analysis.




\section{E2 and E3: conducting charge centroids}\label{sec:slide19}

E2 is the planar equilibrium-charge centroid of either the filled triangular conductor or its boundary: the logarithmic equilibrium measure lives on the boundary. This is why S2 equals E2. E3 is the charge centroid of a thin triangular plate in three-dimensional space. Its projected density fills the plate and becomes singular near its edges. The diagrams show computed distributions in the same 3--4--5 outline. E2 has a certified negative attractivity witness. E3 has refined negative numerical evidence, but its continuum sign is not interval-certified.



Planar logarithmic conductor charge lives on the boundary. Normalize
an exterior conformal map as
\[
g(w)=cw+a_0+a_1w^{-1}+\cdots,\qquad |w|>1,\quad c>0.
\]
Uniform angular measure on the circle maps to equilibrium charge.
Averaging $g(e^{i\theta})$ removes every nonconstant Fourier term,
so E2 is $a_0$. A filled triangle and its skeleton have the same
logarithmic charge centroid. The scale $c$ is capacity; $a_0$ is a
point. Exterior triangle maps may involve Appell functions \cite{finch}.

For a thin plate in three dimensions the kernel is $1/|p-q|$.
Charge occupies the sheet with enhanced edge density. Singular panels
give E3. A planar conformal change transforms this kernel into
$1/|F(z)-F(\zeta)|$, with Jacobians. It does not reduce the problem
to the ordinary disk Coulomb kernel. Holomorphic conformal invariance
applies directly to two-dimensional Laplace problems.



\begin{figure}[htbp]\centering
\begin{tikzpicture}[x=1.7cm,y=1.7cm]
\fill[teal!75!white] (3.142857,0.6428571) -- (2.857143,0.8571429) -- (2.857143,0.6428571) -- cycle;
\fill[teal!83!white] (1.714286,1.714286) -- (1.428571,1.928571) -- (1.428571,1.714286) -- cycle;
\fill[teal!63!white] (1.714286,1.714286) -- (1.714286,1.5) -- (2,1.5) -- cycle;
\fill[teal!94!white] (0.5714286,2.571429) -- (0.2857143,2.785714) -- (0.2857143,2.571429) -- cycle;
\fill[teal!83!white] (2.571429,1.071429) -- (2.285714,1.285714) -- (2.285714,1.071429) -- cycle;
\fill[teal!94!white] (3.714286,0.2142857) -- (3.428571,0.4285714) -- (3.428571,0.2142857) -- cycle;
\fill[teal!94!white] (3.714286,0.2142857) -- (3.714286,0) -- (4,0) -- cycle;
\fill[teal!71!white] (0.8571429,2.142857) -- (1.142857,2.142857) -- (0.8571429,2.357143) -- cycle;
\fill[teal!74!white] (1.142857,2.142857) -- (1.142857,1.928571) -- (1.428571,1.928571) -- cycle;
\fill[teal!94!white] (0,2.785714) -- (0.2857143,2.785714) -- (0,3) -- cycle;
\fill[teal!88!white] (0.5714286,2.571429) -- (0.5714286,2.357143) -- (0.8571429,2.357143) -- cycle;
\fill[teal!64!white] (2,1.285714) -- (2.285714,1.285714) -- (2,1.5) -- cycle;
\fill[teal!79!white] (2.571429,1.071429) -- (2.571429,0.8571429) -- (2.857143,0.8571429) -- cycle;
\fill[teal!94!white] (3.142857,0.4285714) -- (3.428571,0.4285714) -- (3.142857,0.6428571) -- cycle;
\fill[teal!94!white] (0.2857143,0) -- (0,0.2142857) -- (0,0) -- cycle;
\fill[teal!65!white] (0,0.2142857) -- (0.2857143,0) -- (0.2857143,0.2142857) -- cycle;
\fill[teal!70!white] (0.5714286,0.2142857) -- (0.2857143,0) -- (0.5714286,0) -- cycle;
\fill[teal!50!white] (0.2857143,0) -- (0.5714286,0.2142857) -- (0.2857143,0.2142857) -- cycle;
\fill[teal!41!white] (1.714286,1.5) -- (1.714286,1.714286) -- (1.428571,1.714286) -- cycle;
\fill[teal!23!white] (1.428571,1.5) -- (1.714286,1.5) -- (1.428571,1.714286) -- cycle;
\fill[teal!55!white] (0,1.928571) -- (0,1.714286) -- (0.2857143,1.928571) -- cycle;
\fill[teal!33!white] (0,1.714286) -- (0.2857143,1.714286) -- (0.2857143,1.928571) -- cycle;
\fill[teal!70!white] (0.2857143,2.357143) -- (0,2.571429) -- (0,2.357143) -- cycle;
\fill[teal!53!white] (0,2.571429) -- (0.2857143,2.357143) -- (0.2857143,2.571429) -- cycle;
\fill[teal!64!white] (0,2.142857) -- (0.2857143,2.357143) -- (0,2.357143) -- cycle;
\fill[teal!40!white] (0.2857143,2.357143) -- (0,2.142857) -- (0.2857143,2.142857) -- cycle;
\fill[teal!57!white] (0,2.142857) -- (0,1.928571) -- (0.2857143,1.928571) -- cycle;
\fill[teal!37!white] (0.2857143,2.142857) -- (0,2.142857) -- (0.2857143,1.928571) -- cycle;
\fill[teal!21!white] (0.2857143,1.714286) -- (0.5714286,1.928571) -- (0.2857143,1.928571) -- cycle;
\fill[teal!17!white] (0.5714286,1.714286) -- (0.5714286,1.928571) -- (0.2857143,1.714286) -- cycle;
\fill[teal!52!white] (0.2857143,1.5) -- (0,1.714286) -- (0,1.5) -- cycle;
\fill[teal!32!white] (0,1.714286) -- (0.2857143,1.5) -- (0.2857143,1.714286) -- cycle;
\fill[teal!18!white] (0.2857143,1.5) -- (0.5714286,1.714286) -- (0.2857143,1.714286) -- cycle;
\fill[teal!15!white] (0.5714286,1.714286) -- (0.2857143,1.5) -- (0.5714286,1.5) -- cycle;
\fill[teal!52!white] (0,1.285714) -- (0.2857143,1.5) -- (0,1.5) -- cycle;
\fill[teal!31!white] (0.2857143,1.5) -- (0,1.285714) -- (0.2857143,1.285714) -- cycle;
\fill[teal!52!white] (0,1.285714) -- (0,1.071429) -- (0.2857143,1.071429) -- cycle;
\fill[teal!31!white] (0.2857143,1.285714) -- (0,1.285714) -- (0.2857143,1.071429) -- cycle;
\fill[teal!94!white] (3.428571,0) -- (3.714286,0) -- (3.428571,0.2142857) -- cycle;
\fill[teal!94!white] (3.714286,0) -- (3.714286,0.2142857) -- (3.428571,0.2142857) -- cycle;
\fill[teal!62!white] (1.428571,0) -- (1.714286,0) -- (1.428571,0.2142857) -- cycle;
\fill[teal!36!white] (1.714286,0) -- (1.714286,0.2142857) -- (1.428571,0.2142857) -- cycle;
\fill[teal!67!white] (0.8571429,0) -- (0.5714286,0.2142857) -- (0.5714286,0) -- cycle;
\fill[teal!40!white] (0.5714286,0.2142857) -- (0.8571429,0) -- (0.8571429,0.2142857) -- cycle;
\fill[teal!63!white] (1.142857,0.2142857) -- (1.142857,0) -- (1.428571,0) -- cycle;
\fill[teal!36!white] (1.142857,0.2142857) -- (1.428571,0) -- (1.428571,0.2142857) -- cycle;
\fill[teal!67!white] (1.142857,0.2142857) -- (0.8571429,0) -- (1.142857,0) -- cycle;
\fill[teal!38!white] (0.8571429,0) -- (1.142857,0.2142857) -- (0.8571429,0.2142857) -- cycle;
\fill[teal!24!white] (1.142857,1.928571) -- (1.142857,1.714286) -- (1.428571,1.714286) -- cycle;
\fill[teal!42!white] (1.428571,1.928571) -- (1.142857,1.928571) -- (1.428571,1.714286) -- cycle;
\fill[teal!45!white] (0.8571429,2.142857) -- (1.142857,1.928571) -- (1.142857,2.142857) -- cycle;
\fill[teal!26!white] (1.142857,1.928571) -- (0.8571429,2.142857) -- (0.8571429,1.928571) -- cycle;
\fill[teal!87!white] (0,2.785714) -- (0,2.571429) -- (0.2857143,2.571429) -- cycle;
\fill[teal!86!white] (0.2857143,2.785714) -- (0,2.785714) -- (0.2857143,2.571429) -- cycle;
\fill[teal!18!white] (1.142857,1.928571) -- (0.8571429,1.714286) -- (1.142857,1.714286) -- cycle;
\fill[teal!18!white] (0.8571429,1.714286) -- (1.142857,1.928571) -- (0.8571429,1.928571) -- cycle;
\fill[teal!15!white] (0.8571429,1.714286) -- (0.5714286,1.928571) -- (0.5714286,1.714286) -- cycle;
\fill[teal!16!white] (0.5714286,1.928571) -- (0.8571429,1.714286) -- (0.8571429,1.928571) -- cycle;
\fill[teal!41!white] (0.2857143,2.357143) -- (0.5714286,2.357143) -- (0.2857143,2.571429) -- cycle;
\fill[teal!59!white] (0.5714286,2.357143) -- (0.5714286,2.571429) -- (0.2857143,2.571429) -- cycle;
\fill[teal!58!white] (0,0.6428571) -- (0,0.4285714) -- (0.2857143,0.6428571) -- cycle;
\fill[teal!37!white] (0,0.4285714) -- (0.2857143,0.4285714) -- (0.2857143,0.6428571) -- cycle;
\fill[teal!71!white] (0,0.4285714) -- (0,0.2142857) -- (0.2857143,0.2142857) -- cycle;
\fill[teal!43!white] (0.2857143,0.4285714) -- (0,0.4285714) -- (0.2857143,0.2142857) -- cycle;
\fill[teal!42!white] (2.285714,1.285714) -- (2,1.285714) -- (2.285714,1.071429) -- cycle;
\fill[teal!24!white] (2,1.285714) -- (2,1.071429) -- (2.285714,1.071429) -- cycle;
\fill[teal!17!white] (2,1.071429) -- (2,0.8571429) -- (2.285714,1.071429) -- cycle;
\fill[teal!18!white] (2,0.8571429) -- (2.285714,0.8571429) -- (2.285714,1.071429) -- cycle;
\fill[teal!26!white] (2.285714,0.8571429) -- (2.571429,0.8571429) -- (2.285714,1.071429) -- cycle;
\fill[teal!44!white] (2.571429,0.8571429) -- (2.571429,1.071429) -- (2.285714,1.071429) -- cycle;
\fill[teal!16!white] (1.428571,0.4285714) -- (1.714286,0.2142857) -- (1.714286,0.4285714) -- cycle;
\fill[teal!20!white] (1.714286,0.2142857) -- (1.428571,0.4285714) -- (1.428571,0.2142857) -- cycle;
\fill[teal!64!white] (3.142857,0.2142857) -- (3.428571,0) -- (3.428571,0.2142857) -- cycle;
\fill[teal!88!white] (3.142857,0.2142857) -- (3.142857,0) -- (3.428571,0) -- cycle;
\fill[teal!83!white] (3.142857,0) -- (3.142857,0.2142857) -- (2.857143,0) -- cycle;
\fill[teal!49!white] (3.142857,0.2142857) -- (2.857143,0.2142857) -- (2.857143,0) -- cycle;
\fill[teal!36!white] (1.714286,0) -- (2,0.2142857) -- (1.714286,0.2142857) -- cycle;
\fill[teal!63!white] (2,0.2142857) -- (1.714286,0) -- (2,0) -- cycle;
\fill[teal!29!white] (0.5714286,0.2142857) -- (0.5714286,0.4285714) -- (0.2857143,0.2142857) -- cycle;
\fill[teal!28!white] (0.5714286,0.4285714) -- (0.2857143,0.4285714) -- (0.2857143,0.2142857) -- cycle;
\fill[teal!23!white] (0.8571429,0.4285714) -- (0.5714286,0.2142857) -- (0.8571429,0.2142857) -- cycle;
\fill[teal!21!white] (0.8571429,0.4285714) -- (0.5714286,0.4285714) -- (0.5714286,0.2142857) -- cycle;
\fill[teal!16!white] (1.142857,1.5) -- (1.428571,1.5) -- (1.428571,1.714286) -- cycle;
\fill[teal!17!white] (1.142857,1.714286) -- (1.142857,1.5) -- (1.428571,1.714286) -- cycle;
\fill[teal!24!white] (0.5714286,2.142857) -- (0.2857143,2.142857) -- (0.2857143,1.928571) -- cycle;
\fill[teal!21!white] (0.5714286,1.928571) -- (0.5714286,2.142857) -- (0.2857143,1.928571) -- cycle;
\fill[teal!28!white] (0.5714286,2.142857) -- (0.2857143,2.357143) -- (0.2857143,2.142857) -- cycle;
\fill[teal!31!white] (0.5714286,2.142857) -- (0.5714286,2.357143) -- (0.2857143,2.357143) -- cycle;
\fill[teal!24!white] (0.8571429,2.142857) -- (0.5714286,2.142857) -- (0.8571429,1.928571) -- cycle;
\fill[teal!20!white] (0.5714286,2.142857) -- (0.5714286,1.928571) -- (0.8571429,1.928571) -- cycle;
\fill[teal!36!white] (0.5714286,2.142857) -- (0.8571429,2.142857) -- (0.8571429,2.357143) -- cycle;
\fill[teal!38!white] (0.5714286,2.357143) -- (0.5714286,2.142857) -- (0.8571429,2.357143) -- cycle;
\fill[teal!30!white] (1.714286,1.285714) -- (2,1.285714) -- (2,1.5) -- cycle;
\fill[teal!29!white] (1.714286,1.5) -- (1.714286,1.285714) -- (2,1.5) -- cycle;
\fill[teal!8!white] (1.142857,1.071429) -- (1.428571,0.8571429) -- (1.428571,1.071429) -- cycle;
\fill[teal!8!white] (1.428571,0.8571429) -- (1.142857,1.071429) -- (1.142857,0.8571429) -- cycle;
\fill[teal!47!white] (3.142857,0.4285714) -- (3.142857,0.2142857) -- (3.428571,0.2142857) -- cycle;
\fill[teal!66!white] (3.428571,0.4285714) -- (3.142857,0.4285714) -- (3.428571,0.2142857) -- cycle;
\fill[teal!35!white] (2.857143,0.4285714) -- (3.142857,0.4285714) -- (2.857143,0.6428571) -- cycle;
\fill[teal!54!white] (3.142857,0.4285714) -- (3.142857,0.6428571) -- (2.857143,0.6428571) -- cycle;
\fill[teal!37!white] (3.142857,0.2142857) -- (3.142857,0.4285714) -- (2.857143,0.2142857) -- cycle;
\fill[teal!32!white] (3.142857,0.4285714) -- (2.857143,0.4285714) -- (2.857143,0.2142857) -- cycle;
\fill[teal!35!white] (2.571429,0.8571429) -- (2.571429,0.6428571) -- (2.857143,0.8571429) -- cycle;
\fill[teal!36!white] (2.857143,0.8571429) -- (2.571429,0.6428571) -- (2.857143,0.6428571) -- cycle;
\fill[teal!65!white] (2.285714,0) -- (2.285714,0.2142857) -- (2,0) -- cycle;
\fill[teal!38!white] (2.285714,0.2142857) -- (2,0.2142857) -- (2,0) -- cycle;
\fill[teal!22!white] (2.285714,0.2142857) -- (2,0.4285714) -- (2,0.2142857) -- cycle;
\fill[teal!19!white] (2,0.4285714) -- (2.285714,0.2142857) -- (2.285714,0.4285714) -- cycle;
\fill[teal!13!white] (2,0.4285714) -- (1.714286,0.6428571) -- (1.714286,0.4285714) -- cycle;
\fill[teal!12!white] (1.714286,0.6428571) -- (2,0.4285714) -- (2,0.6428571) -- cycle;
\fill[teal!17!white] (1.714286,0.2142857) -- (2,0.4285714) -- (1.714286,0.4285714) -- cycle;
\fill[teal!21!white] (2,0.2142857) -- (2,0.4285714) -- (1.714286,0.2142857) -- cycle;
\fill[teal!54!white] (0.2857143,0.8571429) -- (0,0.8571429) -- (0,0.6428571) -- cycle;
\fill[teal!33!white] (0.2857143,0.8571429) -- (0,0.6428571) -- (0.2857143,0.6428571) -- cycle;
\fill[teal!31!white] (0,1.071429) -- (0.2857143,0.8571429) -- (0.2857143,1.071429) -- cycle;
\fill[teal!53!white] (0,0.8571429) -- (0.2857143,0.8571429) -- (0,1.071429) -- cycle;
\fill[teal!14!white] (0.2857143,1.5) -- (0.5714286,1.285714) -- (0.5714286,1.5) -- cycle;
\fill[teal!17!white] (0.5714286,1.285714) -- (0.2857143,1.5) -- (0.2857143,1.285714) -- cycle;
\fill[teal!19!white] (0.5714286,0.6428571) -- (0.2857143,0.8571429) -- (0.2857143,0.6428571) -- cycle;
\fill[teal!15!white] (0.2857143,0.8571429) -- (0.5714286,0.6428571) -- (0.5714286,0.8571429) -- cycle;
\fill[teal!21!white] (0.2857143,0.4285714) -- (0.5714286,0.6428571) -- (0.2857143,0.6428571) -- cycle;
\fill[teal!20!white] (0.5714286,0.4285714) -- (0.5714286,0.6428571) -- (0.2857143,0.4285714) -- cycle;
\fill[teal!16!white] (0.8571429,0.4285714) -- (0.5714286,0.6428571) -- (0.5714286,0.4285714) -- cycle;
\fill[teal!13!white] (0.5714286,0.6428571) -- (0.8571429,0.4285714) -- (0.8571429,0.6428571) -- cycle;
\fill[teal!20!white] (1.142857,0.4285714) -- (1.142857,0.2142857) -- (1.428571,0.2142857) -- cycle;
\fill[teal!16!white] (1.428571,0.4285714) -- (1.142857,0.4285714) -- (1.428571,0.2142857) -- cycle;
\fill[teal!21!white] (1.142857,0.2142857) -- (1.142857,0.4285714) -- (0.8571429,0.2142857) -- cycle;
\fill[teal!18!white] (1.142857,0.4285714) -- (0.8571429,0.4285714) -- (0.8571429,0.2142857) -- cycle;
\fill[teal!13!white] (0.8571429,0.4285714) -- (1.142857,0.4285714) -- (0.8571429,0.6428571) -- cycle;
\fill[teal!11!white] (1.142857,0.4285714) -- (1.142857,0.6428571) -- (0.8571429,0.6428571) -- cycle;
\fill[teal!16!white] (1.428571,1.285714) -- (1.714286,1.5) -- (1.428571,1.5) -- cycle;
\fill[teal!16!white] (1.428571,1.285714) -- (1.714286,1.285714) -- (1.714286,1.5) -- cycle;
\fill[teal!12!white] (1.142857,1.5) -- (1.428571,1.285714) -- (1.428571,1.5) -- cycle;
\fill[teal!10!white] (1.428571,1.285714) -- (1.142857,1.5) -- (1.142857,1.285714) -- cycle;
\fill[teal!9!white] (1.428571,1.285714) -- (1.142857,1.071429) -- (1.428571,1.071429) -- cycle;
\fill[teal!9!white] (1.142857,1.071429) -- (1.428571,1.285714) -- (1.142857,1.285714) -- cycle;
\fill[teal!12!white] (1.428571,0.6428571) -- (1.428571,0.4285714) -- (1.714286,0.4285714) -- cycle;
\fill[teal!11!white] (1.714286,0.6428571) -- (1.428571,0.6428571) -- (1.714286,0.4285714) -- cycle;
\fill[teal!12!white] (1.428571,0.6428571) -- (1.142857,0.4285714) -- (1.428571,0.4285714) -- cycle;
\fill[teal!11!white] (1.142857,0.4285714) -- (1.428571,0.6428571) -- (1.142857,0.6428571) -- cycle;
\fill[teal!9!white] (1.428571,0.6428571) -- (1.428571,0.8571429) -- (1.142857,0.8571429) -- cycle;
\fill[teal!9!white] (1.142857,0.6428571) -- (1.428571,0.6428571) -- (1.142857,0.8571429) -- cycle;
\fill[teal!16!white] (2,1.285714) -- (1.714286,1.071429) -- (2,1.071429) -- cycle;
\fill[teal!16!white] (1.714286,1.285714) -- (1.714286,1.071429) -- (2,1.285714) -- cycle;
\fill[teal!12!white] (1.428571,1.285714) -- (1.714286,1.071429) -- (1.714286,1.285714) -- cycle;
\fill[teal!10!white] (1.714286,1.071429) -- (1.428571,1.285714) -- (1.428571,1.071429) -- cycle;
\fill[teal!20!white] (2.285714,0.6428571) -- (2.571429,0.8571429) -- (2.285714,0.8571429) -- cycle;
\fill[teal!21!white] (2.285714,0.6428571) -- (2.571429,0.6428571) -- (2.571429,0.8571429) -- cycle;
\fill[teal!15!white] (2,0.8571429) -- (2.285714,0.6428571) -- (2.285714,0.8571429) -- cycle;
\fill[teal!13!white] (2.285714,0.6428571) -- (2,0.8571429) -- (2,0.6428571) -- cycle;
\fill[teal!15!white] (2.285714,0.6428571) -- (2,0.4285714) -- (2.285714,0.4285714) -- cycle;
\fill[teal!14!white] (2,0.4285714) -- (2.285714,0.6428571) -- (2,0.6428571) -- cycle;
\fill[teal!26!white] (2.571429,0.4285714) -- (2.857143,0.4285714) -- (2.857143,0.6428571) -- cycle;
\fill[teal!24!white] (2.571429,0.6428571) -- (2.571429,0.4285714) -- (2.857143,0.6428571) -- cycle;
\fill[teal!19!white] (2.285714,0.6428571) -- (2.571429,0.4285714) -- (2.571429,0.6428571) -- cycle;
\fill[teal!18!white] (2.571429,0.4285714) -- (2.285714,0.6428571) -- (2.285714,0.4285714) -- cycle;
\fill[teal!17!white] (0.2857143,0.8571429) -- (0.5714286,1.071429) -- (0.2857143,1.071429) -- cycle;
\fill[teal!15!white] (0.5714286,1.071429) -- (0.2857143,0.8571429) -- (0.5714286,0.8571429) -- cycle;
\fill[teal!17!white] (0.5714286,1.071429) -- (0.2857143,1.285714) -- (0.2857143,1.071429) -- cycle;
\fill[teal!14!white] (0.5714286,1.071429) -- (0.5714286,1.285714) -- (0.2857143,1.285714) -- cycle;
\fill[teal!12!white] (0.8571429,1.5) -- (0.5714286,1.714286) -- (0.5714286,1.5) -- cycle;
\fill[teal!12!white] (0.8571429,1.5) -- (0.8571429,1.714286) -- (0.5714286,1.714286) -- cycle;
\fill[teal!13!white] (0.8571429,1.714286) -- (0.8571429,1.5) -- (1.142857,1.714286) -- cycle;
\fill[teal!12!white] (0.8571429,1.5) -- (1.142857,1.5) -- (1.142857,1.714286) -- cycle;
\fill[teal!11!white] (1.714286,0.8571429) -- (1.714286,0.6428571) -- (2,0.6428571) -- cycle;
\fill[teal!11!white] (2,0.8571429) -- (1.714286,0.8571429) -- (2,0.6428571) -- cycle;
\fill[teal!12!white] (1.714286,0.8571429) -- (2,0.8571429) -- (2,1.071429) -- cycle;
\fill[teal!12!white] (1.714286,1.071429) -- (1.714286,0.8571429) -- (2,1.071429) -- cycle;
\fill[teal!10!white] (1.714286,0.8571429) -- (1.428571,0.6428571) -- (1.714286,0.6428571) -- cycle;
\fill[teal!9!white] (1.428571,0.6428571) -- (1.714286,0.8571429) -- (1.428571,0.8571429) -- cycle;
\fill[teal!9!white] (1.428571,0.8571429) -- (1.714286,0.8571429) -- (1.428571,1.071429) -- cycle;
\fill[teal!10!white] (1.714286,0.8571429) -- (1.714286,1.071429) -- (1.428571,1.071429) -- cycle;
\fill[teal!62!white] (2.571429,0) -- (2.571429,0.2142857) -- (2.285714,0) -- cycle;
\fill[teal!40!white] (2.571429,0.2142857) -- (2.285714,0.2142857) -- (2.285714,0) -- cycle;
\fill[teal!24!white] (2.285714,0.2142857) -- (2.571429,0.2142857) -- (2.285714,0.4285714) -- cycle;
\fill[teal!22!white] (2.571429,0.2142857) -- (2.571429,0.4285714) -- (2.285714,0.4285714) -- cycle;
\fill[teal!45!white] (2.857143,0.2142857) -- (2.571429,0.2142857) -- (2.857143,0) -- cycle;
\fill[teal!68!white] (2.571429,0.2142857) -- (2.571429,0) -- (2.857143,0) -- cycle;
\fill[teal!29!white] (2.857143,0.4285714) -- (2.571429,0.2142857) -- (2.857143,0.2142857) -- cycle;
\fill[teal!25!white] (2.571429,0.4285714) -- (2.571429,0.2142857) -- (2.857143,0.4285714) -- cycle;
\fill[teal!10!white] (1.142857,0.6428571) -- (0.8571429,0.8571429) -- (0.8571429,0.6428571) -- cycle;
\fill[teal!9!white] (0.8571429,0.8571429) -- (1.142857,0.6428571) -- (1.142857,0.8571429) -- cycle;
\fill[teal!12!white] (0.5714286,0.6428571) -- (0.8571429,0.8571429) -- (0.5714286,0.8571429) -- cycle;
\fill[teal!11!white] (0.8571429,0.8571429) -- (0.5714286,0.6428571) -- (0.8571429,0.6428571) -- cycle;
\fill[teal!9!white] (1.142857,1.5) -- (0.8571429,1.285714) -- (1.142857,1.285714) -- cycle;
\fill[teal!10!white] (0.8571429,1.5) -- (0.8571429,1.285714) -- (1.142857,1.5) -- cycle;
\fill[teal!11!white] (0.5714286,1.285714) -- (0.8571429,1.285714) -- (0.5714286,1.5) -- cycle;
\fill[teal!10!white] (0.8571429,1.285714) -- (0.8571429,1.5) -- (0.5714286,1.5) -- cycle;
\fill[teal!8!white] (1.142857,1.071429) -- (0.8571429,1.071429) -- (1.142857,0.8571429) -- cycle;
\fill[teal!9!white] (0.8571429,1.071429) -- (0.8571429,0.8571429) -- (1.142857,0.8571429) -- cycle;
\fill[teal!8!white] (0.8571429,1.071429) -- (1.142857,1.071429) -- (1.142857,1.285714) -- cycle;
\fill[teal!9!white] (0.8571429,1.285714) -- (0.8571429,1.071429) -- (1.142857,1.285714) -- cycle;
\fill[teal!11!white] (0.8571429,1.071429) -- (0.5714286,1.071429) -- (0.5714286,0.8571429) -- cycle;
\fill[teal!10!white] (0.8571429,0.8571429) -- (0.8571429,1.071429) -- (0.5714286,0.8571429) -- cycle;
\fill[teal!10!white] (0.5714286,1.071429) -- (0.8571429,1.071429) -- (0.5714286,1.285714) -- cycle;
\fill[teal!10!white] (0.8571429,1.071429) -- (0.8571429,1.285714) -- (0.5714286,1.285714) -- cycle;
\draw[teal,thick] (0,0) -- (4,0) -- (0,3) -- cycle;
\draw[magenta!70,thin] (1.417142,1.002754) -- (3.6,2.35) -- (4.1,2.35);\filldraw[fill=magenta,draw=white,line width=.45pt] (1.417142,1.002754) circle(1.5pt) node[above right,font=\scriptsize,text=magenta,fill=white,inner sep=.6pt] {};
\node[anchor=west,text=magenta,font=\small] at(4.12,2.35){E3};\draw[gold!70,thin] (1,1) -- (3.6,1.8) -- (4.1,1.8);\filldraw[fill=gold,draw=white,line width=.45pt] (1,1) circle(1.5pt) node[above right,font=\scriptsize,text=gold,fill=white,inner sep=.6pt] {};
\node[anchor=west,text=gold,font=\small] at(4.12,1.8){$I$};\draw[navy!70,thin] (1.333333,1) -- (3.6,1.25) -- (4.1,1.25);\filldraw[fill=navy,draw=white,line width=.45pt] (1.333333,1) circle(1.5pt) node[above right,font=\scriptsize,text=navy,fill=white,inner sep=.6pt] {};
\node[anchor=west,text=navy,font=\small] at(4.12,1.25){$G$};\end{tikzpicture}
\caption[Projected equilibrium plate-charge density]{Projected equilibrium plate-charge density. Edge singularities are unresolved by this display mesh. Fields are archived numerical illustrations; teal saturation increases with the displayed value.}\label{fig:field-E3}
\end{figure}

\paragraph{Demonstration and investigation.} Compare the marked charge means with G and W. Emphasize that color scales differ between charge per unit length and charge per unit area.


\section{E0: equal probability of first exit}\label{sec:slide20}

The equal-exit center E0 has an intuitive stochastic definition. Start an absorbing Brownian particle at a point in the triangle and record the side it first reaches. At E0 all three exit probabilities agree. In two dimensions these probabilities are harmonic measures and are conformally invariant. A Schwarz--Christoffel map sends the upper half-plane to the triangle. The point one half plus i times square root of three over two sees the three boundary intervals equally. Mapping it gives E0. Equal exit probabilities do not imply equal distances to the sides.


\[\Pr_{\rm E0}\{\hbox{first exit through side }i\}=1/3.\]



Solve $\Delta h_i=0$ with boundary value one on side $i$ and zero on
the others. Brownian exit probability is $h_i(p)$ and their sum is
one. E0 solves $h_i=1/3$.
Map the triangle to a disk with the trial point at the center.
Harmonic measure becomes normalized arc length, so the three sides
must correspond to arcs of angular length $2\pi/3$.
In the upper half-plane a real interval's measure is its viewing
angle divided by $\pi$. This gives a small conformal-coordinate
problem followed by a Schwarz--Christoffel evaluation
\cite{iannaccone}.




\paragraph{Demonstration and investigation.} Reveal three equal-probability exit fans. Save an actual Brownian simulation for lecture 4 or the companion demo. Mention now that the new certificate disproves E0 attractivity.


\section{E1: least boundary influence}\label{sec:slide21}

Map the unit disk conformally onto the triangle with its origin sent to p. The derivative magnitude at the origin is the conformal radius. It measures a local scale relative to the surrounding boundary. Maximizing it gives E1. Equivalently, a tiny electrode of fixed radius placed inside a grounded triangular enclosure has an asymptotic capacity governed by the regular part of the Green function. This is a point-probe limit, with equal probe footprints. A freely moving test charge sees an induced self-energy with the opposite sign, so E1 is not a stable electrostatic trap.


\[r_T(p)=|F_p^\prime(0)|,\quad F_p:\mathbb D\to T,\quad F_p(0)=p;\quad {\rm E1}=\argmax_T r_T.\]



Let $F:\mathbb H\to T$ send $0,1,\infty$ to $A,B,C$. Put
$\alpha=\theta_A/\pi$, $\beta=\theta_B/\pi$, $\gamma=\theta_C/\pi$.
Then
\[
F(z)=A+(B-A)\frac{B_z(\alpha,\beta)}{B(\alpha,\beta)},\qquad
B_z(\alpha,\beta)=\int_0^z t^{\alpha-1}(1-t)^{\beta-1}\dd t,
\]
with branches continued through $\mathbb H$.
The radius is $r_T(F(z))=2\Im z\,|F'(z)|$.
Differentiating its logarithm gives the unique stationary maximum
\[
z_*=\frac{\alpha}{\alpha+\beta}
+i\frac{\sqrt{\alpha\beta/\gamma}}{\alpha+\beta},
\qquad \mathrm{E1}=F(z_*).
\]
The radius tends to zero at the boundary. Accurate incomplete beta
functions are essential for thin triangles.

A small circular electrode of radius $\varepsilon$ inside a grounded
triangle has relative capacity asymptotic to
$2\pi/\log(r_T(p)/\varepsilon)$ in convenient units. The least-capacity
point maximizes radius.
For $G_T(p,q)=-(2\pi)^{-1}\log|p-q|+H_T(p,q)$,
$H_T(p,p)=(2\pi)^{-1}\log r_T(p)$.
Robin functions sometimes have the opposite sign; declare conventions.
In higher dimensions the analogous correction uses the regular part
of the Newtonian Dirichlet Green function, with an explicitly chosen
domain and electrode regularization.



\begin{figure}[htbp]\centering
\begin{tikzpicture}[x=1.05cm,y=1.05cm]
\draw[teal,thick](0,0)--(2,0)--(0,1.5)--cycle;\draw[teal,thick](4,1)circle(1);\draw[->,thick](2.2,.8)--(2.8,.8)node[midway,above]{$f$};\draw[<-,magenta,thick](2.2,.4)--(2.8,.4)node[midway,below]{$F$};\fill[magenta](4,1)circle(1.7pt);\node at(4,-.3){unit disk};\node at(1,-.3){triangle};\end{tikzpicture}
\caption[A conformal map takes the trial point to the disk center]{A conformal map takes the trial point to the disk center. The inverse derivative supplies conformal radius.}\label{fig:conformal}
\end{figure}

\section{A current around the perimeter}\label{sec:slide22}

Apply Biot--Savart's law to each straight side. Resolve the
observer position into perpendicular side distance and along-side
coordinate. All three contributions have the same normal sign
inside an oriented loop; near a side the field diverges as $1/d$.
The next section evaluates the segment integral explicitly.


\[\dd\mathbf B=\frac{\mu_0I}{4\pi}\frac{\dd\boldsymbol\ell\times\mathbf r}{|\mathbf r|^3}.\]


Integrating $d/[d^2+(s-t)^2]^{3/2}$ along a side yields
$(t-s)/[d\sqrt{d^2+(t-s)^2}]$ as antiderivative. Restore the common
factor $\mu_0I/(4\pi)$. The next section writes the complete side
contribution explicitly.



\paragraph{Demonstration and investigation.} Trace a current arrow around the editable wire. Indicate a small normal probe above the plane rather than proposing contact with an ideal zero-thickness wire.


\section{M1: the weakest interior loop field}\label{sec:slide23}

The loop field becomes large near the boundary, so its minimum occurs inside. For this triangular problem the known convexity argument gives a unique minimum. Numerically, sum exact segment fields, differentiate, and find the interior stationary point. Compare its Hessian and global behavior instead of trusting the first optimizer result. M1 is a natural single point but it is not globally vertex-attractive: an interval-certified finite-motion counterexample already exists. The point marks low field, not a generic point-charge mechanical trap.


\[\mathrm{M1}=\argmin_{p\in T^\circ}|B_z(p)|.\]



Orient the loop counterclockwise. Each interior side contribution is
positive. If $d>0$ is side distance and $t_\pm$ are signed endpoint
coordinates measured from the perpendicular foot,
\[
B_{z,\rm side}=\frac{\mu_0I}{4\pi d}
\left[\frac{t_+}{\sqrt{d^2+t_+^2}}-
\frac{t_-}{\sqrt{d^2+t_-^2}}\right].
\]
Integrate $d\dd t/(d^2+t^2)^{3/2}$ in Biot--Savart's law to obtain
this expression. Summing three terms gives a smooth interior field
diverging at the sides. Since $B_z>0$, its minimum also minimizes
$B_z^2$.
Use analytic gradients, an interior bracket and a Hessian or global
uniqueness argument, rather than one optimizer run.
The unique loop minimum is established in the project analysis.
Its fixed-triangle stability coexists with a certified negative
vertex motion.



\begin{figure}[htbp]\centering
\begin{tikzpicture}[x=1.7cm,y=1.7cm]
\fill[teal!94!white] (3.142857,0.6428571) -- (2.857143,0.8571429) -- (2.857143,0.6428571) -- cycle;
\fill[teal!91!white] (1.714286,1.714286) -- (1.428571,1.928571) -- (1.428571,1.714286) -- cycle;
\fill[teal!91!white] (1.714286,1.714286) -- (1.714286,1.5) -- (2,1.5) -- cycle;
\fill[teal!94!white] (0.5714286,2.571429) -- (0.2857143,2.785714) -- (0.2857143,2.571429) -- cycle;
\fill[teal!92!white] (2.571429,1.071429) -- (2.285714,1.285714) -- (2.285714,1.071429) -- cycle;
\fill[teal!94!white] (3.714286,0.2142857) -- (3.428571,0.4285714) -- (3.428571,0.2142857) -- cycle;
\fill[teal!94!white] (3.714286,0.2142857) -- (3.714286,0) -- (4,0) -- cycle;
\fill[teal!92!white] (0.8571429,2.142857) -- (1.142857,2.142857) -- (0.8571429,2.357143) -- cycle;
\fill[teal!92!white] (1.142857,2.142857) -- (1.142857,1.928571) -- (1.428571,1.928571) -- cycle;
\fill[teal!94!white] (0,2.785714) -- (0.2857143,2.785714) -- (0,3) -- cycle;
\fill[teal!94!white] (0.5714286,2.571429) -- (0.5714286,2.357143) -- (0.8571429,2.357143) -- cycle;
\fill[teal!91!white] (2,1.285714) -- (2.285714,1.285714) -- (2,1.5) -- cycle;
\fill[teal!93!white] (2.571429,1.071429) -- (2.571429,0.8571429) -- (2.857143,0.8571429) -- cycle;
\fill[teal!94!white] (3.142857,0.4285714) -- (3.428571,0.4285714) -- (3.142857,0.6428571) -- cycle;
\fill[teal!94!white] (0.2857143,0) -- (0,0.2142857) -- (0,0) -- cycle;
\fill[teal!52!white] (0,0.2142857) -- (0.2857143,0) -- (0.2857143,0.2142857) -- cycle;
\fill[teal!76!white] (0.5714286,0.2142857) -- (0.2857143,0) -- (0.5714286,0) -- cycle;
\fill[teal!41!white] (0.2857143,0) -- (0.5714286,0.2142857) -- (0.2857143,0.2142857) -- cycle;
\fill[teal!45!white] (1.714286,1.5) -- (1.714286,1.714286) -- (1.428571,1.714286) -- cycle;
\fill[teal!23!white] (1.428571,1.5) -- (1.714286,1.5) -- (1.428571,1.714286) -- cycle;
\fill[teal!56!white] (0,1.928571) -- (0,1.714286) -- (0.2857143,1.928571) -- cycle;
\fill[teal!29!white] (0,1.714286) -- (0.2857143,1.714286) -- (0.2857143,1.928571) -- cycle;
\fill[teal!62!white] (0.2857143,2.357143) -- (0,2.571429) -- (0,2.357143) -- cycle;
\fill[teal!39!white] (0,2.571429) -- (0.2857143,2.357143) -- (0.2857143,2.571429) -- cycle;
\fill[teal!59!white] (0,2.142857) -- (0.2857143,2.357143) -- (0,2.357143) -- cycle;
\fill[teal!32!white] (0.2857143,2.357143) -- (0,2.142857) -- (0.2857143,2.142857) -- cycle;
\fill[teal!57!white] (0,2.142857) -- (0,1.928571) -- (0.2857143,1.928571) -- cycle;
\fill[teal!30!white] (0.2857143,2.142857) -- (0,2.142857) -- (0.2857143,1.928571) -- cycle;
\fill[teal!16!white] (0.2857143,1.714286) -- (0.5714286,1.928571) -- (0.2857143,1.928571) -- cycle;
\fill[teal!14!white] (0.5714286,1.714286) -- (0.5714286,1.928571) -- (0.2857143,1.714286) -- cycle;
\fill[teal!55!white] (0.2857143,1.5) -- (0,1.714286) -- (0,1.5) -- cycle;
\fill[teal!28!white] (0,1.714286) -- (0.2857143,1.5) -- (0.2857143,1.714286) -- cycle;
\fill[teal!15!white] (0.2857143,1.5) -- (0.5714286,1.714286) -- (0.2857143,1.714286) -- cycle;
\fill[teal!13!white] (0.5714286,1.714286) -- (0.2857143,1.5) -- (0.5714286,1.5) -- cycle;
\fill[teal!55!white] (0,1.285714) -- (0.2857143,1.5) -- (0,1.5) -- cycle;
\fill[teal!28!white] (0.2857143,1.5) -- (0,1.285714) -- (0.2857143,1.285714) -- cycle;
\fill[teal!55!white] (0,1.285714) -- (0,1.071429) -- (0.2857143,1.071429) -- cycle;
\fill[teal!28!white] (0.2857143,1.285714) -- (0,1.285714) -- (0.2857143,1.071429) -- cycle;
\fill[teal!90!white] (3.428571,0) -- (3.714286,0) -- (3.428571,0.2142857) -- cycle;
\fill[teal!76!white] (3.714286,0) -- (3.714286,0.2142857) -- (3.428571,0.2142857) -- cycle;
\fill[teal!73!white] (1.428571,0) -- (1.714286,0) -- (1.428571,0.2142857) -- cycle;
\fill[teal!36!white] (1.714286,0) -- (1.714286,0.2142857) -- (1.428571,0.2142857) -- cycle;
\fill[teal!74!white] (0.8571429,0) -- (0.5714286,0.2142857) -- (0.5714286,0) -- cycle;
\fill[teal!37!white] (0.5714286,0.2142857) -- (0.8571429,0) -- (0.8571429,0.2142857) -- cycle;
\fill[teal!73!white] (1.142857,0.2142857) -- (1.142857,0) -- (1.428571,0) -- cycle;
\fill[teal!36!white] (1.142857,0.2142857) -- (1.428571,0) -- (1.428571,0.2142857) -- cycle;
\fill[teal!73!white] (1.142857,0.2142857) -- (0.8571429,0) -- (1.142857,0) -- cycle;
\fill[teal!37!white] (0.8571429,0) -- (1.142857,0.2142857) -- (0.8571429,0.2142857) -- cycle;
\fill[teal!23!white] (1.142857,1.928571) -- (1.142857,1.714286) -- (1.428571,1.714286) -- cycle;
\fill[teal!45!white] (1.428571,1.928571) -- (1.142857,1.928571) -- (1.428571,1.714286) -- cycle;
\fill[teal!46!white] (0.8571429,2.142857) -- (1.142857,1.928571) -- (1.142857,2.142857) -- cycle;
\fill[teal!24!white] (1.142857,1.928571) -- (0.8571429,2.142857) -- (0.8571429,1.928571) -- cycle;
\fill[teal!70!white] (0,2.785714) -- (0,2.571429) -- (0.2857143,2.571429) -- cycle;
\fill[teal!65!white] (0.2857143,2.785714) -- (0,2.785714) -- (0.2857143,2.571429) -- cycle;
\fill[teal!16!white] (1.142857,1.928571) -- (0.8571429,1.714286) -- (1.142857,1.714286) -- cycle;
\fill[teal!16!white] (0.8571429,1.714286) -- (1.142857,1.928571) -- (0.8571429,1.928571) -- cycle;
\fill[teal!12!white] (0.8571429,1.714286) -- (0.5714286,1.928571) -- (0.5714286,1.714286) -- cycle;
\fill[teal!14!white] (0.5714286,1.928571) -- (0.8571429,1.714286) -- (0.8571429,1.928571) -- cycle;
\fill[teal!31!white] (0.2857143,2.357143) -- (0.5714286,2.357143) -- (0.2857143,2.571429) -- cycle;
\fill[teal!51!white] (0.5714286,2.357143) -- (0.5714286,2.571429) -- (0.2857143,2.571429) -- cycle;
\fill[teal!57!white] (0,0.6428571) -- (0,0.4285714) -- (0.2857143,0.6428571) -- cycle;
\fill[teal!31!white] (0,0.4285714) -- (0.2857143,0.4285714) -- (0.2857143,0.6428571) -- cycle;
\fill[teal!62!white] (0,0.4285714) -- (0,0.2142857) -- (0.2857143,0.2142857) -- cycle;
\fill[teal!33!white] (0.2857143,0.4285714) -- (0,0.4285714) -- (0.2857143,0.2142857) -- cycle;
\fill[teal!46!white] (2.285714,1.285714) -- (2,1.285714) -- (2.285714,1.071429) -- cycle;
\fill[teal!23!white] (2,1.285714) -- (2,1.071429) -- (2.285714,1.071429) -- cycle;
\fill[teal!16!white] (2,1.071429) -- (2,0.8571429) -- (2.285714,1.071429) -- cycle;
\fill[teal!16!white] (2,0.8571429) -- (2.285714,0.8571429) -- (2.285714,1.071429) -- cycle;
\fill[teal!24!white] (2.285714,0.8571429) -- (2.571429,0.8571429) -- (2.285714,1.071429) -- cycle;
\fill[teal!46!white] (2.571429,0.8571429) -- (2.571429,1.071429) -- (2.285714,1.071429) -- cycle;
\fill[teal!15!white] (1.428571,0.4285714) -- (1.714286,0.2142857) -- (1.714286,0.4285714) -- cycle;
\fill[teal!18!white] (1.714286,0.2142857) -- (1.428571,0.4285714) -- (1.428571,0.2142857) -- cycle;
\fill[teal!49!white] (3.142857,0.2142857) -- (3.428571,0) -- (3.428571,0.2142857) -- cycle;
\fill[teal!81!white] (3.142857,0.2142857) -- (3.142857,0) -- (3.428571,0) -- cycle;
\fill[teal!78!white] (3.142857,0) -- (3.142857,0.2142857) -- (2.857143,0) -- cycle;
\fill[teal!41!white] (3.142857,0.2142857) -- (2.857143,0.2142857) -- (2.857143,0) -- cycle;
\fill[teal!37!white] (1.714286,0) -- (2,0.2142857) -- (1.714286,0.2142857) -- cycle;
\fill[teal!73!white] (2,0.2142857) -- (1.714286,0) -- (2,0) -- cycle;
\fill[teal!22!white] (0.5714286,0.2142857) -- (0.5714286,0.4285714) -- (0.2857143,0.2142857) -- cycle;
\fill[teal!21!white] (0.5714286,0.4285714) -- (0.2857143,0.4285714) -- (0.2857143,0.2142857) -- cycle;
\fill[teal!19!white] (0.8571429,0.4285714) -- (0.5714286,0.2142857) -- (0.8571429,0.2142857) -- cycle;
\fill[teal!17!white] (0.8571429,0.4285714) -- (0.5714286,0.4285714) -- (0.5714286,0.2142857) -- cycle;
\fill[teal!15!white] (1.142857,1.5) -- (1.428571,1.5) -- (1.428571,1.714286) -- cycle;
\fill[teal!15!white] (1.142857,1.714286) -- (1.142857,1.5) -- (1.428571,1.714286) -- cycle;
\fill[teal!19!white] (0.5714286,2.142857) -- (0.2857143,2.142857) -- (0.2857143,1.928571) -- cycle;
\fill[teal!16!white] (0.5714286,1.928571) -- (0.5714286,2.142857) -- (0.2857143,1.928571) -- cycle;
\fill[teal!21!white] (0.5714286,2.142857) -- (0.2857143,2.357143) -- (0.2857143,2.142857) -- cycle;
\fill[teal!24!white] (0.5714286,2.142857) -- (0.5714286,2.357143) -- (0.2857143,2.357143) -- cycle;
\fill[teal!20!white] (0.8571429,2.142857) -- (0.5714286,2.142857) -- (0.8571429,1.928571) -- cycle;
\fill[teal!16!white] (0.5714286,2.142857) -- (0.5714286,1.928571) -- (0.8571429,1.928571) -- cycle;
\fill[teal!32!white] (0.5714286,2.142857) -- (0.8571429,2.142857) -- (0.8571429,2.357143) -- cycle;
\fill[teal!33!white] (0.5714286,2.357143) -- (0.5714286,2.142857) -- (0.8571429,2.357143) -- cycle;
\fill[teal!30!white] (1.714286,1.285714) -- (2,1.285714) -- (2,1.5) -- cycle;
\fill[teal!30!white] (1.714286,1.5) -- (1.714286,1.285714) -- (2,1.5) -- cycle;
\fill[teal!8!white] (1.142857,1.071429) -- (1.428571,0.8571429) -- (1.428571,1.071429) -- cycle;
\fill[teal!8!white] (1.428571,0.8571429) -- (1.142857,1.071429) -- (1.142857,0.8571429) -- cycle;
\fill[teal!36!white] (3.142857,0.4285714) -- (3.142857,0.2142857) -- (3.428571,0.2142857) -- cycle;
\fill[teal!55!white] (3.428571,0.4285714) -- (3.142857,0.4285714) -- (3.428571,0.2142857) -- cycle;
\fill[teal!28!white] (2.857143,0.4285714) -- (3.142857,0.4285714) -- (2.857143,0.6428571) -- cycle;
\fill[teal!50!white] (3.142857,0.4285714) -- (3.142857,0.6428571) -- (2.857143,0.6428571) -- cycle;
\fill[teal!28!white] (3.142857,0.2142857) -- (3.142857,0.4285714) -- (2.857143,0.2142857) -- cycle;
\fill[teal!25!white] (3.142857,0.4285714) -- (2.857143,0.4285714) -- (2.857143,0.2142857) -- cycle;
\fill[teal!32!white] (2.571429,0.8571429) -- (2.571429,0.6428571) -- (2.857143,0.8571429) -- cycle;
\fill[teal!33!white] (2.857143,0.8571429) -- (2.571429,0.6428571) -- (2.857143,0.6428571) -- cycle;
\fill[teal!74!white] (2.285714,0) -- (2.285714,0.2142857) -- (2,0) -- cycle;
\fill[teal!37!white] (2.285714,0.2142857) -- (2,0.2142857) -- (2,0) -- cycle;
\fill[teal!19!white] (2.285714,0.2142857) -- (2,0.4285714) -- (2,0.2142857) -- cycle;
\fill[teal!16!white] (2,0.4285714) -- (2.285714,0.2142857) -- (2.285714,0.4285714) -- cycle;
\fill[teal!12!white] (2,0.4285714) -- (1.714286,0.6428571) -- (1.714286,0.4285714) -- cycle;
\fill[teal!11!white] (1.714286,0.6428571) -- (2,0.4285714) -- (2,0.6428571) -- cycle;
\fill[teal!15!white] (1.714286,0.2142857) -- (2,0.4285714) -- (1.714286,0.4285714) -- cycle;
\fill[teal!19!white] (2,0.2142857) -- (2,0.4285714) -- (1.714286,0.2142857) -- cycle;
\fill[teal!56!white] (0.2857143,0.8571429) -- (0,0.8571429) -- (0,0.6428571) -- cycle;
\fill[teal!29!white] (0.2857143,0.8571429) -- (0,0.6428571) -- (0.2857143,0.6428571) -- cycle;
\fill[teal!28!white] (0,1.071429) -- (0.2857143,0.8571429) -- (0.2857143,1.071429) -- cycle;
\fill[teal!55!white] (0,0.8571429) -- (0.2857143,0.8571429) -- (0,1.071429) -- cycle;
\fill[teal!12!white] (0.2857143,1.5) -- (0.5714286,1.285714) -- (0.5714286,1.5) -- cycle;
\fill[teal!15!white] (0.5714286,1.285714) -- (0.2857143,1.5) -- (0.2857143,1.285714) -- cycle;
\fill[teal!16!white] (0.5714286,0.6428571) -- (0.2857143,0.8571429) -- (0.2857143,0.6428571) -- cycle;
\fill[teal!13!white] (0.2857143,0.8571429) -- (0.5714286,0.6428571) -- (0.5714286,0.8571429) -- cycle;
\fill[teal!17!white] (0.2857143,0.4285714) -- (0.5714286,0.6428571) -- (0.2857143,0.6428571) -- cycle;
\fill[teal!15!white] (0.5714286,0.4285714) -- (0.5714286,0.6428571) -- (0.2857143,0.4285714) -- cycle;
\fill[teal!13!white] (0.8571429,0.4285714) -- (0.5714286,0.6428571) -- (0.5714286,0.4285714) -- cycle;
\fill[teal!11!white] (0.5714286,0.6428571) -- (0.8571429,0.4285714) -- (0.8571429,0.6428571) -- cycle;
\fill[teal!18!white] (1.142857,0.4285714) -- (1.142857,0.2142857) -- (1.428571,0.2142857) -- cycle;
\fill[teal!15!white] (1.428571,0.4285714) -- (1.142857,0.4285714) -- (1.428571,0.2142857) -- cycle;
\fill[teal!19!white] (1.142857,0.2142857) -- (1.142857,0.4285714) -- (0.8571429,0.2142857) -- cycle;
\fill[teal!15!white] (1.142857,0.4285714) -- (0.8571429,0.4285714) -- (0.8571429,0.2142857) -- cycle;
\fill[teal!12!white] (0.8571429,0.4285714) -- (1.142857,0.4285714) -- (0.8571429,0.6428571) -- cycle;
\fill[teal!10!white] (1.142857,0.4285714) -- (1.142857,0.6428571) -- (0.8571429,0.6428571) -- cycle;
\fill[teal!15!white] (1.428571,1.285714) -- (1.714286,1.5) -- (1.428571,1.5) -- cycle;
\fill[teal!15!white] (1.428571,1.285714) -- (1.714286,1.285714) -- (1.714286,1.5) -- cycle;
\fill[teal!12!white] (1.142857,1.5) -- (1.428571,1.285714) -- (1.428571,1.5) -- cycle;
\fill[teal!10!white] (1.428571,1.285714) -- (1.142857,1.5) -- (1.142857,1.285714) -- cycle;
\fill[teal!9!white] (1.428571,1.285714) -- (1.142857,1.071429) -- (1.428571,1.071429) -- cycle;
\fill[teal!9!white] (1.142857,1.071429) -- (1.428571,1.285714) -- (1.142857,1.285714) -- cycle;
\fill[teal!11!white] (1.428571,0.6428571) -- (1.428571,0.4285714) -- (1.714286,0.4285714) -- cycle;
\fill[teal!10!white] (1.714286,0.6428571) -- (1.428571,0.6428571) -- (1.714286,0.4285714) -- cycle;
\fill[teal!11!white] (1.428571,0.6428571) -- (1.142857,0.4285714) -- (1.428571,0.4285714) -- cycle;
\fill[teal!10!white] (1.142857,0.4285714) -- (1.428571,0.6428571) -- (1.142857,0.6428571) -- cycle;
\fill[teal!8!white] (1.428571,0.6428571) -- (1.428571,0.8571429) -- (1.142857,0.8571429) -- cycle;
\fill[teal!9!white] (1.142857,0.6428571) -- (1.428571,0.6428571) -- (1.142857,0.8571429) -- cycle;
\fill[teal!15!white] (2,1.285714) -- (1.714286,1.071429) -- (2,1.071429) -- cycle;
\fill[teal!15!white] (1.714286,1.285714) -- (1.714286,1.071429) -- (2,1.285714) -- cycle;
\fill[teal!12!white] (1.428571,1.285714) -- (1.714286,1.071429) -- (1.714286,1.285714) -- cycle;
\fill[teal!10!white] (1.714286,1.071429) -- (1.428571,1.285714) -- (1.428571,1.071429) -- cycle;
\fill[teal!17!white] (2.285714,0.6428571) -- (2.571429,0.8571429) -- (2.285714,0.8571429) -- cycle;
\fill[teal!18!white] (2.285714,0.6428571) -- (2.571429,0.6428571) -- (2.571429,0.8571429) -- cycle;
\fill[teal!13!white] (2,0.8571429) -- (2.285714,0.6428571) -- (2.285714,0.8571429) -- cycle;
\fill[teal!12!white] (2.285714,0.6428571) -- (2,0.8571429) -- (2,0.6428571) -- cycle;
\fill[teal!13!white] (2.285714,0.6428571) -- (2,0.4285714) -- (2.285714,0.4285714) -- cycle;
\fill[teal!12!white] (2,0.4285714) -- (2.285714,0.6428571) -- (2,0.6428571) -- cycle;
\fill[teal!21!white] (2.571429,0.4285714) -- (2.857143,0.4285714) -- (2.857143,0.6428571) -- cycle;
\fill[teal!19!white] (2.571429,0.6428571) -- (2.571429,0.4285714) -- (2.857143,0.6428571) -- cycle;
\fill[teal!16!white] (2.285714,0.6428571) -- (2.571429,0.4285714) -- (2.571429,0.6428571) -- cycle;
\fill[teal!15!white] (2.571429,0.4285714) -- (2.285714,0.6428571) -- (2.285714,0.4285714) -- cycle;
\fill[teal!15!white] (0.2857143,0.8571429) -- (0.5714286,1.071429) -- (0.2857143,1.071429) -- cycle;
\fill[teal!12!white] (0.5714286,1.071429) -- (0.2857143,0.8571429) -- (0.5714286,0.8571429) -- cycle;
\fill[teal!14!white] (0.5714286,1.071429) -- (0.2857143,1.285714) -- (0.2857143,1.071429) -- cycle;
\fill[teal!12!white] (0.5714286,1.071429) -- (0.5714286,1.285714) -- (0.2857143,1.285714) -- cycle;
\fill[teal!11!white] (0.8571429,1.5) -- (0.5714286,1.714286) -- (0.5714286,1.5) -- cycle;
\fill[teal!11!white] (0.8571429,1.5) -- (0.8571429,1.714286) -- (0.5714286,1.714286) -- cycle;
\fill[teal!11!white] (0.8571429,1.714286) -- (0.8571429,1.5) -- (1.142857,1.714286) -- cycle;
\fill[teal!11!white] (0.8571429,1.5) -- (1.142857,1.5) -- (1.142857,1.714286) -- cycle;
\fill[teal!10!white] (1.714286,0.8571429) -- (1.714286,0.6428571) -- (2,0.6428571) -- cycle;
\fill[teal!11!white] (2,0.8571429) -- (1.714286,0.8571429) -- (2,0.6428571) -- cycle;
\fill[teal!11!white] (1.714286,0.8571429) -- (2,0.8571429) -- (2,1.071429) -- cycle;
\fill[teal!11!white] (1.714286,1.071429) -- (1.714286,0.8571429) -- (2,1.071429) -- cycle;
\fill[teal!9!white] (1.714286,0.8571429) -- (1.428571,0.6428571) -- (1.714286,0.6428571) -- cycle;
\fill[teal!9!white] (1.428571,0.6428571) -- (1.714286,0.8571429) -- (1.428571,0.8571429) -- cycle;
\fill[teal!9!white] (1.428571,0.8571429) -- (1.714286,0.8571429) -- (1.428571,1.071429) -- cycle;
\fill[teal!10!white] (1.714286,0.8571429) -- (1.714286,1.071429) -- (1.428571,1.071429) -- cycle;
\fill[teal!75!white] (2.571429,0) -- (2.571429,0.2142857) -- (2.285714,0) -- cycle;
\fill[teal!38!white] (2.571429,0.2142857) -- (2.285714,0.2142857) -- (2.285714,0) -- cycle;
\fill[teal!20!white] (2.285714,0.2142857) -- (2.571429,0.2142857) -- (2.285714,0.4285714) -- cycle;
\fill[teal!18!white] (2.571429,0.2142857) -- (2.571429,0.4285714) -- (2.285714,0.4285714) -- cycle;
\fill[teal!40!white] (2.857143,0.2142857) -- (2.571429,0.2142857) -- (2.857143,0) -- cycle;
\fill[teal!75!white] (2.571429,0.2142857) -- (2.571429,0) -- (2.857143,0) -- cycle;
\fill[teal!23!white] (2.857143,0.4285714) -- (2.571429,0.2142857) -- (2.857143,0.2142857) -- cycle;
\fill[teal!20!white] (2.571429,0.4285714) -- (2.571429,0.2142857) -- (2.857143,0.4285714) -- cycle;
\fill[teal!9!white] (1.142857,0.6428571) -- (0.8571429,0.8571429) -- (0.8571429,0.6428571) -- cycle;
\fill[teal!8!white] (0.8571429,0.8571429) -- (1.142857,0.6428571) -- (1.142857,0.8571429) -- cycle;
\fill[teal!10!white] (0.5714286,0.6428571) -- (0.8571429,0.8571429) -- (0.5714286,0.8571429) -- cycle;
\fill[teal!10!white] (0.8571429,0.8571429) -- (0.5714286,0.6428571) -- (0.8571429,0.6428571) -- cycle;
\fill[teal!9!white] (1.142857,1.5) -- (0.8571429,1.285714) -- (1.142857,1.285714) -- cycle;
\fill[teal!9!white] (0.8571429,1.5) -- (0.8571429,1.285714) -- (1.142857,1.5) -- cycle;
\fill[teal!10!white] (0.5714286,1.285714) -- (0.8571429,1.285714) -- (0.5714286,1.5) -- cycle;
\fill[teal!9!white] (0.8571429,1.285714) -- (0.8571429,1.5) -- (0.5714286,1.5) -- cycle;
\fill[teal!8!white] (1.142857,1.071429) -- (0.8571429,1.071429) -- (1.142857,0.8571429) -- cycle;
\fill[teal!8!white] (0.8571429,1.071429) -- (0.8571429,0.8571429) -- (1.142857,0.8571429) -- cycle;
\fill[teal!8!white] (0.8571429,1.071429) -- (1.142857,1.071429) -- (1.142857,1.285714) -- cycle;
\fill[teal!8!white] (0.8571429,1.285714) -- (0.8571429,1.071429) -- (1.142857,1.285714) -- cycle;
\fill[teal!9!white] (0.8571429,1.071429) -- (0.5714286,1.071429) -- (0.5714286,0.8571429) -- cycle;
\fill[teal!9!white] (0.8571429,0.8571429) -- (0.8571429,1.071429) -- (0.5714286,0.8571429) -- cycle;
\fill[teal!9!white] (0.5714286,1.071429) -- (0.8571429,1.071429) -- (0.5714286,1.285714) -- cycle;
\fill[teal!9!white] (0.8571429,1.071429) -- (0.8571429,1.285714) -- (0.5714286,1.285714) -- cycle;
\draw[teal,thick] (0,0) -- (4,0) -- (0,3) -- cycle;
\draw[magenta!70,thin] (1.103977,0.9985762) -- (3.6,2.35) -- (4.1,2.35);\filldraw[fill=magenta,draw=white,line width=.45pt] (1.103977,0.9985762) circle(1.5pt) node[above right,font=\scriptsize,text=magenta,fill=white,inner sep=.6pt] {};
\node[anchor=west,text=magenta,font=\small] at(4.12,2.35){M1};\draw[gold!70,thin] (1,1) -- (3.6,1.8) -- (4.1,1.8);\filldraw[fill=gold,draw=white,line width=.45pt] (1,1) circle(1.5pt) node[above right,font=\scriptsize,text=gold,fill=white,inner sep=.6pt] {};
\node[anchor=west,text=gold,font=\small] at(4.12,1.8){$I$};\draw[navy!70,thin] (1.333333,1) -- (3.6,1.25) -- (4.1,1.25);\filldraw[fill=navy,draw=white,line width=.45pt] (1.333333,1) circle(1.5pt) node[above right,font=\scriptsize,text=navy,fill=white,inner sep=.6pt] {};
\node[anchor=west,text=navy,font=\small] at(4.12,1.25){$G$};\end{tikzpicture}
\caption[Normal field of a perimeter current]{Normal field of a perimeter current. M1 selects the minimum; high values accumulate near the wire. Fields are archived numerical illustrations; teal saturation increases with the displayed value.}\label{fig:field-M1}
\end{figure}

\section{A cutoff selects the loop-field centroid}\label{sec:slide24}

The unregularized interior field integral diverges near the wire, so an ordinary centroid does not exist without specifying a limit. Exclude a tube of uniform radius epsilon. Near a side, Bz behaves like a universal constant divided by perpendicular distance. Integrating Bz gives a logarithmic divergence per unit wire length. Integrating its square gives a stronger inverse-epsilon divergence. Their leading weights are both uniform along the perimeter, so their normalized centroids tend to W. Unequal cutoff thicknesses could change the answer.


\[\frac{\int_{T_\varepsilon}pB_z\,\dd A}{\int_{T_\varepsilon}B_z\,\dd A}\to X_{10},\qquad\frac{\int_{T_\varepsilon}pB_z^2\,\dd A}{\int_{T_\varepsilon}B_z^2\,\dd A}\to X_{10}.\]



Away from vertices $B_z\sim\mu_0I/(2\pi d)$ near a side.
Equal perpendicular clearance $d\ge\varepsilon$ yields a common
$\log(1/\varepsilon)$ coefficient per unit side length for $B_z$,
and a common $1/\varepsilon$ coefficient for $B_z^2$.
The dominating normalized measure is uniform arclength; corner
terms are lower order. Both limits are $W$.

Neither uncut integral defines a finite weighted mean. Squaring
the ideal field in the whole plane gives nonintegrable positive
singularities. Signed whole-plane means need explicit cutoffs and
cancellation conventions. Far-field decay does not cure a local
wire singularity. A finite-radius wire is a different well-defined
model, with limits that can depend on the core prescription.




\section{Three vertex currents reveal the Steiner foci}\label{sec:slide25}

Represent the plane by complex coordinates. Each axial wire contributes a rotated reciprocal-distance field. Equal currents make its zero condition the sum of reciprocals shown. The product P(z)=(z\ensuremath{-}zA)(z\ensuremath{-}zB)(z\ensuremath{-}zC) has logarithmic derivative P\ensuremath{^\prime}/P. The two roots of P\ensuremath{^\prime} are the Steiner inellipse foci. This is a striking link: mass covariance determines a conic, and a different physical experiment recovers its focal pair. The pair is intrinsic even when no preferred label distinguishes its two points.


\[P(z)=\prod_{j=1}^3(z-z_j),\quad\sum_j\frac1{z-z_j}=\frac{P^\prime(z)}{P(z)}.\]



With equal axial currents at complex positions $z_j$,
$B_x+iB_y$ is proportional to
$i\,\overline{\sum_j1/(z-z_j)}$. Zeros solve $P'(z)=0$.
For the reference triangle,
\[
z_\pm=\frac{4+3i\pm\sqrt{7-12i}}3
\simeq(2.4107,0.3812),\ (0.2560,1.6188).
\]
The cubic Marden relation identifies these as Steiner foci
\cite{marden}. Affine maps preserve the ellipse but not its foci:
diagonalize covariance after transforming.
There are two zero minima. Selecting one requires an additional
invariant rule; merely choosing left or right violates reflection
equivariance.



\begin{figure}[htbp]\centering
\begin{tikzpicture}[x=1.45cm,y=1.45cm]
\draw[teal,thick] (0,0) -- (4,0) -- (0,3) -- cycle;
\draw[magenta,thick] (2.607971,0.2678861) -- (2.630831,0.314005) -- (2.648135,0.3630615) -- (2.659808,0.4148455) -- (2.665802,0.4691352) -- (2.666089,0.5256981) -- (2.66067,0.5842921) -- (2.649567,0.6446662) -- (2.632827,0.7065619) -- (2.610523,0.7697141) -- (2.58275,0.8338524) -- (2.549626,0.8987023) -- (2.511294,0.9639858) -- (2.467918,1.029424) -- (2.419684,1.094735) -- (2.366797,1.159642) -- (2.309485,1.223864) -- (2.247993,1.287128) -- (2.182585,1.349162) -- (2.113539,1.409702) -- (2.041153,1.468486) -- (1.965736,1.525265) -- (1.887611,1.579795) -- (1.807112,1.631841) -- (1.724584,1.681182) -- (1.640381,1.727606) -- (1.554863,1.770915) -- (1.468397,1.810922) -- (1.381352,1.847457) -- (1.294102,1.880362) -- (1.207019,1.909498) -- (1.120478,1.934739) -- (1.034848,1.955978) -- (0.9504959,1.973123) -- (0.8677834,1.986101) -- (0.7870645,1.994856) -- (0.7086847,1.999351) -- (0.6329798,1.999567) -- (0.5602739,1.995502) -- (0.4908784,1.987175) -- (0.4250904,1.97462) -- (0.3631916,1.957892) -- (0.3054471,1.937062) -- (0.2521042,1.91222) -- (0.2033913,1.883471) -- (0.159517,1.850939) -- (0.1206691,1.814763) -- (0.08701404,1.775098) -- (0.05869592,1.732114) -- (0.03583598,1.685995) -- (0.01853213,1.636938) -- (0.006858461,1.585154) -- (0.0008649599,1.530865) -- (0.0005772928,1.474302) -- (0.005996692,1.415708) -- (0.01709995,1.355334) -- (0.03383952,1.293438) -- (0.05614372,1.230286) -- (0.08391705,1.166148) -- (0.1170406,1.101298) -- (0.1553724,1.036014) -- (0.1987485,0.9705763) -- (0.2469831,0.9052645) -- (0.2998695,0.8403584) -- (0.3571814,0.7761359) -- (0.4186734,0.7128719) -- (0.4840821,0.6508376) -- (0.5531273,0.5902983) -- (0.6255136,0.5315135) -- (0.7009308,0.4747348) -- (0.7790561,0.4202054) -- (0.8595549,0.3681588) -- (0.9420825,0.3188178) -- (1.026285,0.2723937) -- (1.111803,0.2290853) -- (1.19827,0.1890782) -- (1.285314,0.1525435) -- (1.372565,0.1196377) -- (1.459647,0.09050183) -- (1.546189,0.06526053) -- (1.631819,0.04402194) -- (1.716171,0.02687699) -- (1.798883,0.0138991) -- (1.879602,0.005143846) -- (1.957982,0.0006487199) -- (2.033687,0.0004329696) -- (2.106393,0.004497519) -- (2.175788,0.01282496) -- (2.241576,0.02537964) -- (2.303475,0.04210779) -- (2.36122,0.06293779) -- (2.414562,0.08778043) -- (2.463275,0.1165293) -- (2.50715,0.1490614) -- (2.545998,0.1852373) -- (2.579653,0.2249021) -- (2.607971,0.2678861);
\draw[blue,thick] (0.1193343,1.697285) -- (2.547332,0.302715);
\draw[gold,thick] (0.6360483,-0.213999) -- (2.030618,2.213999);
\filldraw[fill=navy,draw=white,line width=.45pt] (1.333333,1) circle(1.5pt) node[below,font=\scriptsize,text=navy,fill=white,inner sep=.6pt] {$G$};
\filldraw[fill=magenta,draw=white,line width=.45pt] (0.2559795,1.6188) circle(1.5pt) node[above,font=\scriptsize,text=magenta,fill=white,inner sep=.6pt] {$F_-$};
\filldraw[fill=magenta,draw=white,line width=.45pt] (2.410687,0.3811999) circle(1.5pt) node[above,font=\scriptsize,text=magenta,fill=white,inner sep=.6pt] {$F_+$};
\node[below left] at(0,0){$A$};\node[below right] at(4,0){$B$};\node[above] at(0,3){$C$};\end{tikzpicture}
\caption[The two field zeros of equal axial vertex currents are the Steiner foci]{The two field zeros of equal axial vertex currents are the Steiner foci.}\label{fig:geometry-25}
\end{figure}

\paragraph{Demonstration and investigation.} Reveal the ellipse, then its two foci, then the zero-field markers. Ask students to predict the equilateral limit.


\section{M3: angle-weighted magnetic energy}\label{sec:slide26}

For equal point-like vertex currents, the magnetic energy density diverges like 1/r\ensuremath{^2} near each source. Inside the triangle, a small annulus around a vertex occupies only its interior angular sector. Its divergent energy is therefore proportional to that angle. Equal circular core cutoffs yield weights \ensuremath{\alpha}, \ensuremath{\beta}, \ensuremath{\gamma} and the angle-weighted mean M3. This is the ETC center X360, whose center function contains an angle. Its ETC membership directly refutes the idea that every ETC center must have an algebraic formula.


\[\mathrm{M3}=\frac{\theta_AA+\theta_BB+\theta_CC}{\pi}=X_{360}.\]



Near vertex $A$, $|\mathbf B|^2\sim c/r^2$.
Integration over its interior wedge gives
$c\theta_A\log(1/\varepsilon)$ for equal circular core cutoffs.
Cross terms from other wires are lower order. The normalized first
moment therefore weights the vertices by angles.
Unequal currents instead give dominant weights $I_A^2\theta_A$.

Here $\theta_A=\pi/2$, $\theta_B=\arctan(3/4)$,
$\theta_C=\arctan(4/3)$, so M3 is approximately $(0.8193,0.8855)$.
It differs from $G$ and $W$. Its strict vertex attractivity is proved
in the appendix rather than inferred from a sample.




\section{M6: a uniformly current-filled prism}\label{sec:slide27}

Take an infinitely long prism with uniform axial current density over its triangular cross section. Integrate the two-dimensional reciprocal-distance kernel over that area. A boundary log-distance formula evaluates the field without a singular area-panel self term. The field magnitude is continuous inside, including at its zeros. Its positive first moment defines M6. The source current itself has centroid G, but the centroid of field magnitude generally differs. This is currently our strongest positive attractivity conjecture beyond the four proved points.


\[\mathbf B(p)=\frac{\mu_0J_z}{2\pi}\int_T\frac{\mathsf R(p-q)}{|p-q|^2}\dd A_q,\quad {\rm M6}=\frac{\int_Tp|\mathbf B|\dd A}{\int_T|\mathbf B|\dd A}.\]



A volume current is a continuum of line currents with
$J_z=I/\mathcal A$. A uniform prism-surface current instead has
line density $I/P$ on $\partial T$, giving a different field.
The rotation $\mathsf R(x,y)=(-y,x)$ enforces the right-hand rule.
The area singularity $1/|p-q|$ is integrable. Use singularity
subtraction, exact edge formulas or a logarithmic-potential boundary
reduction for quadrature.

A vector has no positive scalar centroid without a chosen weight.
M6 uses $|\mathbf B|$; the $|\mathbf B|^2$ mean is another selector.
Zeros and extrema are separate problems, with uniqueness requiring
its own argument.
Far away $|\mathbf B|\sim\mu_0I/(2\pi r)$.
Whole-plane mass of $|\mathbf B|$ diverges linearly and that of
$|\mathbf B|^2$ logarithmically. Interior moments are the clean
finite quantities. M6 remains a strong numerical attractivity
conjecture.



\begin{figure}[htbp]\centering
\begin{tikzpicture}[x=1.7cm,y=1.7cm]
\fill[teal!84!white] (3.142857,0.6428571) -- (2.857143,0.8571429) -- (2.857143,0.6428571) -- cycle;
\fill[teal!90!white] (1.714286,1.714286) -- (1.428571,1.928571) -- (1.428571,1.714286) -- cycle;
\fill[teal!90!white] (1.714286,1.714286) -- (1.714286,1.5) -- (2,1.5) -- cycle;
\fill[teal!89!white] (0.5714286,2.571429) -- (0.2857143,2.785714) -- (0.2857143,2.571429) -- cycle;
\fill[teal!87!white] (2.571429,1.071429) -- (2.285714,1.285714) -- (2.285714,1.071429) -- cycle;
\fill[teal!80!white] (3.714286,0.2142857) -- (3.428571,0.4285714) -- (3.428571,0.2142857) -- cycle;
\fill[teal!76!white] (3.714286,0.2142857) -- (3.714286,0) -- (4,0) -- cycle;
\fill[teal!90!white] (0.8571429,2.142857) -- (1.142857,2.142857) -- (0.8571429,2.357143) -- cycle;
\fill[teal!91!white] (1.142857,2.142857) -- (1.142857,1.928571) -- (1.428571,1.928571) -- cycle;
\fill[teal!86!white] (0,2.785714) -- (0.2857143,2.785714) -- (0,3) -- cycle;
\fill[teal!90!white] (0.5714286,2.571429) -- (0.5714286,2.357143) -- (0.8571429,2.357143) -- cycle;
\fill[teal!89!white] (2,1.285714) -- (2.285714,1.285714) -- (2,1.5) -- cycle;
\fill[teal!86!white] (2.571429,1.071429) -- (2.571429,0.8571429) -- (2.857143,0.8571429) -- cycle;
\fill[teal!82!white] (3.142857,0.4285714) -- (3.428571,0.4285714) -- (3.142857,0.6428571) -- cycle;
\fill[teal!94!white] (0.2857143,0) -- (0,0.2142857) -- (0,0) -- cycle;
\fill[teal!94!white] (0,0.2142857) -- (0.2857143,0) -- (0.2857143,0.2142857) -- cycle;
\fill[teal!94!white] (0.5714286,0.2142857) -- (0.2857143,0) -- (0.5714286,0) -- cycle;
\fill[teal!92!white] (0.2857143,0) -- (0.5714286,0.2142857) -- (0.2857143,0.2142857) -- cycle;
\fill[teal!83!white] (1.714286,1.5) -- (1.714286,1.714286) -- (1.428571,1.714286) -- cycle;
\fill[teal!71!white] (1.428571,1.5) -- (1.714286,1.5) -- (1.428571,1.714286) -- cycle;
\fill[teal!91!white] (0,1.928571) -- (0,1.714286) -- (0.2857143,1.928571) -- cycle;
\fill[teal!83!white] (0,1.714286) -- (0.2857143,1.714286) -- (0.2857143,1.928571) -- cycle;
\fill[teal!88!white] (0.2857143,2.357143) -- (0,2.571429) -- (0,2.357143) -- cycle;
\fill[teal!85!white] (0,2.571429) -- (0.2857143,2.357143) -- (0.2857143,2.571429) -- cycle;
\fill[teal!89!white] (0,2.142857) -- (0.2857143,2.357143) -- (0,2.357143) -- cycle;
\fill[teal!84!white] (0.2857143,2.357143) -- (0,2.142857) -- (0.2857143,2.142857) -- cycle;
\fill[teal!90!white] (0,2.142857) -- (0,1.928571) -- (0.2857143,1.928571) -- cycle;
\fill[teal!83!white] (0.2857143,2.142857) -- (0,2.142857) -- (0.2857143,1.928571) -- cycle;
\fill[teal!70!white] (0.2857143,1.714286) -- (0.5714286,1.928571) -- (0.2857143,1.928571) -- cycle;
\fill[teal!63!white] (0.5714286,1.714286) -- (0.5714286,1.928571) -- (0.2857143,1.714286) -- cycle;
\fill[teal!92!white] (0.2857143,1.5) -- (0,1.714286) -- (0,1.5) -- cycle;
\fill[teal!83!white] (0,1.714286) -- (0.2857143,1.5) -- (0.2857143,1.714286) -- cycle;
\fill[teal!67!white] (0.2857143,1.5) -- (0.5714286,1.714286) -- (0.2857143,1.714286) -- cycle;
\fill[teal!60!white] (0.5714286,1.714286) -- (0.2857143,1.5) -- (0.5714286,1.5) -- cycle;
\fill[teal!92!white] (0,1.285714) -- (0.2857143,1.5) -- (0,1.5) -- cycle;
\fill[teal!83!white] (0.2857143,1.5) -- (0,1.285714) -- (0.2857143,1.285714) -- cycle;
\fill[teal!93!white] (0,1.285714) -- (0,1.071429) -- (0.2857143,1.071429) -- cycle;
\fill[teal!84!white] (0.2857143,1.285714) -- (0,1.285714) -- (0.2857143,1.071429) -- cycle;
\fill[teal!79!white] (3.428571,0) -- (3.714286,0) -- (3.428571,0.2142857) -- cycle;
\fill[teal!77!white] (3.714286,0) -- (3.714286,0.2142857) -- (3.428571,0.2142857) -- cycle;
\fill[teal!91!white] (1.428571,0) -- (1.714286,0) -- (1.428571,0.2142857) -- cycle;
\fill[teal!82!white] (1.714286,0) -- (1.714286,0.2142857) -- (1.428571,0.2142857) -- cycle;
\fill[teal!94!white] (0.8571429,0) -- (0.5714286,0.2142857) -- (0.5714286,0) -- cycle;
\fill[teal!88!white] (0.5714286,0.2142857) -- (0.8571429,0) -- (0.8571429,0.2142857) -- cycle;
\fill[teal!92!white] (1.142857,0.2142857) -- (1.142857,0) -- (1.428571,0) -- cycle;
\fill[teal!84!white] (1.142857,0.2142857) -- (1.428571,0) -- (1.428571,0.2142857) -- cycle;
\fill[teal!93!white] (1.142857,0.2142857) -- (0.8571429,0) -- (1.142857,0) -- cycle;
\fill[teal!87!white] (0.8571429,0) -- (1.142857,0.2142857) -- (0.8571429,0.2142857) -- cycle;
\fill[teal!72!white] (1.142857,1.928571) -- (1.142857,1.714286) -- (1.428571,1.714286) -- cycle;
\fill[teal!84!white] (1.428571,1.928571) -- (1.142857,1.928571) -- (1.428571,1.714286) -- cycle;
\fill[teal!85!white] (0.8571429,2.142857) -- (1.142857,1.928571) -- (1.142857,2.142857) -- cycle;
\fill[teal!75!white] (1.142857,1.928571) -- (0.8571429,2.142857) -- (0.8571429,1.928571) -- cycle;
\fill[teal!87!white] (0,2.785714) -- (0,2.571429) -- (0.2857143,2.571429) -- cycle;
\fill[teal!86!white] (0.2857143,2.785714) -- (0,2.785714) -- (0.2857143,2.571429) -- cycle;
\fill[teal!63!white] (1.142857,1.928571) -- (0.8571429,1.714286) -- (1.142857,1.714286) -- cycle;
\fill[teal!64!white] (0.8571429,1.714286) -- (1.142857,1.928571) -- (0.8571429,1.928571) -- cycle;
\fill[teal!57!white] (0.8571429,1.714286) -- (0.5714286,1.928571) -- (0.5714286,1.714286) -- cycle;
\fill[teal!60!white] (0.5714286,1.928571) -- (0.8571429,1.714286) -- (0.8571429,1.928571) -- cycle;
\fill[teal!82!white] (0.2857143,2.357143) -- (0.5714286,2.357143) -- (0.2857143,2.571429) -- cycle;
\fill[teal!86!white] (0.5714286,2.357143) -- (0.5714286,2.571429) -- (0.2857143,2.571429) -- cycle;
\fill[teal!94!white] (0,0.6428571) -- (0,0.4285714) -- (0.2857143,0.6428571) -- cycle;
\fill[teal!89!white] (0,0.4285714) -- (0.2857143,0.4285714) -- (0.2857143,0.6428571) -- cycle;
\fill[teal!94!white] (0,0.4285714) -- (0,0.2142857) -- (0.2857143,0.2142857) -- cycle;
\fill[teal!91!white] (0.2857143,0.4285714) -- (0,0.4285714) -- (0.2857143,0.2142857) -- cycle;
\fill[teal!82!white] (2.285714,1.285714) -- (2,1.285714) -- (2.285714,1.071429) -- cycle;
\fill[teal!70!white] (2,1.285714) -- (2,1.071429) -- (2.285714,1.071429) -- cycle;
\fill[teal!59!white] (2,1.071429) -- (2,0.8571429) -- (2.285714,1.071429) -- cycle;
\fill[teal!60!white] (2,0.8571429) -- (2.285714,0.8571429) -- (2.285714,1.071429) -- cycle;
\fill[teal!70!white] (2.285714,0.8571429) -- (2.571429,0.8571429) -- (2.285714,1.071429) -- cycle;
\fill[teal!81!white] (2.571429,0.8571429) -- (2.571429,1.071429) -- (2.285714,1.071429) -- cycle;
\fill[teal!60!white] (1.428571,0.4285714) -- (1.714286,0.2142857) -- (1.714286,0.4285714) -- cycle;
\fill[teal!68!white] (1.714286,0.2142857) -- (1.428571,0.4285714) -- (1.428571,0.2142857) -- cycle;
\fill[teal!78!white] (3.142857,0.2142857) -- (3.428571,0) -- (3.428571,0.2142857) -- cycle;
\fill[teal!81!white] (3.142857,0.2142857) -- (3.142857,0) -- (3.428571,0) -- cycle;
\fill[teal!82!white] (3.142857,0) -- (3.142857,0.2142857) -- (2.857143,0) -- cycle;
\fill[teal!78!white] (3.142857,0.2142857) -- (2.857143,0.2142857) -- (2.857143,0) -- cycle;
\fill[teal!82!white] (1.714286,0) -- (2,0.2142857) -- (1.714286,0.2142857) -- cycle;
\fill[teal!89!white] (2,0.2142857) -- (1.714286,0) -- (2,0) -- cycle;
\fill[teal!82!white] (0.5714286,0.2142857) -- (0.5714286,0.4285714) -- (0.2857143,0.2142857) -- cycle;
\fill[teal!82!white] (0.5714286,0.4285714) -- (0.2857143,0.4285714) -- (0.2857143,0.2142857) -- cycle;
\fill[teal!76!white] (0.8571429,0.4285714) -- (0.5714286,0.2142857) -- (0.8571429,0.2142857) -- cycle;
\fill[teal!73!white] (0.8571429,0.4285714) -- (0.5714286,0.4285714) -- (0.5714286,0.2142857) -- cycle;
\fill[teal!59!white] (1.142857,1.5) -- (1.428571,1.5) -- (1.428571,1.714286) -- cycle;
\fill[teal!60!white] (1.142857,1.714286) -- (1.142857,1.5) -- (1.428571,1.714286) -- cycle;
\fill[teal!73!white] (0.5714286,2.142857) -- (0.2857143,2.142857) -- (0.2857143,1.928571) -- cycle;
\fill[teal!69!white] (0.5714286,1.928571) -- (0.5714286,2.142857) -- (0.2857143,1.928571) -- cycle;
\fill[teal!76!white] (0.5714286,2.142857) -- (0.2857143,2.357143) -- (0.2857143,2.142857) -- cycle;
\fill[teal!78!white] (0.5714286,2.142857) -- (0.5714286,2.357143) -- (0.2857143,2.357143) -- cycle;
\fill[teal!72!white] (0.8571429,2.142857) -- (0.5714286,2.142857) -- (0.8571429,1.928571) -- cycle;
\fill[teal!67!white] (0.5714286,2.142857) -- (0.5714286,1.928571) -- (0.8571429,1.928571) -- cycle;
\fill[teal!81!white] (0.5714286,2.142857) -- (0.8571429,2.142857) -- (0.8571429,2.357143) -- cycle;
\fill[teal!82!white] (0.5714286,2.357143) -- (0.5714286,2.142857) -- (0.8571429,2.357143) -- cycle;
\fill[teal!76!white] (1.714286,1.285714) -- (2,1.285714) -- (2,1.5) -- cycle;
\fill[teal!76!white] (1.714286,1.5) -- (1.714286,1.285714) -- (2,1.5) -- cycle;
\fill[teal!11!white] (1.142857,1.071429) -- (1.428571,0.8571429) -- (1.428571,1.071429) -- cycle;
\fill[teal!8!white] (1.428571,0.8571429) -- (1.142857,1.071429) -- (1.142857,0.8571429) -- cycle;
\fill[teal!76!white] (3.142857,0.4285714) -- (3.142857,0.2142857) -- (3.428571,0.2142857) -- cycle;
\fill[teal!79!white] (3.428571,0.4285714) -- (3.142857,0.4285714) -- (3.428571,0.2142857) -- cycle;
\fill[teal!73!white] (2.857143,0.4285714) -- (3.142857,0.4285714) -- (2.857143,0.6428571) -- cycle;
\fill[teal!79!white] (3.142857,0.4285714) -- (3.142857,0.6428571) -- (2.857143,0.6428571) -- cycle;
\fill[teal!73!white] (3.142857,0.2142857) -- (3.142857,0.4285714) -- (2.857143,0.2142857) -- cycle;
\fill[teal!71!white] (3.142857,0.4285714) -- (2.857143,0.4285714) -- (2.857143,0.2142857) -- cycle;
\fill[teal!76!white] (2.571429,0.8571429) -- (2.571429,0.6428571) -- (2.857143,0.8571429) -- cycle;
\fill[teal!76!white] (2.857143,0.8571429) -- (2.571429,0.6428571) -- (2.857143,0.6428571) -- cycle;
\fill[teal!87!white] (2.285714,0) -- (2.285714,0.2142857) -- (2,0) -- cycle;
\fill[teal!80!white] (2.285714,0.2142857) -- (2,0.2142857) -- (2,0) -- cycle;
\fill[teal!67!white] (2.285714,0.2142857) -- (2,0.4285714) -- (2,0.2142857) -- cycle;
\fill[teal!62!white] (2,0.4285714) -- (2.285714,0.2142857) -- (2.285714,0.4285714) -- cycle;
\fill[teal!48!white] (2,0.4285714) -- (1.714286,0.6428571) -- (1.714286,0.4285714) -- cycle;
\fill[teal!44!white] (1.714286,0.6428571) -- (2,0.4285714) -- (2,0.6428571) -- cycle;
\fill[teal!60!white] (1.714286,0.2142857) -- (2,0.4285714) -- (1.714286,0.4285714) -- cycle;
\fill[teal!67!white] (2,0.2142857) -- (2,0.4285714) -- (1.714286,0.2142857) -- cycle;
\fill[teal!94!white] (0.2857143,0.8571429) -- (0,0.8571429) -- (0,0.6428571) -- cycle;
\fill[teal!87!white] (0.2857143,0.8571429) -- (0,0.6428571) -- (0.2857143,0.6428571) -- cycle;
\fill[teal!85!white] (0,1.071429) -- (0.2857143,0.8571429) -- (0.2857143,1.071429) -- cycle;
\fill[teal!94!white] (0,0.8571429) -- (0.2857143,0.8571429) -- (0,1.071429) -- cycle;
\fill[teal!58!white] (0.2857143,1.5) -- (0.5714286,1.285714) -- (0.5714286,1.5) -- cycle;
\fill[teal!66!white] (0.5714286,1.285714) -- (0.2857143,1.5) -- (0.2857143,1.285714) -- cycle;
\fill[teal!72!white] (0.5714286,0.6428571) -- (0.2857143,0.8571429) -- (0.2857143,0.6428571) -- cycle;
\fill[teal!63!white] (0.2857143,0.8571429) -- (0.5714286,0.6428571) -- (0.5714286,0.8571429) -- cycle;
\fill[teal!75!white] (0.2857143,0.4285714) -- (0.5714286,0.6428571) -- (0.2857143,0.6428571) -- cycle;
\fill[teal!72!white] (0.5714286,0.4285714) -- (0.5714286,0.6428571) -- (0.2857143,0.4285714) -- cycle;
\fill[teal!64!white] (0.8571429,0.4285714) -- (0.5714286,0.6428571) -- (0.5714286,0.4285714) -- cycle;
\fill[teal!55!white] (0.5714286,0.6428571) -- (0.8571429,0.4285714) -- (0.8571429,0.6428571) -- cycle;
\fill[teal!70!white] (1.142857,0.4285714) -- (1.142857,0.2142857) -- (1.428571,0.2142857) -- cycle;
\fill[teal!62!white] (1.428571,0.4285714) -- (1.142857,0.4285714) -- (1.428571,0.2142857) -- cycle;
\fill[teal!72!white] (1.142857,0.2142857) -- (1.142857,0.4285714) -- (0.8571429,0.2142857) -- cycle;
\fill[teal!66!white] (1.142857,0.4285714) -- (0.8571429,0.4285714) -- (0.8571429,0.2142857) -- cycle;
\fill[teal!54!white] (0.8571429,0.4285714) -- (1.142857,0.4285714) -- (0.8571429,0.6428571) -- cycle;
\fill[teal!45!white] (1.142857,0.4285714) -- (1.142857,0.6428571) -- (0.8571429,0.6428571) -- cycle;
\fill[teal!58!white] (1.428571,1.285714) -- (1.714286,1.5) -- (1.428571,1.5) -- cycle;
\fill[teal!58!white] (1.428571,1.285714) -- (1.714286,1.285714) -- (1.714286,1.5) -- cycle;
\fill[teal!47!white] (1.142857,1.5) -- (1.428571,1.285714) -- (1.428571,1.5) -- cycle;
\fill[teal!36!white] (1.428571,1.285714) -- (1.142857,1.5) -- (1.142857,1.285714) -- cycle;
\fill[teal!22!white] (1.428571,1.285714) -- (1.142857,1.071429) -- (1.428571,1.071429) -- cycle;
\fill[teal!23!white] (1.142857,1.071429) -- (1.428571,1.285714) -- (1.142857,1.285714) -- cycle;
\fill[teal!46!white] (1.428571,0.6428571) -- (1.428571,0.4285714) -- (1.714286,0.4285714) -- cycle;
\fill[teal!40!white] (1.714286,0.6428571) -- (1.428571,0.6428571) -- (1.714286,0.4285714) -- cycle;
\fill[teal!47!white] (1.428571,0.6428571) -- (1.142857,0.4285714) -- (1.428571,0.4285714) -- cycle;
\fill[teal!41!white] (1.142857,0.4285714) -- (1.428571,0.6428571) -- (1.142857,0.6428571) -- cycle;
\fill[teal!20!white] (1.428571,0.6428571) -- (1.428571,0.8571429) -- (1.142857,0.8571429) -- cycle;
\fill[teal!28!white] (1.142857,0.6428571) -- (1.428571,0.6428571) -- (1.142857,0.8571429) -- cycle;
\fill[teal!58!white] (2,1.285714) -- (1.714286,1.071429) -- (2,1.071429) -- cycle;
\fill[teal!57!white] (1.714286,1.285714) -- (1.714286,1.071429) -- (2,1.285714) -- cycle;
\fill[teal!45!white] (1.428571,1.285714) -- (1.714286,1.071429) -- (1.714286,1.285714) -- cycle;
\fill[teal!33!white] (1.714286,1.071429) -- (1.428571,1.285714) -- (1.428571,1.071429) -- cycle;
\fill[teal!62!white] (2.285714,0.6428571) -- (2.571429,0.8571429) -- (2.285714,0.8571429) -- cycle;
\fill[teal!63!white] (2.285714,0.6428571) -- (2.571429,0.6428571) -- (2.571429,0.8571429) -- cycle;
\fill[teal!53!white] (2,0.8571429) -- (2.285714,0.6428571) -- (2.285714,0.8571429) -- cycle;
\fill[teal!47!white] (2.285714,0.6428571) -- (2,0.8571429) -- (2,0.6428571) -- cycle;
\fill[teal!54!white] (2.285714,0.6428571) -- (2,0.4285714) -- (2.285714,0.4285714) -- cycle;
\fill[teal!49!white] (2,0.4285714) -- (2.285714,0.6428571) -- (2,0.6428571) -- cycle;
\fill[teal!68!white] (2.571429,0.4285714) -- (2.857143,0.4285714) -- (2.857143,0.6428571) -- cycle;
\fill[teal!66!white] (2.571429,0.6428571) -- (2.571429,0.4285714) -- (2.857143,0.6428571) -- cycle;
\fill[teal!60!white] (2.285714,0.6428571) -- (2.571429,0.4285714) -- (2.571429,0.6428571) -- cycle;
\fill[teal!58!white] (2.571429,0.4285714) -- (2.285714,0.6428571) -- (2.285714,0.4285714) -- cycle;
\fill[teal!67!white] (0.2857143,0.8571429) -- (0.5714286,1.071429) -- (0.2857143,1.071429) -- cycle;
\fill[teal!60!white] (0.5714286,1.071429) -- (0.2857143,0.8571429) -- (0.5714286,0.8571429) -- cycle;
\fill[teal!66!white] (0.5714286,1.071429) -- (0.2857143,1.285714) -- (0.2857143,1.071429) -- cycle;
\fill[teal!57!white] (0.5714286,1.071429) -- (0.5714286,1.285714) -- (0.2857143,1.285714) -- cycle;
\fill[teal!49!white] (0.8571429,1.5) -- (0.5714286,1.714286) -- (0.5714286,1.5) -- cycle;
\fill[teal!49!white] (0.8571429,1.5) -- (0.8571429,1.714286) -- (0.5714286,1.714286) -- cycle;
\fill[teal!49!white] (0.8571429,1.714286) -- (0.8571429,1.5) -- (1.142857,1.714286) -- cycle;
\fill[teal!46!white] (0.8571429,1.5) -- (1.142857,1.5) -- (1.142857,1.714286) -- cycle;
\fill[teal!36!white] (1.714286,0.8571429) -- (1.714286,0.6428571) -- (2,0.6428571) -- cycle;
\fill[teal!39!white] (2,0.8571429) -- (1.714286,0.8571429) -- (2,0.6428571) -- cycle;
\fill[teal!43!white] (1.714286,0.8571429) -- (2,0.8571429) -- (2,1.071429) -- cycle;
\fill[teal!42!white] (1.714286,1.071429) -- (1.714286,0.8571429) -- (2,1.071429) -- cycle;
\fill[teal!30!white] (1.714286,0.8571429) -- (1.428571,0.6428571) -- (1.714286,0.6428571) -- cycle;
\fill[teal!23!white] (1.428571,0.6428571) -- (1.714286,0.8571429) -- (1.428571,0.8571429) -- cycle;
\fill[teal!21!white] (1.428571,0.8571429) -- (1.714286,0.8571429) -- (1.428571,1.071429) -- cycle;
\fill[teal!30!white] (1.714286,0.8571429) -- (1.714286,1.071429) -- (1.428571,1.071429) -- cycle;
\fill[teal!86!white] (2.571429,0) -- (2.571429,0.2142857) -- (2.285714,0) -- cycle;
\fill[teal!80!white] (2.571429,0.2142857) -- (2.285714,0.2142857) -- (2.285714,0) -- cycle;
\fill[teal!68!white] (2.285714,0.2142857) -- (2.571429,0.2142857) -- (2.285714,0.4285714) -- cycle;
\fill[teal!65!white] (2.571429,0.2142857) -- (2.571429,0.4285714) -- (2.285714,0.4285714) -- cycle;
\fill[teal!79!white] (2.857143,0.2142857) -- (2.571429,0.2142857) -- (2.857143,0) -- cycle;
\fill[teal!85!white] (2.571429,0.2142857) -- (2.571429,0) -- (2.857143,0) -- cycle;
\fill[teal!71!white] (2.857143,0.4285714) -- (2.571429,0.2142857) -- (2.857143,0.2142857) -- cycle;
\fill[teal!67!white] (2.571429,0.4285714) -- (2.571429,0.2142857) -- (2.857143,0.4285714) -- cycle;
\fill[teal!37!white] (1.142857,0.6428571) -- (0.8571429,0.8571429) -- (0.8571429,0.6428571) -- cycle;
\fill[teal!27!white] (0.8571429,0.8571429) -- (1.142857,0.6428571) -- (1.142857,0.8571429) -- cycle;
\fill[teal!49!white] (0.5714286,0.6428571) -- (0.8571429,0.8571429) -- (0.5714286,0.8571429) -- cycle;
\fill[teal!47!white] (0.8571429,0.8571429) -- (0.5714286,0.6428571) -- (0.8571429,0.6428571) -- cycle;
\fill[teal!30!white] (1.142857,1.5) -- (0.8571429,1.285714) -- (1.142857,1.285714) -- cycle;
\fill[teal!35!white] (0.8571429,1.5) -- (0.8571429,1.285714) -- (1.142857,1.5) -- cycle;
\fill[teal!43!white] (0.5714286,1.285714) -- (0.8571429,1.285714) -- (0.5714286,1.5) -- cycle;
\fill[teal!40!white] (0.8571429,1.285714) -- (0.8571429,1.5) -- (0.5714286,1.5) -- cycle;
\fill[teal!14!white] (1.142857,1.071429) -- (0.8571429,1.071429) -- (1.142857,0.8571429) -- cycle;
\fill[teal!24!white] (0.8571429,1.071429) -- (0.8571429,0.8571429) -- (1.142857,0.8571429) -- cycle;
\fill[teal!16!white] (0.8571429,1.071429) -- (1.142857,1.071429) -- (1.142857,1.285714) -- cycle;
\fill[teal!23!white] (0.8571429,1.285714) -- (0.8571429,1.071429) -- (1.142857,1.285714) -- cycle;
\fill[teal!43!white] (0.8571429,1.071429) -- (0.5714286,1.071429) -- (0.5714286,0.8571429) -- cycle;
\fill[teal!38!white] (0.8571429,0.8571429) -- (0.8571429,1.071429) -- (0.5714286,0.8571429) -- cycle;
\fill[teal!42!white] (0.5714286,1.071429) -- (0.8571429,1.071429) -- (0.5714286,1.285714) -- cycle;
\fill[teal!35!white] (0.8571429,1.071429) -- (0.8571429,1.285714) -- (0.5714286,1.285714) -- cycle;
\draw[teal,thick] (0,0) -- (4,0) -- (0,3) -- cycle;
\draw[magenta!70,thin] (1.310681,1.006307) -- (3.6,2.35) -- (4.1,2.35);\filldraw[fill=magenta,draw=white,line width=.45pt] (1.310681,1.006307) circle(1.5pt) node[above right,font=\scriptsize,text=magenta,fill=white,inner sep=.6pt] {};
\node[anchor=west,text=magenta,font=\small] at(4.12,2.35){M6};\draw[gold!70,thin] (1,1) -- (3.6,1.8) -- (4.1,1.8);\filldraw[fill=gold,draw=white,line width=.45pt] (1,1) circle(1.5pt) node[above right,font=\scriptsize,text=gold,fill=white,inner sep=.6pt] {};
\node[anchor=west,text=gold,font=\small] at(4.12,1.8){$I$};\draw[navy!70,thin] (1.333333,1) -- (3.6,1.25) -- (4.1,1.25);\filldraw[fill=navy,draw=white,line width=.45pt] (1.333333,1) circle(1.5pt) node[above right,font=\scriptsize,text=navy,fill=white,inner sep=.6pt] {};
\node[anchor=west,text=navy,font=\small] at(4.12,1.25){$G$};\end{tikzpicture}
\caption[Uniform-volume-current field magnitude]{Uniform-volume-current field magnitude. M6 is its first moment, not its minimum. Fields are archived numerical illustrations; teal saturation increases with the displayed value.}\label{fig:field-M6}
\end{figure}

\section{Field stability and vertex response}\label{sec:slide28}

A field extremum describes variation within one fixed triangle.
Vertex response changes the domain, field and selected point together.
A negative response can occur even when each equilibrium is stable
and unique. The certified M1 motion differs qualitatively from the
exact centroid response in Chapter 1. An illustrative arrow must not
replace evaluation of the actual dot product with controlled errors.



\paragraph{Demonstration and investigation.} Run the before-and-after center motions. State the sign first, then let students interpret the directions. The exact magnitude shown in the notes is the evidence; the arrow is a visualization.


\section{A certified failure for E0}\label{sec:slide29}

This is a correction in the current atlas. The test triangle has A=(0,0), B=(1,0) and C approximately (0.996364,0.001599). Move A by ten to the minus five times the stated nearly horizontal vector. The normalized dot product is negative by a narrow outward-rounded interval enclosure. The effect is too small for a reliable ordinary difference of double-precision positions. The certificate computes the incomplete-beta map in two bounded stages, using a power series inside the unit disk and an analytic continuation with an explicit remainder. Its proof is in the companion addendum.


\[\frac{h\cdot\Delta\mathrm{E0}}{|h|^2}\in[-3.13100322614,-3.13100322613]\times10^{-7}.\]


The exact decimal witness inputs are listed in Appendix A.
Independent 45-, 75- and 105-digit evaluations agree, while the interval
record supplies the actual sign proof. Agreement alone is not a proof.



\section{Properties of a physical center}\label{sec:slide30}

Review four examples. E2 is a well-defined charge mean with a certified failure of attractivity. M1 is a unique minimum and is nonattractive. M3 is exactly X360 and strictly attractive. M6 is an interior mean with strong positive screens and an open global theorem. These differences guide numerical work. A laboratory demonstration can measure a field map and then compute its digital weighted mean. It need not imply that a freely moving test particle seeks that mean.



\paragraph{Demonstration and investigation.} For a practical loop demonstration, use low current, finite wire radius, a three-axis Hall probe and a fixed-height scan. Record lead geometry and probe averaging. For equal long axial wires, search for the focal zero pair rather than an energy centroid.


\section*{Questions and exercises}\addcontentsline{toc}{section}{Questions and exercises}

\begin{enumerate}

\item\label{ex:17} What does positive charge density guarantee about its centroid?

\item\label{ex:18} Why can the two-dimensional and three-dimensional charge centroids differ?

\item\label{ex:19} Would a denser panel mesh by itself prove the continuum centroid?

\item\label{ex:20} Which familiar center has equal side distances?

\item\label{ex:21} Why must every trial electrode have the same size?

\item\label{ex:23} Does a unique minimum establish geometric attractivity?

\item\label{ex:24} Can we simply write \ensuremath{\int}pBz/\ensuremath{\int}Bz without mentioning epsilon?

\item\label{ex:25} What happens to the two foci at equilateral?

\item\label{ex:26} Why does integrating over the whole plane give different vertex weights?

\item\label{ex:27} Would positivity of |B| alone prove attractivity?

\item\label{ex:28} Which calculation could settle a failure globally?

\item\label{ex:29} What would 612 positive numerical cases have established?

\item\label{ex:30} Which new positive theorem would be most useful now?

\end{enumerate}

\chapter{Membranes, elasticity and fluids}\label{ch:3}

The triangle now supports a displacement field. A taut membrane resists changes in slope, while an elastic plate resists curvature. Their mathematical difference is physical: tension and bending store different energies. The linear membrane equation will reappear in viscous duct flow and Brownian lifetime. We will also see an exact incenter experiment, but only after carefully restricting what shapes the sheet may form.


\section{A stretched membrane under gravity}\label{sec:slide32}

Assume a massless or lightly massive membrane under a prescribed uniform downward load. The edge is fixed at zero displacement and the tension is isotropic and constant. For small slopes, the vertical components of tension across a small patch depend on the curvature of its graph. Their net upward force balances q times the patch area. This approximation ignores stretching changes of tension and bending stiffness. It is a clean physical model rather than a claim about every metallic sheet.


\[-\tau\Delta w=q,\quad w|_{\partial T}=0.\]


A rectangular patch's net vertical tension is
$\tau(w_{xx}+w_{yy})\dd A$. With downward $w$ and positive downward
load, equilibrium gives $-\tau\Delta w=q$.



\paragraph{Demonstration and investigation.} Reveal edge support, downward load arrows and the displacement field in successive steps.


\section{Force balance gives a Poisson equation}\label{sec:slide33}

The Laplacian is the sum of second derivatives in the two in-plane directions. It measures how the local slope changes around a point. Dividing by q/\ensuremath{\tau} rescales the displacement but leaves its maximizing position unchanged. The normalized torsion equation is therefore \ensuremath{-}\ensuremath{\Delta}u=1 with zero boundary values. This is the central boundary-value problem of the lecture. The word torsion originally refers to the equivalent Prandtl stress function of a twisted rod, but it is equally useful here for a membrane.


\[\int_{\partial U}\tau\,\partial_nw\dd s+\int_Uq\dd A=0.\]


Tension $\tau$ has units force per length, load $q$ force per area,
and $\Delta w$ inverse length. Thus $u=\tau w/q$ has units length
squared and solves $-\Delta u=1$.



\section{The energy formulation}\label{sec:slide34}

Small slopes stretch a membrane locally by an area proportional to one half the squared slope. At fixed imposed tension this gives the displayed quadratic energy. The load lowers the potential energy when the membrane moves downward. Vary w by a test function that vanishes at the boundary. Integrate the gradient term by parts. The stationary equation is exactly the force-balance equation. The positive quadratic term makes the linear field solution unique. Energy stability of this field still says nothing about how the maximum moves when the boundary triangle changes.


\[\mathcal J[w]=\frac{\tau}{2}\int_T|\nabla w|^2\dd A-\int_Tqw\dd A.\]



Vary $w$ to $w+\varepsilon v$, with $v=0$ on the boundary.
The stationary energy condition is
\[
\tau\int_T\nabla w\cdot\nabla v\dd A=\int_Tqv\dd A
\quad\hbox{for every }v\in H_0^1(T).
\]
Integration by parts recovers the PDE. Conversely this identity
minimizes the strictly convex quadratic functional.
$H_0^1$ consists of square-integrable functions with square-integrable
first derivatives and zero boundary trace. The Poincare inequality
gives coercivity, hence existence and uniqueness of the field.
The maximum principle gives interior positivity for $q>0$.
These statements do not alone prove a unique maximum.
Torsion in a convex planar domain has concavity properties sufficient
to give its unique maximum; related results are discussed in
\cite{torsion}. Its mean and maximum remain different rules.




\paragraph{Demonstration and investigation.} Pause before the integration by parts and let students identify the disappearing boundary term.


\section{F1: the deepest membrane point}\label{sec:slide35}

The torsion field is positive in the interior and zero on the boundary. For a convex triangle its maximum is unique. The red marker is F1, while I and G are comparison points. F1 generally differs from both. An independent theorem supports uniqueness, but the attractivity question remains open in this project. Earlier negative numerical witnesses became sensitive to degree refinement, so we do not present them as established continuum counterexamples.


\[-\Delta u=1,\quad u|_{\partial T}=0;\quad {\rm F1}=\argmax_Tu.\]


The maximum principle gives positivity. Stronger concavity and critical-point results give uniqueness for this convex planar torsion problem. The numerical task is to solve a linear system for the field, then locate its maximum accurately inside an element rather than only among mesh vertices.


\begin{figure}[htbp]\centering
\begin{tikzpicture}[x=1.7cm,y=1.7cm]
\fill[teal!14!white] (3.142857,0.6428571) -- (2.857143,0.8571429) -- (2.857143,0.6428571) -- cycle;
\fill[teal!19!white] (1.714286,1.714286) -- (1.428571,1.928571) -- (1.428571,1.714286) -- cycle;
\fill[teal!19!white] (1.714286,1.714286) -- (1.714286,1.5) -- (2,1.5) -- cycle;
\fill[teal!12!white] (0.5714286,2.571429) -- (0.2857143,2.785714) -- (0.2857143,2.571429) -- cycle;
\fill[teal!17!white] (2.571429,1.071429) -- (2.285714,1.285714) -- (2.285714,1.071429) -- cycle;
\fill[teal!10!white] (3.714286,0.2142857) -- (3.428571,0.4285714) -- (3.428571,0.2142857) -- cycle;
\fill[teal!8!white] (3.714286,0.2142857) -- (3.714286,0) -- (4,0) -- cycle;
\fill[teal!17!white] (0.8571429,2.142857) -- (1.142857,2.142857) -- (0.8571429,2.357143) -- cycle;
\fill[teal!18!white] (1.142857,2.142857) -- (1.142857,1.928571) -- (1.428571,1.928571) -- cycle;
\fill[teal!9!white] (0,2.785714) -- (0.2857143,2.785714) -- (0,3) -- cycle;
\fill[teal!14!white] (0.5714286,2.571429) -- (0.5714286,2.357143) -- (0.8571429,2.357143) -- cycle;
\fill[teal!18!white] (2,1.285714) -- (2.285714,1.285714) -- (2,1.5) -- cycle;
\fill[teal!16!white] (2.571429,1.071429) -- (2.571429,0.8571429) -- (2.857143,0.8571429) -- cycle;
\fill[teal!12!white] (3.142857,0.4285714) -- (3.428571,0.4285714) -- (3.142857,0.6428571) -- cycle;
\fill[teal!11!white] (0.2857143,0) -- (0,0.2142857) -- (0,0) -- cycle;
\fill[teal!18!white] (0,0.2142857) -- (0.2857143,0) -- (0.2857143,0.2142857) -- cycle;
\fill[teal!18!white] (0.5714286,0.2142857) -- (0.2857143,0) -- (0.5714286,0) -- cycle;
\fill[teal!24!white] (0.2857143,0) -- (0.5714286,0.2142857) -- (0.2857143,0.2142857) -- cycle;
\fill[teal!29!white] (1.714286,1.5) -- (1.714286,1.714286) -- (1.428571,1.714286) -- cycle;
\fill[teal!48!white] (1.428571,1.5) -- (1.714286,1.5) -- (1.428571,1.714286) -- cycle;
\fill[teal!21!white] (0,1.928571) -- (0,1.714286) -- (0.2857143,1.928571) -- cycle;
\fill[teal!35!white] (0,1.714286) -- (0.2857143,1.714286) -- (0.2857143,1.928571) -- cycle;
\fill[teal!15!white] (0.2857143,2.357143) -- (0,2.571429) -- (0,2.357143) -- cycle;
\fill[teal!18!white] (0,2.571429) -- (0.2857143,2.357143) -- (0.2857143,2.571429) -- cycle;
\fill[teal!17!white] (0,2.142857) -- (0.2857143,2.357143) -- (0,2.357143) -- cycle;
\fill[teal!26!white] (0.2857143,2.357143) -- (0,2.142857) -- (0.2857143,2.142857) -- cycle;
\fill[teal!20!white] (0,2.142857) -- (0,1.928571) -- (0.2857143,1.928571) -- cycle;
\fill[teal!29!white] (0.2857143,2.142857) -- (0,2.142857) -- (0.2857143,1.928571) -- cycle;
\fill[teal!50!white] (0.2857143,1.714286) -- (0.5714286,1.928571) -- (0.2857143,1.928571) -- cycle;
\fill[teal!59!white] (0.5714286,1.714286) -- (0.5714286,1.928571) -- (0.2857143,1.714286) -- cycle;
\fill[teal!24!white] (0.2857143,1.5) -- (0,1.714286) -- (0,1.5) -- cycle;
\fill[teal!37!white] (0,1.714286) -- (0.2857143,1.5) -- (0.2857143,1.714286) -- cycle;
\fill[teal!57!white] (0.2857143,1.5) -- (0.5714286,1.714286) -- (0.2857143,1.714286) -- cycle;
\fill[teal!67!white] (0.5714286,1.714286) -- (0.2857143,1.5) -- (0.5714286,1.5) -- cycle;
\fill[teal!24!white] (0,1.285714) -- (0.2857143,1.5) -- (0,1.5) -- cycle;
\fill[teal!39!white] (0.2857143,1.5) -- (0,1.285714) -- (0.2857143,1.285714) -- cycle;
\fill[teal!25!white] (0,1.285714) -- (0,1.071429) -- (0.2857143,1.071429) -- cycle;
\fill[teal!40!white] (0.2857143,1.285714) -- (0,1.285714) -- (0.2857143,1.071429) -- cycle;
\fill[teal!10!white] (3.428571,0) -- (3.714286,0) -- (3.428571,0.2142857) -- cycle;
\fill[teal!10!white] (3.714286,0) -- (3.714286,0.2142857) -- (3.428571,0.2142857) -- cycle;
\fill[teal!21!white] (1.428571,0) -- (1.714286,0) -- (1.428571,0.2142857) -- cycle;
\fill[teal!33!white] (1.714286,0) -- (1.714286,0.2142857) -- (1.428571,0.2142857) -- cycle;
\fill[teal!19!white] (0.8571429,0) -- (0.5714286,0.2142857) -- (0.5714286,0) -- cycle;
\fill[teal!31!white] (0.5714286,0.2142857) -- (0.8571429,0) -- (0.8571429,0.2142857) -- cycle;
\fill[teal!21!white] (1.142857,0.2142857) -- (1.142857,0) -- (1.428571,0) -- cycle;
\fill[teal!34!white] (1.142857,0.2142857) -- (1.428571,0) -- (1.428571,0.2142857) -- cycle;
\fill[teal!21!white] (1.142857,0.2142857) -- (0.8571429,0) -- (1.142857,0) -- cycle;
\fill[teal!33!white] (0.8571429,0) -- (1.142857,0.2142857) -- (0.8571429,0.2142857) -- cycle;
\fill[teal!46!white] (1.142857,1.928571) -- (1.142857,1.714286) -- (1.428571,1.714286) -- cycle;
\fill[teal!29!white] (1.428571,1.928571) -- (1.142857,1.928571) -- (1.428571,1.714286) -- cycle;
\fill[teal!26!white] (0.8571429,2.142857) -- (1.142857,1.928571) -- (1.142857,2.142857) -- cycle;
\fill[teal!42!white] (1.142857,1.928571) -- (0.8571429,2.142857) -- (0.8571429,1.928571) -- cycle;
\fill[teal!12!white] (0,2.785714) -- (0,2.571429) -- (0.2857143,2.571429) -- cycle;
\fill[teal!12!white] (0.2857143,2.785714) -- (0,2.785714) -- (0.2857143,2.571429) -- cycle;
\fill[teal!59!white] (1.142857,1.928571) -- (0.8571429,1.714286) -- (1.142857,1.714286) -- cycle;
\fill[teal!56!white] (0.8571429,1.714286) -- (1.142857,1.928571) -- (0.8571429,1.928571) -- cycle;
\fill[teal!66!white] (0.8571429,1.714286) -- (0.5714286,1.928571) -- (0.5714286,1.714286) -- cycle;
\fill[teal!61!white] (0.5714286,1.928571) -- (0.8571429,1.714286) -- (0.8571429,1.928571) -- cycle;
\fill[teal!24!white] (0.2857143,2.357143) -- (0.5714286,2.357143) -- (0.2857143,2.571429) -- cycle;
\fill[teal!18!white] (0.5714286,2.357143) -- (0.5714286,2.571429) -- (0.2857143,2.571429) -- cycle;
\fill[teal!22!white] (0,0.6428571) -- (0,0.4285714) -- (0.2857143,0.6428571) -- cycle;
\fill[teal!33!white] (0,0.4285714) -- (0.2857143,0.4285714) -- (0.2857143,0.6428571) -- cycle;
\fill[teal!17!white] (0,0.4285714) -- (0,0.2142857) -- (0.2857143,0.2142857) -- cycle;
\fill[teal!28!white] (0.2857143,0.4285714) -- (0,0.4285714) -- (0.2857143,0.2142857) -- cycle;
\fill[teal!27!white] (2.285714,1.285714) -- (2,1.285714) -- (2.285714,1.071429) -- cycle;
\fill[teal!44!white] (2,1.285714) -- (2,1.071429) -- (2.285714,1.071429) -- cycle;
\fill[teal!56!white] (2,1.071429) -- (2,0.8571429) -- (2.285714,1.071429) -- cycle;
\fill[teal!53!white] (2,0.8571429) -- (2.285714,0.8571429) -- (2.285714,1.071429) -- cycle;
\fill[teal!39!white] (2.285714,0.8571429) -- (2.571429,0.8571429) -- (2.285714,1.071429) -- cycle;
\fill[teal!25!white] (2.571429,0.8571429) -- (2.571429,1.071429) -- (2.285714,1.071429) -- cycle;
\fill[teal!61!white] (1.428571,0.4285714) -- (1.714286,0.2142857) -- (1.714286,0.4285714) -- cycle;
\fill[teal!54!white] (1.714286,0.2142857) -- (1.428571,0.4285714) -- (1.428571,0.2142857) -- cycle;
\fill[teal!14!white] (3.142857,0.2142857) -- (3.428571,0) -- (3.428571,0.2142857) -- cycle;
\fill[teal!12!white] (3.142857,0.2142857) -- (3.142857,0) -- (3.428571,0) -- cycle;
\fill[teal!13!white] (3.142857,0) -- (3.142857,0.2142857) -- (2.857143,0) -- cycle;
\fill[teal!19!white] (3.142857,0.2142857) -- (2.857143,0.2142857) -- (2.857143,0) -- cycle;
\fill[teal!32!white] (1.714286,0) -- (2,0.2142857) -- (1.714286,0.2142857) -- cycle;
\fill[teal!20!white] (2,0.2142857) -- (1.714286,0) -- (2,0) -- cycle;
\fill[teal!42!white] (0.5714286,0.2142857) -- (0.5714286,0.4285714) -- (0.2857143,0.2142857) -- cycle;
\fill[teal!42!white] (0.5714286,0.4285714) -- (0.2857143,0.4285714) -- (0.2857143,0.2142857) -- cycle;
\fill[teal!50!white] (0.8571429,0.4285714) -- (0.5714286,0.2142857) -- (0.8571429,0.2142857) -- cycle;
\fill[teal!56!white] (0.8571429,0.4285714) -- (0.5714286,0.4285714) -- (0.5714286,0.2142857) -- cycle;
\fill[teal!63!white] (1.142857,1.5) -- (1.428571,1.5) -- (1.428571,1.714286) -- cycle;
\fill[teal!62!white] (1.142857,1.714286) -- (1.142857,1.5) -- (1.428571,1.714286) -- cycle;
\fill[teal!42!white] (0.5714286,2.142857) -- (0.2857143,2.142857) -- (0.2857143,1.928571) -- cycle;
\fill[teal!49!white] (0.5714286,1.928571) -- (0.5714286,2.142857) -- (0.2857143,1.928571) -- cycle;
\fill[teal!35!white] (0.5714286,2.142857) -- (0.2857143,2.357143) -- (0.2857143,2.142857) -- cycle;
\fill[teal!33!white] (0.5714286,2.142857) -- (0.5714286,2.357143) -- (0.2857143,2.357143) -- cycle;
\fill[teal!44!white] (0.8571429,2.142857) -- (0.5714286,2.142857) -- (0.8571429,1.928571) -- cycle;
\fill[teal!51!white] (0.5714286,2.142857) -- (0.5714286,1.928571) -- (0.8571429,1.928571) -- cycle;
\fill[teal!30!white] (0.5714286,2.142857) -- (0.8571429,2.142857) -- (0.8571429,2.357143) -- cycle;
\fill[teal!28!white] (0.5714286,2.357143) -- (0.5714286,2.142857) -- (0.8571429,2.357143) -- cycle;
\fill[teal!38!white] (1.714286,1.285714) -- (2,1.285714) -- (2,1.5) -- cycle;
\fill[teal!39!white] (1.714286,1.5) -- (1.714286,1.285714) -- (2,1.5) -- cycle;
\fill[teal!94!white] (1.142857,1.071429) -- (1.428571,0.8571429) -- (1.428571,1.071429) -- cycle;
\fill[teal!94!white] (1.428571,0.8571429) -- (1.142857,1.071429) -- (1.142857,0.8571429) -- cycle;
\fill[teal!18!white] (3.142857,0.4285714) -- (3.142857,0.2142857) -- (3.428571,0.2142857) -- cycle;
\fill[teal!14!white] (3.428571,0.4285714) -- (3.142857,0.4285714) -- (3.428571,0.2142857) -- cycle;
\fill[teal!26!white] (2.857143,0.4285714) -- (3.142857,0.4285714) -- (2.857143,0.6428571) -- cycle;
\fill[teal!18!white] (3.142857,0.4285714) -- (3.142857,0.6428571) -- (2.857143,0.6428571) -- cycle;
\fill[teal!23!white] (3.142857,0.2142857) -- (3.142857,0.4285714) -- (2.857143,0.2142857) -- cycle;
\fill[teal!27!white] (3.142857,0.4285714) -- (2.857143,0.4285714) -- (2.857143,0.2142857) -- cycle;
\fill[teal!28!white] (2.571429,0.8571429) -- (2.571429,0.6428571) -- (2.857143,0.8571429) -- cycle;
\fill[teal!27!white] (2.857143,0.8571429) -- (2.571429,0.6428571) -- (2.857143,0.6428571) -- cycle;
\fill[teal!19!white] (2.285714,0) -- (2.285714,0.2142857) -- (2,0) -- cycle;
\fill[teal!29!white] (2.285714,0.2142857) -- (2,0.2142857) -- (2,0) -- cycle;
\fill[teal!46!white] (2.285714,0.2142857) -- (2,0.4285714) -- (2,0.2142857) -- cycle;
\fill[teal!50!white] (2,0.4285714) -- (2.285714,0.2142857) -- (2.285714,0.4285714) -- cycle;
\fill[teal!70!white] (2,0.4285714) -- (1.714286,0.6428571) -- (1.714286,0.4285714) -- cycle;
\fill[teal!71!white] (1.714286,0.6428571) -- (2,0.4285714) -- (2,0.6428571) -- cycle;
\fill[teal!58!white] (1.714286,0.2142857) -- (2,0.4285714) -- (1.714286,0.4285714) -- cycle;
\fill[teal!49!white] (2,0.2142857) -- (2,0.4285714) -- (1.714286,0.2142857) -- cycle;
\fill[teal!24!white] (0.2857143,0.8571429) -- (0,0.8571429) -- (0,0.6428571) -- cycle;
\fill[teal!37!white] (0.2857143,0.8571429) -- (0,0.6428571) -- (0.2857143,0.6428571) -- cycle;
\fill[teal!40!white] (0,1.071429) -- (0.2857143,0.8571429) -- (0.2857143,1.071429) -- cycle;
\fill[teal!25!white] (0,0.8571429) -- (0.2857143,0.8571429) -- (0,1.071429) -- cycle;
\fill[teal!70!white] (0.2857143,1.5) -- (0.5714286,1.285714) -- (0.5714286,1.5) -- cycle;
\fill[teal!63!white] (0.5714286,1.285714) -- (0.2857143,1.5) -- (0.2857143,1.285714) -- cycle;
\fill[teal!59!white] (0.5714286,0.6428571) -- (0.2857143,0.8571429) -- (0.2857143,0.6428571) -- cycle;
\fill[teal!70!white] (0.2857143,0.8571429) -- (0.5714286,0.6428571) -- (0.5714286,0.8571429) -- cycle;
\fill[teal!54!white] (0.2857143,0.4285714) -- (0.5714286,0.6428571) -- (0.2857143,0.6428571) -- cycle;
\fill[teal!58!white] (0.5714286,0.4285714) -- (0.5714286,0.6428571) -- (0.2857143,0.4285714) -- cycle;
\fill[teal!68!white] (0.8571429,0.4285714) -- (0.5714286,0.6428571) -- (0.5714286,0.4285714) -- cycle;
\fill[teal!77!white] (0.5714286,0.6428571) -- (0.8571429,0.4285714) -- (0.8571429,0.6428571) -- cycle;
\fill[teal!55!white] (1.142857,0.4285714) -- (1.142857,0.2142857) -- (1.428571,0.2142857) -- cycle;
\fill[teal!64!white] (1.428571,0.4285714) -- (1.142857,0.4285714) -- (1.428571,0.2142857) -- cycle;
\fill[teal!54!white] (1.142857,0.2142857) -- (1.142857,0.4285714) -- (0.8571429,0.2142857) -- cycle;
\fill[teal!62!white] (1.142857,0.4285714) -- (0.8571429,0.4285714) -- (0.8571429,0.2142857) -- cycle;
\fill[teal!76!white] (0.8571429,0.4285714) -- (1.142857,0.4285714) -- (0.8571429,0.6428571) -- cycle;
\fill[teal!83!white] (1.142857,0.4285714) -- (1.142857,0.6428571) -- (0.8571429,0.6428571) -- cycle;
\fill[teal!64!white] (1.428571,1.285714) -- (1.714286,1.5) -- (1.428571,1.5) -- cycle;
\fill[teal!63!white] (1.428571,1.285714) -- (1.714286,1.285714) -- (1.714286,1.5) -- cycle;
\fill[teal!76!white] (1.142857,1.5) -- (1.428571,1.285714) -- (1.428571,1.5) -- cycle;
\fill[teal!85!white] (1.428571,1.285714) -- (1.142857,1.5) -- (1.142857,1.285714) -- cycle;
\fill[teal!92!white] (1.428571,1.285714) -- (1.142857,1.071429) -- (1.428571,1.071429) -- cycle;
\fill[teal!93!white] (1.142857,1.071429) -- (1.428571,1.285714) -- (1.142857,1.285714) -- cycle;
\fill[teal!76!white] (1.428571,0.6428571) -- (1.428571,0.4285714) -- (1.714286,0.4285714) -- cycle;
\fill[teal!79!white] (1.714286,0.6428571) -- (1.428571,0.6428571) -- (1.714286,0.4285714) -- cycle;
\fill[teal!78!white] (1.428571,0.6428571) -- (1.142857,0.4285714) -- (1.428571,0.4285714) -- cycle;
\fill[teal!84!white] (1.142857,0.4285714) -- (1.428571,0.6428571) -- (1.142857,0.6428571) -- cycle;
\fill[teal!94!white] (1.428571,0.6428571) -- (1.428571,0.8571429) -- (1.142857,0.8571429) -- cycle;
\fill[teal!92!white] (1.142857,0.6428571) -- (1.428571,0.6428571) -- (1.142857,0.8571429) -- cycle;
\fill[teal!60!white] (2,1.285714) -- (1.714286,1.071429) -- (2,1.071429) -- cycle;
\fill[teal!61!white] (1.714286,1.285714) -- (1.714286,1.071429) -- (2,1.285714) -- cycle;
\fill[teal!75!white] (1.428571,1.285714) -- (1.714286,1.071429) -- (1.714286,1.285714) -- cycle;
\fill[teal!85!white] (1.714286,1.071429) -- (1.428571,1.285714) -- (1.428571,1.071429) -- cycle;
\fill[teal!48!white] (2.285714,0.6428571) -- (2.571429,0.8571429) -- (2.285714,0.8571429) -- cycle;
\fill[teal!45!white] (2.285714,0.6428571) -- (2.571429,0.6428571) -- (2.571429,0.8571429) -- cycle;
\fill[teal!60!white] (2,0.8571429) -- (2.285714,0.6428571) -- (2.285714,0.8571429) -- cycle;
\fill[teal!66!white] (2.285714,0.6428571) -- (2,0.8571429) -- (2,0.6428571) -- cycle;
\fill[teal!58!white] (2.285714,0.6428571) -- (2,0.4285714) -- (2.285714,0.4285714) -- cycle;
\fill[teal!64!white] (2,0.4285714) -- (2.285714,0.6428571) -- (2,0.6428571) -- cycle;
\fill[teal!34!white] (2.571429,0.4285714) -- (2.857143,0.4285714) -- (2.857143,0.6428571) -- cycle;
\fill[teal!38!white] (2.571429,0.6428571) -- (2.571429,0.4285714) -- (2.857143,0.6428571) -- cycle;
\fill[teal!48!white] (2.285714,0.6428571) -- (2.571429,0.4285714) -- (2.571429,0.6428571) -- cycle;
\fill[teal!51!white] (2.571429,0.4285714) -- (2.285714,0.6428571) -- (2.285714,0.4285714) -- cycle;
\fill[teal!64!white] (0.2857143,0.8571429) -- (0.5714286,1.071429) -- (0.2857143,1.071429) -- cycle;
\fill[teal!73!white] (0.5714286,1.071429) -- (0.2857143,0.8571429) -- (0.5714286,0.8571429) -- cycle;
\fill[teal!65!white] (0.5714286,1.071429) -- (0.2857143,1.285714) -- (0.2857143,1.071429) -- cycle;
\fill[teal!73!white] (0.5714286,1.071429) -- (0.5714286,1.285714) -- (0.2857143,1.285714) -- cycle;
\fill[teal!77!white] (0.8571429,1.5) -- (0.5714286,1.714286) -- (0.5714286,1.5) -- cycle;
\fill[teal!75!white] (0.8571429,1.5) -- (0.8571429,1.714286) -- (0.5714286,1.714286) -- cycle;
\fill[teal!74!white] (0.8571429,1.714286) -- (0.8571429,1.5) -- (1.142857,1.714286) -- cycle;
\fill[teal!77!white] (0.8571429,1.5) -- (1.142857,1.5) -- (1.142857,1.714286) -- cycle;
\fill[teal!79!white] (1.714286,0.8571429) -- (1.714286,0.6428571) -- (2,0.6428571) -- cycle;
\fill[teal!75!white] (2,0.8571429) -- (1.714286,0.8571429) -- (2,0.6428571) -- cycle;
\fill[teal!72!white] (1.714286,0.8571429) -- (2,0.8571429) -- (2,1.071429) -- cycle;
\fill[teal!75!white] (1.714286,1.071429) -- (1.714286,0.8571429) -- (2,1.071429) -- cycle;
\fill[teal!86!white] (1.714286,0.8571429) -- (1.428571,0.6428571) -- (1.714286,0.6428571) -- cycle;
\fill[teal!90!white] (1.428571,0.6428571) -- (1.714286,0.8571429) -- (1.428571,0.8571429) -- cycle;
\fill[teal!91!white] (1.428571,0.8571429) -- (1.714286,0.8571429) -- (1.428571,1.071429) -- cycle;
\fill[teal!85!white] (1.714286,0.8571429) -- (1.714286,1.071429) -- (1.428571,1.071429) -- cycle;
\fill[teal!17!white] (2.571429,0) -- (2.571429,0.2142857) -- (2.285714,0) -- cycle;
\fill[teal!26!white] (2.571429,0.2142857) -- (2.285714,0.2142857) -- (2.285714,0) -- cycle;
\fill[teal!40!white] (2.285714,0.2142857) -- (2.571429,0.2142857) -- (2.285714,0.4285714) -- cycle;
\fill[teal!42!white] (2.571429,0.2142857) -- (2.571429,0.4285714) -- (2.285714,0.4285714) -- cycle;
\fill[teal!22!white] (2.857143,0.2142857) -- (2.571429,0.2142857) -- (2.857143,0) -- cycle;
\fill[teal!16!white] (2.571429,0.2142857) -- (2.571429,0) -- (2.857143,0) -- cycle;
\fill[teal!31!white] (2.857143,0.4285714) -- (2.571429,0.2142857) -- (2.857143,0.2142857) -- cycle;
\fill[teal!36!white] (2.571429,0.4285714) -- (2.571429,0.2142857) -- (2.857143,0.4285714) -- cycle;
\fill[teal!90!white] (1.142857,0.6428571) -- (0.8571429,0.8571429) -- (0.8571429,0.6428571) -- cycle;
\fill[teal!94!white] (0.8571429,0.8571429) -- (1.142857,0.6428571) -- (1.142857,0.8571429) -- cycle;
\fill[teal!83!white] (0.5714286,0.6428571) -- (0.8571429,0.8571429) -- (0.5714286,0.8571429) -- cycle;
\fill[teal!85!white] (0.8571429,0.8571429) -- (0.5714286,0.6428571) -- (0.8571429,0.6428571) -- cycle;
\fill[teal!90!white] (1.142857,1.5) -- (0.8571429,1.285714) -- (1.142857,1.285714) -- cycle;
\fill[teal!87!white] (0.8571429,1.5) -- (0.8571429,1.285714) -- (1.142857,1.5) -- cycle;
\fill[teal!84!white] (0.5714286,1.285714) -- (0.8571429,1.285714) -- (0.5714286,1.5) -- cycle;
\fill[teal!85!white] (0.8571429,1.285714) -- (0.8571429,1.5) -- (0.5714286,1.5) -- cycle;
\fill[teal!94!white] (1.142857,1.071429) -- (0.8571429,1.071429) -- (1.142857,0.8571429) -- cycle;
\fill[teal!94!white] (0.8571429,1.071429) -- (0.8571429,0.8571429) -- (1.142857,0.8571429) -- cycle;
\fill[teal!94!white] (0.8571429,1.071429) -- (1.142857,1.071429) -- (1.142857,1.285714) -- cycle;
\fill[teal!94!white] (0.8571429,1.285714) -- (0.8571429,1.071429) -- (1.142857,1.285714) -- cycle;
\fill[teal!87!white] (0.8571429,1.071429) -- (0.5714286,1.071429) -- (0.5714286,0.8571429) -- cycle;
\fill[teal!91!white] (0.8571429,0.8571429) -- (0.8571429,1.071429) -- (0.5714286,0.8571429) -- cycle;
\fill[teal!87!white] (0.5714286,1.071429) -- (0.8571429,1.071429) -- (0.5714286,1.285714) -- cycle;
\fill[teal!91!white] (0.8571429,1.071429) -- (0.8571429,1.285714) -- (0.5714286,1.285714) -- cycle;
\draw[teal,thick] (0,0) -- (4,0) -- (0,3) -- cycle;
\draw[magenta!70,thin] (1.094988,0.9982276) -- (3.6,2.35) -- (4.1,2.35);\filldraw[fill=magenta,draw=white,line width=.45pt] (1.094988,0.9982276) circle(1.5pt) node[above right,font=\scriptsize,text=magenta,fill=white,inner sep=.6pt] {};
\node[anchor=west,text=magenta,font=\small] at(4.12,2.35){F1};\draw[gold!70,thin] (1,1) -- (3.6,1.8) -- (4.1,1.8);\filldraw[fill=gold,draw=white,line width=.45pt] (1,1) circle(1.5pt) node[above right,font=\scriptsize,text=gold,fill=white,inner sep=.6pt] {};
\node[anchor=west,text=gold,font=\small] at(4.12,1.8){$I$};\draw[navy!70,thin] (1.333333,1) -- (3.6,1.25) -- (4.1,1.25);\filldraw[fill=navy,draw=white,line width=.45pt] (1.333333,1) circle(1.5pt) node[above right,font=\scriptsize,text=navy,fill=white,inner sep=.6pt] {};
\node[anchor=west,text=navy,font=\small] at(4.12,1.25){$G$};\end{tikzpicture}
\caption[One torsion field represents a tension membrane, a viscous duct and Brownian exit time after rescaling]{One torsion field represents a tension membrane, a viscous duct and Brownian exit time after rescaling. Fields are archived numerical illustrations; teal saturation increases with the displayed value.}\label{fig:field-F1}
\end{figure}

\paragraph{Demonstration and investigation.} Ask students to estimate where a scalene triangle will sag most before revealing the marker. Explain why the maximum need not be at the point farthest from every side.


\section{One field, several centers}\label{sec:slide36}

The scalar field u supplies at least two natural points. Its maximum gives F1. Its first moment gives the center of the volume beneath its graph, called Abs(0) in the reaction family. The energy density is |\ensuremath{\nabla}u|\ensuremath{^2}. A special affine integration identity says that its centroid equals the u-weighted centroid, even though their spatial densities look very different. This is a genuine equality of two physical means, not a visual coincidence.


\[\mathrm{Abs}(0)=\frac{\int_Tpu\dd A}{\int_Tu\dd A}.\]


Multiply $-\Delta u=1$ by $p_i u$ and integrate by parts.
The term $\int u\,\partial_i u=\frac12\int\partial_i(u^2)$ vanishes
because $u=0$ at the boundary. Hence
\[
\int_Tp_i|\nabla u|^2\dd A=\int_Tp_i u\dd A,\qquad
\int_T|\nabla u|^2\dd A=\int_Tu\dd A.
\]
The elastic strain-energy mean and the displacement/flux mean are
therefore exactly the same center, although their densities differ.



\section{Laminar flow through a triangular duct}\label{sec:slide37}

Take a long straight triangular pipe. In steady fully developed laminar flow, velocity has only an axial component v(x,y). The convective term vanishes because the profile does not vary along the pipe. Viscous diffusion balances the constant pressure gradient. No slip at the walls gives v=0. After division by constants this is the same torsion problem. Therefore the fastest streamline passes through F1, while the flux centroid is Abs(0). The dissipation centroid agrees with the flux centroid by the preceding identity.


\[-\eta_{\rm v}\Delta v=-\frac{\dd p_{\rm pressure}}{\dd z},\quad v|_{\partial T}=0.\]



Take a straight duct with $z$-independent velocity
$\mathbf v=(0,0,v(x,y))$. Incompressibility is automatic and the
convective acceleration vanishes. The axial Navier--Stokes equation
becomes $0=-\partial_zp+\eta_{\rm v}\Delta v$.
A positive pressure drop $g=-\partial_zp$ gives
$v=(g/\eta_{\rm v})u$. No slip gives the Dirichlet boundary condition.

The fastest streamline is F1; the flux centroid is Abs(0).
Total flow is $(g/\eta_{\rm v})\int_Tu\dd A$, a scalar torsional
quantity. One solve therefore supplies a maximum point, a moment
center and a transport coefficient.




\paragraph{Demonstration and investigation.} Show the prism cross section and reveal a slow wall region and fast interior region. Compare a maximum marker with a flux-weighted mean.


\section{Brownian mean exit time}\label{sec:slide38}

A Brownian particle with diffusion coefficient D is absorbed at the boundary. Its expected remaining lifetime obeys this equation. To see why, evolve for a small time dt: the expected lifetime at the initial point equals dt plus the expected lifetime after the small Brownian step, unless the boundary has been hit. The Brownian average adds D\ensuremath{\Delta}\ensuremath{\tau} times dt. Cancel the common \ensuremath{\tau} and divide by dt. Consequently F1 is also the point with the largest expected survival time. Lecture 4 will distinguish this lifetime maximum from long-time conditioned particle means.


\[-D_{\rm B}\Delta m=1,\quad m|_{\partial T}=0,\quad m(p)=\mathbb E_p\tau_{\rm exit}.\]



Before exit, the short-time identity is
$m(p)=\delta t+\mathbb E_p m(X_{\delta t})$.
Brownian increments have covariance
$2D_{\rm B}\delta t\,\mathbf1$, so Taylor expansion gives
$\mathbb E m(X_{\delta t})=m(p)+D_{\rm B}\delta t\Delta m+o(\delta t)$.
Cancel and divide to obtain the Poisson problem.
At the boundary the exit time is zero.
F1 maximizes expected survival; E0 equalizes side-exit probabilities.
These are different Brownian questions.




\section{Membrane tension and plate bending}\label{sec:slide39}

A plate has finite bending stiffness D, proportional to Young's modulus times thickness cubed. Its bending energy contains second derivatives of the deflection. Varying a quadratic curvature energy produces a fourth-order equation. The plate needs two boundary conditions per edge. For a clamped edge, both displacement and normal slope vanish. A membrane needs only the displacement condition because its equation is second order. This distinction changes the field and its center, even when the footprint and load are identical.


\[\mathcal J_{\rm bend}[w]=\frac{D_{\rm b}}2\int_T(\Delta w)^2\dd A-\int_Tqw\dd A.\]



For thickness $h_{\rm p}$, Young modulus $E$ and Poisson ratio $\nu$,
$D_{\rm b}=Eh_{\rm p}^3/[12(1-\nu^2)]$.
The full isotropic bending density is
\[
\frac{D_{\rm b}}2
\left[(w_{xx}+w_{yy})^2
-2(1-\nu)(w_{xx}w_{yy}-w_{xy}^2)\right].
\]
For clamped variations the determinant term is a boundary term,
explaining the simplified energy and bulk equation.
A membrane penalizes slope and yields a second-order PDE.
A plate penalizes curvature and yields a fourth-order PDE with
two boundary conditions per side. Simply supported and clamped
are different experiments.
A clamped triangular plate maximum is a numerical candidate unless
global uniqueness is separately established.




\section{Pgrav: a clamped plate under gravity}\label{sec:slide40}

Here the material has uniform thickness, isotropic stiffness and a uniform gravitational load. In linear small-deflection theory q and D multiply the field amplitude and do not move its maximum. The plate maximum is Pgrav in the atlas. Its numerical position differs from F1 because the energy penalizes curvature rather than slope. For general fourth-order problems a maximum principle is not available in the simple Laplacian form. The project keeps global maximum-uniqueness and attractivity qualifications explicit.


\[D_{\rm b}\Delta^2w=q,\quad w=\partial_nw=0\ \hbox{on }\partial T.\]


A conforming Ritz method uses basis functions containing the square of the barycentric bubble, so both displacement and normal derivative vanish at each edge. Assemble curvature energy, solve the resulting positive system, and refine the maximum search. Compare polynomial degree and integration order separately.


\begin{figure}[htbp]\centering
\begin{tikzpicture}[x=1.7cm,y=1.7cm]
\fill[teal!8!white] (3.142857,0.6428571) -- (2.857143,0.8571429) -- (2.857143,0.6428571) -- cycle;
\fill[teal!10!white] (1.714286,1.714286) -- (1.428571,1.928571) -- (1.428571,1.714286) -- cycle;
\fill[teal!10!white] (1.714286,1.714286) -- (1.714286,1.5) -- (2,1.5) -- cycle;
\fill[teal!8!white] (0.5714286,2.571429) -- (0.2857143,2.785714) -- (0.2857143,2.571429) -- cycle;
\fill[teal!9!white] (2.571429,1.071429) -- (2.285714,1.285714) -- (2.285714,1.071429) -- cycle;
\fill[teal!8!white] (3.714286,0.2142857) -- (3.428571,0.4285714) -- (3.428571,0.2142857) -- cycle;
\fill[teal!8!white] (3.714286,0.2142857) -- (3.714286,0) -- (4,0) -- cycle;
\fill[teal!9!white] (0.8571429,2.142857) -- (1.142857,2.142857) -- (0.8571429,2.357143) -- cycle;
\fill[teal!9!white] (1.142857,2.142857) -- (1.142857,1.928571) -- (1.428571,1.928571) -- cycle;
\fill[teal!8!white] (0,2.785714) -- (0.2857143,2.785714) -- (0,3) -- cycle;
\fill[teal!8!white] (0.5714286,2.571429) -- (0.5714286,2.357143) -- (0.8571429,2.357143) -- cycle;
\fill[teal!9!white] (2,1.285714) -- (2.285714,1.285714) -- (2,1.5) -- cycle;
\fill[teal!9!white] (2.571429,1.071429) -- (2.571429,0.8571429) -- (2.857143,0.8571429) -- cycle;
\fill[teal!8!white] (3.142857,0.4285714) -- (3.428571,0.4285714) -- (3.142857,0.6428571) -- cycle;
\fill[teal!8!white] (0.2857143,0) -- (0,0.2142857) -- (0,0) -- cycle;
\fill[teal!9!white] (0,0.2142857) -- (0.2857143,0) -- (0.2857143,0.2142857) -- cycle;
\fill[teal!9!white] (0.5714286,0.2142857) -- (0.2857143,0) -- (0.5714286,0) -- cycle;
\fill[teal!11!white] (0.2857143,0) -- (0.5714286,0.2142857) -- (0.2857143,0.2142857) -- cycle;
\fill[teal!14!white] (1.714286,1.5) -- (1.714286,1.714286) -- (1.428571,1.714286) -- cycle;
\fill[teal!27!white] (1.428571,1.5) -- (1.714286,1.5) -- (1.428571,1.714286) -- cycle;
\fill[teal!10!white] (0,1.928571) -- (0,1.714286) -- (0.2857143,1.928571) -- cycle;
\fill[teal!16!white] (0,1.714286) -- (0.2857143,1.714286) -- (0.2857143,1.928571) -- cycle;
\fill[teal!8!white] (0.2857143,2.357143) -- (0,2.571429) -- (0,2.357143) -- cycle;
\fill[teal!9!white] (0,2.571429) -- (0.2857143,2.357143) -- (0.2857143,2.571429) -- cycle;
\fill[teal!9!white] (0,2.142857) -- (0.2857143,2.357143) -- (0,2.357143) -- cycle;
\fill[teal!11!white] (0.2857143,2.357143) -- (0,2.142857) -- (0.2857143,2.142857) -- cycle;
\fill[teal!10!white] (0,2.142857) -- (0,1.928571) -- (0.2857143,1.928571) -- cycle;
\fill[teal!12!white] (0.2857143,2.142857) -- (0,2.142857) -- (0.2857143,1.928571) -- cycle;
\fill[teal!27!white] (0.2857143,1.714286) -- (0.5714286,1.928571) -- (0.2857143,1.928571) -- cycle;
\fill[teal!36!white] (0.5714286,1.714286) -- (0.5714286,1.928571) -- (0.2857143,1.714286) -- cycle;
\fill[teal!11!white] (0.2857143,1.5) -- (0,1.714286) -- (0,1.5) -- cycle;
\fill[teal!17!white] (0,1.714286) -- (0.2857143,1.5) -- (0.2857143,1.714286) -- cycle;
\fill[teal!35!white] (0.2857143,1.5) -- (0.5714286,1.714286) -- (0.2857143,1.714286) -- cycle;
\fill[teal!47!white] (0.5714286,1.714286) -- (0.2857143,1.5) -- (0.5714286,1.5) -- cycle;
\fill[teal!11!white] (0,1.285714) -- (0.2857143,1.5) -- (0,1.5) -- cycle;
\fill[teal!20!white] (0.2857143,1.5) -- (0,1.285714) -- (0.2857143,1.285714) -- cycle;
\fill[teal!12!white] (0,1.285714) -- (0,1.071429) -- (0.2857143,1.071429) -- cycle;
\fill[teal!20!white] (0.2857143,1.285714) -- (0,1.285714) -- (0.2857143,1.071429) -- cycle;
\fill[teal!8!white] (3.428571,0) -- (3.714286,0) -- (3.428571,0.2142857) -- cycle;
\fill[teal!8!white] (3.714286,0) -- (3.714286,0.2142857) -- (3.428571,0.2142857) -- cycle;
\fill[teal!10!white] (1.428571,0) -- (1.714286,0) -- (1.428571,0.2142857) -- cycle;
\fill[teal!15!white] (1.714286,0) -- (1.714286,0.2142857) -- (1.428571,0.2142857) -- cycle;
\fill[teal!10!white] (0.8571429,0) -- (0.5714286,0.2142857) -- (0.5714286,0) -- cycle;
\fill[teal!15!white] (0.5714286,0.2142857) -- (0.8571429,0) -- (0.8571429,0.2142857) -- cycle;
\fill[teal!10!white] (1.142857,0.2142857) -- (1.142857,0) -- (1.428571,0) -- cycle;
\fill[teal!16!white] (1.142857,0.2142857) -- (1.428571,0) -- (1.428571,0.2142857) -- cycle;
\fill[teal!10!white] (1.142857,0.2142857) -- (0.8571429,0) -- (1.142857,0) -- cycle;
\fill[teal!16!white] (0.8571429,0) -- (1.142857,0.2142857) -- (0.8571429,0.2142857) -- cycle;
\fill[teal!25!white] (1.142857,1.928571) -- (1.142857,1.714286) -- (1.428571,1.714286) -- cycle;
\fill[teal!13!white] (1.428571,1.928571) -- (1.142857,1.928571) -- (1.428571,1.714286) -- cycle;
\fill[teal!12!white] (0.8571429,2.142857) -- (1.142857,1.928571) -- (1.142857,2.142857) -- cycle;
\fill[teal!21!white] (1.142857,1.928571) -- (0.8571429,2.142857) -- (0.8571429,1.928571) -- cycle;
\fill[teal!8!white] (0,2.785714) -- (0,2.571429) -- (0.2857143,2.571429) -- cycle;
\fill[teal!8!white] (0.2857143,2.785714) -- (0,2.785714) -- (0.2857143,2.571429) -- cycle;
\fill[teal!37!white] (1.142857,1.928571) -- (0.8571429,1.714286) -- (1.142857,1.714286) -- cycle;
\fill[teal!34!white] (0.8571429,1.714286) -- (1.142857,1.928571) -- (0.8571429,1.928571) -- cycle;
\fill[teal!45!white] (0.8571429,1.714286) -- (0.5714286,1.928571) -- (0.5714286,1.714286) -- cycle;
\fill[teal!39!white] (0.5714286,1.928571) -- (0.8571429,1.714286) -- (0.8571429,1.928571) -- cycle;
\fill[teal!10!white] (0.2857143,2.357143) -- (0.5714286,2.357143) -- (0.2857143,2.571429) -- cycle;
\fill[teal!9!white] (0.5714286,2.357143) -- (0.5714286,2.571429) -- (0.2857143,2.571429) -- cycle;
\fill[teal!11!white] (0,0.6428571) -- (0,0.4285714) -- (0.2857143,0.6428571) -- cycle;
\fill[teal!15!white] (0,0.4285714) -- (0.2857143,0.4285714) -- (0.2857143,0.6428571) -- cycle;
\fill[teal!9!white] (0,0.4285714) -- (0,0.2142857) -- (0.2857143,0.2142857) -- cycle;
\fill[teal!13!white] (0.2857143,0.4285714) -- (0,0.4285714) -- (0.2857143,0.2142857) -- cycle;
\fill[teal!12!white] (2.285714,1.285714) -- (2,1.285714) -- (2.285714,1.071429) -- cycle;
\fill[teal!22!white] (2,1.285714) -- (2,1.071429) -- (2.285714,1.071429) -- cycle;
\fill[teal!32!white] (2,1.071429) -- (2,0.8571429) -- (2.285714,1.071429) -- cycle;
\fill[teal!29!white] (2,0.8571429) -- (2.285714,0.8571429) -- (2.285714,1.071429) -- cycle;
\fill[teal!18!white] (2.285714,0.8571429) -- (2.571429,0.8571429) -- (2.285714,1.071429) -- cycle;
\fill[teal!11!white] (2.571429,0.8571429) -- (2.571429,1.071429) -- (2.285714,1.071429) -- cycle;
\fill[teal!40!white] (1.428571,0.4285714) -- (1.714286,0.2142857) -- (1.714286,0.4285714) -- cycle;
\fill[teal!32!white] (1.714286,0.2142857) -- (1.428571,0.4285714) -- (1.428571,0.2142857) -- cycle;
\fill[teal!8!white] (3.142857,0.2142857) -- (3.428571,0) -- (3.428571,0.2142857) -- cycle;
\fill[teal!8!white] (3.142857,0.2142857) -- (3.142857,0) -- (3.428571,0) -- cycle;
\fill[teal!8!white] (3.142857,0) -- (3.142857,0.2142857) -- (2.857143,0) -- cycle;
\fill[teal!9!white] (3.142857,0.2142857) -- (2.857143,0.2142857) -- (2.857143,0) -- cycle;
\fill[teal!14!white] (1.714286,0) -- (2,0.2142857) -- (1.714286,0.2142857) -- cycle;
\fill[teal!10!white] (2,0.2142857) -- (1.714286,0) -- (2,0) -- cycle;
\fill[teal!21!white] (0.5714286,0.2142857) -- (0.5714286,0.4285714) -- (0.2857143,0.2142857) -- cycle;
\fill[teal!22!white] (0.5714286,0.4285714) -- (0.2857143,0.4285714) -- (0.2857143,0.2142857) -- cycle;
\fill[teal!30!white] (0.8571429,0.4285714) -- (0.5714286,0.2142857) -- (0.8571429,0.2142857) -- cycle;
\fill[teal!35!white] (0.8571429,0.4285714) -- (0.5714286,0.4285714) -- (0.5714286,0.2142857) -- cycle;
\fill[teal!43!white] (1.142857,1.5) -- (1.428571,1.5) -- (1.428571,1.714286) -- cycle;
\fill[teal!42!white] (1.142857,1.714286) -- (1.142857,1.5) -- (1.428571,1.714286) -- cycle;
\fill[teal!20!white] (0.5714286,2.142857) -- (0.2857143,2.142857) -- (0.2857143,1.928571) -- cycle;
\fill[teal!26!white] (0.5714286,1.928571) -- (0.5714286,2.142857) -- (0.2857143,1.928571) -- cycle;
\fill[teal!15!white] (0.5714286,2.142857) -- (0.2857143,2.357143) -- (0.2857143,2.142857) -- cycle;
\fill[teal!14!white] (0.5714286,2.142857) -- (0.5714286,2.357143) -- (0.2857143,2.357143) -- cycle;
\fill[teal!22!white] (0.8571429,2.142857) -- (0.5714286,2.142857) -- (0.8571429,1.928571) -- cycle;
\fill[teal!28!white] (0.5714286,2.142857) -- (0.5714286,1.928571) -- (0.8571429,1.928571) -- cycle;
\fill[teal!13!white] (0.5714286,2.142857) -- (0.8571429,2.142857) -- (0.8571429,2.357143) -- cycle;
\fill[teal!12!white] (0.5714286,2.357143) -- (0.5714286,2.142857) -- (0.8571429,2.357143) -- cycle;
\fill[teal!18!white] (1.714286,1.285714) -- (2,1.285714) -- (2,1.5) -- cycle;
\fill[teal!19!white] (1.714286,1.5) -- (1.714286,1.285714) -- (2,1.5) -- cycle;
\fill[teal!94!white] (1.142857,1.071429) -- (1.428571,0.8571429) -- (1.428571,1.071429) -- cycle;
\fill[teal!94!white] (1.428571,0.8571429) -- (1.142857,1.071429) -- (1.142857,0.8571429) -- cycle;
\fill[teal!9!white] (3.142857,0.4285714) -- (3.142857,0.2142857) -- (3.428571,0.2142857) -- cycle;
\fill[teal!8!white] (3.428571,0.4285714) -- (3.142857,0.4285714) -- (3.428571,0.2142857) -- cycle;
\fill[teal!11!white] (2.857143,0.4285714) -- (3.142857,0.4285714) -- (2.857143,0.6428571) -- cycle;
\fill[teal!9!white] (3.142857,0.4285714) -- (3.142857,0.6428571) -- (2.857143,0.6428571) -- cycle;
\fill[teal!10!white] (3.142857,0.2142857) -- (3.142857,0.4285714) -- (2.857143,0.2142857) -- cycle;
\fill[teal!11!white] (3.142857,0.4285714) -- (2.857143,0.4285714) -- (2.857143,0.2142857) -- cycle;
\fill[teal!12!white] (2.571429,0.8571429) -- (2.571429,0.6428571) -- (2.857143,0.8571429) -- cycle;
\fill[teal!11!white] (2.857143,0.8571429) -- (2.571429,0.6428571) -- (2.857143,0.6428571) -- cycle;
\fill[teal!9!white] (2.285714,0) -- (2.285714,0.2142857) -- (2,0) -- cycle;
\fill[teal!13!white] (2.285714,0.2142857) -- (2,0.2142857) -- (2,0) -- cycle;
\fill[teal!23!white] (2.285714,0.2142857) -- (2,0.4285714) -- (2,0.2142857) -- cycle;
\fill[teal!26!white] (2,0.4285714) -- (2.285714,0.2142857) -- (2.285714,0.4285714) -- cycle;
\fill[teal!50!white] (2,0.4285714) -- (1.714286,0.6428571) -- (1.714286,0.4285714) -- cycle;
\fill[teal!51!white] (1.714286,0.6428571) -- (2,0.4285714) -- (2,0.6428571) -- cycle;
\fill[teal!35!white] (1.714286,0.2142857) -- (2,0.4285714) -- (1.714286,0.4285714) -- cycle;
\fill[teal!26!white] (2,0.2142857) -- (2,0.4285714) -- (1.714286,0.2142857) -- cycle;
\fill[teal!11!white] (0.2857143,0.8571429) -- (0,0.8571429) -- (0,0.6428571) -- cycle;
\fill[teal!18!white] (0.2857143,0.8571429) -- (0,0.6428571) -- (0.2857143,0.6428571) -- cycle;
\fill[teal!20!white] (0,1.071429) -- (0.2857143,0.8571429) -- (0.2857143,1.071429) -- cycle;
\fill[teal!12!white] (0,0.8571429) -- (0.2857143,0.8571429) -- (0,1.071429) -- cycle;
\fill[teal!53!white] (0.2857143,1.5) -- (0.5714286,1.285714) -- (0.5714286,1.5) -- cycle;
\fill[teal!43!white] (0.5714286,1.285714) -- (0.2857143,1.5) -- (0.2857143,1.285714) -- cycle;
\fill[teal!39!white] (0.5714286,0.6428571) -- (0.2857143,0.8571429) -- (0.2857143,0.6428571) -- cycle;
\fill[teal!54!white] (0.2857143,0.8571429) -- (0.5714286,0.6428571) -- (0.5714286,0.8571429) -- cycle;
\fill[teal!33!white] (0.2857143,0.4285714) -- (0.5714286,0.6428571) -- (0.2857143,0.6428571) -- cycle;
\fill[teal!37!white] (0.5714286,0.4285714) -- (0.5714286,0.6428571) -- (0.2857143,0.4285714) -- cycle;
\fill[teal!51!white] (0.8571429,0.4285714) -- (0.5714286,0.6428571) -- (0.5714286,0.4285714) -- cycle;
\fill[teal!64!white] (0.5714286,0.6428571) -- (0.8571429,0.4285714) -- (0.8571429,0.6428571) -- cycle;
\fill[teal!34!white] (1.142857,0.4285714) -- (1.142857,0.2142857) -- (1.428571,0.2142857) -- cycle;
\fill[teal!44!white] (1.428571,0.4285714) -- (1.142857,0.4285714) -- (1.428571,0.2142857) -- cycle;
\fill[teal!34!white] (1.142857,0.2142857) -- (1.142857,0.4285714) -- (0.8571429,0.2142857) -- cycle;
\fill[teal!43!white] (1.142857,0.4285714) -- (0.8571429,0.4285714) -- (0.8571429,0.2142857) -- cycle;
\fill[teal!64!white] (0.8571429,0.4285714) -- (1.142857,0.4285714) -- (0.8571429,0.6428571) -- cycle;
\fill[teal!75!white] (1.142857,0.4285714) -- (1.142857,0.6428571) -- (0.8571429,0.6428571) -- cycle;
\fill[teal!44!white] (1.428571,1.285714) -- (1.714286,1.5) -- (1.428571,1.5) -- cycle;
\fill[teal!43!white] (1.428571,1.285714) -- (1.714286,1.285714) -- (1.714286,1.5) -- cycle;
\fill[teal!62!white] (1.142857,1.5) -- (1.428571,1.285714) -- (1.428571,1.5) -- cycle;
\fill[teal!78!white] (1.428571,1.285714) -- (1.142857,1.5) -- (1.142857,1.285714) -- cycle;
\fill[teal!91!white] (1.428571,1.285714) -- (1.142857,1.071429) -- (1.428571,1.071429) -- cycle;
\fill[teal!92!white] (1.142857,1.071429) -- (1.428571,1.285714) -- (1.142857,1.285714) -- cycle;
\fill[teal!60!white] (1.428571,0.6428571) -- (1.428571,0.4285714) -- (1.714286,0.4285714) -- cycle;
\fill[teal!65!white] (1.714286,0.6428571) -- (1.428571,0.6428571) -- (1.714286,0.4285714) -- cycle;
\fill[teal!65!white] (1.428571,0.6428571) -- (1.142857,0.4285714) -- (1.428571,0.4285714) -- cycle;
\fill[teal!76!white] (1.142857,0.4285714) -- (1.428571,0.6428571) -- (1.142857,0.6428571) -- cycle;
\fill[teal!94!white] (1.428571,0.6428571) -- (1.428571,0.8571429) -- (1.142857,0.8571429) -- cycle;
\fill[teal!91!white] (1.142857,0.6428571) -- (1.428571,0.6428571) -- (1.142857,0.8571429) -- cycle;
\fill[teal!38!white] (2,1.285714) -- (1.714286,1.071429) -- (2,1.071429) -- cycle;
\fill[teal!40!white] (1.714286,1.285714) -- (1.714286,1.071429) -- (2,1.285714) -- cycle;
\fill[teal!59!white] (1.428571,1.285714) -- (1.714286,1.071429) -- (1.714286,1.285714) -- cycle;
\fill[teal!75!white] (1.714286,1.071429) -- (1.428571,1.285714) -- (1.428571,1.071429) -- cycle;
\fill[teal!24!white] (2.285714,0.6428571) -- (2.571429,0.8571429) -- (2.285714,0.8571429) -- cycle;
\fill[teal!21!white] (2.285714,0.6428571) -- (2.571429,0.6428571) -- (2.571429,0.8571429) -- cycle;
\fill[teal!35!white] (2,0.8571429) -- (2.285714,0.6428571) -- (2.285714,0.8571429) -- cycle;
\fill[teal!43!white] (2.285714,0.6428571) -- (2,0.8571429) -- (2,0.6428571) -- cycle;
\fill[teal!33!white] (2.285714,0.6428571) -- (2,0.4285714) -- (2.285714,0.4285714) -- cycle;
\fill[teal!40!white] (2,0.4285714) -- (2.285714,0.6428571) -- (2,0.6428571) -- cycle;
\fill[teal!14!white] (2.571429,0.4285714) -- (2.857143,0.4285714) -- (2.857143,0.6428571) -- cycle;
\fill[teal!16!white] (2.571429,0.6428571) -- (2.571429,0.4285714) -- (2.857143,0.6428571) -- cycle;
\fill[teal!23!white] (2.285714,0.6428571) -- (2.571429,0.4285714) -- (2.571429,0.6428571) -- cycle;
\fill[teal!26!white] (2.571429,0.4285714) -- (2.285714,0.6428571) -- (2.285714,0.4285714) -- cycle;
\fill[teal!46!white] (0.2857143,0.8571429) -- (0.5714286,1.071429) -- (0.2857143,1.071429) -- cycle;
\fill[teal!58!white] (0.5714286,1.071429) -- (0.2857143,0.8571429) -- (0.5714286,0.8571429) -- cycle;
\fill[teal!46!white] (0.5714286,1.071429) -- (0.2857143,1.285714) -- (0.2857143,1.071429) -- cycle;
\fill[teal!59!white] (0.5714286,1.071429) -- (0.5714286,1.285714) -- (0.2857143,1.285714) -- cycle;
\fill[teal!62!white] (0.8571429,1.5) -- (0.5714286,1.714286) -- (0.5714286,1.5) -- cycle;
\fill[teal!60!white] (0.8571429,1.5) -- (0.8571429,1.714286) -- (0.5714286,1.714286) -- cycle;
\fill[teal!58!white] (0.8571429,1.714286) -- (0.8571429,1.5) -- (1.142857,1.714286) -- cycle;
\fill[teal!63!white] (0.8571429,1.5) -- (1.142857,1.5) -- (1.142857,1.714286) -- cycle;
\fill[teal!63!white] (1.714286,0.8571429) -- (1.714286,0.6428571) -- (2,0.6428571) -- cycle;
\fill[teal!57!white] (2,0.8571429) -- (1.714286,0.8571429) -- (2,0.6428571) -- cycle;
\fill[teal!53!white] (1.714286,0.8571429) -- (2,0.8571429) -- (2,1.071429) -- cycle;
\fill[teal!58!white] (1.714286,1.071429) -- (1.714286,0.8571429) -- (2,1.071429) -- cycle;
\fill[teal!75!white] (1.714286,0.8571429) -- (1.428571,0.6428571) -- (1.714286,0.6428571) -- cycle;
\fill[teal!85!white] (1.428571,0.6428571) -- (1.714286,0.8571429) -- (1.428571,0.8571429) -- cycle;
\fill[teal!86!white] (1.428571,0.8571429) -- (1.714286,0.8571429) -- (1.428571,1.071429) -- cycle;
\fill[teal!76!white] (1.714286,0.8571429) -- (1.714286,1.071429) -- (1.428571,1.071429) -- cycle;
\fill[teal!9!white] (2.571429,0) -- (2.571429,0.2142857) -- (2.285714,0) -- cycle;
\fill[teal!11!white] (2.571429,0.2142857) -- (2.285714,0.2142857) -- (2.285714,0) -- cycle;
\fill[teal!18!white] (2.285714,0.2142857) -- (2.571429,0.2142857) -- (2.285714,0.4285714) -- cycle;
\fill[teal!19!white] (2.571429,0.2142857) -- (2.571429,0.4285714) -- (2.285714,0.4285714) -- cycle;
\fill[teal!10!white] (2.857143,0.2142857) -- (2.571429,0.2142857) -- (2.857143,0) -- cycle;
\fill[teal!9!white] (2.571429,0.2142857) -- (2.571429,0) -- (2.857143,0) -- cycle;
\fill[teal!13!white] (2.857143,0.4285714) -- (2.571429,0.2142857) -- (2.857143,0.2142857) -- cycle;
\fill[teal!15!white] (2.571429,0.4285714) -- (2.571429,0.2142857) -- (2.857143,0.4285714) -- cycle;
\fill[teal!89!white] (1.142857,0.6428571) -- (0.8571429,0.8571429) -- (0.8571429,0.6428571) -- cycle;
\fill[teal!94!white] (0.8571429,0.8571429) -- (1.142857,0.6428571) -- (1.142857,0.8571429) -- cycle;
\fill[teal!75!white] (0.5714286,0.6428571) -- (0.8571429,0.8571429) -- (0.5714286,0.8571429) -- cycle;
\fill[teal!78!white] (0.8571429,0.8571429) -- (0.5714286,0.6428571) -- (0.8571429,0.6428571) -- cycle;
\fill[teal!87!white] (1.142857,1.5) -- (0.8571429,1.285714) -- (1.142857,1.285714) -- cycle;
\fill[teal!82!white] (0.8571429,1.5) -- (0.8571429,1.285714) -- (1.142857,1.5) -- cycle;
\fill[teal!76!white] (0.5714286,1.285714) -- (0.8571429,1.285714) -- (0.5714286,1.5) -- cycle;
\fill[teal!78!white] (0.8571429,1.285714) -- (0.8571429,1.5) -- (0.5714286,1.5) -- cycle;
\fill[teal!94!white] (1.142857,1.071429) -- (0.8571429,1.071429) -- (1.142857,0.8571429) -- cycle;
\fill[teal!94!white] (0.8571429,1.071429) -- (0.8571429,0.8571429) -- (1.142857,0.8571429) -- cycle;
\fill[teal!94!white] (0.8571429,1.071429) -- (1.142857,1.071429) -- (1.142857,1.285714) -- cycle;
\fill[teal!94!white] (0.8571429,1.285714) -- (0.8571429,1.071429) -- (1.142857,1.285714) -- cycle;
\fill[teal!83!white] (0.8571429,1.071429) -- (0.5714286,1.071429) -- (0.5714286,0.8571429) -- cycle;
\fill[teal!91!white] (0.8571429,0.8571429) -- (0.8571429,1.071429) -- (0.5714286,0.8571429) -- cycle;
\fill[teal!83!white] (0.5714286,1.071429) -- (0.8571429,1.071429) -- (0.5714286,1.285714) -- cycle;
\fill[teal!90!white] (0.8571429,1.071429) -- (0.8571429,1.285714) -- (0.5714286,1.285714) -- cycle;
\draw[teal,thick] (0,0) -- (4,0) -- (0,3) -- cycle;
\draw[magenta!70,thin] (1.075314,0.999031) -- (3.6,2.35) -- (4.1,2.35);\filldraw[fill=magenta,draw=white,line width=.45pt] (1.075314,0.999031) circle(1.5pt) node[above right,font=\scriptsize,text=magenta,fill=white,inner sep=.6pt] {};
\node[anchor=west,text=magenta,font=\small] at(4.12,2.35){Pgrav};\draw[gold!70,thin] (1,1) -- (3.6,1.8) -- (4.1,1.8);\filldraw[fill=gold,draw=white,line width=.45pt] (1,1) circle(1.5pt) node[above right,font=\scriptsize,text=gold,fill=white,inner sep=.6pt] {};
\node[anchor=west,text=gold,font=\small] at(4.12,1.8){$I$};\draw[navy!70,thin] (1.333333,1) -- (3.6,1.25) -- (4.1,1.25);\filldraw[fill=navy,draw=white,line width=.45pt] (1.333333,1) circle(1.5pt) node[above right,font=\scriptsize,text=navy,fill=white,inner sep=.6pt] {};
\node[anchor=west,text=navy,font=\small] at(4.12,1.25){$G$};\end{tikzpicture}
\caption[Clamped-plate deflection]{Clamped-plate deflection. The maximum is a numerical candidate; global uniqueness is not assumed. Fields are archived numerical illustrations; teal saturation increases with the displayed value.}\label{fig:field-Pgrav}
\end{figure}

\section{Heating can cause buckling}\label{sec:slide41}

A heated metal lamina wants to expand. Holding its boundary can generate in-plane compression. A linearized buckling model balances bending stiffness against compressive loading, producing an eigenvalue problem. Below threshold, a perfect initially flat plate may remain flat except for gravity-induced bias. Above threshold, a bifurcation can give several symmetry-related shapes. Imperfections, gravity, anisotropy and nonlinear stretching can select a branch. Thus "uniformly heated sheet" alone does not define one unique lowest point or one density extremum.



Heating introduces preferred strain
$\varepsilon_{\rm th}=\alpha_{\rm th}\Delta T\,\mathbf1$.
Fixed boundaries prevent expansion and generate compression.
Below a buckling threshold a linear flat model may stay flat.
Above it, bending and nonlinear stretching couple. A representative
Foppl--von Karman system is
\[
D_{\rm b}\Delta^2w-[\chi,w]=q,\qquad
\Delta^2\chi=-\tfrac E2[w,w]+\hbox{thermal terms},
\]
where $[f,g]=f_{xx}g_{yy}+f_{yy}g_{xx}-2f_{xy}g_{xy}$.
Restraints, load, thickness and imperfections choose a branch.
An up/down symmetric unloaded model has two deflections rather than
an intrinsically lowest point.

Density requires a material model too: at fixed local mass,
areal density is $\rho_A=\rho_{A0}/\det F_\parallel$, whereas volume
density uses the full three-dimensional Jacobian. Neither extremal
density nor a buckled maximum is an unqualified center without
parameters and a branch rule.




\paragraph{Demonstration and investigation.} Reveal a flat state and two opposite buckling branches. Explain that the animation is a branch schematic, not a solved temperature-dependent material model.


\section{A restricted sheet selects the incenter}\label{sec:slide42}

This exact construction uses a restricted geometry. Join a downward apex to the three boundary sides with three planar triangular faces. Let its horizontal projection have signed side distances d\_i and depth h. Each face area is one half its base length times \ensuremath{\sqrt{\vphantom{x}}}(h\ensuremath{^2}+d\_i\ensuremath{^2}). The length-weighted average of the signed distances is the inradius r because \ensuremath{\sum}a\_i d\_i=2K. Convexity of \ensuremath{\sqrt{\vphantom{x}}}(h\ensuremath{^2}+d\ensuremath{^2}) and Jensen's inequality give the displayed bound. Equality requires all side distances equal, so the projection is I.


\[S_{\rm lat}=\frac12\sum_i\ell_i\sqrt{h^2+d_i^2}\ge s\sqrt{h^2+r^2}.\]



An interior apex projection satisfies
$\sum_i\ell_i d_i=2\mathcal A=2sr$.
Each triangular face has area
$\ell_i\sqrt{h^2+d_i^2}/2$.
Jensen's inequality for the strictly convex function
$d\mapsto\sqrt{h^2+d^2}$ proves the stated lower bound.
Equality requires all $d_i=r$, hence $p=I$.
At fixed lateral area this choice admits the greatest depth.
The conclusion is exact in the three-face class.

An unrestricted sheet with fixed area can form increasingly narrow,
deep spikes; a deepest minimizing shape need not exist.
For a shallow smooth membrane,
$\text{excess area}\simeq\frac12\int|\nabla w|^2$.
Balancing this against work $\int w$ leads instead to torsion.
Thus related-looking sheet experiments select $I$ and F1.



\begin{figure}[htbp]\centering
\begin{tikzpicture}[x=1.45cm,y=1.45cm]
\draw[teal,thick] (0,0) -- (4,0) -- (0,3) -- cycle;
\draw[magenta,thick](0,0)--(1,1)--(4,0);\draw[magenta,thick](0,3)--(1,1);\filldraw[fill=magenta,draw=white,line width=.45pt] (1,1) circle(1.5pt) node[above right,font=\scriptsize,text=magenta,fill=white,inner sep=.6pt] {$I$};
\end{tikzpicture}
\caption[Plan view of the three-face sheet: the deepest apex projects to the incenter at fixed lateral area]{Plan view of the three-face sheet: the deepest apex projects to the incenter at fixed lateral area.}\label{fig:sheet}
\end{figure}

\paragraph{Demonstration and investigation.} Show the three faces and the projection to I. Explicitly state that an unrestricted fixed-area sheet can make arbitrarily thin deep spikes. The three-face constraint is essential.


\section{Finite elements as local physics}\label{sec:slide43}

Divide the triangle into small elements. Represent u as a weighted sum of basis functions attached to nodes. The weak equation says that the residual vanishes against every basis function. This gives a stiffness matrix whose entries integrate overlapping basis gradients. Boundary nodes enforce zero displacement. A field solution is then a linear solve. Computing an extremum requires interpolation within an element, while a mean requires coordinate-weighted quadrature. Mesh refinement, quadrature refinement and maximum-location refinement are different checks.


\[K_{ij}=\int_T\nabla\phi_i\cdot\nabla\phi_j\dd A,\quad M_{ij}=\int_T\phi_i\phi_j\dd A,\quad Ku=f.\]



On a mesh element of area $a_e$, affine nodal basis gradients are
constant. Its stiffness entries are
$K_{ij}^{(e)}=a_e\nabla\phi_i\cdot\nabla\phi_j$.
For unit load,
\[
M^{(e)}=\frac{a_e}{12}
\begin{pmatrix}2&1&1\\1&2&1\\1&1&2\end{pmatrix},
\qquad f_i^{(e)}=a_e/3.
\]
Assemble contributions and eliminate boundary degrees of freedom.
$K$ is then positive definite. A piecewise affine field attains its
maximum at a node; this only locates a mesh maximum.
Refine, reconstruct locally or use quadratic elements.
Integrate $p\,u$ with appropriate polynomial quadrature rather than
averaging nodes. Means, eigenvalues, maxima and shape derivatives
have different error behavior.




\paragraph{Demonstration and investigation.} Reveal a coarse triangular mesh, one local basis function, then a refined mesh. Ask why a coarse nodal maximum can be misplaced.


\section{Fingerprints of the torsion center}\label{sec:slide44}

An isosceles triangle has vertices (\ensuremath{-}1,0),(1,0),(0,h). Its center stays on the vertical axis, so one function y(h) describes its position. In a right-triangle family two coordinate functions are needed. Their derivatives reveal sensitivity to changing the height. The displayed plots compare F1 with familiar centers and shortlisted ETC neighbors. A neighbor rank uses RMS position error divided by triangle diameter on the declared 100-shape grid. It is a sample-dependent numerical comparison, not a proof of an identity.


The course appendix shows the independent right-triangle and derivative fingerprints. The live project survey visits X1 through X72807, with 70853 finite supported candidates and 1954 exclusions. A close location match can still have a very different response Jacobian.


\begin{figure}[htbp]\centering
\begin{tikzpicture}\begin{axis}[width=.95\linewidth,height=.56\linewidth,grid=major,grid style={black!10},xlabel={family parameter $x$},ylabel={$c_y$},tick label style={font=\scriptsize},label style={font=\small},legend style={font=\scriptsize,at={(0.5,-0.24)},anchor=north,legend columns=3,draw=none},scaled ticks=false]
\addplot[navy,thick,no marks] coordinates {(0.2,0.09533565) (0.228706,0.1083233) (0.2615321,0.1229526) (0.2990698,0.1393868) (0.3419952,0.1577908) (0.3910817,0.1783309) (0.4472136,0.2009444) (0.5114021,0.2259574) (0.5848035,0.2536648) (0.6687403,0.2841931) (0.7647245,0.3173771) (0.8744853,0.3530622) (1,0.3915748) (1.14353,0.4331087) (1.30766,0.477026) (1.495349,0.523394) (1.709976,0.5725871) (1.732051,0.5773503) (1.955409,0.6239426) (2.236068,0.6777637) (2.55701,0.7336342) (2.924018,0.7908036) (3.343702,0.851589) (3.823622,0.9138039) (4.372426,0.9772452) (5,1.041219)};
\addlegendentry{F1}
\addplot[teal,thick,no marks] coordinates {(0.2,0.09901951) (0.228706,0.1128955) (0.2615321,0.1286033) (0.2990698,0.1463328) (0.3419952,0.1662702) (0.3910817,0.1885865) (0.4472136,0.2134218) (0.5114021,0.2408662) (0.5848035,0.2709374) (0.6687403,0.3035587) (0.7647245,0.33854) (0.8744853,0.3755687) (1,0.4142136) (1.14353,0.4539441) (1.30766,0.4941653) (1.495349,0.5342616) (1.709976,0.5736416) (1.732051,0.5773503) (1.955409,0.6117774) (2.236068,0.6482315) (2.55701,0.682671) (2.924018,0.7148684) (3.343702,0.744694) (3.823622,0.7721018) (4.372426,0.7971139) (5,0.8198039)};
\addlegendentry{X1}
\addplot[magenta,thick,no marks] coordinates {(0.2,0.06666667) (0.228706,0.07623532) (0.2615321,0.08717737) (0.2990698,0.09968992) (0.3419952,0.1139984) (0.3910817,0.1303606) (0.4472136,0.1490712) (0.5114021,0.1704674) (0.5848035,0.1949345) (0.6687403,0.2229134) (0.7647245,0.2549082) (0.8744853,0.2914951) (1,0.3333333) (1.14353,0.3811766) (1.30766,0.4358868) (1.495349,0.4984496) (1.709976,0.569992) (1.732051,0.5773503) (1.955409,0.6518028) (2.236068,0.745356) (2.55701,0.8523368) (2.924018,0.9746726) (3.343702,1.114567) (3.823622,1.274541) (4.372426,1.457475) (5,1.666667)};
\addlegendentry{X2}
\addplot[gold,thick,no marks] coordinates {(0.2,0.05049024) (0.228706,0.05790523) (0.2615321,0.06646438) (0.2990698,0.07636846) (0.3419952,0.08786248) (0.3910817,0.1012476) (0.4472136,0.1168959) (0.5114021,0.135268) (0.5848035,0.1569331) (0.6687403,0.1825908) (0.7647245,0.2130923) (0.8744853,0.2494583) (1,0.2928932) (1.14353,0.3447929) (1.30766,0.4067476) (1.495349,0.4805436) (1.709976,0.5681672) (1.732051,0.5773503) (1.955409,0.6718156) (2.236068,0.7939182) (2.55701,0.9371697) (2.924018,1.104575) (3.343702,1.299504) (3.823622,1.52576) (4.372426,1.787656) (5,2.090098)};
\addlegendentry{X10}
\addplot[blue,thick,no marks] coordinates {(0.2,0.09487137) (0.228706,0.1076444) (0.2615321,0.1219975) (0.2990698,0.1380855) (0.3419952,0.1560679) (0.3910817,0.1761062) (0.4472136,0.1983612) (0.5114021,0.2229905) (0.5848035,0.2501462) (0.6687403,0.2799719) (0.7647245,0.3126008) (0.8744853,0.3481527) (1,0.3867295) (1.14353,0.4284097) (1.30766,0.4732392) (1.495349,0.5212214) (1.709976,0.5723058) (1.732051,0.5773503) (1.955409,0.6263771) (2.236068,0.6832465) (2.55701,0.7426471) (2.924018,0.8042328) (3.343702,0.867584) (3.823622,0.932217) (4.372426,0.9975986) (5,1.063164)};
\addlegendentry{X5393}
\addplot[purple,thick,no marks] coordinates {(0.2,0.09562539) (0.228706,0.1085897) (0.2615321,0.1231749) (0.2990698,0.1395411) (0.3419952,0.1578516) (0.3910817,0.1782691) (0.4472136,0.2009517) (0.5114021,0.2260473) (0.5848035,0.2536881) (0.6687403,0.2839841) (0.7647245,0.3170172) (0.8744853,0.3528347) (1,0.3914451) (1.14353,0.4328148) (1.30766,0.4768685) (1.495349,0.5234923) (1.709976,0.5725388) (1.732051,0.5773503) (1.955409,0.6238338) (2.236068,0.6771805) (2.55701,0.7323627) (2.924018,0.7891444) (3.343702,0.8472672) (3.823622,0.9064475) (4.372426,0.9663736) (5,1.026705)};
\addlegendentry{X71701}
\end{axis}\end{tikzpicture}
\caption[Isosceles fingerprint: the physical selector, two retained ETC location neighbors and the three classical controls]{Isosceles fingerprint: the physical selector, two retained ETC location neighbors and the three classical controls. Similarity of sampled curves is not identity.}\label{fig:fp-F1}
\end{figure}
\begin{figure}[htbp]\centering
\begin{tikzpicture}\begin{axis}[width=.95\linewidth,height=.56\linewidth,grid=major,grid style={black!10},xlabel={apex height $x$},ylabel={$dc_y/dx$},tick label style={font=\scriptsize},label style={font=\small},legend style={font=\scriptsize,at={(0.5,-0.24)},anchor=north,legend columns=3,draw=none},scaled ticks=false]
\addplot[navy,thick,no marks] coordinates {(0.2,0.4555921) (0.228706,0.4464745) (0.2615321,0.4390451) (0.2990698,0.4304627) (0.3419952,0.4231454) (0.3910817,0.4121286) (0.4472136,0.3977715) (0.5114021,0.3837453) (0.5848035,0.3711014) (0.6687403,0.3527772) (0.7647245,0.3374147) (0.8744853,0.3170172) (1,0.2996052) (1.14353,0.2779459) (1.30766,0.2593595) (1.495349,0.2372703) (1.709976,0.2181639) (1.732051,0.2171301) (1.955409,0.202461) (2.236068,0.1819617) (2.55701,0.1638394) (2.924018,0.152571) (3.343702,0.1346967) (3.823622,0.1186128) (4.372426,0.1132283) (5,0.09900385)};
\addlegendentry{F1}
\addplot[teal,thick,no marks] coordinates {(0.2,0.4833837) (0.228706,0.4807868) (0.2615321,0.4752065) (0.2990698,0.4681264) (0.3419952,0.4592189) (0.3910817,0.4481301) (0.4472136,0.434504) (0.5114021,0.4180222) (0.5848035,0.3984569) (0.6687403,0.3757348) (0.7647245,0.3499968) (0.8744853,0.3216382) (1,0.2913098) (1.14353,0.2598701) (1.30766,0.2282923) (1.495349,0.1975472) (1.709976,0.1820376) (1.732051,0.1553818) (1.955409,0.1406326) (2.236068,0.1178415) (2.55701,0.09686305) (2.924018,0.0788403) (3.343702,0.06362058) (3.823622,0.05095623) (4.372426,0.04054999) (5,0.03615511)};
\addlegendentry{X1}
\addplot[magenta,thick,no marks] coordinates {(0.2,0.3333333) (0.228706,0.3333333) (0.2615321,0.3333333) (0.2990698,0.3333333) (0.3419952,0.3333333) (0.3910817,0.3333333) (0.4472136,0.3333333) (0.5114021,0.3333333) (0.5848035,0.3333333) (0.6687403,0.3333333) (0.7647245,0.3333333) (0.8744853,0.3333333) (1,0.3333333) (1.14353,0.3333333) (1.30766,0.3333333) (1.495349,0.3333333) (1.709976,0.3333333) (1.732051,0.3333333) (1.955409,0.3333333) (2.236068,0.3333333) (2.55701,0.3333333) (2.924018,0.3333333) (3.343702,0.3333333) (3.823622,0.3333333) (4.372426,0.3333333) (5,0.3333333)};
\addlegendentry{X2}
\addplot[gold,thick,no marks] coordinates {(0.2,0.2583082) (0.228706,0.2596066) (0.2615321,0.2623967) (0.2990698,0.2659368) (0.3419952,0.2703905) (0.3910817,0.275935) (0.4472136,0.282748) (0.5114021,0.2909889) (0.5848035,0.3007716) (0.6687403,0.3121326) (0.7647245,0.3250016) (0.8744853,0.3391809) (1,0.3543451) (1.14353,0.370065) (1.30766,0.3858539) (1.495349,0.4012264) (1.709976,0.4089812) (1.732051,0.4223091) (1.955409,0.4296837) (2.236068,0.4410793) (2.55701,0.4515685) (2.924018,0.4605799) (3.343702,0.4681897) (3.823622,0.4745219) (4.372426,0.479725) (5,0.4819224)};
\addlegendentry{X10}
\addplot[blue,thick,no marks] coordinates {(0.2,0.4449605) (0.228706,0.4408445) (0.2615321,0.4326248) (0.2990698,0.4234299) (0.3419952,0.4132146) (0.3910817,0.4019566) (0.4472136,0.3896625) (0.5114021,0.3763721) (0.5848035,0.3621582) (0.6687403,0.3471224) (0.7647245,0.3313849) (0.8744853,0.3150721) (1,0.298304) (1.14353,0.2811854) (1.30766,0.2638053) (1.495349,0.2462411) (1.709976,0.2371288) (1.732051,0.2203101) (1.955409,0.2101045) (2.236068,0.1932673) (2.55701,0.175865) (2.924018,0.1588133) (3.343702,0.1422671) (3.823622,0.1263842) (4.372426,0.1113141) (5,0.1044752)};
\addlegendentry{X5393}
\addplot[purple,thick,no marks] coordinates {(0.2,0.4516244) (0.228706,0.4477265) (0.2615321,0.4398773) (0.2990698,0.4309633) (0.3419952,0.4209016) (0.3910817,0.4096255) (0.4472136,0.3970917) (0.5114021,0.3832866) (0.5848035,0.3682308) (0.6687403,0.3519826) (0.7647245,0.3346403) (0.8744853,0.3163436) (1,0.2972744) (1.14353,0.2776549) (1.30766,0.2577391) (1.495349,0.2377994) (1.709976,0.2275351) (1.732051,0.2089982) (1.955409,0.1980691) (2.236068,0.1803999) (2.55701,0.1627501) (2.924018,0.1460605) (3.343702,0.130394) (3.823622,0.1157806) (4.372426,0.1022273) (5,0.09613511)};
\addlegendentry{X71701}
\end{axis}\end{tikzpicture}
\caption[Isosceles derivative comparison]{Isosceles derivative comparison. Neighbor derivatives use sampled differences, and do not replace the full vertex-response matrices.}\label{fig:deriv-F1}
\end{figure}

\section{What the elastic models establish}\label{sec:slide45}

The restricted tetrahedral sheet gives I exactly. The linear membrane, Newtonian duct and Brownian-lifetime maximum give the same F1. Their weighted mean is Abs(0), also equal to a dissipation mean. A clamped plate gives Pgrav, with separate analytical and numerical challenges. The nonlinear p-torsion family gives power-law membrane and fluid analogues, but its parameter and response still need careful study. Unique equilibria and positive energy Hessians do not establish geometric attractivity.



\paragraph{Demonstration and investigation.} Run a final comparison of the two fields and their mean markers. Have students identify which equalities were derived and which locations came from computation.


\section*{Questions and exercises}\addcontentsline{toc}{section}{Questions and exercises}

\begin{enumerate}

\item\label{ex:32} What physical assumption permits a constant coefficient \ensuremath{\tau}?

\item\label{ex:33} Would doubling the uniform load move the linear membrane center?

\item\label{ex:34} Where did the boundary condition enter the derivation?

\item\label{ex:35} Which familiar center maximizes the minimum side distance?

\item\label{ex:36} Must equal centroids imply equal densities?

\item\label{ex:37} Which center would a uniformly chosen position in the duct select?

\item\label{ex:38} Does longest mean lifetime imply equal exit probabilities through the three sides?

\item\label{ex:39} Why is specifying "fixed boundary" insufficient for a plate?

\item\label{ex:40} Why should we not import every membrane theorem into the plate problem?

\item\label{ex:41} What additional data would you ask for before predicting an extremal-density point?

\item\label{ex:42} Does total sheet area alone force the incenter?

\item\label{ex:43} Which center calculation is often smoother numerically, a mean or a maximum?

\item\label{ex:44} Can a center share an isosceles curve with another center and differ on scalene triangles?

\item\label{ex:45} What physical change moves F1 to Pgrav?

\end{enumerate}

\chapter{Quantum confinement, Brownian motion and billiards}\label{ch:4}

Quantum confinement, absorbing Brownian motion and heat share the Dirichlet Laplacian. Their selected means use different weights. Classical billiards adds a different challenge: a single deterministic trajectory may explore its triangle slowly or fail to do so in an exceptional direction. We will distinguish ensemble expectations, stationary distributions and individual time averages.


\section{A particle in a triangular infinite well}\label{sec:slide47}

The stationary Schr\"odinger equation inside a zero-potential infinite well is \ensuremath{-}\ensuremath{\hbar}\ensuremath{^2}\ensuremath{\Delta}\ensuremath{\psi}/(2m)=E\ensuremath{\psi}. Dividing by the common physical factor gives the Dirichlet eigenproblem shown. The energy is E\_j=\ensuremath{\hbar}\ensuremath{^2}\ensuremath{\lambda}\_j/(2m). The shape of the eigenfunction depends on the triangle, while the mass and \ensuremath{\hbar} change the energy scale. The lowest eigenfunction can be chosen positive and normalized. The boundary condition models infinite confinement, not a finite barrier or a leaky cavity.


\[-\Delta\psi_j=\lambda_j\psi_j,\quad\psi_j|_{\partial T}=0,\quad E_j=\hbar^2\lambda_j/(2m).\]



A stationary state has
$\Psi(p,t)=\psi_j(p)e^{-iE_jt/\hbar}$.
The infinite exterior potential gives zero boundary trace.
Substitution in
$i\hbar\partial_t\Psi=-(\hbar^2/2m)\Delta\Psi$
gives the spatial eigenproblem. Normalize
$\int_T|\psi_j|^2=1$ before interpreting probability.
Numbering starts at $j=0$. Mass and size change the energy scale,
but not the normalized location for a fixed shape.

The ground eigenvalue is simple and its eigenfunction is positive.
On convex domains it is log-concave \cite{brascamp}. Analyticity and convex level
sets imply a unique maximum: a segment of maxima would extend
analytically along a line to the zero boundary, a contradiction.
Q0m is this maximum; Q0 is a weighted integral. They agree by symmetry
in the equilateral triangle but generally differ.




\section{The ground state minimizes gradient energy}\label{sec:slide48}

A wavefunction that changes rapidly has high kinetic energy. The zero boundary condition forces it to fall to zero near each side. The ground state makes the smallest possible ratio of gradient energy to probability mass. This Rayleigh quotient is a variational characterization, so the equation can be understood before solving a partial differential equation. Higher states minimize under orthogonality constraints and acquire nodes. The field is not generally a simple polynomial on an arbitrary triangle.


\[\lambda_0=\min_{\psi\in H_0^1(T),\,\int_T|\psi|^2=1}\int_T|\nabla\psi|^2\dd A.\]


Vary the Rayleigh quotient under unit $L^2$ normalization.
Its Lagrange multiplier gives $-\Delta\psi=\lambda\psi$.
In a finite-element basis the quotient is $v^TKv/(v^TMv)$,
yielding the generalized eigenproblem.



\paragraph{Demonstration and investigation.} Reveal a broad single peak and then a rapidly oscillating candidate. Ask which carries larger kinetic energy under the same normalization.


\section{$\mathrm{Q0m}$ and $\mathrm{Q0}$}\label{sec:slide49}

The positive ground amplitude has unique interior maximum Q0m.
Its squared amplitude is probability density, with the same maximum
but a generally different mean Q0. In the reference triangle the two
locations are close: Q0m is approximately $(1.095451,0.997508)$ and
Q0 is $(1.126561,1.001531)$. The later coordinate table and density
plots give the same normalization. Q0 has stable numerical negative
response witnesses. Q0m has a small refinement-sensitive negative
candidate; no certified response result is claimed.


\[{\rm Q0m}=\argmax_T\psi_0,\quad {\rm Q0}=\int_Tp|\psi_0|^2\dd A.\]


Barycentric weights of the expectation are
$\int_T\phi_i|\psi_0|^2\dd A>0$. This proves strict interiority,
without proving a short center formula or vertex attractivity.


\begin{figure}[htbp]\centering
\begin{tikzpicture}[x=1.7cm,y=1.7cm]
\fill[teal!8!white] (3.142857,0.6428571) -- (2.857143,0.8571429) -- (2.857143,0.6428571) -- cycle;
\fill[teal!9!white] (1.714286,1.714286) -- (1.428571,1.928571) -- (1.428571,1.714286) -- cycle;
\fill[teal!9!white] (1.714286,1.714286) -- (1.714286,1.5) -- (2,1.5) -- cycle;
\fill[teal!8!white] (0.5714286,2.571429) -- (0.2857143,2.785714) -- (0.2857143,2.571429) -- cycle;
\fill[teal!8!white] (2.571429,1.071429) -- (2.285714,1.285714) -- (2.285714,1.071429) -- cycle;
\fill[teal!8!white] (3.714286,0.2142857) -- (3.428571,0.4285714) -- (3.428571,0.2142857) -- cycle;
\fill[teal!8!white] (3.714286,0.2142857) -- (3.714286,0) -- (4,0) -- cycle;
\fill[teal!8!white] (0.8571429,2.142857) -- (1.142857,2.142857) -- (0.8571429,2.357143) -- cycle;
\fill[teal!9!white] (1.142857,2.142857) -- (1.142857,1.928571) -- (1.428571,1.928571) -- cycle;
\fill[teal!8!white] (0,2.785714) -- (0.2857143,2.785714) -- (0,3) -- cycle;
\fill[teal!8!white] (0.5714286,2.571429) -- (0.5714286,2.357143) -- (0.8571429,2.357143) -- cycle;
\fill[teal!9!white] (2,1.285714) -- (2.285714,1.285714) -- (2,1.5) -- cycle;
\fill[teal!8!white] (2.571429,1.071429) -- (2.571429,0.8571429) -- (2.857143,0.8571429) -- cycle;
\fill[teal!8!white] (3.142857,0.4285714) -- (3.428571,0.4285714) -- (3.142857,0.6428571) -- cycle;
\fill[teal!8!white] (0.2857143,0) -- (0,0.2142857) -- (0,0) -- cycle;
\fill[teal!8!white] (0,0.2142857) -- (0.2857143,0) -- (0.2857143,0.2142857) -- cycle;
\fill[teal!8!white] (0.5714286,0.2142857) -- (0.2857143,0) -- (0.5714286,0) -- cycle;
\fill[teal!9!white] (0.2857143,0) -- (0.5714286,0.2142857) -- (0.2857143,0.2142857) -- cycle;
\fill[teal!12!white] (1.714286,1.5) -- (1.714286,1.714286) -- (1.428571,1.714286) -- cycle;
\fill[teal!22!white] (1.428571,1.5) -- (1.714286,1.5) -- (1.428571,1.714286) -- cycle;
\fill[teal!9!white] (0,1.928571) -- (0,1.714286) -- (0.2857143,1.928571) -- cycle;
\fill[teal!11!white] (0,1.714286) -- (0.2857143,1.714286) -- (0.2857143,1.928571) -- cycle;
\fill[teal!8!white] (0.2857143,2.357143) -- (0,2.571429) -- (0,2.357143) -- cycle;
\fill[teal!8!white] (0,2.571429) -- (0.2857143,2.357143) -- (0.2857143,2.571429) -- cycle;
\fill[teal!8!white] (0,2.142857) -- (0.2857143,2.357143) -- (0,2.357143) -- cycle;
\fill[teal!9!white] (0.2857143,2.357143) -- (0,2.142857) -- (0.2857143,2.142857) -- cycle;
\fill[teal!8!white] (0,2.142857) -- (0,1.928571) -- (0.2857143,1.928571) -- cycle;
\fill[teal!9!white] (0.2857143,2.142857) -- (0,2.142857) -- (0.2857143,1.928571) -- cycle;
\fill[teal!17!white] (0.2857143,1.714286) -- (0.5714286,1.928571) -- (0.2857143,1.928571) -- cycle;
\fill[teal!24!white] (0.5714286,1.714286) -- (0.5714286,1.928571) -- (0.2857143,1.714286) -- cycle;
\fill[teal!9!white] (0.2857143,1.5) -- (0,1.714286) -- (0,1.5) -- cycle;
\fill[teal!13!white] (0,1.714286) -- (0.2857143,1.5) -- (0.2857143,1.714286) -- cycle;
\fill[teal!24!white] (0.2857143,1.5) -- (0.5714286,1.714286) -- (0.2857143,1.714286) -- cycle;
\fill[teal!35!white] (0.5714286,1.714286) -- (0.2857143,1.5) -- (0.5714286,1.5) -- cycle;
\fill[teal!10!white] (0,1.285714) -- (0.2857143,1.5) -- (0,1.5) -- cycle;
\fill[teal!15!white] (0.2857143,1.5) -- (0,1.285714) -- (0.2857143,1.285714) -- cycle;
\fill[teal!10!white] (0,1.285714) -- (0,1.071429) -- (0.2857143,1.071429) -- cycle;
\fill[teal!16!white] (0.2857143,1.285714) -- (0,1.285714) -- (0.2857143,1.071429) -- cycle;
\fill[teal!8!white] (3.428571,0) -- (3.714286,0) -- (3.428571,0.2142857) -- cycle;
\fill[teal!8!white] (3.714286,0) -- (3.714286,0.2142857) -- (3.428571,0.2142857) -- cycle;
\fill[teal!9!white] (1.428571,0) -- (1.714286,0) -- (1.428571,0.2142857) -- cycle;
\fill[teal!12!white] (1.714286,0) -- (1.714286,0.2142857) -- (1.428571,0.2142857) -- cycle;
\fill[teal!9!white] (0.8571429,0) -- (0.5714286,0.2142857) -- (0.5714286,0) -- cycle;
\fill[teal!12!white] (0.5714286,0.2142857) -- (0.8571429,0) -- (0.8571429,0.2142857) -- cycle;
\fill[teal!9!white] (1.142857,0.2142857) -- (1.142857,0) -- (1.428571,0) -- cycle;
\fill[teal!13!white] (1.142857,0.2142857) -- (1.428571,0) -- (1.428571,0.2142857) -- cycle;
\fill[teal!9!white] (1.142857,0.2142857) -- (0.8571429,0) -- (1.142857,0) -- cycle;
\fill[teal!13!white] (0.8571429,0) -- (1.142857,0.2142857) -- (0.8571429,0.2142857) -- cycle;
\fill[teal!20!white] (1.142857,1.928571) -- (1.142857,1.714286) -- (1.428571,1.714286) -- cycle;
\fill[teal!11!white] (1.428571,1.928571) -- (1.142857,1.928571) -- (1.428571,1.714286) -- cycle;
\fill[teal!10!white] (0.8571429,2.142857) -- (1.142857,1.928571) -- (1.142857,2.142857) -- cycle;
\fill[teal!15!white] (1.142857,1.928571) -- (0.8571429,2.142857) -- (0.8571429,1.928571) -- cycle;
\fill[teal!8!white] (0,2.785714) -- (0,2.571429) -- (0.2857143,2.571429) -- cycle;
\fill[teal!8!white] (0.2857143,2.785714) -- (0,2.785714) -- (0.2857143,2.571429) -- cycle;
\fill[teal!29!white] (1.142857,1.928571) -- (0.8571429,1.714286) -- (1.142857,1.714286) -- cycle;
\fill[teal!25!white] (0.8571429,1.714286) -- (1.142857,1.928571) -- (0.8571429,1.928571) -- cycle;
\fill[teal!31!white] (0.8571429,1.714286) -- (0.5714286,1.928571) -- (0.5714286,1.714286) -- cycle;
\fill[teal!27!white] (0.5714286,1.928571) -- (0.8571429,1.714286) -- (0.8571429,1.928571) -- cycle;
\fill[teal!8!white] (0.2857143,2.357143) -- (0.5714286,2.357143) -- (0.2857143,2.571429) -- cycle;
\fill[teal!8!white] (0.5714286,2.357143) -- (0.5714286,2.571429) -- (0.2857143,2.571429) -- cycle;
\fill[teal!9!white] (0,0.6428571) -- (0,0.4285714) -- (0.2857143,0.6428571) -- cycle;
\fill[teal!12!white] (0,0.4285714) -- (0.2857143,0.4285714) -- (0.2857143,0.6428571) -- cycle;
\fill[teal!8!white] (0,0.4285714) -- (0,0.2142857) -- (0.2857143,0.2142857) -- cycle;
\fill[teal!10!white] (0.2857143,0.4285714) -- (0,0.4285714) -- (0.2857143,0.2142857) -- cycle;
\fill[teal!10!white] (2.285714,1.285714) -- (2,1.285714) -- (2.285714,1.071429) -- cycle;
\fill[teal!16!white] (2,1.285714) -- (2,1.071429) -- (2.285714,1.071429) -- cycle;
\fill[teal!22!white] (2,1.071429) -- (2,0.8571429) -- (2.285714,1.071429) -- cycle;
\fill[teal!19!white] (2,0.8571429) -- (2.285714,0.8571429) -- (2.285714,1.071429) -- cycle;
\fill[teal!12!white] (2.285714,0.8571429) -- (2.571429,0.8571429) -- (2.285714,1.071429) -- cycle;
\fill[teal!9!white] (2.571429,0.8571429) -- (2.571429,1.071429) -- (2.285714,1.071429) -- cycle;
\fill[teal!31!white] (1.428571,0.4285714) -- (1.714286,0.2142857) -- (1.714286,0.4285714) -- cycle;
\fill[teal!25!white] (1.714286,0.2142857) -- (1.428571,0.4285714) -- (1.428571,0.2142857) -- cycle;
\fill[teal!8!white] (3.142857,0.2142857) -- (3.428571,0) -- (3.428571,0.2142857) -- cycle;
\fill[teal!8!white] (3.142857,0.2142857) -- (3.142857,0) -- (3.428571,0) -- cycle;
\fill[teal!8!white] (3.142857,0) -- (3.142857,0.2142857) -- (2.857143,0) -- cycle;
\fill[teal!8!white] (3.142857,0.2142857) -- (2.857143,0.2142857) -- (2.857143,0) -- cycle;
\fill[teal!11!white] (1.714286,0) -- (2,0.2142857) -- (1.714286,0.2142857) -- cycle;
\fill[teal!9!white] (2,0.2142857) -- (1.714286,0) -- (2,0) -- cycle;
\fill[teal!15!white] (0.5714286,0.2142857) -- (0.5714286,0.4285714) -- (0.2857143,0.2142857) -- cycle;
\fill[teal!15!white] (0.5714286,0.4285714) -- (0.2857143,0.4285714) -- (0.2857143,0.2142857) -- cycle;
\fill[teal!23!white] (0.8571429,0.4285714) -- (0.5714286,0.2142857) -- (0.8571429,0.2142857) -- cycle;
\fill[teal!26!white] (0.8571429,0.4285714) -- (0.5714286,0.4285714) -- (0.5714286,0.2142857) -- cycle;
\fill[teal!37!white] (1.142857,1.5) -- (1.428571,1.5) -- (1.428571,1.714286) -- cycle;
\fill[teal!35!white] (1.142857,1.714286) -- (1.142857,1.5) -- (1.428571,1.714286) -- cycle;
\fill[teal!12!white] (0.5714286,2.142857) -- (0.2857143,2.142857) -- (0.2857143,1.928571) -- cycle;
\fill[teal!15!white] (0.5714286,1.928571) -- (0.5714286,2.142857) -- (0.2857143,1.928571) -- cycle;
\fill[teal!10!white] (0.5714286,2.142857) -- (0.2857143,2.357143) -- (0.2857143,2.142857) -- cycle;
\fill[teal!9!white] (0.5714286,2.142857) -- (0.5714286,2.357143) -- (0.2857143,2.357143) -- cycle;
\fill[teal!14!white] (0.8571429,2.142857) -- (0.5714286,2.142857) -- (0.8571429,1.928571) -- cycle;
\fill[teal!18!white] (0.5714286,2.142857) -- (0.5714286,1.928571) -- (0.8571429,1.928571) -- cycle;
\fill[teal!10!white] (0.5714286,2.142857) -- (0.8571429,2.142857) -- (0.8571429,2.357143) -- cycle;
\fill[teal!9!white] (0.5714286,2.357143) -- (0.5714286,2.142857) -- (0.8571429,2.357143) -- cycle;
\fill[teal!15!white] (1.714286,1.285714) -- (2,1.285714) -- (2,1.5) -- cycle;
\fill[teal!16!white] (1.714286,1.5) -- (1.714286,1.285714) -- (2,1.5) -- cycle;
\fill[teal!94!white] (1.142857,1.071429) -- (1.428571,0.8571429) -- (1.428571,1.071429) -- cycle;
\fill[teal!94!white] (1.428571,0.8571429) -- (1.142857,1.071429) -- (1.142857,0.8571429) -- cycle;
\fill[teal!8!white] (3.142857,0.4285714) -- (3.142857,0.2142857) -- (3.428571,0.2142857) -- cycle;
\fill[teal!8!white] (3.428571,0.4285714) -- (3.142857,0.4285714) -- (3.428571,0.2142857) -- cycle;
\fill[teal!8!white] (2.857143,0.4285714) -- (3.142857,0.4285714) -- (2.857143,0.6428571) -- cycle;
\fill[teal!8!white] (3.142857,0.4285714) -- (3.142857,0.6428571) -- (2.857143,0.6428571) -- cycle;
\fill[teal!8!white] (3.142857,0.2142857) -- (3.142857,0.4285714) -- (2.857143,0.2142857) -- cycle;
\fill[teal!8!white] (3.142857,0.4285714) -- (2.857143,0.4285714) -- (2.857143,0.2142857) -- cycle;
\fill[teal!9!white] (2.571429,0.8571429) -- (2.571429,0.6428571) -- (2.857143,0.8571429) -- cycle;
\fill[teal!8!white] (2.857143,0.8571429) -- (2.571429,0.6428571) -- (2.857143,0.6428571) -- cycle;
\fill[teal!8!white] (2.285714,0) -- (2.285714,0.2142857) -- (2,0) -- cycle;
\fill[teal!10!white] (2.285714,0.2142857) -- (2,0.2142857) -- (2,0) -- cycle;
\fill[teal!14!white] (2.285714,0.2142857) -- (2,0.4285714) -- (2,0.2142857) -- cycle;
\fill[teal!15!white] (2,0.4285714) -- (2.285714,0.2142857) -- (2.285714,0.4285714) -- cycle;
\fill[teal!37!white] (2,0.4285714) -- (1.714286,0.6428571) -- (1.714286,0.4285714) -- cycle;
\fill[teal!36!white] (1.714286,0.6428571) -- (2,0.4285714) -- (2,0.6428571) -- cycle;
\fill[teal!25!white] (1.714286,0.2142857) -- (2,0.4285714) -- (1.714286,0.4285714) -- cycle;
\fill[teal!18!white] (2,0.2142857) -- (2,0.4285714) -- (1.714286,0.2142857) -- cycle;
\fill[teal!10!white] (0.2857143,0.8571429) -- (0,0.8571429) -- (0,0.6428571) -- cycle;
\fill[teal!14!white] (0.2857143,0.8571429) -- (0,0.6428571) -- (0.2857143,0.6428571) -- cycle;
\fill[teal!16!white] (0,1.071429) -- (0.2857143,0.8571429) -- (0.2857143,1.071429) -- cycle;
\fill[teal!10!white] (0,0.8571429) -- (0.2857143,0.8571429) -- (0,1.071429) -- cycle;
\fill[teal!42!white] (0.2857143,1.5) -- (0.5714286,1.285714) -- (0.5714286,1.5) -- cycle;
\fill[teal!33!white] (0.5714286,1.285714) -- (0.2857143,1.5) -- (0.2857143,1.285714) -- cycle;
\fill[teal!30!white] (0.5714286,0.6428571) -- (0.2857143,0.8571429) -- (0.2857143,0.6428571) -- cycle;
\fill[teal!44!white] (0.2857143,0.8571429) -- (0.5714286,0.6428571) -- (0.5714286,0.8571429) -- cycle;
\fill[teal!24!white] (0.2857143,0.4285714) -- (0.5714286,0.6428571) -- (0.2857143,0.6428571) -- cycle;
\fill[teal!27!white] (0.5714286,0.4285714) -- (0.5714286,0.6428571) -- (0.2857143,0.4285714) -- cycle;
\fill[teal!41!white] (0.8571429,0.4285714) -- (0.5714286,0.6428571) -- (0.5714286,0.4285714) -- cycle;
\fill[teal!56!white] (0.5714286,0.6428571) -- (0.8571429,0.4285714) -- (0.8571429,0.6428571) -- cycle;
\fill[teal!28!white] (1.142857,0.4285714) -- (1.142857,0.2142857) -- (1.428571,0.2142857) -- cycle;
\fill[teal!37!white] (1.428571,0.4285714) -- (1.142857,0.4285714) -- (1.428571,0.2142857) -- cycle;
\fill[teal!27!white] (1.142857,0.2142857) -- (1.142857,0.4285714) -- (0.8571429,0.2142857) -- cycle;
\fill[teal!36!white] (1.142857,0.4285714) -- (0.8571429,0.4285714) -- (0.8571429,0.2142857) -- cycle;
\fill[teal!56!white] (0.8571429,0.4285714) -- (1.142857,0.4285714) -- (0.8571429,0.6428571) -- cycle;
\fill[teal!70!white] (1.142857,0.4285714) -- (1.142857,0.6428571) -- (0.8571429,0.6428571) -- cycle;
\fill[teal!38!white] (1.428571,1.285714) -- (1.714286,1.5) -- (1.428571,1.5) -- cycle;
\fill[teal!37!white] (1.428571,1.285714) -- (1.714286,1.285714) -- (1.714286,1.5) -- cycle;
\fill[teal!57!white] (1.142857,1.5) -- (1.428571,1.285714) -- (1.428571,1.5) -- cycle;
\fill[teal!75!white] (1.428571,1.285714) -- (1.142857,1.5) -- (1.142857,1.285714) -- cycle;
\fill[teal!92!white] (1.428571,1.285714) -- (1.142857,1.071429) -- (1.428571,1.071429) -- cycle;
\fill[teal!93!white] (1.142857,1.071429) -- (1.428571,1.285714) -- (1.142857,1.285714) -- cycle;
\fill[teal!52!white] (1.428571,0.6428571) -- (1.428571,0.4285714) -- (1.714286,0.4285714) -- cycle;
\fill[teal!55!white] (1.714286,0.6428571) -- (1.428571,0.6428571) -- (1.714286,0.4285714) -- cycle;
\fill[teal!58!white] (1.428571,0.6428571) -- (1.142857,0.4285714) -- (1.428571,0.4285714) -- cycle;
\fill[teal!71!white] (1.142857,0.4285714) -- (1.428571,0.6428571) -- (1.142857,0.6428571) -- cycle;
\fill[teal!94!white] (1.428571,0.6428571) -- (1.428571,0.8571429) -- (1.142857,0.8571429) -- cycle;
\fill[teal!91!white] (1.142857,0.6428571) -- (1.428571,0.6428571) -- (1.142857,0.8571429) -- cycle;
\fill[teal!29!white] (2,1.285714) -- (1.714286,1.071429) -- (2,1.071429) -- cycle;
\fill[teal!33!white] (1.714286,1.285714) -- (1.714286,1.071429) -- (2,1.285714) -- cycle;
\fill[teal!53!white] (1.428571,1.285714) -- (1.714286,1.071429) -- (1.714286,1.285714) -- cycle;
\fill[teal!72!white] (1.714286,1.071429) -- (1.428571,1.285714) -- (1.428571,1.071429) -- cycle;
\fill[teal!14!white] (2.285714,0.6428571) -- (2.571429,0.8571429) -- (2.285714,0.8571429) -- cycle;
\fill[teal!12!white] (2.285714,0.6428571) -- (2.571429,0.6428571) -- (2.571429,0.8571429) -- cycle;
\fill[teal!22!white] (2,0.8571429) -- (2.285714,0.6428571) -- (2.285714,0.8571429) -- cycle;
\fill[teal!28!white] (2.285714,0.6428571) -- (2,0.8571429) -- (2,0.6428571) -- cycle;
\fill[teal!19!white] (2.285714,0.6428571) -- (2,0.4285714) -- (2.285714,0.4285714) -- cycle;
\fill[teal!25!white] (2,0.4285714) -- (2.285714,0.6428571) -- (2,0.6428571) -- cycle;
\fill[teal!9!white] (2.571429,0.4285714) -- (2.857143,0.4285714) -- (2.857143,0.6428571) -- cycle;
\fill[teal!9!white] (2.571429,0.6428571) -- (2.571429,0.4285714) -- (2.857143,0.6428571) -- cycle;
\fill[teal!12!white] (2.285714,0.6428571) -- (2.571429,0.4285714) -- (2.571429,0.6428571) -- cycle;
\fill[teal!14!white] (2.571429,0.4285714) -- (2.285714,0.6428571) -- (2.285714,0.4285714) -- cycle;
\fill[teal!37!white] (0.2857143,0.8571429) -- (0.5714286,1.071429) -- (0.2857143,1.071429) -- cycle;
\fill[teal!49!white] (0.5714286,1.071429) -- (0.2857143,0.8571429) -- (0.5714286,0.8571429) -- cycle;
\fill[teal!37!white] (0.5714286,1.071429) -- (0.2857143,1.285714) -- (0.2857143,1.071429) -- cycle;
\fill[teal!49!white] (0.5714286,1.071429) -- (0.5714286,1.285714) -- (0.2857143,1.285714) -- cycle;
\fill[teal!49!white] (0.8571429,1.5) -- (0.5714286,1.714286) -- (0.5714286,1.5) -- cycle;
\fill[teal!47!white] (0.8571429,1.5) -- (0.8571429,1.714286) -- (0.5714286,1.714286) -- cycle;
\fill[teal!48!white] (0.8571429,1.714286) -- (0.8571429,1.5) -- (1.142857,1.714286) -- cycle;
\fill[teal!55!white] (0.8571429,1.5) -- (1.142857,1.5) -- (1.142857,1.714286) -- cycle;
\fill[teal!50!white] (1.714286,0.8571429) -- (1.714286,0.6428571) -- (2,0.6428571) -- cycle;
\fill[teal!43!white] (2,0.8571429) -- (1.714286,0.8571429) -- (2,0.6428571) -- cycle;
\fill[teal!41!white] (1.714286,0.8571429) -- (2,0.8571429) -- (2,1.071429) -- cycle;
\fill[teal!48!white] (1.714286,1.071429) -- (1.714286,0.8571429) -- (2,1.071429) -- cycle;
\fill[teal!67!white] (1.714286,0.8571429) -- (1.428571,0.6428571) -- (1.714286,0.6428571) -- cycle;
\fill[teal!80!white] (1.428571,0.6428571) -- (1.714286,0.8571429) -- (1.428571,0.8571429) -- cycle;
\fill[teal!83!white] (1.428571,0.8571429) -- (1.714286,0.8571429) -- (1.428571,1.071429) -- cycle;
\fill[teal!70!white] (1.714286,0.8571429) -- (1.714286,1.071429) -- (1.428571,1.071429) -- cycle;
\fill[teal!8!white] (2.571429,0) -- (2.571429,0.2142857) -- (2.285714,0) -- cycle;
\fill[teal!9!white] (2.571429,0.2142857) -- (2.285714,0.2142857) -- (2.285714,0) -- cycle;
\fill[teal!11!white] (2.285714,0.2142857) -- (2.571429,0.2142857) -- (2.285714,0.4285714) -- cycle;
\fill[teal!10!white] (2.571429,0.2142857) -- (2.571429,0.4285714) -- (2.285714,0.4285714) -- cycle;
\fill[teal!8!white] (2.857143,0.2142857) -- (2.571429,0.2142857) -- (2.857143,0) -- cycle;
\fill[teal!8!white] (2.571429,0.2142857) -- (2.571429,0) -- (2.857143,0) -- cycle;
\fill[teal!8!white] (2.857143,0.4285714) -- (2.571429,0.2142857) -- (2.857143,0.2142857) -- cycle;
\fill[teal!9!white] (2.571429,0.4285714) -- (2.571429,0.2142857) -- (2.857143,0.4285714) -- cycle;
\fill[teal!87!white] (1.142857,0.6428571) -- (0.8571429,0.8571429) -- (0.8571429,0.6428571) -- cycle;
\fill[teal!94!white] (0.8571429,0.8571429) -- (1.142857,0.6428571) -- (1.142857,0.8571429) -- cycle;
\fill[teal!68!white] (0.5714286,0.6428571) -- (0.8571429,0.8571429) -- (0.5714286,0.8571429) -- cycle;
\fill[teal!72!white] (0.8571429,0.8571429) -- (0.5714286,0.6428571) -- (0.8571429,0.6428571) -- cycle;
\fill[teal!84!white] (1.142857,1.5) -- (0.8571429,1.285714) -- (1.142857,1.285714) -- cycle;
\fill[teal!76!white] (0.8571429,1.5) -- (0.8571429,1.285714) -- (1.142857,1.5) -- cycle;
\fill[teal!67!white] (0.5714286,1.285714) -- (0.8571429,1.285714) -- (0.5714286,1.5) -- cycle;
\fill[teal!69!white] (0.8571429,1.285714) -- (0.8571429,1.5) -- (0.5714286,1.5) -- cycle;
\fill[teal!94!white] (1.142857,1.071429) -- (0.8571429,1.071429) -- (1.142857,0.8571429) -- cycle;
\fill[teal!94!white] (0.8571429,1.071429) -- (0.8571429,0.8571429) -- (1.142857,0.8571429) -- cycle;
\fill[teal!94!white] (0.8571429,1.071429) -- (1.142857,1.071429) -- (1.142857,1.285714) -- cycle;
\fill[teal!94!white] (0.8571429,1.285714) -- (0.8571429,1.071429) -- (1.142857,1.285714) -- cycle;
\fill[teal!77!white] (0.8571429,1.071429) -- (0.5714286,1.071429) -- (0.5714286,0.8571429) -- cycle;
\fill[teal!87!white] (0.8571429,0.8571429) -- (0.8571429,1.071429) -- (0.5714286,0.8571429) -- cycle;
\fill[teal!77!white] (0.5714286,1.071429) -- (0.8571429,1.071429) -- (0.5714286,1.285714) -- cycle;
\fill[teal!85!white] (0.8571429,1.071429) -- (0.8571429,1.285714) -- (0.5714286,1.285714) -- cycle;
\draw[teal,thick] (0,0) -- (4,0) -- (0,3) -- cycle;
\draw[magenta!70,thin] (1.095451,0.9975085) -- (3.6,2.35) -- (4.1,2.35);\filldraw[fill=magenta,draw=white,line width=.45pt] (1.095451,0.9975085) circle(1.5pt) node[above right,font=\scriptsize,text=magenta,fill=white,inner sep=.6pt] {};
\node[anchor=west,text=magenta,font=\small] at(4.12,2.35){Q0m};\draw[gold!70,thin] (1,1) -- (3.6,1.8) -- (4.1,1.8);\filldraw[fill=gold,draw=white,line width=.45pt] (1,1) circle(1.5pt) node[above right,font=\scriptsize,text=gold,fill=white,inner sep=.6pt] {};
\node[anchor=west,text=gold,font=\small] at(4.12,1.8){$I$};\draw[navy!70,thin] (1.333333,1) -- (3.6,1.25) -- (4.1,1.25);\filldraw[fill=navy,draw=white,line width=.45pt] (1.333333,1) circle(1.5pt) node[above right,font=\scriptsize,text=navy,fill=white,inner sep=.6pt] {};
\node[anchor=west,text=navy,font=\small] at(4.12,1.25){$G$};\end{tikzpicture}
\caption[Positive ground-state amplitude]{Positive ground-state amplitude. The maximum and the probability mean have different definitions. Fields are archived numerical illustrations; teal saturation increases with the displayed value.}\label{fig:field-Q0m}
\end{figure}

\paragraph{Demonstration and investigation.} Show the amplitude field and the measurement-point animation in the companion. Accumulate their sample mean, keeping the maximum marker fixed.


\section{The first energy states}\label{sec:slide50}

For an isolated eigenvalue, the amplitude sign does not change
probability or expectation. Excited states have nodal curves and
separate probability lobes; their mean need not sit at a peak.
The following four archived probability fields show how the lobes
change. The atlas compares their integral means, not arbitrary
lobe maxima. A degenerate eigenspace needs the cluster convention
developed next.


\[\mathrm{Q}j=\int_Tp|\psi_j|^2\dd A,\quad j=0,1,2,3,4.\]



An excited probability density is nonnegative even when its real
amplitude changes sign across nodal lines. A negative amplitude
lobe is not negative mass. The mean averages all lobes and can lie
outside high-density regions. Check mass orthogonality, eigenvalue
residuals and refinement for the first states.
Near crossings, ordered labels can exchange modes abruptly.
Track overlap matrices between neighboring triangles and compare
projectors in a nearly degenerate block. Even low triangular
eigenvalues have subtle shape dependence \cite{endo}.



\begin{figure}[htbp]\centering
\begin{tikzpicture}[x=1.7cm,y=1.7cm]
\fill[teal!8!white] (3.142857,0.6428571) -- (2.857143,0.8571429) -- (2.857143,0.6428571) -- cycle;
\fill[teal!8!white] (1.714286,1.714286) -- (1.428571,1.928571) -- (1.428571,1.714286) -- cycle;
\fill[teal!8!white] (1.714286,1.714286) -- (1.714286,1.5) -- (2,1.5) -- cycle;
\fill[teal!8!white] (0.5714286,2.571429) -- (0.2857143,2.785714) -- (0.2857143,2.571429) -- cycle;
\fill[teal!10!white] (2.571429,1.071429) -- (2.285714,1.285714) -- (2.285714,1.071429) -- cycle;
\fill[teal!8!white] (3.714286,0.2142857) -- (3.428571,0.4285714) -- (3.428571,0.2142857) -- cycle;
\fill[teal!8!white] (3.714286,0.2142857) -- (3.714286,0) -- (4,0) -- cycle;
\fill[teal!9!white] (0.8571429,2.142857) -- (1.142857,2.142857) -- (0.8571429,2.357143) -- cycle;
\fill[teal!9!white] (1.142857,2.142857) -- (1.142857,1.928571) -- (1.428571,1.928571) -- cycle;
\fill[teal!8!white] (0,2.785714) -- (0.2857143,2.785714) -- (0,3) -- cycle;
\fill[teal!8!white] (0.5714286,2.571429) -- (0.5714286,2.357143) -- (0.8571429,2.357143) -- cycle;
\fill[teal!10!white] (2,1.285714) -- (2.285714,1.285714) -- (2,1.5) -- cycle;
\fill[teal!9!white] (2.571429,1.071429) -- (2.571429,0.8571429) -- (2.857143,0.8571429) -- cycle;
\fill[teal!8!white] (3.142857,0.4285714) -- (3.428571,0.4285714) -- (3.142857,0.6428571) -- cycle;
\fill[teal!8!white] (0.2857143,0) -- (0,0.2142857) -- (0,0) -- cycle;
\fill[teal!8!white] (0,0.2142857) -- (0.2857143,0) -- (0.2857143,0.2142857) -- cycle;
\fill[teal!8!white] (0.5714286,0.2142857) -- (0.2857143,0) -- (0.5714286,0) -- cycle;
\fill[teal!8!white] (0.2857143,0) -- (0.5714286,0.2142857) -- (0.2857143,0.2142857) -- cycle;
\fill[teal!8!white] (1.714286,1.5) -- (1.714286,1.714286) -- (1.428571,1.714286) -- cycle;
\fill[teal!8!white] (1.428571,1.5) -- (1.714286,1.5) -- (1.428571,1.714286) -- cycle;
\fill[teal!11!white] (0,1.928571) -- (0,1.714286) -- (0.2857143,1.928571) -- cycle;
\fill[teal!20!white] (0,1.714286) -- (0.2857143,1.714286) -- (0.2857143,1.928571) -- cycle;
\fill[teal!8!white] (0.2857143,2.357143) -- (0,2.571429) -- (0,2.357143) -- cycle;
\fill[teal!8!white] (0,2.571429) -- (0.2857143,2.357143) -- (0.2857143,2.571429) -- cycle;
\fill[teal!9!white] (0,2.142857) -- (0.2857143,2.357143) -- (0,2.357143) -- cycle;
\fill[teal!11!white] (0.2857143,2.357143) -- (0,2.142857) -- (0.2857143,2.142857) -- cycle;
\fill[teal!10!white] (0,2.142857) -- (0,1.928571) -- (0.2857143,1.928571) -- cycle;
\fill[teal!14!white] (0.2857143,2.142857) -- (0,2.142857) -- (0.2857143,1.928571) -- cycle;
\fill[teal!40!white] (0.2857143,1.714286) -- (0.5714286,1.928571) -- (0.2857143,1.928571) -- cycle;
\fill[teal!56!white] (0.5714286,1.714286) -- (0.5714286,1.928571) -- (0.2857143,1.714286) -- cycle;
\fill[teal!12!white] (0.2857143,1.5) -- (0,1.714286) -- (0,1.5) -- cycle;
\fill[teal!22!white] (0,1.714286) -- (0.2857143,1.5) -- (0.2857143,1.714286) -- cycle;
\fill[teal!52!white] (0.2857143,1.5) -- (0.5714286,1.714286) -- (0.2857143,1.714286) -- cycle;
\fill[teal!69!white] (0.5714286,1.714286) -- (0.2857143,1.5) -- (0.5714286,1.5) -- cycle;
\fill[teal!12!white] (0,1.285714) -- (0.2857143,1.5) -- (0,1.5) -- cycle;
\fill[teal!23!white] (0.2857143,1.5) -- (0,1.285714) -- (0.2857143,1.285714) -- cycle;
\fill[teal!11!white] (0,1.285714) -- (0,1.071429) -- (0.2857143,1.071429) -- cycle;
\fill[teal!21!white] (0.2857143,1.285714) -- (0,1.285714) -- (0.2857143,1.071429) -- cycle;
\fill[teal!8!white] (3.428571,0) -- (3.714286,0) -- (3.428571,0.2142857) -- cycle;
\fill[teal!8!white] (3.714286,0) -- (3.714286,0.2142857) -- (3.428571,0.2142857) -- cycle;
\fill[teal!10!white] (1.428571,0) -- (1.714286,0) -- (1.428571,0.2142857) -- cycle;
\fill[teal!18!white] (1.714286,0) -- (1.714286,0.2142857) -- (1.428571,0.2142857) -- cycle;
\fill[teal!8!white] (0.8571429,0) -- (0.5714286,0.2142857) -- (0.5714286,0) -- cycle;
\fill[teal!8!white] (0.5714286,0.2142857) -- (0.8571429,0) -- (0.8571429,0.2142857) -- cycle;
\fill[teal!9!white] (1.142857,0.2142857) -- (1.142857,0) -- (1.428571,0) -- cycle;
\fill[teal!12!white] (1.142857,0.2142857) -- (1.428571,0) -- (1.428571,0.2142857) -- cycle;
\fill[teal!8!white] (1.142857,0.2142857) -- (0.8571429,0) -- (1.142857,0) -- cycle;
\fill[teal!8!white] (0.8571429,0) -- (1.142857,0.2142857) -- (0.8571429,0.2142857) -- cycle;
\fill[teal!16!white] (1.142857,1.928571) -- (1.142857,1.714286) -- (1.428571,1.714286) -- cycle;
\fill[teal!10!white] (1.428571,1.928571) -- (1.142857,1.928571) -- (1.428571,1.714286) -- cycle;
\fill[teal!12!white] (0.8571429,2.142857) -- (1.142857,1.928571) -- (1.142857,2.142857) -- cycle;
\fill[teal!23!white] (1.142857,1.928571) -- (0.8571429,2.142857) -- (0.8571429,1.928571) -- cycle;
\fill[teal!8!white] (0,2.785714) -- (0,2.571429) -- (0.2857143,2.571429) -- cycle;
\fill[teal!8!white] (0.2857143,2.785714) -- (0,2.785714) -- (0.2857143,2.571429) -- cycle;
\fill[teal!33!white] (1.142857,1.928571) -- (0.8571429,1.714286) -- (1.142857,1.714286) -- cycle;
\fill[teal!37!white] (0.8571429,1.714286) -- (1.142857,1.928571) -- (0.8571429,1.928571) -- cycle;
\fill[teal!66!white] (0.8571429,1.714286) -- (0.5714286,1.928571) -- (0.5714286,1.714286) -- cycle;
\fill[teal!54!white] (0.5714286,1.928571) -- (0.8571429,1.714286) -- (0.8571429,1.928571) -- cycle;
\fill[teal!10!white] (0.2857143,2.357143) -- (0.5714286,2.357143) -- (0.2857143,2.571429) -- cycle;
\fill[teal!9!white] (0.5714286,2.357143) -- (0.5714286,2.571429) -- (0.2857143,2.571429) -- cycle;
\fill[teal!9!white] (0,0.6428571) -- (0,0.4285714) -- (0.2857143,0.6428571) -- cycle;
\fill[teal!10!white] (0,0.4285714) -- (0.2857143,0.4285714) -- (0.2857143,0.6428571) -- cycle;
\fill[teal!8!white] (0,0.4285714) -- (0,0.2142857) -- (0.2857143,0.2142857) -- cycle;
\fill[teal!9!white] (0.2857143,0.4285714) -- (0,0.4285714) -- (0.2857143,0.2142857) -- cycle;
\fill[teal!15!white] (2.285714,1.285714) -- (2,1.285714) -- (2.285714,1.071429) -- cycle;
\fill[teal!34!white] (2,1.285714) -- (2,1.071429) -- (2.285714,1.071429) -- cycle;
\fill[teal!64!white] (2,1.071429) -- (2,0.8571429) -- (2.285714,1.071429) -- cycle;
\fill[teal!66!white] (2,0.8571429) -- (2.285714,0.8571429) -- (2.285714,1.071429) -- cycle;
\fill[teal!35!white] (2.285714,0.8571429) -- (2.571429,0.8571429) -- (2.285714,1.071429) -- cycle;
\fill[teal!15!white] (2.571429,0.8571429) -- (2.571429,1.071429) -- (2.285714,1.071429) -- cycle;
\fill[teal!56!white] (1.428571,0.4285714) -- (1.714286,0.2142857) -- (1.714286,0.4285714) -- cycle;
\fill[teal!35!white] (1.714286,0.2142857) -- (1.428571,0.4285714) -- (1.428571,0.2142857) -- cycle;
\fill[teal!8!white] (3.142857,0.2142857) -- (3.428571,0) -- (3.428571,0.2142857) -- cycle;
\fill[teal!8!white] (3.142857,0.2142857) -- (3.142857,0) -- (3.428571,0) -- cycle;
\fill[teal!8!white] (3.142857,0) -- (3.142857,0.2142857) -- (2.857143,0) -- cycle;
\fill[teal!9!white] (3.142857,0.2142857) -- (2.857143,0.2142857) -- (2.857143,0) -- cycle;
\fill[teal!20!white] (1.714286,0) -- (2,0.2142857) -- (1.714286,0.2142857) -- cycle;
\fill[teal!11!white] (2,0.2142857) -- (1.714286,0) -- (2,0) -- cycle;
\fill[teal!10!white] (0.5714286,0.2142857) -- (0.5714286,0.4285714) -- (0.2857143,0.2142857) -- cycle;
\fill[teal!11!white] (0.5714286,0.4285714) -- (0.2857143,0.4285714) -- (0.2857143,0.2142857) -- cycle;
\fill[teal!8!white] (0.8571429,0.4285714) -- (0.5714286,0.2142857) -- (0.8571429,0.2142857) -- cycle;
\fill[teal!10!white] (0.8571429,0.4285714) -- (0.5714286,0.4285714) -- (0.5714286,0.2142857) -- cycle;
\fill[teal!12!white] (1.142857,1.5) -- (1.428571,1.5) -- (1.428571,1.714286) -- cycle;
\fill[teal!19!white] (1.142857,1.714286) -- (1.142857,1.5) -- (1.428571,1.714286) -- cycle;
\fill[teal!25!white] (0.5714286,2.142857) -- (0.2857143,2.142857) -- (0.2857143,1.928571) -- cycle;
\fill[teal!36!white] (0.5714286,1.928571) -- (0.5714286,2.142857) -- (0.2857143,1.928571) -- cycle;
\fill[teal!17!white] (0.5714286,2.142857) -- (0.2857143,2.357143) -- (0.2857143,2.142857) -- cycle;
\fill[teal!15!white] (0.5714286,2.142857) -- (0.5714286,2.357143) -- (0.2857143,2.357143) -- cycle;
\fill[teal!27!white] (0.8571429,2.142857) -- (0.5714286,2.142857) -- (0.8571429,1.928571) -- cycle;
\fill[teal!39!white] (0.5714286,2.142857) -- (0.5714286,1.928571) -- (0.8571429,1.928571) -- cycle;
\fill[teal!15!white] (0.5714286,2.142857) -- (0.8571429,2.142857) -- (0.8571429,2.357143) -- cycle;
\fill[teal!13!white] (0.5714286,2.357143) -- (0.5714286,2.142857) -- (0.8571429,2.357143) -- cycle;
\fill[teal!16!white] (1.714286,1.285714) -- (2,1.285714) -- (2,1.5) -- cycle;
\fill[teal!12!white] (1.714286,1.5) -- (1.714286,1.285714) -- (2,1.5) -- cycle;
\fill[teal!18!white] (1.142857,1.071429) -- (1.428571,0.8571429) -- (1.428571,1.071429) -- cycle;
\fill[teal!12!white] (1.428571,0.8571429) -- (1.142857,1.071429) -- (1.142857,0.8571429) -- cycle;
\fill[teal!8!white] (3.142857,0.4285714) -- (3.142857,0.2142857) -- (3.428571,0.2142857) -- cycle;
\fill[teal!8!white] (3.428571,0.4285714) -- (3.142857,0.4285714) -- (3.428571,0.2142857) -- cycle;
\fill[teal!11!white] (2.857143,0.4285714) -- (3.142857,0.4285714) -- (2.857143,0.6428571) -- cycle;
\fill[teal!9!white] (3.142857,0.4285714) -- (3.142857,0.6428571) -- (2.857143,0.6428571) -- cycle;
\fill[teal!9!white] (3.142857,0.2142857) -- (3.142857,0.4285714) -- (2.857143,0.2142857) -- cycle;
\fill[teal!11!white] (3.142857,0.4285714) -- (2.857143,0.4285714) -- (2.857143,0.2142857) -- cycle;
\fill[teal!17!white] (2.571429,0.8571429) -- (2.571429,0.6428571) -- (2.857143,0.8571429) -- cycle;
\fill[teal!14!white] (2.857143,0.8571429) -- (2.571429,0.6428571) -- (2.857143,0.6428571) -- cycle;
\fill[teal!11!white] (2.285714,0) -- (2.285714,0.2142857) -- (2,0) -- cycle;
\fill[teal!20!white] (2.285714,0.2142857) -- (2,0.2142857) -- (2,0) -- cycle;
\fill[teal!50!white] (2.285714,0.2142857) -- (2,0.4285714) -- (2,0.2142857) -- cycle;
\fill[teal!62!white] (2,0.4285714) -- (2.285714,0.2142857) -- (2.285714,0.4285714) -- cycle;
\fill[teal!94!white] (2,0.4285714) -- (1.714286,0.6428571) -- (1.714286,0.4285714) -- cycle;
\fill[teal!94!white] (1.714286,0.6428571) -- (2,0.4285714) -- (2,0.6428571) -- cycle;
\fill[teal!69!white] (1.714286,0.2142857) -- (2,0.4285714) -- (1.714286,0.4285714) -- cycle;
\fill[teal!53!white] (2,0.2142857) -- (2,0.4285714) -- (1.714286,0.2142857) -- cycle;
\fill[teal!10!white] (0.2857143,0.8571429) -- (0,0.8571429) -- (0,0.6428571) -- cycle;
\fill[teal!13!white] (0.2857143,0.8571429) -- (0,0.6428571) -- (0.2857143,0.6428571) -- cycle;
\fill[teal!18!white] (0,1.071429) -- (0.2857143,0.8571429) -- (0.2857143,1.071429) -- cycle;
\fill[teal!10!white] (0,0.8571429) -- (0.2857143,0.8571429) -- (0,1.071429) -- cycle;
\fill[teal!71!white] (0.2857143,1.5) -- (0.5714286,1.285714) -- (0.5714286,1.5) -- cycle;
\fill[teal!55!white] (0.5714286,1.285714) -- (0.2857143,1.5) -- (0.2857143,1.285714) -- cycle;
\fill[teal!21!white] (0.5714286,0.6428571) -- (0.2857143,0.8571429) -- (0.2857143,0.6428571) -- cycle;
\fill[teal!29!white] (0.2857143,0.8571429) -- (0.5714286,0.6428571) -- (0.5714286,0.8571429) -- cycle;
\fill[teal!16!white] (0.2857143,0.4285714) -- (0.5714286,0.6428571) -- (0.2857143,0.6428571) -- cycle;
\fill[teal!15!white] (0.5714286,0.4285714) -- (0.5714286,0.6428571) -- (0.2857143,0.4285714) -- cycle;
\fill[teal!13!white] (0.8571429,0.4285714) -- (0.5714286,0.6428571) -- (0.5714286,0.4285714) -- cycle;
\fill[teal!12!white] (0.5714286,0.6428571) -- (0.8571429,0.4285714) -- (0.8571429,0.6428571) -- cycle;
\fill[teal!16!white] (1.142857,0.4285714) -- (1.142857,0.2142857) -- (1.428571,0.2142857) -- cycle;
\fill[teal!27!white] (1.428571,0.4285714) -- (1.142857,0.4285714) -- (1.428571,0.2142857) -- cycle;
\fill[teal!10!white] (1.142857,0.2142857) -- (1.142857,0.4285714) -- (0.8571429,0.2142857) -- cycle;
\fill[teal!8!white] (1.142857,0.4285714) -- (0.8571429,0.4285714) -- (0.8571429,0.2142857) -- cycle;
\fill[teal!8!white] (0.8571429,0.4285714) -- (1.142857,0.4285714) -- (0.8571429,0.6428571) -- cycle;
\fill[teal!9!white] (1.142857,0.4285714) -- (1.142857,0.6428571) -- (0.8571429,0.6428571) -- cycle;
\fill[teal!9!white] (1.428571,1.285714) -- (1.714286,1.5) -- (1.428571,1.5) -- cycle;
\fill[teal!15!white] (1.428571,1.285714) -- (1.714286,1.285714) -- (1.714286,1.5) -- cycle;
\fill[teal!10!white] (1.142857,1.5) -- (1.428571,1.285714) -- (1.428571,1.5) -- cycle;
\fill[teal!13!white] (1.428571,1.285714) -- (1.142857,1.5) -- (1.142857,1.285714) -- cycle;
\fill[teal!11!white] (1.428571,1.285714) -- (1.142857,1.071429) -- (1.428571,1.071429) -- cycle;
\fill[teal!9!white] (1.142857,1.071429) -- (1.428571,1.285714) -- (1.142857,1.285714) -- cycle;
\fill[teal!66!white] (1.428571,0.6428571) -- (1.428571,0.4285714) -- (1.714286,0.4285714) -- cycle;
\fill[teal!89!white] (1.714286,0.6428571) -- (1.428571,0.6428571) -- (1.714286,0.4285714) -- cycle;
\fill[teal!36!white] (1.428571,0.6428571) -- (1.142857,0.4285714) -- (1.428571,0.4285714) -- cycle;
\fill[teal!24!white] (1.142857,0.4285714) -- (1.428571,0.6428571) -- (1.142857,0.6428571) -- cycle;
\fill[teal!32!white] (1.428571,0.6428571) -- (1.428571,0.8571429) -- (1.142857,0.8571429) -- cycle;
\fill[teal!21!white] (1.142857,0.6428571) -- (1.428571,0.6428571) -- (1.142857,0.8571429) -- cycle;
\fill[teal!48!white] (2,1.285714) -- (1.714286,1.071429) -- (2,1.071429) -- cycle;
\fill[teal!36!white] (1.714286,1.285714) -- (1.714286,1.071429) -- (2,1.285714) -- cycle;
\fill[teal!28!white] (1.428571,1.285714) -- (1.714286,1.071429) -- (1.714286,1.285714) -- cycle;
\fill[teal!28!white] (1.714286,1.071429) -- (1.428571,1.285714) -- (1.428571,1.071429) -- cycle;
\fill[teal!55!white] (2.285714,0.6428571) -- (2.571429,0.8571429) -- (2.285714,0.8571429) -- cycle;
\fill[teal!46!white] (2.285714,0.6428571) -- (2.571429,0.6428571) -- (2.571429,0.8571429) -- cycle;
\fill[teal!91!white] (2,0.8571429) -- (2.285714,0.6428571) -- (2.285714,0.8571429) -- cycle;
\fill[teal!94!white] (2.285714,0.6428571) -- (2,0.8571429) -- (2,0.6428571) -- cycle;
\fill[teal!89!white] (2.285714,0.6428571) -- (2,0.4285714) -- (2.285714,0.4285714) -- cycle;
\fill[teal!94!white] (2,0.4285714) -- (2.285714,0.6428571) -- (2,0.6428571) -- cycle;
\fill[teal!20!white] (2.571429,0.4285714) -- (2.857143,0.4285714) -- (2.857143,0.6428571) -- cycle;
\fill[teal!27!white] (2.571429,0.6428571) -- (2.571429,0.4285714) -- (2.857143,0.6428571) -- cycle;
\fill[teal!52!white] (2.285714,0.6428571) -- (2.571429,0.4285714) -- (2.571429,0.6428571) -- cycle;
\fill[teal!63!white] (2.571429,0.4285714) -- (2.285714,0.6428571) -- (2.285714,0.4285714) -- cycle;
\fill[teal!37!white] (0.2857143,0.8571429) -- (0.5714286,1.071429) -- (0.2857143,1.071429) -- cycle;
\fill[teal!39!white] (0.5714286,1.071429) -- (0.2857143,0.8571429) -- (0.5714286,0.8571429) -- cycle;
\fill[teal!46!white] (0.5714286,1.071429) -- (0.2857143,1.285714) -- (0.2857143,1.071429) -- cycle;
\fill[teal!62!white] (0.5714286,1.071429) -- (0.5714286,1.285714) -- (0.2857143,1.285714) -- cycle;
\fill[teal!83!white] (0.8571429,1.5) -- (0.5714286,1.714286) -- (0.5714286,1.5) -- cycle;
\fill[teal!76!white] (0.8571429,1.5) -- (0.8571429,1.714286) -- (0.5714286,1.714286) -- cycle;
\fill[teal!54!white] (0.8571429,1.714286) -- (0.8571429,1.5) -- (1.142857,1.714286) -- cycle;
\fill[teal!42!white] (0.8571429,1.5) -- (1.142857,1.5) -- (1.142857,1.714286) -- cycle;
\fill[teal!94!white] (1.714286,0.8571429) -- (1.714286,0.6428571) -- (2,0.6428571) -- cycle;
\fill[teal!94!white] (2,0.8571429) -- (1.714286,0.8571429) -- (2,0.6428571) -- cycle;
\fill[teal!94!white] (1.714286,0.8571429) -- (2,0.8571429) -- (2,1.071429) -- cycle;
\fill[teal!80!white] (1.714286,1.071429) -- (1.714286,0.8571429) -- (2,1.071429) -- cycle;
\fill[teal!94!white] (1.714286,0.8571429) -- (1.428571,0.6428571) -- (1.714286,0.6428571) -- cycle;
\fill[teal!70!white] (1.428571,0.6428571) -- (1.714286,0.8571429) -- (1.428571,0.8571429) -- cycle;
\fill[teal!55!white] (1.428571,0.8571429) -- (1.714286,0.8571429) -- (1.428571,1.071429) -- cycle;
\fill[teal!60!white] (1.714286,0.8571429) -- (1.714286,1.071429) -- (1.428571,1.071429) -- cycle;
\fill[teal!10!white] (2.571429,0) -- (2.571429,0.2142857) -- (2.285714,0) -- cycle;
\fill[teal!16!white] (2.571429,0.2142857) -- (2.285714,0.2142857) -- (2.285714,0) -- cycle;
\fill[teal!35!white] (2.285714,0.2142857) -- (2.571429,0.2142857) -- (2.285714,0.4285714) -- cycle;
\fill[teal!38!white] (2.571429,0.2142857) -- (2.571429,0.4285714) -- (2.285714,0.4285714) -- cycle;
\fill[teal!10!white] (2.857143,0.2142857) -- (2.571429,0.2142857) -- (2.857143,0) -- cycle;
\fill[teal!9!white] (2.571429,0.2142857) -- (2.571429,0) -- (2.857143,0) -- cycle;
\fill[teal!15!white] (2.857143,0.4285714) -- (2.571429,0.2142857) -- (2.857143,0.2142857) -- cycle;
\fill[teal!23!white] (2.571429,0.4285714) -- (2.571429,0.2142857) -- (2.857143,0.4285714) -- cycle;
\fill[teal!9!white] (1.142857,0.6428571) -- (0.8571429,0.8571429) -- (0.8571429,0.6428571) -- cycle;
\fill[teal!8!white] (0.8571429,0.8571429) -- (1.142857,0.6428571) -- (1.142857,0.8571429) -- cycle;
\fill[teal!27!white] (0.5714286,0.6428571) -- (0.8571429,0.8571429) -- (0.5714286,0.8571429) -- cycle;
\fill[teal!18!white] (0.8571429,0.8571429) -- (0.5714286,0.6428571) -- (0.8571429,0.6428571) -- cycle;
\fill[teal!35!white] (1.142857,1.5) -- (0.8571429,1.285714) -- (1.142857,1.285714) -- cycle;
\fill[teal!54!white] (0.8571429,1.5) -- (0.8571429,1.285714) -- (1.142857,1.5) -- cycle;
\fill[teal!80!white] (0.5714286,1.285714) -- (0.8571429,1.285714) -- (0.5714286,1.5) -- cycle;
\fill[teal!79!white] (0.8571429,1.285714) -- (0.8571429,1.5) -- (0.5714286,1.5) -- cycle;
\fill[teal!12!white] (1.142857,1.071429) -- (0.8571429,1.071429) -- (1.142857,0.8571429) -- cycle;
\fill[teal!16!white] (0.8571429,1.071429) -- (0.8571429,0.8571429) -- (1.142857,0.8571429) -- cycle;
\fill[teal!20!white] (0.8571429,1.071429) -- (1.142857,1.071429) -- (1.142857,1.285714) -- cycle;
\fill[teal!39!white] (0.8571429,1.285714) -- (0.8571429,1.071429) -- (1.142857,1.285714) -- cycle;
\fill[teal!47!white] (0.8571429,1.071429) -- (0.5714286,1.071429) -- (0.5714286,0.8571429) -- cycle;
\fill[teal!34!white] (0.8571429,0.8571429) -- (0.8571429,1.071429) -- (0.5714286,0.8571429) -- cycle;
\fill[teal!62!white] (0.5714286,1.071429) -- (0.8571429,1.071429) -- (0.5714286,1.285714) -- cycle;
\fill[teal!64!white] (0.8571429,1.071429) -- (0.8571429,1.285714) -- (0.5714286,1.285714) -- cycle;
\draw[teal,thick] (0,0) -- (4,0) -- (0,3) -- cycle;
\draw[magenta!70,thin] (1.413362,0.9877257) -- (3.6,2.35) -- (4.1,2.35);\filldraw[fill=magenta,draw=white,line width=.45pt] (1.413362,0.9877257) circle(1.5pt) node[above right,font=\scriptsize,text=magenta,fill=white,inner sep=.6pt] {};
\node[anchor=west,text=magenta,font=\small] at(4.12,2.35){Q1};\draw[gold!70,thin] (1,1) -- (3.6,1.8) -- (4.1,1.8);\filldraw[fill=gold,draw=white,line width=.45pt] (1,1) circle(1.5pt) node[above right,font=\scriptsize,text=gold,fill=white,inner sep=.6pt] {};
\node[anchor=west,text=gold,font=\small] at(4.12,1.8){$I$};\draw[navy!70,thin] (1.333333,1) -- (3.6,1.25) -- (4.1,1.25);\filldraw[fill=navy,draw=white,line width=.45pt] (1.333333,1) circle(1.5pt) node[above right,font=\scriptsize,text=navy,fill=white,inner sep=.6pt] {};
\node[anchor=west,text=navy,font=\small] at(4.12,1.25){$G$};\end{tikzpicture}
\caption[Probability density for Q1]{Probability density for Q1. The first moment averages all probability, including separate lobes. Numerical illustration.}\label{fig:state-Q1}
\end{figure}
\begin{figure}[htbp]\centering
\begin{tikzpicture}[x=1.7cm,y=1.7cm]
\fill[teal!8!white] (3.142857,0.6428571) -- (2.857143,0.8571429) -- (2.857143,0.6428571) -- cycle;
\fill[teal!9!white] (1.714286,1.714286) -- (1.428571,1.928571) -- (1.428571,1.714286) -- cycle;
\fill[teal!9!white] (1.714286,1.714286) -- (1.714286,1.5) -- (2,1.5) -- cycle;
\fill[teal!8!white] (0.5714286,2.571429) -- (0.2857143,2.785714) -- (0.2857143,2.571429) -- cycle;
\fill[teal!9!white] (2.571429,1.071429) -- (2.285714,1.285714) -- (2.285714,1.071429) -- cycle;
\fill[teal!8!white] (3.714286,0.2142857) -- (3.428571,0.4285714) -- (3.428571,0.2142857) -- cycle;
\fill[teal!8!white] (3.714286,0.2142857) -- (3.714286,0) -- (4,0) -- cycle;
\fill[teal!10!white] (0.8571429,2.142857) -- (1.142857,2.142857) -- (0.8571429,2.357143) -- cycle;
\fill[teal!10!white] (1.142857,2.142857) -- (1.142857,1.928571) -- (1.428571,1.928571) -- cycle;
\fill[teal!8!white] (0,2.785714) -- (0.2857143,2.785714) -- (0,3) -- cycle;
\fill[teal!9!white] (0.5714286,2.571429) -- (0.5714286,2.357143) -- (0.8571429,2.357143) -- cycle;
\fill[teal!9!white] (2,1.285714) -- (2.285714,1.285714) -- (2,1.5) -- cycle;
\fill[teal!9!white] (2.571429,1.071429) -- (2.571429,0.8571429) -- (2.857143,0.8571429) -- cycle;
\fill[teal!8!white] (3.142857,0.4285714) -- (3.428571,0.4285714) -- (3.142857,0.6428571) -- cycle;
\fill[teal!8!white] (0.2857143,0) -- (0,0.2142857) -- (0,0) -- cycle;
\fill[teal!10!white] (0,0.2142857) -- (0.2857143,0) -- (0.2857143,0.2142857) -- cycle;
\fill[teal!11!white] (0.5714286,0.2142857) -- (0.2857143,0) -- (0.5714286,0) -- cycle;
\fill[teal!16!white] (0.2857143,0) -- (0.5714286,0.2142857) -- (0.2857143,0.2142857) -- cycle;
\fill[teal!11!white] (1.714286,1.5) -- (1.714286,1.714286) -- (1.428571,1.714286) -- cycle;
\fill[teal!19!white] (1.428571,1.5) -- (1.714286,1.5) -- (1.428571,1.714286) -- cycle;
\fill[teal!12!white] (0,1.928571) -- (0,1.714286) -- (0.2857143,1.928571) -- cycle;
\fill[teal!24!white] (0,1.714286) -- (0.2857143,1.714286) -- (0.2857143,1.928571) -- cycle;
\fill[teal!9!white] (0.2857143,2.357143) -- (0,2.571429) -- (0,2.357143) -- cycle;
\fill[teal!9!white] (0,2.571429) -- (0.2857143,2.357143) -- (0.2857143,2.571429) -- cycle;
\fill[teal!10!white] (0,2.142857) -- (0.2857143,2.357143) -- (0,2.357143) -- cycle;
\fill[teal!16!white] (0.2857143,2.357143) -- (0,2.142857) -- (0.2857143,2.142857) -- cycle;
\fill[teal!12!white] (0,2.142857) -- (0,1.928571) -- (0.2857143,1.928571) -- cycle;
\fill[teal!21!white] (0.2857143,2.142857) -- (0,2.142857) -- (0.2857143,1.928571) -- cycle;
\fill[teal!60!white] (0.2857143,1.714286) -- (0.5714286,1.928571) -- (0.2857143,1.928571) -- cycle;
\fill[teal!78!white] (0.5714286,1.714286) -- (0.5714286,1.928571) -- (0.2857143,1.714286) -- cycle;
\fill[teal!11!white] (0.2857143,1.5) -- (0,1.714286) -- (0,1.5) -- cycle;
\fill[teal!21!white] (0,1.714286) -- (0.2857143,1.5) -- (0.2857143,1.714286) -- cycle;
\fill[teal!51!white] (0.2857143,1.5) -- (0.5714286,1.714286) -- (0.2857143,1.714286) -- cycle;
\fill[teal!56!white] (0.5714286,1.714286) -- (0.2857143,1.5) -- (0.5714286,1.5) -- cycle;
\fill[teal!9!white] (0,1.285714) -- (0.2857143,1.5) -- (0,1.5) -- cycle;
\fill[teal!10!white] (0.2857143,1.5) -- (0,1.285714) -- (0.2857143,1.285714) -- cycle;
\fill[teal!8!white] (0,1.285714) -- (0,1.071429) -- (0.2857143,1.071429) -- cycle;
\fill[teal!8!white] (0.2857143,1.285714) -- (0,1.285714) -- (0.2857143,1.071429) -- cycle;
\fill[teal!8!white] (3.428571,0) -- (3.714286,0) -- (3.428571,0.2142857) -- cycle;
\fill[teal!8!white] (3.714286,0) -- (3.714286,0.2142857) -- (3.428571,0.2142857) -- cycle;
\fill[teal!8!white] (1.428571,0) -- (1.714286,0) -- (1.428571,0.2142857) -- cycle;
\fill[teal!8!white] (1.714286,0) -- (1.714286,0.2142857) -- (1.428571,0.2142857) -- cycle;
\fill[teal!12!white] (0.8571429,0) -- (0.5714286,0.2142857) -- (0.5714286,0) -- cycle;
\fill[teal!25!white] (0.5714286,0.2142857) -- (0.8571429,0) -- (0.8571429,0.2142857) -- cycle;
\fill[teal!10!white] (1.142857,0.2142857) -- (1.142857,0) -- (1.428571,0) -- cycle;
\fill[teal!13!white] (1.142857,0.2142857) -- (1.428571,0) -- (1.428571,0.2142857) -- cycle;
\fill[teal!12!white] (1.142857,0.2142857) -- (0.8571429,0) -- (1.142857,0) -- cycle;
\fill[teal!24!white] (0.8571429,0) -- (1.142857,0.2142857) -- (0.8571429,0.2142857) -- cycle;
\fill[teal!31!white] (1.142857,1.928571) -- (1.142857,1.714286) -- (1.428571,1.714286) -- cycle;
\fill[teal!14!white] (1.428571,1.928571) -- (1.142857,1.928571) -- (1.428571,1.714286) -- cycle;
\fill[teal!17!white] (0.8571429,2.142857) -- (1.142857,1.928571) -- (1.142857,2.142857) -- cycle;
\fill[teal!41!white] (1.142857,1.928571) -- (0.8571429,2.142857) -- (0.8571429,1.928571) -- cycle;
\fill[teal!8!white] (0,2.785714) -- (0,2.571429) -- (0.2857143,2.571429) -- cycle;
\fill[teal!8!white] (0.2857143,2.785714) -- (0,2.785714) -- (0.2857143,2.571429) -- cycle;
\fill[teal!60!white] (1.142857,1.928571) -- (0.8571429,1.714286) -- (1.142857,1.714286) -- cycle;
\fill[teal!66!white] (0.8571429,1.714286) -- (1.142857,1.928571) -- (0.8571429,1.928571) -- cycle;
\fill[teal!94!white] (0.8571429,1.714286) -- (0.5714286,1.928571) -- (0.5714286,1.714286) -- cycle;
\fill[teal!90!white] (0.5714286,1.928571) -- (0.8571429,1.714286) -- (0.8571429,1.928571) -- cycle;
\fill[teal!14!white] (0.2857143,2.357143) -- (0.5714286,2.357143) -- (0.2857143,2.571429) -- cycle;
\fill[teal!10!white] (0.5714286,2.357143) -- (0.5714286,2.571429) -- (0.2857143,2.571429) -- cycle;
\fill[teal!13!white] (0,0.6428571) -- (0,0.4285714) -- (0.2857143,0.6428571) -- cycle;
\fill[teal!24!white] (0,0.4285714) -- (0.2857143,0.4285714) -- (0.2857143,0.6428571) -- cycle;
\fill[teal!10!white] (0,0.4285714) -- (0,0.2142857) -- (0.2857143,0.2142857) -- cycle;
\fill[teal!20!white] (0.2857143,0.4285714) -- (0,0.4285714) -- (0.2857143,0.2142857) -- cycle;
\fill[teal!11!white] (2.285714,1.285714) -- (2,1.285714) -- (2.285714,1.071429) -- cycle;
\fill[teal!17!white] (2,1.285714) -- (2,1.071429) -- (2.285714,1.071429) -- cycle;
\fill[teal!26!white] (2,1.071429) -- (2,0.8571429) -- (2.285714,1.071429) -- cycle;
\fill[teal!28!white] (2,0.8571429) -- (2.285714,0.8571429) -- (2.285714,1.071429) -- cycle;
\fill[teal!19!white] (2.285714,0.8571429) -- (2.571429,0.8571429) -- (2.285714,1.071429) -- cycle;
\fill[teal!11!white] (2.571429,0.8571429) -- (2.571429,1.071429) -- (2.285714,1.071429) -- cycle;
\fill[teal!9!white] (1.428571,0.4285714) -- (1.714286,0.2142857) -- (1.714286,0.4285714) -- cycle;
\fill[teal!11!white] (1.714286,0.2142857) -- (1.428571,0.4285714) -- (1.428571,0.2142857) -- cycle;
\fill[teal!8!white] (3.142857,0.2142857) -- (3.428571,0) -- (3.428571,0.2142857) -- cycle;
\fill[teal!8!white] (3.142857,0.2142857) -- (3.142857,0) -- (3.428571,0) -- cycle;
\fill[teal!8!white] (3.142857,0) -- (3.142857,0.2142857) -- (2.857143,0) -- cycle;
\fill[teal!9!white] (3.142857,0.2142857) -- (2.857143,0.2142857) -- (2.857143,0) -- cycle;
\fill[teal!8!white] (1.714286,0) -- (2,0.2142857) -- (1.714286,0.2142857) -- cycle;
\fill[teal!8!white] (2,0.2142857) -- (1.714286,0) -- (2,0) -- cycle;
\fill[teal!47!white] (0.5714286,0.2142857) -- (0.5714286,0.4285714) -- (0.2857143,0.2142857) -- cycle;
\fill[teal!47!white] (0.5714286,0.4285714) -- (0.2857143,0.4285714) -- (0.2857143,0.2142857) -- cycle;
\fill[teal!66!white] (0.8571429,0.4285714) -- (0.5714286,0.2142857) -- (0.8571429,0.2142857) -- cycle;
\fill[teal!85!white] (0.8571429,0.4285714) -- (0.5714286,0.4285714) -- (0.5714286,0.2142857) -- cycle;
\fill[teal!33!white] (1.142857,1.5) -- (1.428571,1.5) -- (1.428571,1.714286) -- cycle;
\fill[teal!41!white] (1.142857,1.714286) -- (1.142857,1.5) -- (1.428571,1.714286) -- cycle;
\fill[teal!46!white] (0.5714286,2.142857) -- (0.2857143,2.142857) -- (0.2857143,1.928571) -- cycle;
\fill[teal!65!white] (0.5714286,1.928571) -- (0.5714286,2.142857) -- (0.2857143,1.928571) -- cycle;
\fill[teal!31!white] (0.5714286,2.142857) -- (0.2857143,2.357143) -- (0.2857143,2.142857) -- cycle;
\fill[teal!26!white] (0.5714286,2.142857) -- (0.5714286,2.357143) -- (0.2857143,2.357143) -- cycle;
\fill[teal!51!white] (0.8571429,2.142857) -- (0.5714286,2.142857) -- (0.8571429,1.928571) -- cycle;
\fill[teal!72!white] (0.5714286,2.142857) -- (0.5714286,1.928571) -- (0.8571429,1.928571) -- cycle;
\fill[teal!24!white] (0.5714286,2.142857) -- (0.8571429,2.142857) -- (0.8571429,2.357143) -- cycle;
\fill[teal!20!white] (0.5714286,2.357143) -- (0.5714286,2.142857) -- (0.8571429,2.357143) -- cycle;
\fill[teal!13!white] (1.714286,1.285714) -- (2,1.285714) -- (2,1.5) -- cycle;
\fill[teal!12!white] (1.714286,1.5) -- (1.714286,1.285714) -- (2,1.5) -- cycle;
\fill[teal!9!white] (1.142857,1.071429) -- (1.428571,0.8571429) -- (1.428571,1.071429) -- cycle;
\fill[teal!18!white] (1.428571,0.8571429) -- (1.142857,1.071429) -- (1.142857,0.8571429) -- cycle;
\fill[teal!8!white] (3.142857,0.4285714) -- (3.142857,0.2142857) -- (3.428571,0.2142857) -- cycle;
\fill[teal!8!white] (3.428571,0.4285714) -- (3.142857,0.4285714) -- (3.428571,0.2142857) -- cycle;
\fill[teal!10!white] (2.857143,0.4285714) -- (3.142857,0.4285714) -- (2.857143,0.6428571) -- cycle;
\fill[teal!9!white] (3.142857,0.4285714) -- (3.142857,0.6428571) -- (2.857143,0.6428571) -- cycle;
\fill[teal!9!white] (3.142857,0.2142857) -- (3.142857,0.4285714) -- (2.857143,0.2142857) -- cycle;
\fill[teal!10!white] (3.142857,0.4285714) -- (2.857143,0.4285714) -- (2.857143,0.2142857) -- cycle;
\fill[teal!13!white] (2.571429,0.8571429) -- (2.571429,0.6428571) -- (2.857143,0.8571429) -- cycle;
\fill[teal!12!white] (2.857143,0.8571429) -- (2.571429,0.6428571) -- (2.857143,0.6428571) -- cycle;
\fill[teal!9!white] (2.285714,0) -- (2.285714,0.2142857) -- (2,0) -- cycle;
\fill[teal!11!white] (2.285714,0.2142857) -- (2,0.2142857) -- (2,0) -- cycle;
\fill[teal!18!white] (2.285714,0.2142857) -- (2,0.4285714) -- (2,0.2142857) -- cycle;
\fill[teal!24!white] (2,0.4285714) -- (2.285714,0.2142857) -- (2.285714,0.4285714) -- cycle;
\fill[teal!13!white] (2,0.4285714) -- (1.714286,0.6428571) -- (1.714286,0.4285714) -- cycle;
\fill[teal!21!white] (1.714286,0.6428571) -- (2,0.4285714) -- (2,0.6428571) -- cycle;
\fill[teal!10!white] (1.714286,0.2142857) -- (2,0.4285714) -- (1.714286,0.4285714) -- cycle;
\fill[teal!12!white] (2,0.2142857) -- (2,0.4285714) -- (1.714286,0.2142857) -- cycle;
\fill[teal!12!white] (0.2857143,0.8571429) -- (0,0.8571429) -- (0,0.6428571) -- cycle;
\fill[teal!24!white] (0.2857143,0.8571429) -- (0,0.6428571) -- (0.2857143,0.6428571) -- cycle;
\fill[teal!12!white] (0,1.071429) -- (0.2857143,0.8571429) -- (0.2857143,1.071429) -- cycle;
\fill[teal!10!white] (0,0.8571429) -- (0.2857143,0.8571429) -- (0,1.071429) -- cycle;
\fill[teal!31!white] (0.2857143,1.5) -- (0.5714286,1.285714) -- (0.5714286,1.5) -- cycle;
\fill[teal!16!white] (0.5714286,1.285714) -- (0.2857143,1.5) -- (0.2857143,1.285714) -- cycle;
\fill[teal!62!white] (0.5714286,0.6428571) -- (0.2857143,0.8571429) -- (0.2857143,0.6428571) -- cycle;
\fill[teal!71!white] (0.2857143,0.8571429) -- (0.5714286,0.6428571) -- (0.5714286,0.8571429) -- cycle;
\fill[teal!67!white] (0.2857143,0.4285714) -- (0.5714286,0.6428571) -- (0.2857143,0.6428571) -- cycle;
\fill[teal!85!white] (0.5714286,0.4285714) -- (0.5714286,0.6428571) -- (0.2857143,0.4285714) -- cycle;
\fill[teal!94!white] (0.8571429,0.4285714) -- (0.5714286,0.6428571) -- (0.5714286,0.4285714) -- cycle;
\fill[teal!94!white] (0.5714286,0.6428571) -- (0.8571429,0.4285714) -- (0.8571429,0.6428571) -- cycle;
\fill[teal!34!white] (1.142857,0.4285714) -- (1.142857,0.2142857) -- (1.428571,0.2142857) -- cycle;
\fill[teal!30!white] (1.428571,0.4285714) -- (1.142857,0.4285714) -- (1.428571,0.2142857) -- cycle;
\fill[teal!54!white] (1.142857,0.2142857) -- (1.142857,0.4285714) -- (0.8571429,0.2142857) -- cycle;
\fill[teal!81!white] (1.142857,0.4285714) -- (0.8571429,0.4285714) -- (0.8571429,0.2142857) -- cycle;
\fill[teal!94!white] (0.8571429,0.4285714) -- (1.142857,0.4285714) -- (0.8571429,0.6428571) -- cycle;
\fill[teal!94!white] (1.142857,0.4285714) -- (1.142857,0.6428571) -- (0.8571429,0.6428571) -- cycle;
\fill[teal!22!white] (1.428571,1.285714) -- (1.714286,1.5) -- (1.428571,1.5) -- cycle;
\fill[teal!20!white] (1.428571,1.285714) -- (1.714286,1.285714) -- (1.714286,1.5) -- cycle;
\fill[teal!30!white] (1.142857,1.5) -- (1.428571,1.285714) -- (1.428571,1.5) -- cycle;
\fill[teal!27!white] (1.428571,1.285714) -- (1.142857,1.5) -- (1.142857,1.285714) -- cycle;
\fill[teal!11!white] (1.428571,1.285714) -- (1.142857,1.071429) -- (1.428571,1.071429) -- cycle;
\fill[teal!13!white] (1.142857,1.071429) -- (1.428571,1.285714) -- (1.142857,1.285714) -- cycle;
\fill[teal!13!white] (1.428571,0.6428571) -- (1.428571,0.4285714) -- (1.714286,0.4285714) -- cycle;
\fill[teal!8!white] (1.714286,0.6428571) -- (1.428571,0.6428571) -- (1.714286,0.4285714) -- cycle;
\fill[teal!37!white] (1.428571,0.6428571) -- (1.142857,0.4285714) -- (1.428571,0.4285714) -- cycle;
\fill[teal!55!white] (1.142857,0.4285714) -- (1.428571,0.6428571) -- (1.142857,0.6428571) -- cycle;
\fill[teal!22!white] (1.428571,0.6428571) -- (1.428571,0.8571429) -- (1.142857,0.8571429) -- cycle;
\fill[teal!45!white] (1.142857,0.6428571) -- (1.428571,0.6428571) -- (1.142857,0.8571429) -- cycle;
\fill[teal!21!white] (2,1.285714) -- (1.714286,1.071429) -- (2,1.071429) -- cycle;
\fill[teal!20!white] (1.714286,1.285714) -- (1.714286,1.071429) -- (2,1.285714) -- cycle;
\fill[teal!20!white] (1.428571,1.285714) -- (1.714286,1.071429) -- (1.714286,1.285714) -- cycle;
\fill[teal!15!white] (1.714286,1.071429) -- (1.428571,1.285714) -- (1.428571,1.071429) -- cycle;
\fill[teal!28!white] (2.285714,0.6428571) -- (2.571429,0.8571429) -- (2.285714,0.8571429) -- cycle;
\fill[teal!26!white] (2.285714,0.6428571) -- (2.571429,0.6428571) -- (2.571429,0.8571429) -- cycle;
\fill[teal!35!white] (2,0.8571429) -- (2.285714,0.6428571) -- (2.285714,0.8571429) -- cycle;
\fill[teal!36!white] (2.285714,0.6428571) -- (2,0.8571429) -- (2,0.6428571) -- cycle;
\fill[teal!33!white] (2.285714,0.6428571) -- (2,0.4285714) -- (2.285714,0.4285714) -- cycle;
\fill[teal!33!white] (2,0.4285714) -- (2.285714,0.6428571) -- (2,0.6428571) -- cycle;
\fill[teal!16!white] (2.571429,0.4285714) -- (2.857143,0.4285714) -- (2.857143,0.6428571) -- cycle;
\fill[teal!20!white] (2.571429,0.6428571) -- (2.571429,0.4285714) -- (2.857143,0.6428571) -- cycle;
\fill[teal!29!white] (2.285714,0.6428571) -- (2.571429,0.4285714) -- (2.571429,0.6428571) -- cycle;
\fill[teal!32!white] (2.571429,0.4285714) -- (2.285714,0.6428571) -- (2.285714,0.4285714) -- cycle;
\fill[teal!23!white] (0.2857143,0.8571429) -- (0.5714286,1.071429) -- (0.2857143,1.071429) -- cycle;
\fill[teal!42!white] (0.5714286,1.071429) -- (0.2857143,0.8571429) -- (0.5714286,0.8571429) -- cycle;
\fill[teal!10!white] (0.5714286,1.071429) -- (0.2857143,1.285714) -- (0.2857143,1.071429) -- cycle;
\fill[teal!8!white] (0.5714286,1.071429) -- (0.5714286,1.285714) -- (0.2857143,1.285714) -- cycle;
\fill[teal!74!white] (0.8571429,1.5) -- (0.5714286,1.714286) -- (0.5714286,1.5) -- cycle;
\fill[teal!88!white] (0.8571429,1.5) -- (0.8571429,1.714286) -- (0.5714286,1.714286) -- cycle;
\fill[teal!77!white] (0.8571429,1.714286) -- (0.8571429,1.5) -- (1.142857,1.714286) -- cycle;
\fill[teal!62!white] (0.8571429,1.5) -- (1.142857,1.5) -- (1.142857,1.714286) -- cycle;
\fill[teal!18!white] (1.714286,0.8571429) -- (1.714286,0.6428571) -- (2,0.6428571) -- cycle;
\fill[teal!26!white] (2,0.8571429) -- (1.714286,0.8571429) -- (2,0.6428571) -- cycle;
\fill[teal!26!white] (1.714286,0.8571429) -- (2,0.8571429) -- (2,1.071429) -- cycle;
\fill[teal!22!white] (1.714286,1.071429) -- (1.714286,0.8571429) -- (2,1.071429) -- cycle;
\fill[teal!8!white] (1.714286,0.8571429) -- (1.428571,0.6428571) -- (1.714286,0.6428571) -- cycle;
\fill[teal!9!white] (1.428571,0.6428571) -- (1.714286,0.8571429) -- (1.428571,0.8571429) -- cycle;
\fill[teal!9!white] (1.428571,0.8571429) -- (1.714286,0.8571429) -- (1.428571,1.071429) -- cycle;
\fill[teal!14!white] (1.714286,0.8571429) -- (1.714286,1.071429) -- (1.428571,1.071429) -- cycle;
\fill[teal!9!white] (2.571429,0) -- (2.571429,0.2142857) -- (2.285714,0) -- cycle;
\fill[teal!11!white] (2.571429,0.2142857) -- (2.285714,0.2142857) -- (2.285714,0) -- cycle;
\fill[teal!20!white] (2.285714,0.2142857) -- (2.571429,0.2142857) -- (2.285714,0.4285714) -- cycle;
\fill[teal!23!white] (2.571429,0.2142857) -- (2.571429,0.4285714) -- (2.285714,0.4285714) -- cycle;
\fill[teal!10!white] (2.857143,0.2142857) -- (2.571429,0.2142857) -- (2.857143,0) -- cycle;
\fill[teal!9!white] (2.571429,0.2142857) -- (2.571429,0) -- (2.857143,0) -- cycle;
\fill[teal!13!white] (2.857143,0.4285714) -- (2.571429,0.2142857) -- (2.857143,0.2142857) -- cycle;
\fill[teal!17!white] (2.571429,0.4285714) -- (2.571429,0.2142857) -- (2.857143,0.4285714) -- cycle;
\fill[teal!94!white] (1.142857,0.6428571) -- (0.8571429,0.8571429) -- (0.8571429,0.6428571) -- cycle;
\fill[teal!68!white] (0.8571429,0.8571429) -- (1.142857,0.6428571) -- (1.142857,0.8571429) -- cycle;
\fill[teal!94!white] (0.5714286,0.6428571) -- (0.8571429,0.8571429) -- (0.5714286,0.8571429) -- cycle;
\fill[teal!94!white] (0.8571429,0.8571429) -- (0.5714286,0.6428571) -- (0.8571429,0.6428571) -- cycle;
\fill[teal!29!white] (1.142857,1.5) -- (0.8571429,1.285714) -- (1.142857,1.285714) -- cycle;
\fill[teal!43!white] (0.8571429,1.5) -- (0.8571429,1.285714) -- (1.142857,1.5) -- cycle;
\fill[teal!25!white] (0.5714286,1.285714) -- (0.8571429,1.285714) -- (0.5714286,1.5) -- cycle;
\fill[teal!43!white] (0.8571429,1.285714) -- (0.8571429,1.5) -- (0.5714286,1.5) -- cycle;
\fill[teal!21!white] (1.142857,1.071429) -- (0.8571429,1.071429) -- (1.142857,0.8571429) -- cycle;
\fill[teal!43!white] (0.8571429,1.071429) -- (0.8571429,0.8571429) -- (1.142857,0.8571429) -- cycle;
\fill[teal!8!white] (0.8571429,1.071429) -- (1.142857,1.071429) -- (1.142857,1.285714) -- cycle;
\fill[teal!10!white] (0.8571429,1.285714) -- (0.8571429,1.071429) -- (1.142857,1.285714) -- cycle;
\fill[teal!34!white] (0.8571429,1.071429) -- (0.5714286,1.071429) -- (0.5714286,0.8571429) -- cycle;
\fill[teal!53!white] (0.8571429,0.8571429) -- (0.8571429,1.071429) -- (0.5714286,0.8571429) -- cycle;
\fill[teal!10!white] (0.5714286,1.071429) -- (0.8571429,1.071429) -- (0.5714286,1.285714) -- cycle;
\fill[teal!9!white] (0.8571429,1.071429) -- (0.8571429,1.285714) -- (0.5714286,1.285714) -- cycle;
\draw[teal,thick] (0,0) -- (4,0) -- (0,3) -- cycle;
\draw[magenta!70,thin] (0.9581951,1.050956) -- (3.6,2.35) -- (4.1,2.35);\filldraw[fill=magenta,draw=white,line width=.45pt] (0.9581951,1.050956) circle(1.5pt) node[above right,font=\scriptsize,text=magenta,fill=white,inner sep=.6pt] {};
\node[anchor=west,text=magenta,font=\small] at(4.12,2.35){Q2};\draw[gold!70,thin] (1,1) -- (3.6,1.8) -- (4.1,1.8);\filldraw[fill=gold,draw=white,line width=.45pt] (1,1) circle(1.5pt) node[above right,font=\scriptsize,text=gold,fill=white,inner sep=.6pt] {};
\node[anchor=west,text=gold,font=\small] at(4.12,1.8){$I$};\draw[navy!70,thin] (1.333333,1) -- (3.6,1.25) -- (4.1,1.25);\filldraw[fill=navy,draw=white,line width=.45pt] (1.333333,1) circle(1.5pt) node[above right,font=\scriptsize,text=navy,fill=white,inner sep=.6pt] {};
\node[anchor=west,text=navy,font=\small] at(4.12,1.25){$G$};\end{tikzpicture}
\caption[Probability density for Q2]{Probability density for Q2. The first moment averages all probability, including separate lobes. Numerical illustration.}\label{fig:state-Q2}
\end{figure}
\begin{figure}[htbp]\centering
\begin{tikzpicture}[x=1.7cm,y=1.7cm]
\fill[teal!10!white] (3.142857,0.6428571) -- (2.857143,0.8571429) -- (2.857143,0.6428571) -- cycle;
\fill[teal!11!white] (1.714286,1.714286) -- (1.428571,1.928571) -- (1.428571,1.714286) -- cycle;
\fill[teal!11!white] (1.714286,1.714286) -- (1.714286,1.5) -- (2,1.5) -- cycle;
\fill[teal!8!white] (0.5714286,2.571429) -- (0.2857143,2.785714) -- (0.2857143,2.571429) -- cycle;
\fill[teal!9!white] (2.571429,1.071429) -- (2.285714,1.285714) -- (2.285714,1.071429) -- cycle;
\fill[teal!8!white] (3.714286,0.2142857) -- (3.428571,0.4285714) -- (3.428571,0.2142857) -- cycle;
\fill[teal!8!white] (3.714286,0.2142857) -- (3.714286,0) -- (4,0) -- cycle;
\fill[teal!8!white] (0.8571429,2.142857) -- (1.142857,2.142857) -- (0.8571429,2.357143) -- cycle;
\fill[teal!8!white] (1.142857,2.142857) -- (1.142857,1.928571) -- (1.428571,1.928571) -- cycle;
\fill[teal!8!white] (0,2.785714) -- (0.2857143,2.785714) -- (0,3) -- cycle;
\fill[teal!8!white] (0.5714286,2.571429) -- (0.5714286,2.357143) -- (0.8571429,2.357143) -- cycle;
\fill[teal!8!white] (2,1.285714) -- (2.285714,1.285714) -- (2,1.5) -- cycle;
\fill[teal!11!white] (2.571429,1.071429) -- (2.571429,0.8571429) -- (2.857143,0.8571429) -- cycle;
\fill[teal!9!white] (3.142857,0.4285714) -- (3.428571,0.4285714) -- (3.142857,0.6428571) -- cycle;
\fill[teal!8!white] (0.2857143,0) -- (0,0.2142857) -- (0,0) -- cycle;
\fill[teal!10!white] (0,0.2142857) -- (0.2857143,0) -- (0.2857143,0.2142857) -- cycle;
\fill[teal!9!white] (0.5714286,0.2142857) -- (0.2857143,0) -- (0.5714286,0) -- cycle;
\fill[teal!12!white] (0.2857143,0) -- (0.5714286,0.2142857) -- (0.2857143,0.2142857) -- cycle;
\fill[teal!19!white] (1.714286,1.5) -- (1.714286,1.714286) -- (1.428571,1.714286) -- cycle;
\fill[teal!46!white] (1.428571,1.5) -- (1.714286,1.5) -- (1.428571,1.714286) -- cycle;
\fill[teal!10!white] (0,1.928571) -- (0,1.714286) -- (0.2857143,1.928571) -- cycle;
\fill[teal!17!white] (0,1.714286) -- (0.2857143,1.714286) -- (0.2857143,1.928571) -- cycle;
\fill[teal!8!white] (0.2857143,2.357143) -- (0,2.571429) -- (0,2.357143) -- cycle;
\fill[teal!9!white] (0,2.571429) -- (0.2857143,2.357143) -- (0.2857143,2.571429) -- cycle;
\fill[teal!9!white] (0,2.142857) -- (0.2857143,2.357143) -- (0,2.357143) -- cycle;
\fill[teal!11!white] (0.2857143,2.357143) -- (0,2.142857) -- (0.2857143,2.142857) -- cycle;
\fill[teal!10!white] (0,2.142857) -- (0,1.928571) -- (0.2857143,1.928571) -- cycle;
\fill[teal!13!white] (0.2857143,2.142857) -- (0,2.142857) -- (0.2857143,1.928571) -- cycle;
\fill[teal!28!white] (0.2857143,1.714286) -- (0.5714286,1.928571) -- (0.2857143,1.928571) -- cycle;
\fill[teal!31!white] (0.5714286,1.714286) -- (0.5714286,1.928571) -- (0.2857143,1.714286) -- cycle;
\fill[teal!11!white] (0.2857143,1.5) -- (0,1.714286) -- (0,1.5) -- cycle;
\fill[teal!19!white] (0,1.714286) -- (0.2857143,1.5) -- (0.2857143,1.714286) -- cycle;
\fill[teal!34!white] (0.2857143,1.5) -- (0.5714286,1.714286) -- (0.2857143,1.714286) -- cycle;
\fill[teal!36!white] (0.5714286,1.714286) -- (0.2857143,1.5) -- (0.5714286,1.5) -- cycle;
\fill[teal!12!white] (0,1.285714) -- (0.2857143,1.5) -- (0,1.5) -- cycle;
\fill[teal!23!white] (0.2857143,1.5) -- (0,1.285714) -- (0.2857143,1.285714) -- cycle;
\fill[teal!13!white] (0,1.285714) -- (0,1.071429) -- (0.2857143,1.071429) -- cycle;
\fill[teal!25!white] (0.2857143,1.285714) -- (0,1.285714) -- (0.2857143,1.071429) -- cycle;
\fill[teal!8!white] (3.428571,0) -- (3.714286,0) -- (3.428571,0.2142857) -- cycle;
\fill[teal!8!white] (3.714286,0) -- (3.714286,0.2142857) -- (3.428571,0.2142857) -- cycle;
\fill[teal!8!white] (1.428571,0) -- (1.714286,0) -- (1.428571,0.2142857) -- cycle;
\fill[teal!9!white] (1.714286,0) -- (1.714286,0.2142857) -- (1.428571,0.2142857) -- cycle;
\fill[teal!9!white] (0.8571429,0) -- (0.5714286,0.2142857) -- (0.5714286,0) -- cycle;
\fill[teal!12!white] (0.5714286,0.2142857) -- (0.8571429,0) -- (0.8571429,0.2142857) -- cycle;
\fill[teal!8!white] (1.142857,0.2142857) -- (1.142857,0) -- (1.428571,0) -- cycle;
\fill[teal!8!white] (1.142857,0.2142857) -- (1.428571,0) -- (1.428571,0.2142857) -- cycle;
\fill[teal!8!white] (1.142857,0.2142857) -- (0.8571429,0) -- (1.142857,0) -- cycle;
\fill[teal!9!white] (0.8571429,0) -- (1.142857,0.2142857) -- (0.8571429,0.2142857) -- cycle;
\fill[teal!19!white] (1.142857,1.928571) -- (1.142857,1.714286) -- (1.428571,1.714286) -- cycle;
\fill[teal!12!white] (1.428571,1.928571) -- (1.142857,1.928571) -- (1.428571,1.714286) -- cycle;
\fill[teal!8!white] (0.8571429,2.142857) -- (1.142857,1.928571) -- (1.142857,2.142857) -- cycle;
\fill[teal!8!white] (1.142857,1.928571) -- (0.8571429,2.142857) -- (0.8571429,1.928571) -- cycle;
\fill[teal!8!white] (0,2.785714) -- (0,2.571429) -- (0.2857143,2.571429) -- cycle;
\fill[teal!8!white] (0.2857143,2.785714) -- (0,2.785714) -- (0.2857143,2.571429) -- cycle;
\fill[teal!12!white] (1.142857,1.928571) -- (0.8571429,1.714286) -- (1.142857,1.714286) -- cycle;
\fill[teal!8!white] (0.8571429,1.714286) -- (1.142857,1.928571) -- (0.8571429,1.928571) -- cycle;
\fill[teal!20!white] (0.8571429,1.714286) -- (0.5714286,1.928571) -- (0.5714286,1.714286) -- cycle;
\fill[teal!13!white] (0.5714286,1.928571) -- (0.8571429,1.714286) -- (0.8571429,1.928571) -- cycle;
\fill[teal!10!white] (0.2857143,2.357143) -- (0.5714286,2.357143) -- (0.2857143,2.571429) -- cycle;
\fill[teal!9!white] (0.5714286,2.357143) -- (0.5714286,2.571429) -- (0.2857143,2.571429) -- cycle;
\fill[teal!12!white] (0,0.6428571) -- (0,0.4285714) -- (0.2857143,0.6428571) -- cycle;
\fill[teal!21!white] (0,0.4285714) -- (0.2857143,0.4285714) -- (0.2857143,0.6428571) -- cycle;
\fill[teal!10!white] (0,0.4285714) -- (0,0.2142857) -- (0.2857143,0.2142857) -- cycle;
\fill[teal!16!white] (0.2857143,0.4285714) -- (0,0.4285714) -- (0.2857143,0.2142857) -- cycle;
\fill[teal!8!white] (2.285714,1.285714) -- (2,1.285714) -- (2.285714,1.071429) -- cycle;
\fill[teal!8!white] (2,1.285714) -- (2,1.071429) -- (2.285714,1.071429) -- cycle;
\fill[teal!13!white] (2,1.071429) -- (2,0.8571429) -- (2.285714,1.071429) -- cycle;
\fill[teal!30!white] (2,0.8571429) -- (2.285714,0.8571429) -- (2.285714,1.071429) -- cycle;
\fill[teal!36!white] (2.285714,0.8571429) -- (2.571429,0.8571429) -- (2.285714,1.071429) -- cycle;
\fill[teal!17!white] (2.571429,0.8571429) -- (2.571429,1.071429) -- (2.285714,1.071429) -- cycle;
\fill[teal!10!white] (1.428571,0.4285714) -- (1.714286,0.2142857) -- (1.714286,0.4285714) -- cycle;
\fill[teal!8!white] (1.714286,0.2142857) -- (1.428571,0.4285714) -- (1.428571,0.2142857) -- cycle;
\fill[teal!8!white] (3.142857,0.2142857) -- (3.428571,0) -- (3.428571,0.2142857) -- cycle;
\fill[teal!8!white] (3.142857,0.2142857) -- (3.142857,0) -- (3.428571,0) -- cycle;
\fill[teal!9!white] (3.142857,0) -- (3.142857,0.2142857) -- (2.857143,0) -- cycle;
\fill[teal!14!white] (3.142857,0.2142857) -- (2.857143,0.2142857) -- (2.857143,0) -- cycle;
\fill[teal!12!white] (1.714286,0) -- (2,0.2142857) -- (1.714286,0.2142857) -- cycle;
\fill[teal!10!white] (2,0.2142857) -- (1.714286,0) -- (2,0) -- cycle;
\fill[teal!26!white] (0.5714286,0.2142857) -- (0.5714286,0.4285714) -- (0.2857143,0.2142857) -- cycle;
\fill[teal!30!white] (0.5714286,0.4285714) -- (0.2857143,0.4285714) -- (0.2857143,0.2142857) -- cycle;
\fill[teal!20!white] (0.8571429,0.4285714) -- (0.5714286,0.2142857) -- (0.8571429,0.2142857) -- cycle;
\fill[teal!30!white] (0.8571429,0.4285714) -- (0.5714286,0.4285714) -- (0.5714286,0.2142857) -- cycle;
\fill[teal!60!white] (1.142857,1.5) -- (1.428571,1.5) -- (1.428571,1.714286) -- cycle;
\fill[teal!40!white] (1.142857,1.714286) -- (1.142857,1.5) -- (1.428571,1.714286) -- cycle;
\fill[teal!20!white] (0.5714286,2.142857) -- (0.2857143,2.142857) -- (0.2857143,1.928571) -- cycle;
\fill[teal!24!white] (0.5714286,1.928571) -- (0.5714286,2.142857) -- (0.2857143,1.928571) -- cycle;
\fill[teal!15!white] (0.5714286,2.142857) -- (0.2857143,2.357143) -- (0.2857143,2.142857) -- cycle;
\fill[teal!13!white] (0.5714286,2.142857) -- (0.5714286,2.357143) -- (0.2857143,2.357143) -- cycle;
\fill[teal!12!white] (0.8571429,2.142857) -- (0.5714286,2.142857) -- (0.8571429,1.928571) -- cycle;
\fill[teal!18!white] (0.5714286,2.142857) -- (0.5714286,1.928571) -- (0.8571429,1.928571) -- cycle;
\fill[teal!10!white] (0.5714286,2.142857) -- (0.8571429,2.142857) -- (0.8571429,2.357143) -- cycle;
\fill[teal!11!white] (0.5714286,2.357143) -- (0.5714286,2.142857) -- (0.8571429,2.357143) -- cycle;
\fill[teal!18!white] (1.714286,1.285714) -- (2,1.285714) -- (2,1.5) -- cycle;
\fill[teal!25!white] (1.714286,1.5) -- (1.714286,1.285714) -- (2,1.5) -- cycle;
\fill[teal!91!white] (1.142857,1.071429) -- (1.428571,0.8571429) -- (1.428571,1.071429) -- cycle;
\fill[teal!62!white] (1.428571,0.8571429) -- (1.142857,1.071429) -- (1.142857,0.8571429) -- cycle;
\fill[teal!10!white] (3.142857,0.4285714) -- (3.142857,0.2142857) -- (3.428571,0.2142857) -- cycle;
\fill[teal!9!white] (3.428571,0.4285714) -- (3.142857,0.4285714) -- (3.428571,0.2142857) -- cycle;
\fill[teal!27!white] (2.857143,0.4285714) -- (3.142857,0.4285714) -- (2.857143,0.6428571) -- cycle;
\fill[teal!13!white] (3.142857,0.4285714) -- (3.142857,0.6428571) -- (2.857143,0.6428571) -- cycle;
\fill[teal!18!white] (3.142857,0.2142857) -- (3.142857,0.4285714) -- (2.857143,0.2142857) -- cycle;
\fill[teal!29!white] (3.142857,0.4285714) -- (2.857143,0.4285714) -- (2.857143,0.2142857) -- cycle;
\fill[teal!35!white] (2.571429,0.8571429) -- (2.571429,0.6428571) -- (2.857143,0.8571429) -- cycle;
\fill[teal!33!white] (2.857143,0.8571429) -- (2.571429,0.6428571) -- (2.857143,0.6428571) -- cycle;
\fill[teal!13!white] (2.285714,0) -- (2.285714,0.2142857) -- (2,0) -- cycle;
\fill[teal!23!white] (2.285714,0.2142857) -- (2,0.2142857) -- (2,0) -- cycle;
\fill[teal!55!white] (2.285714,0.2142857) -- (2,0.4285714) -- (2,0.2142857) -- cycle;
\fill[teal!85!white] (2,0.4285714) -- (2.285714,0.2142857) -- (2.285714,0.4285714) -- cycle;
\fill[teal!20!white] (2,0.4285714) -- (1.714286,0.6428571) -- (1.714286,0.4285714) -- cycle;
\fill[teal!32!white] (1.714286,0.6428571) -- (2,0.4285714) -- (2,0.6428571) -- cycle;
\fill[teal!22!white] (1.714286,0.2142857) -- (2,0.4285714) -- (1.714286,0.4285714) -- cycle;
\fill[teal!30!white] (2,0.2142857) -- (2,0.4285714) -- (1.714286,0.2142857) -- cycle;
\fill[teal!13!white] (0.2857143,0.8571429) -- (0,0.8571429) -- (0,0.6428571) -- cycle;
\fill[teal!26!white] (0.2857143,0.8571429) -- (0,0.6428571) -- (0.2857143,0.6428571) -- cycle;
\fill[teal!27!white] (0,1.071429) -- (0.2857143,0.8571429) -- (0.2857143,1.071429) -- cycle;
\fill[teal!14!white] (0,0.8571429) -- (0.2857143,0.8571429) -- (0,1.071429) -- cycle;
\fill[teal!40!white] (0.2857143,1.5) -- (0.5714286,1.285714) -- (0.5714286,1.5) -- cycle;
\fill[teal!42!white] (0.5714286,1.285714) -- (0.2857143,1.5) -- (0.2857143,1.285714) -- cycle;
\fill[teal!56!white] (0.5714286,0.6428571) -- (0.2857143,0.8571429) -- (0.2857143,0.6428571) -- cycle;
\fill[teal!62!white] (0.2857143,0.8571429) -- (0.5714286,0.6428571) -- (0.5714286,0.8571429) -- cycle;
\fill[teal!49!white] (0.2857143,0.4285714) -- (0.5714286,0.6428571) -- (0.2857143,0.6428571) -- cycle;
\fill[teal!49!white] (0.5714286,0.4285714) -- (0.5714286,0.6428571) -- (0.2857143,0.4285714) -- cycle;
\fill[teal!41!white] (0.8571429,0.4285714) -- (0.5714286,0.6428571) -- (0.5714286,0.4285714) -- cycle;
\fill[teal!31!white] (0.5714286,0.6428571) -- (0.8571429,0.4285714) -- (0.8571429,0.6428571) -- cycle;
\fill[teal!8!white] (1.142857,0.4285714) -- (1.142857,0.2142857) -- (1.428571,0.2142857) -- cycle;
\fill[teal!8!white] (1.428571,0.4285714) -- (1.142857,0.4285714) -- (1.428571,0.2142857) -- cycle;
\fill[teal!9!white] (1.142857,0.2142857) -- (1.142857,0.4285714) -- (0.8571429,0.2142857) -- cycle;
\fill[teal!13!white] (1.142857,0.4285714) -- (0.8571429,0.4285714) -- (0.8571429,0.2142857) -- cycle;
\fill[teal!12!white] (0.8571429,0.4285714) -- (1.142857,0.4285714) -- (0.8571429,0.6428571) -- cycle;
\fill[teal!8!white] (1.142857,0.4285714) -- (1.142857,0.6428571) -- (0.8571429,0.6428571) -- cycle;
\fill[teal!82!white] (1.428571,1.285714) -- (1.714286,1.5) -- (1.428571,1.5) -- cycle;
\fill[teal!76!white] (1.428571,1.285714) -- (1.714286,1.285714) -- (1.714286,1.5) -- cycle;
\fill[teal!94!white] (1.142857,1.5) -- (1.428571,1.285714) -- (1.428571,1.5) -- cycle;
\fill[teal!92!white] (1.428571,1.285714) -- (1.142857,1.5) -- (1.142857,1.285714) -- cycle;
\fill[teal!94!white] (1.428571,1.285714) -- (1.142857,1.071429) -- (1.428571,1.071429) -- cycle;
\fill[teal!94!white] (1.142857,1.071429) -- (1.428571,1.285714) -- (1.142857,1.285714) -- cycle;
\fill[teal!9!white] (1.428571,0.6428571) -- (1.428571,0.4285714) -- (1.714286,0.4285714) -- cycle;
\fill[teal!9!white] (1.714286,0.6428571) -- (1.428571,0.6428571) -- (1.714286,0.4285714) -- cycle;
\fill[teal!11!white] (1.428571,0.6428571) -- (1.142857,0.4285714) -- (1.428571,0.4285714) -- cycle;
\fill[teal!13!white] (1.142857,0.4285714) -- (1.428571,0.6428571) -- (1.142857,0.6428571) -- cycle;
\fill[teal!43!white] (1.428571,0.6428571) -- (1.428571,0.8571429) -- (1.142857,0.8571429) -- cycle;
\fill[teal!25!white] (1.142857,0.6428571) -- (1.428571,0.6428571) -- (1.142857,0.8571429) -- cycle;
\fill[teal!18!white] (2,1.285714) -- (1.714286,1.071429) -- (2,1.071429) -- cycle;
\fill[teal!37!white] (1.714286,1.285714) -- (1.714286,1.071429) -- (2,1.285714) -- cycle;
\fill[teal!88!white] (1.428571,1.285714) -- (1.714286,1.071429) -- (1.714286,1.285714) -- cycle;
\fill[teal!94!white] (1.714286,1.071429) -- (1.428571,1.285714) -- (1.428571,1.071429) -- cycle;
\fill[teal!75!white] (2.285714,0.6428571) -- (2.571429,0.8571429) -- (2.285714,0.8571429) -- cycle;
\fill[teal!89!white] (2.285714,0.6428571) -- (2.571429,0.6428571) -- (2.571429,0.8571429) -- cycle;
\fill[teal!62!white] (2,0.8571429) -- (2.285714,0.6428571) -- (2.285714,0.8571429) -- cycle;
\fill[teal!59!white] (2.285714,0.6428571) -- (2,0.8571429) -- (2,0.6428571) -- cycle;
\fill[teal!94!white] (2.285714,0.6428571) -- (2,0.4285714) -- (2.285714,0.4285714) -- cycle;
\fill[teal!81!white] (2,0.4285714) -- (2.285714,0.6428571) -- (2,0.6428571) -- cycle;
\fill[teal!63!white] (2.571429,0.4285714) -- (2.857143,0.4285714) -- (2.857143,0.6428571) -- cycle;
\fill[teal!80!white] (2.571429,0.6428571) -- (2.571429,0.4285714) -- (2.857143,0.6428571) -- cycle;
\fill[teal!94!white] (2.285714,0.6428571) -- (2.571429,0.4285714) -- (2.571429,0.6428571) -- cycle;
\fill[teal!94!white] (2.571429,0.4285714) -- (2.285714,0.6428571) -- (2.285714,0.4285714) -- cycle;
\fill[teal!55!white] (0.2857143,0.8571429) -- (0.5714286,1.071429) -- (0.2857143,1.071429) -- cycle;
\fill[teal!59!white] (0.5714286,1.071429) -- (0.2857143,0.8571429) -- (0.5714286,0.8571429) -- cycle;
\fill[teal!50!white] (0.5714286,1.071429) -- (0.2857143,1.285714) -- (0.2857143,1.071429) -- cycle;
\fill[teal!47!white] (0.5714286,1.071429) -- (0.5714286,1.285714) -- (0.2857143,1.285714) -- cycle;
\fill[teal!20!white] (0.8571429,1.5) -- (0.5714286,1.714286) -- (0.5714286,1.5) -- cycle;
\fill[teal!12!white] (0.8571429,1.5) -- (0.8571429,1.714286) -- (0.5714286,1.714286) -- cycle;
\fill[teal!11!white] (0.8571429,1.714286) -- (0.8571429,1.5) -- (1.142857,1.714286) -- cycle;
\fill[teal!26!white] (0.8571429,1.5) -- (1.142857,1.5) -- (1.142857,1.714286) -- cycle;
\fill[teal!9!white] (1.714286,0.8571429) -- (1.714286,0.6428571) -- (2,0.6428571) -- cycle;
\fill[teal!13!white] (2,0.8571429) -- (1.714286,0.8571429) -- (2,0.6428571) -- cycle;
\fill[teal!8!white] (1.714286,0.8571429) -- (2,0.8571429) -- (2,1.071429) -- cycle;
\fill[teal!21!white] (1.714286,1.071429) -- (1.714286,0.8571429) -- (2,1.071429) -- cycle;
\fill[teal!16!white] (1.714286,0.8571429) -- (1.428571,0.6428571) -- (1.714286,0.6428571) -- cycle;
\fill[teal!36!white] (1.428571,0.6428571) -- (1.714286,0.8571429) -- (1.428571,0.8571429) -- cycle;
\fill[teal!68!white] (1.428571,0.8571429) -- (1.714286,0.8571429) -- (1.428571,1.071429) -- cycle;
\fill[teal!64!white] (1.714286,0.8571429) -- (1.714286,1.071429) -- (1.428571,1.071429) -- cycle;
\fill[teal!14!white] (2.571429,0) -- (2.571429,0.2142857) -- (2.285714,0) -- cycle;
\fill[teal!29!white] (2.571429,0.2142857) -- (2.285714,0.2142857) -- (2.285714,0) -- cycle;
\fill[teal!77!white] (2.285714,0.2142857) -- (2.571429,0.2142857) -- (2.285714,0.4285714) -- cycle;
\fill[teal!94!white] (2.571429,0.2142857) -- (2.571429,0.4285714) -- (2.285714,0.4285714) -- cycle;
\fill[teal!20!white] (2.857143,0.2142857) -- (2.571429,0.2142857) -- (2.857143,0) -- cycle;
\fill[teal!12!white] (2.571429,0.2142857) -- (2.571429,0) -- (2.857143,0) -- cycle;
\fill[teal!46!white] (2.857143,0.4285714) -- (2.571429,0.2142857) -- (2.857143,0.2142857) -- cycle;
\fill[teal!74!white] (2.571429,0.4285714) -- (2.571429,0.2142857) -- (2.857143,0.4285714) -- cycle;
\fill[teal!9!white] (1.142857,0.6428571) -- (0.8571429,0.8571429) -- (0.8571429,0.6428571) -- cycle;
\fill[teal!12!white] (0.8571429,0.8571429) -- (1.142857,0.6428571) -- (1.142857,0.8571429) -- cycle;
\fill[teal!43!white] (0.5714286,0.6428571) -- (0.8571429,0.8571429) -- (0.5714286,0.8571429) -- cycle;
\fill[teal!29!white] (0.8571429,0.8571429) -- (0.5714286,0.6428571) -- (0.8571429,0.6428571) -- cycle;
\fill[teal!40!white] (1.142857,1.5) -- (0.8571429,1.285714) -- (1.142857,1.285714) -- cycle;
\fill[teal!18!white] (0.8571429,1.5) -- (0.8571429,1.285714) -- (1.142857,1.5) -- cycle;
\fill[teal!21!white] (0.5714286,1.285714) -- (0.8571429,1.285714) -- (0.5714286,1.5) -- cycle;
\fill[teal!11!white] (0.8571429,1.285714) -- (0.8571429,1.5) -- (0.5714286,1.5) -- cycle;
\fill[teal!28!white] (1.142857,1.071429) -- (0.8571429,1.071429) -- (1.142857,0.8571429) -- cycle;
\fill[teal!10!white] (0.8571429,1.071429) -- (0.8571429,0.8571429) -- (1.142857,0.8571429) -- cycle;
\fill[teal!38!white] (0.8571429,1.071429) -- (1.142857,1.071429) -- (1.142857,1.285714) -- cycle;
\fill[teal!19!white] (0.8571429,1.285714) -- (0.8571429,1.071429) -- (1.142857,1.285714) -- cycle;
\fill[teal!33!white] (0.8571429,1.071429) -- (0.5714286,1.071429) -- (0.5714286,0.8571429) -- cycle;
\fill[teal!20!white] (0.8571429,0.8571429) -- (0.8571429,1.071429) -- (0.5714286,0.8571429) -- cycle;
\fill[teal!27!white] (0.5714286,1.071429) -- (0.8571429,1.071429) -- (0.5714286,1.285714) -- cycle;
\fill[teal!12!white] (0.8571429,1.071429) -- (0.8571429,1.285714) -- (0.5714286,1.285714) -- cycle;
\draw[teal,thick] (0,0) -- (4,0) -- (0,3) -- cycle;
\draw[magenta!70,thin] (1.52641,0.9056026) -- (3.6,2.35) -- (4.1,2.35);\filldraw[fill=magenta,draw=white,line width=.45pt] (1.52641,0.9056026) circle(1.5pt) node[above right,font=\scriptsize,text=magenta,fill=white,inner sep=.6pt] {};
\node[anchor=west,text=magenta,font=\small] at(4.12,2.35){Q3};\draw[gold!70,thin] (1,1) -- (3.6,1.8) -- (4.1,1.8);\filldraw[fill=gold,draw=white,line width=.45pt] (1,1) circle(1.5pt) node[above right,font=\scriptsize,text=gold,fill=white,inner sep=.6pt] {};
\node[anchor=west,text=gold,font=\small] at(4.12,1.8){$I$};\draw[navy!70,thin] (1.333333,1) -- (3.6,1.25) -- (4.1,1.25);\filldraw[fill=navy,draw=white,line width=.45pt] (1.333333,1) circle(1.5pt) node[above right,font=\scriptsize,text=navy,fill=white,inner sep=.6pt] {};
\node[anchor=west,text=navy,font=\small] at(4.12,1.25){$G$};\end{tikzpicture}
\caption[Probability density for Q3]{Probability density for Q3. The first moment averages all probability, including separate lobes. Numerical illustration.}\label{fig:state-Q3}
\end{figure}
\begin{figure}[htbp]\centering
\begin{tikzpicture}[x=1.7cm,y=1.7cm]
\fill[teal!10!white] (3.142857,0.6428571) -- (2.857143,0.8571429) -- (2.857143,0.6428571) -- cycle;
\fill[teal!9!white] (1.714286,1.714286) -- (1.428571,1.928571) -- (1.428571,1.714286) -- cycle;
\fill[teal!9!white] (1.714286,1.714286) -- (1.714286,1.5) -- (2,1.5) -- cycle;
\fill[teal!9!white] (0.5714286,2.571429) -- (0.2857143,2.785714) -- (0.2857143,2.571429) -- cycle;
\fill[teal!9!white] (2.571429,1.071429) -- (2.285714,1.285714) -- (2.285714,1.071429) -- cycle;
\fill[teal!8!white] (3.714286,0.2142857) -- (3.428571,0.4285714) -- (3.428571,0.2142857) -- cycle;
\fill[teal!8!white] (3.714286,0.2142857) -- (3.714286,0) -- (4,0) -- cycle;
\fill[teal!11!white] (0.8571429,2.142857) -- (1.142857,2.142857) -- (0.8571429,2.357143) -- cycle;
\fill[teal!8!white] (1.142857,2.142857) -- (1.142857,1.928571) -- (1.428571,1.928571) -- cycle;
\fill[teal!8!white] (0,2.785714) -- (0.2857143,2.785714) -- (0,3) -- cycle;
\fill[teal!12!white] (0.5714286,2.571429) -- (0.5714286,2.357143) -- (0.8571429,2.357143) -- cycle;
\fill[teal!9!white] (2,1.285714) -- (2.285714,1.285714) -- (2,1.5) -- cycle;
\fill[teal!10!white] (2.571429,1.071429) -- (2.571429,0.8571429) -- (2.857143,0.8571429) -- cycle;
\fill[teal!9!white] (3.142857,0.4285714) -- (3.428571,0.4285714) -- (3.142857,0.6428571) -- cycle;
\fill[teal!8!white] (0.2857143,0) -- (0,0.2142857) -- (0,0) -- cycle;
\fill[teal!8!white] (0,0.2142857) -- (0.2857143,0) -- (0.2857143,0.2142857) -- cycle;
\fill[teal!8!white] (0.5714286,0.2142857) -- (0.2857143,0) -- (0.5714286,0) -- cycle;
\fill[teal!8!white] (0.2857143,0) -- (0.5714286,0.2142857) -- (0.2857143,0.2142857) -- cycle;
\fill[teal!13!white] (1.714286,1.5) -- (1.714286,1.714286) -- (1.428571,1.714286) -- cycle;
\fill[teal!27!white] (1.428571,1.5) -- (1.714286,1.5) -- (1.428571,1.714286) -- cycle;
\fill[teal!18!white] (0,1.928571) -- (0,1.714286) -- (0.2857143,1.928571) -- cycle;
\fill[teal!38!white] (0,1.714286) -- (0.2857143,1.714286) -- (0.2857143,1.928571) -- cycle;
\fill[teal!11!white] (0.2857143,2.357143) -- (0,2.571429) -- (0,2.357143) -- cycle;
\fill[teal!15!white] (0,2.571429) -- (0.2857143,2.357143) -- (0.2857143,2.571429) -- cycle;
\fill[teal!15!white] (0,2.142857) -- (0.2857143,2.357143) -- (0,2.357143) -- cycle;
\fill[teal!38!white] (0.2857143,2.357143) -- (0,2.142857) -- (0.2857143,2.142857) -- cycle;
\fill[teal!19!white] (0,2.142857) -- (0,1.928571) -- (0.2857143,1.928571) -- cycle;
\fill[teal!47!white] (0.2857143,2.142857) -- (0,2.142857) -- (0.2857143,1.928571) -- cycle;
\fill[teal!94!white] (0.2857143,1.714286) -- (0.5714286,1.928571) -- (0.2857143,1.928571) -- cycle;
\fill[teal!94!white] (0.5714286,1.714286) -- (0.5714286,1.928571) -- (0.2857143,1.714286) -- cycle;
\fill[teal!11!white] (0.2857143,1.5) -- (0,1.714286) -- (0,1.5) -- cycle;
\fill[teal!24!white] (0,1.714286) -- (0.2857143,1.5) -- (0.2857143,1.714286) -- cycle;
\fill[teal!46!white] (0.2857143,1.5) -- (0.5714286,1.714286) -- (0.2857143,1.714286) -- cycle;
\fill[teal!29!white] (0.5714286,1.714286) -- (0.2857143,1.5) -- (0.5714286,1.5) -- cycle;
\fill[teal!8!white] (0,1.285714) -- (0.2857143,1.5) -- (0,1.5) -- cycle;
\fill[teal!8!white] (0.2857143,1.5) -- (0,1.285714) -- (0.2857143,1.285714) -- cycle;
\fill[teal!9!white] (0,1.285714) -- (0,1.071429) -- (0.2857143,1.071429) -- cycle;
\fill[teal!10!white] (0.2857143,1.285714) -- (0,1.285714) -- (0.2857143,1.071429) -- cycle;
\fill[teal!8!white] (3.428571,0) -- (3.714286,0) -- (3.428571,0.2142857) -- cycle;
\fill[teal!8!white] (3.714286,0) -- (3.714286,0.2142857) -- (3.428571,0.2142857) -- cycle;
\fill[teal!14!white] (1.428571,0) -- (1.714286,0) -- (1.428571,0.2142857) -- cycle;
\fill[teal!27!white] (1.714286,0) -- (1.714286,0.2142857) -- (1.428571,0.2142857) -- cycle;
\fill[teal!8!white] (0.8571429,0) -- (0.5714286,0.2142857) -- (0.5714286,0) -- cycle;
\fill[teal!10!white] (0.5714286,0.2142857) -- (0.8571429,0) -- (0.8571429,0.2142857) -- cycle;
\fill[teal!14!white] (1.142857,0.2142857) -- (1.142857,0) -- (1.428571,0) -- cycle;
\fill[teal!33!white] (1.142857,0.2142857) -- (1.428571,0) -- (1.428571,0.2142857) -- cycle;
\fill[teal!12!white] (1.142857,0.2142857) -- (0.8571429,0) -- (1.142857,0) -- cycle;
\fill[teal!17!white] (0.8571429,0) -- (1.142857,0.2142857) -- (0.8571429,0.2142857) -- cycle;
\fill[teal!11!white] (1.142857,1.928571) -- (1.142857,1.714286) -- (1.428571,1.714286) -- cycle;
\fill[teal!9!white] (1.428571,1.928571) -- (1.142857,1.928571) -- (1.428571,1.714286) -- cycle;
\fill[teal!12!white] (0.8571429,2.142857) -- (1.142857,1.928571) -- (1.142857,2.142857) -- cycle;
\fill[teal!24!white] (1.142857,1.928571) -- (0.8571429,2.142857) -- (0.8571429,1.928571) -- cycle;
\fill[teal!9!white] (0,2.785714) -- (0,2.571429) -- (0.2857143,2.571429) -- cycle;
\fill[teal!8!white] (0.2857143,2.785714) -- (0,2.785714) -- (0.2857143,2.571429) -- cycle;
\fill[teal!8!white] (1.142857,1.928571) -- (0.8571429,1.714286) -- (1.142857,1.714286) -- cycle;
\fill[teal!19!white] (0.8571429,1.714286) -- (1.142857,1.928571) -- (0.8571429,1.928571) -- cycle;
\fill[teal!70!white] (0.8571429,1.714286) -- (0.5714286,1.928571) -- (0.5714286,1.714286) -- cycle;
\fill[teal!62!white] (0.5714286,1.928571) -- (0.8571429,1.714286) -- (0.8571429,1.928571) -- cycle;
\fill[teal!31!white] (0.2857143,2.357143) -- (0.5714286,2.357143) -- (0.2857143,2.571429) -- cycle;
\fill[teal!16!white] (0.5714286,2.357143) -- (0.5714286,2.571429) -- (0.2857143,2.571429) -- cycle;
\fill[teal!9!white] (0,0.6428571) -- (0,0.4285714) -- (0.2857143,0.6428571) -- cycle;
\fill[teal!10!white] (0,0.4285714) -- (0.2857143,0.4285714) -- (0.2857143,0.6428571) -- cycle;
\fill[teal!8!white] (0,0.4285714) -- (0,0.2142857) -- (0.2857143,0.2142857) -- cycle;
\fill[teal!9!white] (0.2857143,0.4285714) -- (0,0.4285714) -- (0.2857143,0.2142857) -- cycle;
\fill[teal!10!white] (2.285714,1.285714) -- (2,1.285714) -- (2.285714,1.071429) -- cycle;
\fill[teal!13!white] (2,1.285714) -- (2,1.071429) -- (2.285714,1.071429) -- cycle;
\fill[teal!15!white] (2,1.071429) -- (2,0.8571429) -- (2.285714,1.071429) -- cycle;
\fill[teal!19!white] (2,0.8571429) -- (2.285714,0.8571429) -- (2.285714,1.071429) -- cycle;
\fill[teal!22!white] (2.285714,0.8571429) -- (2.571429,0.8571429) -- (2.285714,1.071429) -- cycle;
\fill[teal!13!white] (2.571429,0.8571429) -- (2.571429,1.071429) -- (2.285714,1.071429) -- cycle;
\fill[teal!83!white] (1.428571,0.4285714) -- (1.714286,0.2142857) -- (1.714286,0.4285714) -- cycle;
\fill[teal!79!white] (1.714286,0.2142857) -- (1.428571,0.4285714) -- (1.428571,0.2142857) -- cycle;
\fill[teal!8!white] (3.142857,0.2142857) -- (3.428571,0) -- (3.428571,0.2142857) -- cycle;
\fill[teal!8!white] (3.142857,0.2142857) -- (3.142857,0) -- (3.428571,0) -- cycle;
\fill[teal!9!white] (3.142857,0) -- (3.142857,0.2142857) -- (2.857143,0) -- cycle;
\fill[teal!13!white] (3.142857,0.2142857) -- (2.857143,0.2142857) -- (2.857143,0) -- cycle;
\fill[teal!16!white] (1.714286,0) -- (2,0.2142857) -- (1.714286,0.2142857) -- cycle;
\fill[teal!9!white] (2,0.2142857) -- (1.714286,0) -- (2,0) -- cycle;
\fill[teal!8!white] (0.5714286,0.2142857) -- (0.5714286,0.4285714) -- (0.2857143,0.2142857) -- cycle;
\fill[teal!9!white] (0.5714286,0.4285714) -- (0.2857143,0.4285714) -- (0.2857143,0.2142857) -- cycle;
\fill[teal!14!white] (0.8571429,0.4285714) -- (0.5714286,0.2142857) -- (0.8571429,0.2142857) -- cycle;
\fill[teal!10!white] (0.8571429,0.4285714) -- (0.5714286,0.4285714) -- (0.5714286,0.2142857) -- cycle;
\fill[teal!36!white] (1.142857,1.5) -- (1.428571,1.5) -- (1.428571,1.714286) -- cycle;
\fill[teal!24!white] (1.142857,1.714286) -- (1.142857,1.5) -- (1.428571,1.714286) -- cycle;
\fill[teal!94!white] (0.5714286,2.142857) -- (0.2857143,2.142857) -- (0.2857143,1.928571) -- cycle;
\fill[teal!94!white] (0.5714286,1.928571) -- (0.5714286,2.142857) -- (0.2857143,1.928571) -- cycle;
\fill[teal!84!white] (0.5714286,2.142857) -- (0.2857143,2.357143) -- (0.2857143,2.142857) -- cycle;
\fill[teal!67!white] (0.5714286,2.142857) -- (0.5714286,2.357143) -- (0.2857143,2.357143) -- cycle;
\fill[teal!67!white] (0.8571429,2.142857) -- (0.5714286,2.142857) -- (0.8571429,1.928571) -- cycle;
\fill[teal!94!white] (0.5714286,2.142857) -- (0.5714286,1.928571) -- (0.8571429,1.928571) -- cycle;
\fill[teal!39!white] (0.5714286,2.142857) -- (0.8571429,2.142857) -- (0.8571429,2.357143) -- cycle;
\fill[teal!40!white] (0.5714286,2.357143) -- (0.5714286,2.142857) -- (0.8571429,2.357143) -- cycle;
\fill[teal!13!white] (1.714286,1.285714) -- (2,1.285714) -- (2,1.5) -- cycle;
\fill[teal!15!white] (1.714286,1.5) -- (1.714286,1.285714) -- (2,1.5) -- cycle;
\fill[teal!10!white] (1.142857,1.071429) -- (1.428571,0.8571429) -- (1.428571,1.071429) -- cycle;
\fill[teal!8!white] (1.428571,0.8571429) -- (1.142857,1.071429) -- (1.142857,0.8571429) -- cycle;
\fill[teal!10!white] (3.142857,0.4285714) -- (3.142857,0.2142857) -- (3.428571,0.2142857) -- cycle;
\fill[teal!9!white] (3.428571,0.4285714) -- (3.142857,0.4285714) -- (3.428571,0.2142857) -- cycle;
\fill[teal!23!white] (2.857143,0.4285714) -- (3.142857,0.4285714) -- (2.857143,0.6428571) -- cycle;
\fill[teal!12!white] (3.142857,0.4285714) -- (3.142857,0.6428571) -- (2.857143,0.6428571) -- cycle;
\fill[teal!17!white] (3.142857,0.2142857) -- (3.142857,0.4285714) -- (2.857143,0.2142857) -- cycle;
\fill[teal!25!white] (3.142857,0.4285714) -- (2.857143,0.4285714) -- (2.857143,0.2142857) -- cycle;
\fill[teal!24!white] (2.571429,0.8571429) -- (2.571429,0.6428571) -- (2.857143,0.8571429) -- cycle;
\fill[teal!24!white] (2.857143,0.8571429) -- (2.571429,0.6428571) -- (2.857143,0.6428571) -- cycle;
\fill[teal!8!white] (2.285714,0) -- (2.285714,0.2142857) -- (2,0) -- cycle;
\fill[teal!8!white] (2.285714,0.2142857) -- (2,0.2142857) -- (2,0) -- cycle;
\fill[teal!8!white] (2.285714,0.2142857) -- (2,0.4285714) -- (2,0.2142857) -- cycle;
\fill[teal!14!white] (2,0.4285714) -- (2.285714,0.2142857) -- (2.285714,0.4285714) -- cycle;
\fill[teal!42!white] (2,0.4285714) -- (1.714286,0.6428571) -- (1.714286,0.4285714) -- cycle;
\fill[teal!19!white] (1.714286,0.6428571) -- (2,0.4285714) -- (2,0.6428571) -- cycle;
\fill[teal!40!white] (1.714286,0.2142857) -- (2,0.4285714) -- (1.714286,0.4285714) -- cycle;
\fill[teal!19!white] (2,0.2142857) -- (2,0.4285714) -- (1.714286,0.2142857) -- cycle;
\fill[teal!10!white] (0.2857143,0.8571429) -- (0,0.8571429) -- (0,0.6428571) -- cycle;
\fill[teal!14!white] (0.2857143,0.8571429) -- (0,0.6428571) -- (0.2857143,0.6428571) -- cycle;
\fill[teal!15!white] (0,1.071429) -- (0.2857143,0.8571429) -- (0.2857143,1.071429) -- cycle;
\fill[teal!10!white] (0,0.8571429) -- (0.2857143,0.8571429) -- (0,1.071429) -- cycle;
\fill[teal!9!white] (0.2857143,1.5) -- (0.5714286,1.285714) -- (0.5714286,1.5) -- cycle;
\fill[teal!8!white] (0.5714286,1.285714) -- (0.2857143,1.5) -- (0.2857143,1.285714) -- cycle;
\fill[teal!22!white] (0.5714286,0.6428571) -- (0.2857143,0.8571429) -- (0.2857143,0.6428571) -- cycle;
\fill[teal!28!white] (0.2857143,0.8571429) -- (0.5714286,0.6428571) -- (0.5714286,0.8571429) -- cycle;
\fill[teal!14!white] (0.2857143,0.4285714) -- (0.5714286,0.6428571) -- (0.2857143,0.6428571) -- cycle;
\fill[teal!11!white] (0.5714286,0.4285714) -- (0.5714286,0.6428571) -- (0.2857143,0.4285714) -- cycle;
\fill[teal!8!white] (0.8571429,0.4285714) -- (0.5714286,0.6428571) -- (0.5714286,0.4285714) -- cycle;
\fill[teal!9!white] (0.5714286,0.6428571) -- (0.8571429,0.4285714) -- (0.8571429,0.6428571) -- cycle;
\fill[teal!76!white] (1.142857,0.4285714) -- (1.142857,0.2142857) -- (1.428571,0.2142857) -- cycle;
\fill[teal!94!white] (1.428571,0.4285714) -- (1.142857,0.4285714) -- (1.428571,0.2142857) -- cycle;
\fill[teal!48!white] (1.142857,0.2142857) -- (1.142857,0.4285714) -- (0.8571429,0.2142857) -- cycle;
\fill[teal!39!white] (1.142857,0.4285714) -- (0.8571429,0.4285714) -- (0.8571429,0.2142857) -- cycle;
\fill[teal!34!white] (0.8571429,0.4285714) -- (1.142857,0.4285714) -- (0.8571429,0.6428571) -- cycle;
\fill[teal!45!white] (1.142857,0.4285714) -- (1.142857,0.6428571) -- (0.8571429,0.6428571) -- cycle;
\fill[teal!40!white] (1.428571,1.285714) -- (1.714286,1.5) -- (1.428571,1.5) -- cycle;
\fill[teal!32!white] (1.428571,1.285714) -- (1.714286,1.285714) -- (1.714286,1.5) -- cycle;
\fill[teal!53!white] (1.142857,1.5) -- (1.428571,1.285714) -- (1.428571,1.5) -- cycle;
\fill[teal!58!white] (1.428571,1.285714) -- (1.142857,1.5) -- (1.142857,1.285714) -- cycle;
\fill[teal!31!white] (1.428571,1.285714) -- (1.142857,1.071429) -- (1.428571,1.071429) -- cycle;
\fill[teal!49!white] (1.142857,1.071429) -- (1.428571,1.285714) -- (1.142857,1.285714) -- cycle;
\fill[teal!94!white] (1.428571,0.6428571) -- (1.428571,0.4285714) -- (1.714286,0.4285714) -- cycle;
\fill[teal!85!white] (1.714286,0.6428571) -- (1.428571,0.6428571) -- (1.714286,0.4285714) -- cycle;
\fill[teal!94!white] (1.428571,0.6428571) -- (1.142857,0.4285714) -- (1.428571,0.4285714) -- cycle;
\fill[teal!89!white] (1.142857,0.4285714) -- (1.428571,0.6428571) -- (1.142857,0.6428571) -- cycle;
\fill[teal!34!white] (1.428571,0.6428571) -- (1.428571,0.8571429) -- (1.142857,0.8571429) -- cycle;
\fill[teal!46!white] (1.142857,0.6428571) -- (1.428571,0.6428571) -- (1.142857,0.8571429) -- cycle;
\fill[teal!14!white] (2,1.285714) -- (1.714286,1.071429) -- (2,1.071429) -- cycle;
\fill[teal!18!white] (1.714286,1.285714) -- (1.714286,1.071429) -- (2,1.285714) -- cycle;
\fill[teal!26!white] (1.428571,1.285714) -- (1.714286,1.071429) -- (1.714286,1.285714) -- cycle;
\fill[teal!22!white] (1.714286,1.071429) -- (1.428571,1.285714) -- (1.428571,1.071429) -- cycle;
\fill[teal!35!white] (2.285714,0.6428571) -- (2.571429,0.8571429) -- (2.285714,0.8571429) -- cycle;
\fill[teal!45!white] (2.285714,0.6428571) -- (2.571429,0.6428571) -- (2.571429,0.8571429) -- cycle;
\fill[teal!21!white] (2,0.8571429) -- (2.285714,0.6428571) -- (2.285714,0.8571429) -- cycle;
\fill[teal!12!white] (2.285714,0.6428571) -- (2,0.8571429) -- (2,0.6428571) -- cycle;
\fill[teal!18!white] (2.285714,0.6428571) -- (2,0.4285714) -- (2.285714,0.4285714) -- cycle;
\fill[teal!10!white] (2,0.4285714) -- (2.285714,0.6428571) -- (2,0.6428571) -- cycle;
\fill[teal!45!white] (2.571429,0.4285714) -- (2.857143,0.4285714) -- (2.857143,0.6428571) -- cycle;
\fill[teal!51!white] (2.571429,0.6428571) -- (2.571429,0.4285714) -- (2.857143,0.6428571) -- cycle;
\fill[teal!55!white] (2.285714,0.6428571) -- (2.571429,0.4285714) -- (2.571429,0.6428571) -- cycle;
\fill[teal!46!white] (2.571429,0.4285714) -- (2.285714,0.6428571) -- (2.285714,0.4285714) -- cycle;
\fill[teal!29!white] (0.2857143,0.8571429) -- (0.5714286,1.071429) -- (0.2857143,1.071429) -- cycle;
\fill[teal!36!white] (0.5714286,1.071429) -- (0.2857143,0.8571429) -- (0.5714286,0.8571429) -- cycle;
\fill[teal!21!white] (0.5714286,1.071429) -- (0.2857143,1.285714) -- (0.2857143,1.071429) -- cycle;
\fill[teal!20!white] (0.5714286,1.071429) -- (0.5714286,1.285714) -- (0.2857143,1.285714) -- cycle;
\fill[teal!16!white] (0.8571429,1.5) -- (0.5714286,1.714286) -- (0.5714286,1.5) -- cycle;
\fill[teal!19!white] (0.8571429,1.5) -- (0.8571429,1.714286) -- (0.5714286,1.714286) -- cycle;
\fill[teal!8!white] (0.8571429,1.714286) -- (0.8571429,1.5) -- (1.142857,1.714286) -- cycle;
\fill[teal!19!white] (0.8571429,1.5) -- (1.142857,1.5) -- (1.142857,1.714286) -- cycle;
\fill[teal!22!white] (1.714286,0.8571429) -- (1.714286,0.6428571) -- (2,0.6428571) -- cycle;
\fill[teal!9!white] (2,0.8571429) -- (1.714286,0.8571429) -- (2,0.6428571) -- cycle;
\fill[teal!9!white] (1.714286,0.8571429) -- (2,0.8571429) -- (2,1.071429) -- cycle;
\fill[teal!9!white] (1.714286,1.071429) -- (1.714286,0.8571429) -- (2,1.071429) -- cycle;
\fill[teal!51!white] (1.714286,0.8571429) -- (1.428571,0.6428571) -- (1.714286,0.6428571) -- cycle;
\fill[teal!38!white] (1.428571,0.6428571) -- (1.714286,0.8571429) -- (1.428571,0.8571429) -- cycle;
\fill[teal!11!white] (1.428571,0.8571429) -- (1.714286,0.8571429) -- (1.428571,1.071429) -- cycle;
\fill[teal!8!white] (1.714286,0.8571429) -- (1.714286,1.071429) -- (1.428571,1.071429) -- cycle;
\fill[teal!10!white] (2.571429,0) -- (2.571429,0.2142857) -- (2.285714,0) -- cycle;
\fill[teal!14!white] (2.571429,0.2142857) -- (2.285714,0.2142857) -- (2.285714,0) -- cycle;
\fill[teal!28!white] (2.285714,0.2142857) -- (2.571429,0.2142857) -- (2.285714,0.4285714) -- cycle;
\fill[teal!44!white] (2.571429,0.2142857) -- (2.571429,0.4285714) -- (2.285714,0.4285714) -- cycle;
\fill[teal!17!white] (2.857143,0.2142857) -- (2.571429,0.2142857) -- (2.857143,0) -- cycle;
\fill[teal!11!white] (2.571429,0.2142857) -- (2.571429,0) -- (2.857143,0) -- cycle;
\fill[teal!34!white] (2.857143,0.4285714) -- (2.571429,0.2142857) -- (2.857143,0.2142857) -- cycle;
\fill[teal!47!white] (2.571429,0.4285714) -- (2.571429,0.2142857) -- (2.857143,0.4285714) -- cycle;
\fill[teal!11!white] (1.142857,0.6428571) -- (0.8571429,0.8571429) -- (0.8571429,0.6428571) -- cycle;
\fill[teal!10!white] (0.8571429,0.8571429) -- (1.142857,0.6428571) -- (1.142857,0.8571429) -- cycle;
\fill[teal!22!white] (0.5714286,0.6428571) -- (0.8571429,0.8571429) -- (0.5714286,0.8571429) -- cycle;
\fill[teal!10!white] (0.8571429,0.8571429) -- (0.5714286,0.6428571) -- (0.8571429,0.6428571) -- cycle;
\fill[teal!50!white] (1.142857,1.5) -- (0.8571429,1.285714) -- (1.142857,1.285714) -- cycle;
\fill[teal!30!white] (0.8571429,1.5) -- (0.8571429,1.285714) -- (1.142857,1.5) -- cycle;
\fill[teal!15!white] (0.5714286,1.285714) -- (0.8571429,1.285714) -- (0.5714286,1.5) -- cycle;
\fill[teal!13!white] (0.8571429,1.285714) -- (0.8571429,1.5) -- (0.5714286,1.5) -- cycle;
\fill[teal!23!white] (1.142857,1.071429) -- (0.8571429,1.071429) -- (1.142857,0.8571429) -- cycle;
\fill[teal!19!white] (0.8571429,1.071429) -- (0.8571429,0.8571429) -- (1.142857,0.8571429) -- cycle;
\fill[teal!46!white] (0.8571429,1.071429) -- (1.142857,1.071429) -- (1.142857,1.285714) -- cycle;
\fill[teal!51!white] (0.8571429,1.285714) -- (0.8571429,1.071429) -- (1.142857,1.285714) -- cycle;
\fill[teal!42!white] (0.8571429,1.071429) -- (0.5714286,1.071429) -- (0.5714286,0.8571429) -- cycle;
\fill[teal!32!white] (0.8571429,0.8571429) -- (0.8571429,1.071429) -- (0.5714286,0.8571429) -- cycle;
\fill[teal!40!white] (0.5714286,1.071429) -- (0.8571429,1.071429) -- (0.5714286,1.285714) -- cycle;
\fill[teal!40!white] (0.8571429,1.071429) -- (0.8571429,1.285714) -- (0.5714286,1.285714) -- cycle;
\draw[teal,thick] (0,0) -- (4,0) -- (0,3) -- cycle;
\draw[magenta!70,thin] (1.191916,1.127577) -- (3.6,2.35) -- (4.1,2.35);\filldraw[fill=magenta,draw=white,line width=.45pt] (1.191916,1.127577) circle(1.5pt) node[above right,font=\scriptsize,text=magenta,fill=white,inner sep=.6pt] {};
\node[anchor=west,text=magenta,font=\small] at(4.12,2.35){Q4};\draw[gold!70,thin] (1,1) -- (3.6,1.8) -- (4.1,1.8);\filldraw[fill=gold,draw=white,line width=.45pt] (1,1) circle(1.5pt) node[above right,font=\scriptsize,text=gold,fill=white,inner sep=.6pt] {};
\node[anchor=west,text=gold,font=\small] at(4.12,1.8){$I$};\draw[navy!70,thin] (1.333333,1) -- (3.6,1.25) -- (4.1,1.25);\filldraw[fill=navy,draw=white,line width=.45pt] (1.333333,1) circle(1.5pt) node[above right,font=\scriptsize,text=navy,fill=white,inner sep=.6pt] {};
\node[anchor=west,text=navy,font=\small] at(4.12,1.25){$G$};\end{tikzpicture}
\caption[Probability density for Q4]{Probability density for Q4. The first moment averages all probability, including separate lobes. Numerical illustration.}\label{fig:state-Q4}
\end{figure}
\begin{center}\small\begin{tabular}{p{0.3\linewidth}p{0.3\linewidth}p{0.3\linewidth}}\toprule
Rule & $c_x$ & $c_y$\\\midrule
$\mathrm{Q0m}$ & 1.095451 & 0.997508\\
$\mathrm{Q0}$ & 1.126561 & 1.001531\\
Q1 & 1.413362 & 0.987726\\
Q2 & 0.958195 & 1.050956\\
Q3 & 1.526410 & 0.905603\\
Q4 & 1.191916 & 1.127577\\\bottomrule\end{tabular}\end{center}

\section{Degeneracy and symmetry}\label{sec:slide51}

At an equilateral triangle, some energy levels form symmetry-related doublets. Any orthonormal rotation within a degenerate eigenspace is still a valid eigenbasis. An individual member's density and moment can change under that rotation. Therefore its index alone does not define a canonical center at the degeneracy. Summing the squared amplitudes over a complete eigenspace gives the diagonal of the spectral projector. Dividing by its dimension normalizes it and removes the basis choice. Its mean is intrinsic.


\[\rho_{\rm cluster}=\frac1m\sum_{j=1}^m|\psi_j|^2.\]



For an orthonormal basis of an eigenspace, another is
$\widetilde\psi_a=\sum_bU_{ab}\psi_b$ with $U$ unitary. Then
\[
\sum_a|\widetilde\psi_a(p)|^2
=\sum_{bc}\left(\sum_aU_{ab}\overline U_{ac}\right)
\psi_b(p)\overline{\psi_c(p)}
=\sum_b|\psi_b(p)|^2.
\]
The equal mixture is basis independent and respects all geometric
symmetries. An individual degenerate eigenvector can break them.
A center based on an arbitrary solver basis is not single valued
at the degeneracy.




\paragraph{Demonstration and investigation.} Use the equilateral diagram to explain why the complete projector density respects triangle symmetry and has mean G. The basis labels describe freedom inside the doublet, not computed nodal patterns.


\section{Repeated position measurements}\label{sec:slide52}

Prepare one normalized state repeatedly and measure position.
Independent outcomes follow $|\psi_j|^2$ and their sample mean
fluctuates around Q$j$. This is an estimator of an expectation,
not the time path of a particle. It differs both from a classical
billiard trajectory and from a localized evolving wave packet.
The sampling experiment can be repeated with the offline companion.


\[\overline p_N=N^{-1}\sum_{k=1}^Np_k,\quad \operatorname{Cov}(\overline p_N)=\Sigma_{\rm quantum}/N.\]


For independent measurements the mean covariance is
$\Sigma_\psi/N$, so directional standard deviation scales as
$N^{-1/2}$. This is not a claim about correlated billiard samples.



\paragraph{Demonstration and investigation.} Reveal twelve illustrative sample points and the accumulating mean. State that the animation is schematic and the companion app can draw an actual normalized archived ground-state sample.


\section{Numerical eigenstates}\label{sec:slide53}

Use the same local basis and stiffness matrix introduced in the membrane lecture. The mass matrix contains \ensuremath{\int}\ensuremath{\varphi}\_i\ensuremath{\varphi}\_j. Solve the smallest generalized eigenpairs, normalize in M, and integrate coordinate-weighted probability moments. Quadratic elements give an interpolated maximum inside an element. Check eigenvalue gaps before naming an individual state. A right-isosceles triangle has exact separated eigenstates that provide an independent test of this solver.


\[Kv_j=\lambda_jMv_j,\quad v_i^TMv_j=\delta_{ij}.\]



Weak integration yields $Kv_j=\lambda_jMv_j$.
Euclidean normalization $v_j^Tv_j=1$ does not normalize continuum
probability; use $v_j^TMv_j=1$.
For each element, exact moments of $p\,\psi^2$ follow from
\[
\int_e\phi_1^{a_1}\phi_2^{a_2}\phi_3^{a_3}\dd A
=2a_e\frac{a_1!a_2!a_3!}{(a_1+a_2+a_3+2)!}.
\]
This cubic integration avoids treating coarse midpoint colors as
probability masses. Convergence tables need actual solver refinements;
the low-resolution teaching illustrations are not those meshes.




\section{Quantum fingerprints and ETC neighbors}\label{sec:slide54}

The same fingerprint machinery applies to the ground-state expectation. The isosceles curve is easy to view, but the right-triangle curve and derivative slices supply independent information. The neighbor labels and errors come from the frozen ETC location survey. No algebraic identity follows from small RMS error. Individual excited-state fingerprints may have gaps where the state branch is undefined or unresolved. Those gaps should remain visible rather than being interpolated into a false center.


A similarity-normalized location error is |Q\ensuremath{-}X| divided by triangle diameter. Derivative-fingerprint RMS and symmetric-Jacobian response discrepancies are separate diagnostics. A center can match location well but behave differently under vertex motion.


\begin{figure}[htbp]\centering
\begin{tikzpicture}\begin{axis}[width=.95\linewidth,height=.56\linewidth,grid=major,grid style={black!10},xlabel={family parameter $x$},ylabel={$c_y$},tick label style={font=\scriptsize},label style={font=\small},legend style={font=\scriptsize,at={(0.5,-0.24)},anchor=north,legend columns=3,draw=none},scaled ticks=false]
\addplot[navy,thick,no marks] coordinates {(0.2,0.09200876) (0.228706,0.104382) (0.2615321,0.1183134) (0.2990698,0.1339837) (0.3419952,0.1515642) (0.3910817,0.1712399) (0.4472136,0.193192) (0.5114021,0.2176239) (0.5848035,0.2447062) (0.6687403,0.2746066) (0.7647245,0.3074743) (0.8744853,0.3434347) (1,0.3825863) (1.14353,0.4249999) (1.30766,0.4707208) (1.495349,0.5197728) (1.709976,0.5721646) (1.732051,0.5773503) (1.955409,0.6278962) (2.236068,0.6869642) (2.55701,0.7493674) (2.924018,0.8151102) (3.343702,0.8842047) (3.823622,0.9566731) (4.372426,1.032548) (5,1.111869)};
\addlegendentry{$\mathrm{Q0}$}
\addplot[teal,thick,no marks] coordinates {(0.2,0.09901951) (0.228706,0.1128955) (0.2615321,0.1286033) (0.2990698,0.1463328) (0.3419952,0.1662702) (0.3910817,0.1885865) (0.4472136,0.2134218) (0.5114021,0.2408662) (0.5848035,0.2709374) (0.6687403,0.3035587) (0.7647245,0.33854) (0.8744853,0.3755687) (1,0.4142136) (1.14353,0.4539441) (1.30766,0.4941653) (1.495349,0.5342616) (1.709976,0.5736416) (1.732051,0.5773503) (1.955409,0.6117774) (2.236068,0.6482315) (2.55701,0.682671) (2.924018,0.7148684) (3.343702,0.744694) (3.823622,0.7721018) (4.372426,0.7971139) (5,0.8198039)};
\addlegendentry{X1}
\addplot[magenta,thick,no marks] coordinates {(0.2,0.06666667) (0.228706,0.07623532) (0.2615321,0.08717737) (0.2990698,0.09968992) (0.3419952,0.1139984) (0.3910817,0.1303606) (0.4472136,0.1490712) (0.5114021,0.1704674) (0.5848035,0.1949345) (0.6687403,0.2229134) (0.7647245,0.2549082) (0.8744853,0.2914951) (1,0.3333333) (1.14353,0.3811766) (1.30766,0.4358868) (1.495349,0.4984496) (1.709976,0.569992) (1.732051,0.5773503) (1.955409,0.6518028) (2.236068,0.745356) (2.55701,0.8523368) (2.924018,0.9746726) (3.343702,1.114567) (3.823622,1.274541) (4.372426,1.457475) (5,1.666667)};
\addlegendentry{X2}
\addplot[gold,thick,no marks] coordinates {(0.2,0.05049024) (0.228706,0.05790523) (0.2615321,0.06646438) (0.2990698,0.07636846) (0.3419952,0.08786248) (0.3910817,0.1012476) (0.4472136,0.1168959) (0.5114021,0.135268) (0.5848035,0.1569331) (0.6687403,0.1825908) (0.7647245,0.2130923) (0.8744853,0.2494583) (1,0.2928932) (1.14353,0.3447929) (1.30766,0.4067476) (1.495349,0.4805436) (1.709976,0.5681672) (1.732051,0.5773503) (1.955409,0.6718156) (2.236068,0.7939182) (2.55701,0.9371697) (2.924018,1.104575) (3.343702,1.299504) (3.823622,1.52576) (4.372426,1.787656) (5,2.090098)};
\addlegendentry{X10}
\addplot[blue,thick,no marks] coordinates {(0.2,0.08800741) (0.228706,0.1003429) (0.2615321,0.1143128) (0.2990698,0.1300933) (0.3419952,0.1478641) (0.3910817,0.1678039) (0.4472136,0.1900843) (0.5114021,0.2148655) (0.5848035,0.2422934) (0.6687403,0.2724997) (0.7647245,0.3056067) (0.8744853,0.3417333) (1,0.3810025) (1.14353,0.4235435) (1.30766,0.4694891) (1.495349,0.5189656) (1.709976,0.5720775) (1.732051,0.5773503) (1.955409,0.6288901) (2.236068,0.6894118) (2.55701,0.7535791) (2.924018,0.821244) (3.343702,0.8921665) (3.823622,0.9660108) (4.372426,1.042347) (5,1.12066)};
\addlegendentry{X24936}
\addplot[purple,thick,no marks] coordinates {(0.2,0.09052478) (0.228706,0.1023182) (0.2615321,0.1155502) (0.2990698,0.1303858) (0.3419952,0.147009) (0.3910817,0.1656234) (0.4472136,0.1864526) (0.5114021,0.2097398) (0.5848035,0.2357453) (0.6687403,0.2647409) (0.7647245,0.2970028) (0.8744853,0.3327993) (1,0.3723769) (1.14353,0.415944) (1.30766,0.4636541) (1.495349,0.5155914) (1.709976,0.5717592) (1.732051,0.5773503) (1.955409,0.6320731) (2.236068,0.6963574) (2.55701,0.7643466) (2.924018,0.8356881) (3.343702,0.9099495) (3.823622,0.9866261) (4.372426,1.065152) (5,1.144915)};
\addlegendentry{X39622}
\end{axis}\end{tikzpicture}
\caption[Isosceles fingerprint: the physical selector, two retained ETC location neighbors and the three classical controls]{Isosceles fingerprint: the physical selector, two retained ETC location neighbors and the three classical controls. Similarity of sampled curves is not identity.}\label{fig:fp-Q0}
\end{figure}
\begin{figure}[htbp]\centering
\begin{tikzpicture}\begin{axis}[width=.95\linewidth,height=.56\linewidth,grid=major,grid style={black!10},xlabel={apex height $x$},ylabel={$dc_y/dx$},tick label style={font=\scriptsize},label style={font=\small},legend style={font=\scriptsize,at={(0.5,-0.24)},anchor=north,legend columns=3,draw=none},scaled ticks=false]
\addplot[navy,thick,no marks] coordinates {(0.2,0.4341277) (0.228706,0.4279763) (0.2615321,0.4212254) (0.2990698,0.4137261) (0.3419952,0.4054403) (0.3910817,0.3962865) (0.4472136,0.3861861) (0.5114021,0.3750715) (0.5848035,0.3628957) (0.6687403,0.3496442) (0.7647245,0.3353468) (0.8744853,0.3200868) (1,0.3040038) (1.14353,0.2872894) (1.30766,0.2701732) (1.495349,0.2529039) (1.709976,0.2357272) (1.732051,0.2340978) (1.955409,0.2188683) (2.236068,0.2025179) (2.55701,0.1868268) (2.924018,0.1719044) (3.343702,0.1578224) (3.823622,0.14462) (4.372426,0.1323103) (5,0.1208859)};
\addlegendentry{$\mathrm{Q0}$}
\addplot[teal,thick,no marks] coordinates {(0.2,0.4833837) (0.228706,0.4807868) (0.2615321,0.4752065) (0.2990698,0.4681264) (0.3419952,0.4592189) (0.3910817,0.4481301) (0.4472136,0.434504) (0.5114021,0.4180222) (0.5848035,0.3984569) (0.6687403,0.3757348) (0.7647245,0.3499968) (0.8744853,0.3216382) (1,0.2913098) (1.14353,0.2598701) (1.30766,0.2282923) (1.495349,0.1975472) (1.709976,0.1820376) (1.732051,0.1553818) (1.955409,0.1406326) (2.236068,0.1178415) (2.55701,0.09686305) (2.924018,0.0788403) (3.343702,0.06362058) (3.823622,0.05095623) (4.372426,0.04054999) (5,0.03615511)};
\addlegendentry{X1}
\addplot[magenta,thick,no marks] coordinates {(0.2,0.3333333) (0.228706,0.3333333) (0.2615321,0.3333333) (0.2990698,0.3333333) (0.3419952,0.3333333) (0.3910817,0.3333333) (0.4472136,0.3333333) (0.5114021,0.3333333) (0.5848035,0.3333333) (0.6687403,0.3333333) (0.7647245,0.3333333) (0.8744853,0.3333333) (1,0.3333333) (1.14353,0.3333333) (1.30766,0.3333333) (1.495349,0.3333333) (1.709976,0.3333333) (1.732051,0.3333333) (1.955409,0.3333333) (2.236068,0.3333333) (2.55701,0.3333333) (2.924018,0.3333333) (3.343702,0.3333333) (3.823622,0.3333333) (4.372426,0.3333333) (5,0.3333333)};
\addlegendentry{X2}
\addplot[gold,thick,no marks] coordinates {(0.2,0.2583082) (0.228706,0.2596066) (0.2615321,0.2623967) (0.2990698,0.2659368) (0.3419952,0.2703905) (0.3910817,0.275935) (0.4472136,0.282748) (0.5114021,0.2909889) (0.5848035,0.3007716) (0.6687403,0.3121326) (0.7647245,0.3250016) (0.8744853,0.3391809) (1,0.3543451) (1.14353,0.370065) (1.30766,0.3858539) (1.495349,0.4012264) (1.709976,0.4089812) (1.732051,0.4223091) (1.955409,0.4296837) (2.236068,0.4410793) (2.55701,0.4515685) (2.924018,0.4605799) (3.343702,0.4681897) (3.823622,0.4745219) (4.372426,0.479725) (5,0.4819224)};
\addlegendentry{X10}
\addplot[blue,thick,no marks] coordinates {(0.2,0.4297202) (0.228706,0.4275074) (0.2615321,0.4228075) (0.2990698,0.4169774) (0.3419952,0.4098445) (0.3910817,0.401262) (0.4472136,0.391136) (0.5114021,0.3794542) (0.5848035,0.3663077) (0.6687403,0.351895) (0.7647245,0.3365022) (0.8744853,0.3204577) (1,0.3040767) (1.14353,0.2876113) (1.30766,0.2712251) (1.495349,0.2549949) (1.709976,0.2466588) (1.732051,0.2314793) (1.955409,0.2223368) (2.236068,0.2072617) (2.55701,0.1916306) (2.924018,0.176165) (3.343702,0.1609227) (3.823622,0.1459874) (4.372426,0.1314625) (5,0.1247869)};
\addlegendentry{X24936}
\addplot[purple,thick,no marks] coordinates {(0.2,0.4108368) (0.228706,0.4067043) (0.2615321,0.3988918) (0.2990698,0.3909725) (0.3419952,0.3829673) (0.3910817,0.3748733) (0.4472136,0.3666584) (0.5114021,0.3582579) (0.5848035,0.3495722) (0.6687403,0.340469) (0.7647245,0.3307898) (0.8744853,0.3203655) (1,0.3090372) (1.14353,0.2966816) (1.30766,0.2832348) (1.495349,0.2687073) (1.709976,0.260914) (1.732051,0.2457451) (1.955409,0.2361173) (2.236068,0.2198688) (2.55701,0.2025303) (2.924018,0.1850827) (3.343702,0.1677826) (3.823622,0.1508693) (4.372426,0.1345559) (5,0.1270962)};
\addlegendentry{X39622}
\end{axis}\end{tikzpicture}
\caption[Isosceles derivative comparison]{Isosceles derivative comparison. Neighbor derivatives use sampled differences, and do not replace the full vertex-response matrices.}\label{fig:deriv-Q0}
\end{figure}

\section{Four ways to probe triangle shape}\label{sec:slide55}

The atlas uses four complementary special classes. Isosceles triangles reduce the point to one coordinate. Right triangles retain two coordinates with a simple boundary. Thin triangles expose singular limits and boundary layers. Nearly equilateral triangles expose symmetry constraints and local Taylor coefficients. A separate two-dimensional shape sample is necessary because none of these one-parameter families covers every scalene triangle. Excited-state degeneracies make the equilateral limit especially delicate.


\[(-1,0),(1,0),(x,\varepsilon);\quad(-1,0),(1,0),(x,\sqrt3+y).\]



For isosceles triangles take $A=(-1,0),B=(1,0),C=(0,x)$, $x>0$.
Symmetry forces $c_x=0$. For right triangles take
$A=(0,0),B=(1,0),C=(0,x)$ and record both coordinates.
In the thin family, $0\le x\le1$ and $\varepsilon\downarrow0$;
the order of geometric and physical limits matters.

At the equilateral shape every single-valued equivariant point is
$G=(0,\sqrt3/3)$. Symmetry constrains the two-variable Taylor expansion
in $x,y$ but does not determine every coefficient.
An individual degenerate quantum mode may lack an analytic expansion;
a cluster rule avoids that ambiguity.
One-dimensional fingerprints cannot uniquely determine a function
on two-dimensional shape space.




\section{Reflecting Brownian motion}\label{sec:slide56}

An ideal reflecting wall returns diffusing particles rather than absorbing them. The stationary diffusion equation with zero normal flux admits a constant density. On a connected triangle it is the normalized invariant distribution. Its mean is therefore G. Long-time ensemble density and ergodic time averages recover this area mean under the standard reflecting process. The implementation of reflection near corners needs care, but the invariant measure statement does not require our particular animation algorithm.


\[\partial_t\rho=D_{\rm B}\Delta\rho,\quad\partial_n\rho=0,\quad\rho_\infty=\mathcal A^{-1}.\]


Integrating the diffusion equation gives zero boundary flux
and conserved probability. At stationarity integration by parts yields
$\int_T|\nabla\rho|^2=0$, hence a constant density.



\paragraph{Demonstration and investigation.} Run the offline reflecting-walk illustration or use its archived path excerpt. Explain that the short path is a schematic sample, not proof of uniform exploration.


\section{Absorption changes the conditioning}\label{sec:slide57}

An absorbing particle is removed at its first boundary hit. At long times, the lowest Dirichlet mode dominates the transition density. Condition only on surviving up to the observation time, and the endpoint density is proportional to $\psi_0$. Its mean differs from the quantum expectation because the amplitude is not squared. A different construction conditions the process never to be killed, using the ground-state Doob transform. Its invariant density is proportional to $\psi_0^2$, so its mean is $\mathrm{Q0}$. Meanwhile the longest expected lifetime still selects F1, and equal exit probabilities select E0.


\[\rho_{\rm endpoint}\propto\psi_0,\qquad\rho_{\rm forever}\propto\psi_0^2.\]



The killed heat kernel expands as
$p_t(x,y)=\sum_j e^{-D_{\rm B}\lambda_jt}\psi_j(x)\psi_j(y)$.
Divide by survival probability $\int_Tp_t(x,y)\dd y$.
The ground term dominates, giving endpoint density
$\psi_0(y)/\int\psi_0$, the quasistationary distribution.

Condition instead on survival indefinitely. The Doob-transformed
transition kernel is
\[
p_t^Q(x,y)=e^{D_{\rm B}\lambda_0t}
p_t(x,y)\frac{\psi_0(y)}{\psi_0(x)}.
\]
It has invariant density $\psi_0^2$, by integration or detailed
balance, and generator
$D_{\rm B}[\Delta+2\nabla\log\psi_0\cdot\nabla]$.
Its mean is Q0. The late surviving endpoint mean equals the late
temperature mean, generally a different point.




\section{A classical billiard at speed one}\label{sec:slide58}

A point particle travels in straight lines and reflects at a side with equal incidence and reflection angles. Its speed stays one. The position mean must weight continuous time, not give every collision equal weight. Each straight flight contributes duration times its midpoint, including a final partial flight. A rational-angle polygon has strong almost-every-direction exploration results, while exceptional periodic directions can fail. The 3--4--5 triangle is not a rational-angle triangle, so that theorem does not settle every trajectory here.


\[\overline p(t)=t^{-1}\int_0^tp(s)\dd s,\quad\int_{\rm flight}p(s)\dd s=\frac{\ell}{2}(p_{\rm start}+p_{\rm end}).\]


Reflection is $v'=v-2(v\cdot n)n$. A flight of duration
$\delta t$ beginning at $p$ contributes
$\delta t[p+v\delta t/2]$ to the time integral.



\paragraph{Demonstration and investigation.} Play the moving particle along a precomputed exact flight path. The companion app reports its time mean. Ask why equal averaging of collision locations would bias the result.


\section{Convergence of billiard means}\label{sec:slide59}

The archived experiment uses 1024 random boundary launches per law, unit diameter and unit speed, through time 30000. It reports mean distance from G, the standard deviation of that distance and RMS distance. At the last time the RMS is about 0.0204 triangle diameters. The simulations show slow finite-time behavior and do not establish a vanishing asymptotic RMS for this triangle. A small bias of the ensemble-average vector would not imply that each individual trajectory mean is close to G. The choice of uniform inward angle or cosine-flux angle is also recorded.


\[\mathrm{RMS}^2=(\mathbb E|\overline p(t)-G|)^2+\operatorname{Var}|\overline p(t)-G|.\]



Store time-weighted flight integrals. Collision endpoint averages
overweight short flights. Reflection is
$v\mapsto v-2(v\cdot n)n$; exactly vertex-hitting trajectories are
a measure-zero ideal-model ambiguity.
For $N$ launches define
$d_i(t)=|\overline p_i(t)-G|/D_T$ and compute mean distance,
standard deviation, RMS and vector ensemble bias separately.
The RMS identity uses population variance; a sample variance with
denominator $N-1$ requires a corresponding correction.

The archived 1,024-launch run reached $t=30{,}000$ at unit speed and
unit diameter, with RMS about $0.0204$.
Uniform boundary location and uniform inward angle are launch
conventions. Invariant collision measure has a cosine angular weight.
The reference triangle is not a rational-angle polygon.
Known almost-everywhere direction results \cite{kms} and
density-one quantum equidistribution for rational polygons
\cite{marklof} have hypotheses. Neither proves that all $Q_j$ in
every triangle approach $G$.



\begin{figure}[htbp]\centering
\begin{tikzpicture}\begin{axis}[width=.95\linewidth,height=.56\linewidth,grid=major,grid style={black!10},xlabel={$t$; speed and diameter $=1$},ylabel={distance / $D$},tick label style={font=\scriptsize},label style={font=\small},legend style={font=\scriptsize,at={(0.5,-0.24)},anchor=north,legend columns=3,draw=none},scaled ticks=false,xmode=log]
\addplot[navy,thick,no marks] coordinates {(0.1,0.3192981) (0.1233915,0.3147119) (0.1522547,0.3089831) (0.1878694,0.3017961) (0.2318149,0.2926799) (0.28604,0.2810961) (0.3529491,0.2663581) (0.4355093,0.2480027) (0.5373817,0.225977) (0.6630834,0.2010844) (0.8181888,0.1755115) (1,0.1533367) (1.009576,0.1523946) (1.245731,0.1348914) (1.537126,0.1210677) (1.896684,0.1089524) (2.340347,0.09490182) (2.88779,0.08081031) (3.563289,0.0780296) (4.396797,0.07029419) (5.425275,0.06601469) (6.69433,0.06030054) (8.260236,0.05656261) (10,0.05320303) (10.19243,0.05268738) (12.5766,0.05011105) (15.51846,0.04868586) (19.14846,0.04735904) (23.62758,0.04538589) (29.15443,0.04309874) (35.9741,0.04198975) (44.38899,0.04062645) (54.77226,0.03933369) (67.58433,0.03784679) (83.39333,0.03676873) (100,0.03575898) (102.9003,0.03564728) (126.9703,0.03468947) (156.6706,0.03365755) (193.3182,0.03254668) (238.5383,0.03158588) (294.3361,0.03085117) (363.1858,0.03007687) (448.1405,0.02954992) (552.9674,0.02883767) (682.3149,0.02852384) (841.9188,0.02809202) (1000,0.02760755) (1038.857,0.0274702) (1281.861,0.02710385) (1581.708,0.0264858) (1951.694,0.02595546) (2408.225,0.02535879) (2971.545,0.0246456) (3000,0.02461636) (3666.635,0.02417024) (4524.317,0.02363008) (5582.625,0.02310872) (6888.486,0.02285814) (8499.808,0.02239649) (10000,0.02226308) (10488.04,0.0222406) (12941.36,0.02196998) (15968.54,0.0214912) (19703.83,0.02109931) (24312.85,0.02075769) (30000,0.02041631)};
\addlegendentry{RMS}
\addplot[teal,thick,no marks] coordinates {(0.1,0.1084825) (0.1233915,0.1095966) (0.1522547,0.1108017) (0.1878694,0.112025) (0.2318149,0.1130898) (0.28604,0.1136268) (0.3529491,0.1126083) (0.4355093,0.1082373) (0.5373817,0.1009242) (0.6630834,0.09130985) (0.8181888,0.07964985) (1,0.06847023) (1.009576,0.0679662) (1.245731,0.05804676) (1.537126,0.05359218) (1.896684,0.05231751) (2.340347,0.04549745) (2.88779,0.03761229) (3.563289,0.03781603) (4.396797,0.03422519) (5.425275,0.03393798) (6.69433,0.03308676) (8.260236,0.03286836) (10,0.03073262) (10.19243,0.03054214) (12.5766,0.02950364) (15.51846,0.02963598) (19.14846,0.02953506) (23.62758,0.0288071) (29.15443,0.0277857) (35.9741,0.02680149) (44.38899,0.02600277) (54.77226,0.02522386) (67.58433,0.02407669) (83.39333,0.02336869) (100,0.02280279) (102.9003,0.02270746) (126.9703,0.02213069) (156.6706,0.02145221) (193.3182,0.02073943) (238.5383,0.02009072) (294.3361,0.01943643) (363.1858,0.018968) (448.1405,0.0188203) (552.9674,0.01832115) (682.3149,0.01817953) (841.9188,0.01796954) (1000,0.01759878) (1038.857,0.01751758) (1281.861,0.01715346) (1581.708,0.01680022) (1951.694,0.01643251) (2408.225,0.01588984) (2971.545,0.01528387) (3000,0.0152634) (3666.635,0.01499447) (4524.317,0.01457941) (5582.625,0.01419833) (6888.486,0.01416694) (8499.808,0.01394391) (10000,0.01386459) (10488.04,0.01383611) (12941.36,0.01372975) (15968.54,0.01342358) (19703.83,0.01319894) (24312.85,0.01303186) (30000,0.01289504)};
\addlegendentry{SD of distance}
\addplot[magenta,thick,no marks] coordinates {(0.1,0.0481274) (0.1233915,0.04796202) (0.1522547,0.04767098) (0.1878694,0.04713304) (0.2318149,0.04617215) (0.28604,0.04449985) (0.3529491,0.04195839) (0.4355093,0.03826826) (0.5373817,0.03293481) (0.6630834,0.02547408) (0.8181888,0.01641734) (1,0.009433875) (1.009576,0.009230989) (1.245731,0.007899144) (1.537126,0.005628525) (1.896684,0.005722334) (2.340347,0.003951745) (2.88779,0.003751989) (3.563289,0.00819651) (4.396797,0.005735032) (5.425275,0.005247861) (6.69433,0.003946115) (8.260236,0.004333325) (10,0.004650546) (10.19243,0.004339938) (12.5766,0.004377176) (15.51846,0.003464022) (19.14846,0.004034008) (23.62758,0.003784687) (29.15443,0.00301958) (35.9741,0.002799961) (44.38899,0.002290242) (54.77226,0.001925868) (67.58433,0.001823472) (83.39333,0.001810109) (100,0.00177967) (102.9003,0.001753192) (126.9703,0.00172975) (156.6706,0.001727772) (193.3182,0.001703865) (238.5383,0.001452027) (294.3361,0.001399991) (363.1858,0.001287501) (448.1405,0.001097423) (552.9674,0.001005678) (682.3149,0.000828099) (841.9188,0.000480684) (1000,0.000398719) (1038.857,0.0003560706) (1281.861,0.0002156173) (1581.708,0.0001166584) (1951.694,0.0002709158) (2408.225,0.0002386872) (2971.545,0.0001731672) (3000,0.0001726055) (3666.635,6.104232e-05) (4524.317,0.0001183162) (5582.625,0.0002223564) (6888.486,0.000129103) (8499.808,3.360565e-05) (10000,9.204582e-05) (10488.04,0.000102535) (12941.36,0.0001699513) (15968.54,0.0001341773) (19703.83,0.0001954704) (24312.85,0.0003300798) (30000,0.0003936469)};
\addlegendentry{ensemble mean bias}
\end{axis}\end{tikzpicture}
\caption[Reference billiard: 1,024 launches with uniform inward angles]{Reference billiard: 1,024 launches with uniform inward angles. Small ensemble bias does not prove concentration of individual trajectory means.}\label{fig:billiard}
\end{figure}

\section{Quantum and classical limits}\label{sec:slide60}

A classical uniform spatial measure suggests G, but passing from classical motion to quantum eigenstates needs a theorem. For rational polygons, quantum spatial equidistribution results concern a density-one subsequence of eigenstates, not necessarily every state. Individual eigenstates can have special structure, and degeneracies require a projector convention. For a general triangular shape, do not turn the heuristic Qj\ensuremath{\to}G into a universal identity. A thermal quantum mean has a different, well-defined high-temperature route toward the area mean.


A thermal quantum mixture uses
\[
\rho_\tau(p)=
\frac{\sum_j e^{-\tau\lambda_j}|\psi_j(p)|^2}
{\sum_j e^{-\tau\lambda_j}},\qquad \tau>0.
\]
Large $\tau$ selects the ground probability density.
As $\tau\downarrow0$, heat-kernel asymptotics give normalized uniform
area measure on a fixed triangle, hence mean $G$. Sufficiently many
modes are needed. This averaged high-temperature limit does not prove
a limit for every individually ordered high-energy state.



\paragraph{Demonstration and investigation.} End with three labels: independent measurements, one billiard orbit, and thermal mixture. Let students match the averaging procedure to each one.


\section*{Questions and exercises}\addcontentsline{toc}{section}{Questions and exercises}

\begin{enumerate}

\item\label{ex:47} Does changing the particle mass move the ground-state expectation in this ideal well?

\item\label{ex:48} Why is normalization necessary in the minimization?

\item\label{ex:49} Could a wavefunction have its expectation on a nodal line?

\item\label{ex:50} Do two orthogonal states necessarily have different expectation positions?

\item\label{ex:51} What is the smallest canonical unit at a double eigenvalue?

\item\label{ex:52} Why is a quantum expectation not a prediction of every measured position?

\item\label{ex:53} Why does ordinary Euclidean normalization of the coefficient vector fail?

\item\label{ex:54} What would persuade us that $\mathrm{Q0}$ equals an ETC entry?

\item\label{ex:55} Which family is most likely to reveal singular numerical conditioning?

\item\label{ex:56} What changes if the walls absorb particles?

\item\label{ex:57} How many different Brownian centers have we met?

\item\label{ex:58} Does speed one imply the particle must visit every region equally?

\item\label{ex:59} Why can the vector average of the cloud look good while the RMS stays large?

\item\label{ex:60} Which quantity is basis-independent at a degenerate energy?

\end{enumerate}

\chapter{Heat, reaction and the numerical atlas}\label{ch:5}

Heat diffusion gives a family of centers indexed by time. Production and removal give another family indexed by a dimensionless rate. These examples show why a physically meaningful center needs a parameter convention, a numerical method and a statement of evidence. The second half of this lecture develops the atlas tools that let us compare implicit centers with classical ones without inventing an exact formula.


\section{A uniformly hot triangular plate}\label{sec:slide62}

Assume isotropic, constant thermal diffusivity \ensuremath{\kappa} and no internal source. The plate begins with uniform excess temperature one in its interior. At time zero its boundary is clamped to the bath temperature zero. At positive times heat flows out. The ideal initial boundary jump is allowed in the weak heat equation. The spatial temperature maximum is a natural center, but its position can move even though the initial temperature was uniform. All lengths must be scaled consistently when comparing different triangles.


\[\partial_tu=\kappa\Delta u,\quad u|_{\partial T}=0,\quad u(p,0)=1.\]


Energy balance says heat-capacity density times $\partial_t T$ equals the divergence of thermal conductivity times the gradient. With uniform coefficients, \ensuremath{\kappa} is conductivity divided by heat-capacity density.



\section{The same eigenmodes describe cooling}\label{sec:slide63}

Expand the initial temperature in the Dirichlet eigenfunctions from the quantum lecture. Substituting one spatial eigenmode into the heat equation gives an exponential time factor. High eigenvalues decay faster, leaving the positive ground mode at late times. The coefficients are integrals of $\psi_j$ against the initial condition, so for uniform initial temperature $c_j$=$\int_T\psi_j$. They differ from the thermal quantum weights, which involve the diagonal heat kernel. We use \ensuremath{\Delta}=$\kappa t$/$L^2$ with $L^2$=(a\ensuremath{^2}+b\ensuremath{^2}+c\ensuremath{^2})/3.


\[u(p,t)=\sum_{j=0}^\infty a_je^{-\kappa\lambda_jt}\psi_j(p),\quad a_j=\int_T\psi_j\dd A.\]



Substitution of the mode expansion into the heat equation gives
$\dot b_j=-\kappa\lambda_jb_j$. Initial value one is understood in
$L^2$, since it has nonzero boundary trace before instantaneous
edge cooling. This is Dirichlet cooling, not insulating Neumann flow.

At long times
$u\sim a_0e^{-\kappa\lambda_0t}\psi_0$.
The hottest point tends to Q0m, but the normalized mean weights
$\psi_0$ rather than $\psi_0^2$. Heat content is
$\sum_j a_j^2e^{-\kappa\lambda_jt}$.
Concavity questions for heat flow need care \cite{heatconc}.




\section{H(\ensuremath{\Delta}): the hottest point}\label{sec:slide64}

At early time the temperature is nearly flat except near the sides, making maximum localization numerically difficult. In the vanishing-time limit, the point with greatest boundary distance is the incenter. At late time the lowest eigenmode dominates and the maximum tends to $\mathrm{Q0m}$. The finite-time path between these endpoints is a new family of centers. The atlas uses \ensuremath{\Delta}=.01,.03,.1,.3,1. Its retained refined witnesses have negative numerical responses, but they lack continuum interval certificates.


\[\Delta_h=\kappa t/L^2,\quad {\rm H}(\Delta_h)=\argmax_Tu(p,\Delta_hL^2/\kappa).\]



At short times the point farthest from the cold boundary loses
heat least. Distance-to-boundary asymptotics select $I$:
$\mathrm H(\Delta_h)\to I$ as $\Delta_h\downarrow0$.
The early interior temperature is nearly flat, so maximizing raw
floating-point values is ill conditioned. Use boundary-layer
asymptotics, high precision or a heat-loss variable.
At long times factor out the leading exponential to avoid underflow.
Hold dimensionless $\Delta_h$ fixed when vertices move; holding
dimensional time fixed mixes shape and scale.




\paragraph{Demonstration and investigation.} Run the center motion across the five archived heat times. Keep the visible center path separate from the normalized field colors.


\section{Cooling changes the whole field}\label{sec:slide65}

The cold boundary influences the field before the interior
loses most of its heat. The reference maximum drifts only slightly.
The two archived spectral snapshots below show early and late
spatial structure. Each is normalized to its own peak, so the
color comparison removes the cooling amplitude.
The companion slider interpolates selected snapshots for display;
interpolated frames are not fresh PDE solves.


\begin{figure}[htbp]\centering
\begin{tikzpicture}[x=1.7cm,y=1.7cm]
\fill[teal!11!white] (3.142857,0.6428571) -- (2.857143,0.8571429) -- (2.857143,0.6428571) -- cycle;
\fill[teal!17!white] (1.714286,1.714286) -- (1.428571,1.928571) -- (1.428571,1.714286) -- cycle;
\fill[teal!17!white] (1.714286,1.714286) -- (1.714286,1.5) -- (2,1.5) -- cycle;
\fill[teal!10!white] (0.5714286,2.571429) -- (0.2857143,2.785714) -- (0.2857143,2.571429) -- cycle;
\fill[teal!15!white] (2.571429,1.071429) -- (2.285714,1.285714) -- (2.285714,1.071429) -- cycle;
\fill[teal!8!white] (3.714286,0.2142857) -- (3.428571,0.4285714) -- (3.428571,0.2142857) -- cycle;
\fill[teal!8!white] (3.714286,0.2142857) -- (3.714286,0) -- (4,0) -- cycle;
\fill[teal!14!white] (0.8571429,2.142857) -- (1.142857,2.142857) -- (0.8571429,2.357143) -- cycle;
\fill[teal!16!white] (1.142857,2.142857) -- (1.142857,1.928571) -- (1.428571,1.928571) -- cycle;
\fill[teal!8!white] (0,2.785714) -- (0.2857143,2.785714) -- (0,3) -- cycle;
\fill[teal!12!white] (0.5714286,2.571429) -- (0.5714286,2.357143) -- (0.8571429,2.357143) -- cycle;
\fill[teal!16!white] (2,1.285714) -- (2.285714,1.285714) -- (2,1.5) -- cycle;
\fill[teal!14!white] (2.571429,1.071429) -- (2.571429,0.8571429) -- (2.857143,0.8571429) -- cycle;
\fill[teal!9!white] (3.142857,0.4285714) -- (3.428571,0.4285714) -- (3.142857,0.6428571) -- cycle;
\fill[teal!9!white] (0.2857143,0) -- (0,0.2142857) -- (0,0) -- cycle;
\fill[teal!14!white] (0,0.2142857) -- (0.2857143,0) -- (0.2857143,0.2142857) -- cycle;
\fill[teal!15!white] (0.5714286,0.2142857) -- (0.2857143,0) -- (0.5714286,0) -- cycle;
\fill[teal!19!white] (0.2857143,0) -- (0.5714286,0.2142857) -- (0.2857143,0.2142857) -- cycle;
\fill[teal!25!white] (1.714286,1.5) -- (1.714286,1.714286) -- (1.428571,1.714286) -- cycle;
\fill[teal!41!white] (1.428571,1.5) -- (1.714286,1.5) -- (1.428571,1.714286) -- cycle;
\fill[teal!18!white] (0,1.928571) -- (0,1.714286) -- (0.2857143,1.928571) -- cycle;
\fill[teal!29!white] (0,1.714286) -- (0.2857143,1.714286) -- (0.2857143,1.928571) -- cycle;
\fill[teal!11!white] (0.2857143,2.357143) -- (0,2.571429) -- (0,2.357143) -- cycle;
\fill[teal!12!white] (0,2.571429) -- (0.2857143,2.357143) -- (0.2857143,2.571429) -- cycle;
\fill[teal!13!white] (0,2.142857) -- (0.2857143,2.357143) -- (0,2.357143) -- cycle;
\fill[teal!19!white] (0.2857143,2.357143) -- (0,2.142857) -- (0.2857143,2.142857) -- cycle;
\fill[teal!16!white] (0,2.142857) -- (0,1.928571) -- (0.2857143,1.928571) -- cycle;
\fill[teal!22!white] (0.2857143,2.142857) -- (0,2.142857) -- (0.2857143,1.928571) -- cycle;
\fill[teal!43!white] (0.2857143,1.714286) -- (0.5714286,1.928571) -- (0.2857143,1.928571) -- cycle;
\fill[teal!52!white] (0.5714286,1.714286) -- (0.5714286,1.928571) -- (0.2857143,1.714286) -- cycle;
\fill[teal!20!white] (0.2857143,1.5) -- (0,1.714286) -- (0,1.5) -- cycle;
\fill[teal!31!white] (0,1.714286) -- (0.2857143,1.5) -- (0.2857143,1.714286) -- cycle;
\fill[teal!50!white] (0.2857143,1.5) -- (0.5714286,1.714286) -- (0.2857143,1.714286) -- cycle;
\fill[teal!60!white] (0.5714286,1.714286) -- (0.2857143,1.5) -- (0.5714286,1.5) -- cycle;
\fill[teal!21!white] (0,1.285714) -- (0.2857143,1.5) -- (0,1.5) -- cycle;
\fill[teal!34!white] (0.2857143,1.5) -- (0,1.285714) -- (0.2857143,1.285714) -- cycle;
\fill[teal!21!white] (0,1.285714) -- (0,1.071429) -- (0.2857143,1.071429) -- cycle;
\fill[teal!34!white] (0.2857143,1.285714) -- (0,1.285714) -- (0.2857143,1.071429) -- cycle;
\fill[teal!8!white] (3.428571,0) -- (3.714286,0) -- (3.428571,0.2142857) -- cycle;
\fill[teal!8!white] (3.714286,0) -- (3.714286,0.2142857) -- (3.428571,0.2142857) -- cycle;
\fill[teal!18!white] (1.428571,0) -- (1.714286,0) -- (1.428571,0.2142857) -- cycle;
\fill[teal!28!white] (1.714286,0) -- (1.714286,0.2142857) -- (1.428571,0.2142857) -- cycle;
\fill[teal!16!white] (0.8571429,0) -- (0.5714286,0.2142857) -- (0.5714286,0) -- cycle;
\fill[teal!26!white] (0.5714286,0.2142857) -- (0.8571429,0) -- (0.8571429,0.2142857) -- cycle;
\fill[teal!18!white] (1.142857,0.2142857) -- (1.142857,0) -- (1.428571,0) -- cycle;
\fill[teal!29!white] (1.142857,0.2142857) -- (1.428571,0) -- (1.428571,0.2142857) -- cycle;
\fill[teal!18!white] (1.142857,0.2142857) -- (0.8571429,0) -- (1.142857,0) -- cycle;
\fill[teal!28!white] (0.8571429,0) -- (1.142857,0.2142857) -- (0.8571429,0.2142857) -- cycle;
\fill[teal!40!white] (1.142857,1.928571) -- (1.142857,1.714286) -- (1.428571,1.714286) -- cycle;
\fill[teal!24!white] (1.428571,1.928571) -- (1.142857,1.928571) -- (1.428571,1.714286) -- cycle;
\fill[teal!22!white] (0.8571429,2.142857) -- (1.142857,1.928571) -- (1.142857,2.142857) -- cycle;
\fill[teal!36!white] (1.142857,1.928571) -- (0.8571429,2.142857) -- (0.8571429,1.928571) -- cycle;
\fill[teal!9!white] (0,2.785714) -- (0,2.571429) -- (0.2857143,2.571429) -- cycle;
\fill[teal!9!white] (0.2857143,2.785714) -- (0,2.785714) -- (0.2857143,2.571429) -- cycle;
\fill[teal!52!white] (1.142857,1.928571) -- (0.8571429,1.714286) -- (1.142857,1.714286) -- cycle;
\fill[teal!49!white] (0.8571429,1.714286) -- (1.142857,1.928571) -- (0.8571429,1.928571) -- cycle;
\fill[teal!59!white] (0.8571429,1.714286) -- (0.5714286,1.928571) -- (0.5714286,1.714286) -- cycle;
\fill[teal!54!white] (0.5714286,1.928571) -- (0.8571429,1.714286) -- (0.8571429,1.928571) -- cycle;
\fill[teal!17!white] (0.2857143,2.357143) -- (0.5714286,2.357143) -- (0.2857143,2.571429) -- cycle;
\fill[teal!13!white] (0.5714286,2.357143) -- (0.5714286,2.571429) -- (0.2857143,2.571429) -- cycle;
\fill[teal!18!white] (0,0.6428571) -- (0,0.4285714) -- (0.2857143,0.6428571) -- cycle;
\fill[teal!26!white] (0,0.4285714) -- (0.2857143,0.4285714) -- (0.2857143,0.6428571) -- cycle;
\fill[teal!14!white] (0,0.4285714) -- (0,0.2142857) -- (0.2857143,0.2142857) -- cycle;
\fill[teal!22!white] (0.2857143,0.4285714) -- (0,0.4285714) -- (0.2857143,0.2142857) -- cycle;
\fill[teal!23!white] (2.285714,1.285714) -- (2,1.285714) -- (2.285714,1.071429) -- cycle;
\fill[teal!38!white] (2,1.285714) -- (2,1.071429) -- (2.285714,1.071429) -- cycle;
\fill[teal!50!white] (2,1.071429) -- (2,0.8571429) -- (2.285714,1.071429) -- cycle;
\fill[teal!48!white] (2,0.8571429) -- (2.285714,0.8571429) -- (2.285714,1.071429) -- cycle;
\fill[teal!34!white] (2.285714,0.8571429) -- (2.571429,0.8571429) -- (2.285714,1.071429) -- cycle;
\fill[teal!21!white] (2.571429,0.8571429) -- (2.571429,1.071429) -- (2.285714,1.071429) -- cycle;
\fill[teal!55!white] (1.428571,0.4285714) -- (1.714286,0.2142857) -- (1.714286,0.4285714) -- cycle;
\fill[teal!48!white] (1.714286,0.2142857) -- (1.428571,0.4285714) -- (1.428571,0.2142857) -- cycle;
\fill[teal!9!white] (3.142857,0.2142857) -- (3.428571,0) -- (3.428571,0.2142857) -- cycle;
\fill[teal!9!white] (3.142857,0.2142857) -- (3.142857,0) -- (3.428571,0) -- cycle;
\fill[teal!10!white] (3.142857,0) -- (3.142857,0.2142857) -- (2.857143,0) -- cycle;
\fill[teal!13!white] (3.142857,0.2142857) -- (2.857143,0.2142857) -- (2.857143,0) -- cycle;
\fill[teal!27!white] (1.714286,0) -- (2,0.2142857) -- (1.714286,0.2142857) -- cycle;
\fill[teal!17!white] (2,0.2142857) -- (1.714286,0) -- (2,0) -- cycle;
\fill[teal!33!white] (0.5714286,0.2142857) -- (0.5714286,0.4285714) -- (0.2857143,0.2142857) -- cycle;
\fill[teal!34!white] (0.5714286,0.4285714) -- (0.2857143,0.4285714) -- (0.2857143,0.2142857) -- cycle;
\fill[teal!43!white] (0.8571429,0.4285714) -- (0.5714286,0.2142857) -- (0.8571429,0.2142857) -- cycle;
\fill[teal!47!white] (0.8571429,0.4285714) -- (0.5714286,0.4285714) -- (0.5714286,0.2142857) -- cycle;
\fill[teal!56!white] (1.142857,1.5) -- (1.428571,1.5) -- (1.428571,1.714286) -- cycle;
\fill[teal!55!white] (1.142857,1.714286) -- (1.142857,1.5) -- (1.428571,1.714286) -- cycle;
\fill[teal!33!white] (0.5714286,2.142857) -- (0.2857143,2.142857) -- (0.2857143,1.928571) -- cycle;
\fill[teal!41!white] (0.5714286,1.928571) -- (0.5714286,2.142857) -- (0.2857143,1.928571) -- cycle;
\fill[teal!26!white] (0.5714286,2.142857) -- (0.2857143,2.357143) -- (0.2857143,2.142857) -- cycle;
\fill[teal!24!white] (0.5714286,2.142857) -- (0.5714286,2.357143) -- (0.2857143,2.357143) -- cycle;
\fill[teal!36!white] (0.8571429,2.142857) -- (0.5714286,2.142857) -- (0.8571429,1.928571) -- cycle;
\fill[teal!44!white] (0.5714286,2.142857) -- (0.5714286,1.928571) -- (0.8571429,1.928571) -- cycle;
\fill[teal!24!white] (0.5714286,2.142857) -- (0.8571429,2.142857) -- (0.8571429,2.357143) -- cycle;
\fill[teal!22!white] (0.5714286,2.357143) -- (0.5714286,2.142857) -- (0.8571429,2.357143) -- cycle;
\fill[teal!33!white] (1.714286,1.285714) -- (2,1.285714) -- (2,1.5) -- cycle;
\fill[teal!33!white] (1.714286,1.5) -- (1.714286,1.285714) -- (2,1.5) -- cycle;
\fill[teal!91!white] (1.142857,1.071429) -- (1.428571,0.8571429) -- (1.428571,1.071429) -- cycle;
\fill[teal!93!white] (1.428571,0.8571429) -- (1.142857,1.071429) -- (1.142857,0.8571429) -- cycle;
\fill[teal!11!white] (3.142857,0.4285714) -- (3.142857,0.2142857) -- (3.428571,0.2142857) -- cycle;
\fill[teal!10!white] (3.428571,0.4285714) -- (3.142857,0.4285714) -- (3.428571,0.2142857) -- cycle;
\fill[teal!18!white] (2.857143,0.4285714) -- (3.142857,0.4285714) -- (2.857143,0.6428571) -- cycle;
\fill[teal!13!white] (3.142857,0.4285714) -- (3.142857,0.6428571) -- (2.857143,0.6428571) -- cycle;
\fill[teal!15!white] (3.142857,0.2142857) -- (3.142857,0.4285714) -- (2.857143,0.2142857) -- cycle;
\fill[teal!18!white] (3.142857,0.4285714) -- (2.857143,0.4285714) -- (2.857143,0.2142857) -- cycle;
\fill[teal!23!white] (2.571429,0.8571429) -- (2.571429,0.6428571) -- (2.857143,0.8571429) -- cycle;
\fill[teal!21!white] (2.857143,0.8571429) -- (2.571429,0.6428571) -- (2.857143,0.6428571) -- cycle;
\fill[teal!16!white] (2.285714,0) -- (2.285714,0.2142857) -- (2,0) -- cycle;
\fill[teal!25!white] (2.285714,0.2142857) -- (2,0.2142857) -- (2,0) -- cycle;
\fill[teal!40!white] (2.285714,0.2142857) -- (2,0.4285714) -- (2,0.2142857) -- cycle;
\fill[teal!44!white] (2,0.4285714) -- (2.285714,0.2142857) -- (2.285714,0.4285714) -- cycle;
\fill[teal!65!white] (2,0.4285714) -- (1.714286,0.6428571) -- (1.714286,0.4285714) -- cycle;
\fill[teal!66!white] (1.714286,0.6428571) -- (2,0.4285714) -- (2,0.6428571) -- cycle;
\fill[teal!53!white] (1.714286,0.2142857) -- (2,0.4285714) -- (1.714286,0.4285714) -- cycle;
\fill[teal!44!white] (2,0.2142857) -- (2,0.4285714) -- (1.714286,0.2142857) -- cycle;
\fill[teal!20!white] (0.2857143,0.8571429) -- (0,0.8571429) -- (0,0.6428571) -- cycle;
\fill[teal!31!white] (0.2857143,0.8571429) -- (0,0.6428571) -- (0.2857143,0.6428571) -- cycle;
\fill[teal!34!white] (0,1.071429) -- (0.2857143,0.8571429) -- (0.2857143,1.071429) -- cycle;
\fill[teal!21!white] (0,0.8571429) -- (0.2857143,0.8571429) -- (0,1.071429) -- cycle;
\fill[teal!64!white] (0.2857143,1.5) -- (0.5714286,1.285714) -- (0.5714286,1.5) -- cycle;
\fill[teal!57!white] (0.5714286,1.285714) -- (0.2857143,1.5) -- (0.2857143,1.285714) -- cycle;
\fill[teal!51!white] (0.5714286,0.6428571) -- (0.2857143,0.8571429) -- (0.2857143,0.6428571) -- cycle;
\fill[teal!62!white] (0.2857143,0.8571429) -- (0.5714286,0.6428571) -- (0.5714286,0.8571429) -- cycle;
\fill[teal!45!white] (0.2857143,0.4285714) -- (0.5714286,0.6428571) -- (0.2857143,0.6428571) -- cycle;
\fill[teal!49!white] (0.5714286,0.4285714) -- (0.5714286,0.6428571) -- (0.2857143,0.4285714) -- cycle;
\fill[teal!59!white] (0.8571429,0.4285714) -- (0.5714286,0.6428571) -- (0.5714286,0.4285714) -- cycle;
\fill[teal!69!white] (0.5714286,0.6428571) -- (0.8571429,0.4285714) -- (0.8571429,0.6428571) -- cycle;
\fill[teal!48!white] (1.142857,0.4285714) -- (1.142857,0.2142857) -- (1.428571,0.2142857) -- cycle;
\fill[teal!57!white] (1.428571,0.4285714) -- (1.142857,0.4285714) -- (1.428571,0.2142857) -- cycle;
\fill[teal!47!white] (1.142857,0.2142857) -- (1.142857,0.4285714) -- (0.8571429,0.2142857) -- cycle;
\fill[teal!54!white] (1.142857,0.4285714) -- (0.8571429,0.4285714) -- (0.8571429,0.2142857) -- cycle;
\fill[teal!69!white] (0.8571429,0.4285714) -- (1.142857,0.4285714) -- (0.8571429,0.6428571) -- cycle;
\fill[teal!77!white] (1.142857,0.4285714) -- (1.142857,0.6428571) -- (0.8571429,0.6428571) -- cycle;
\fill[teal!56!white] (1.428571,1.285714) -- (1.714286,1.5) -- (1.428571,1.5) -- cycle;
\fill[teal!56!white] (1.428571,1.285714) -- (1.714286,1.285714) -- (1.714286,1.5) -- cycle;
\fill[teal!69!white] (1.142857,1.5) -- (1.428571,1.285714) -- (1.428571,1.5) -- cycle;
\fill[teal!79!white] (1.428571,1.285714) -- (1.142857,1.5) -- (1.142857,1.285714) -- cycle;
\fill[teal!87!white] (1.428571,1.285714) -- (1.142857,1.071429) -- (1.428571,1.071429) -- cycle;
\fill[teal!87!white] (1.142857,1.071429) -- (1.428571,1.285714) -- (1.142857,1.285714) -- cycle;
\fill[teal!70!white] (1.428571,0.6428571) -- (1.428571,0.4285714) -- (1.714286,0.4285714) -- cycle;
\fill[teal!74!white] (1.714286,0.6428571) -- (1.428571,0.6428571) -- (1.714286,0.4285714) -- cycle;
\fill[teal!72!white] (1.428571,0.6428571) -- (1.142857,0.4285714) -- (1.428571,0.4285714) -- cycle;
\fill[teal!78!white] (1.142857,0.4285714) -- (1.428571,0.6428571) -- (1.142857,0.6428571) -- cycle;
\fill[teal!89!white] (1.428571,0.6428571) -- (1.428571,0.8571429) -- (1.142857,0.8571429) -- cycle;
\fill[teal!87!white] (1.142857,0.6428571) -- (1.428571,0.6428571) -- (1.142857,0.8571429) -- cycle;
\fill[teal!53!white] (2,1.285714) -- (1.714286,1.071429) -- (2,1.071429) -- cycle;
\fill[teal!55!white] (1.714286,1.285714) -- (1.714286,1.071429) -- (2,1.285714) -- cycle;
\fill[teal!68!white] (1.428571,1.285714) -- (1.714286,1.071429) -- (1.714286,1.285714) -- cycle;
\fill[teal!78!white] (1.714286,1.071429) -- (1.428571,1.285714) -- (1.428571,1.071429) -- cycle;
\fill[teal!42!white] (2.285714,0.6428571) -- (2.571429,0.8571429) -- (2.285714,0.8571429) -- cycle;
\fill[teal!38!white] (2.285714,0.6428571) -- (2.571429,0.6428571) -- (2.571429,0.8571429) -- cycle;
\fill[teal!54!white] (2,0.8571429) -- (2.285714,0.6428571) -- (2.285714,0.8571429) -- cycle;
\fill[teal!61!white] (2.285714,0.6428571) -- (2,0.8571429) -- (2,0.6428571) -- cycle;
\fill[teal!52!white] (2.285714,0.6428571) -- (2,0.4285714) -- (2.285714,0.4285714) -- cycle;
\fill[teal!59!white] (2,0.4285714) -- (2.285714,0.6428571) -- (2,0.6428571) -- cycle;
\fill[teal!26!white] (2.571429,0.4285714) -- (2.857143,0.4285714) -- (2.857143,0.6428571) -- cycle;
\fill[teal!30!white] (2.571429,0.6428571) -- (2.571429,0.4285714) -- (2.857143,0.6428571) -- cycle;
\fill[teal!40!white] (2.285714,0.6428571) -- (2.571429,0.4285714) -- (2.571429,0.6428571) -- cycle;
\fill[teal!44!white] (2.571429,0.4285714) -- (2.285714,0.6428571) -- (2.285714,0.4285714) -- cycle;
\fill[teal!57!white] (0.2857143,0.8571429) -- (0.5714286,1.071429) -- (0.2857143,1.071429) -- cycle;
\fill[teal!66!white] (0.5714286,1.071429) -- (0.2857143,0.8571429) -- (0.5714286,0.8571429) -- cycle;
\fill[teal!58!white] (0.5714286,1.071429) -- (0.2857143,1.285714) -- (0.2857143,1.071429) -- cycle;
\fill[teal!67!white] (0.5714286,1.071429) -- (0.5714286,1.285714) -- (0.2857143,1.285714) -- cycle;
\fill[teal!71!white] (0.8571429,1.5) -- (0.5714286,1.714286) -- (0.5714286,1.5) -- cycle;
\fill[teal!69!white] (0.8571429,1.5) -- (0.8571429,1.714286) -- (0.5714286,1.714286) -- cycle;
\fill[teal!68!white] (0.8571429,1.714286) -- (0.8571429,1.5) -- (1.142857,1.714286) -- cycle;
\fill[teal!71!white] (0.8571429,1.5) -- (1.142857,1.5) -- (1.142857,1.714286) -- cycle;
\fill[teal!74!white] (1.714286,0.8571429) -- (1.714286,0.6428571) -- (2,0.6428571) -- cycle;
\fill[teal!70!white] (2,0.8571429) -- (1.714286,0.8571429) -- (2,0.6428571) -- cycle;
\fill[teal!67!white] (1.714286,0.8571429) -- (2,0.8571429) -- (2,1.071429) -- cycle;
\fill[teal!69!white] (1.714286,1.071429) -- (1.714286,0.8571429) -- (2,1.071429) -- cycle;
\fill[teal!81!white] (1.714286,0.8571429) -- (1.428571,0.6428571) -- (1.714286,0.6428571) -- cycle;
\fill[teal!85!white] (1.428571,0.6428571) -- (1.714286,0.8571429) -- (1.428571,0.8571429) -- cycle;
\fill[teal!86!white] (1.428571,0.8571429) -- (1.714286,0.8571429) -- (1.428571,1.071429) -- cycle;
\fill[teal!80!white] (1.714286,0.8571429) -- (1.714286,1.071429) -- (1.428571,1.071429) -- cycle;
\fill[teal!14!white] (2.571429,0) -- (2.571429,0.2142857) -- (2.285714,0) -- cycle;
\fill[teal!22!white] (2.571429,0.2142857) -- (2.285714,0.2142857) -- (2.285714,0) -- cycle;
\fill[teal!33!white] (2.285714,0.2142857) -- (2.571429,0.2142857) -- (2.285714,0.4285714) -- cycle;
\fill[teal!35!white] (2.571429,0.2142857) -- (2.571429,0.4285714) -- (2.285714,0.4285714) -- cycle;
\fill[teal!16!white] (2.857143,0.2142857) -- (2.571429,0.2142857) -- (2.857143,0) -- cycle;
\fill[teal!13!white] (2.571429,0.2142857) -- (2.571429,0) -- (2.857143,0) -- cycle;
\fill[teal!22!white] (2.857143,0.4285714) -- (2.571429,0.2142857) -- (2.857143,0.2142857) -- cycle;
\fill[teal!28!white] (2.571429,0.4285714) -- (2.571429,0.2142857) -- (2.857143,0.4285714) -- cycle;
\fill[teal!85!white] (1.142857,0.6428571) -- (0.8571429,0.8571429) -- (0.8571429,0.6428571) -- cycle;
\fill[teal!90!white] (0.8571429,0.8571429) -- (1.142857,0.6428571) -- (1.142857,0.8571429) -- cycle;
\fill[teal!76!white] (0.5714286,0.6428571) -- (0.8571429,0.8571429) -- (0.5714286,0.8571429) -- cycle;
\fill[teal!78!white] (0.8571429,0.8571429) -- (0.5714286,0.6428571) -- (0.8571429,0.6428571) -- cycle;
\fill[teal!85!white] (1.142857,1.5) -- (0.8571429,1.285714) -- (1.142857,1.285714) -- cycle;
\fill[teal!82!white] (0.8571429,1.5) -- (0.8571429,1.285714) -- (1.142857,1.5) -- cycle;
\fill[teal!79!white] (0.5714286,1.285714) -- (0.8571429,1.285714) -- (0.5714286,1.5) -- cycle;
\fill[teal!80!white] (0.8571429,1.285714) -- (0.8571429,1.5) -- (0.5714286,1.5) -- cycle;
\fill[teal!94!white] (1.142857,1.071429) -- (0.8571429,1.071429) -- (1.142857,0.8571429) -- cycle;
\fill[teal!92!white] (0.8571429,1.071429) -- (0.8571429,0.8571429) -- (1.142857,0.8571429) -- cycle;
\fill[teal!93!white] (0.8571429,1.071429) -- (1.142857,1.071429) -- (1.142857,1.285714) -- cycle;
\fill[teal!91!white] (0.8571429,1.285714) -- (0.8571429,1.071429) -- (1.142857,1.285714) -- cycle;
\fill[teal!82!white] (0.8571429,1.071429) -- (0.5714286,1.071429) -- (0.5714286,0.8571429) -- cycle;
\fill[teal!86!white] (0.8571429,0.8571429) -- (0.8571429,1.071429) -- (0.5714286,0.8571429) -- cycle;
\fill[teal!82!white] (0.5714286,1.071429) -- (0.8571429,1.071429) -- (0.5714286,1.285714) -- cycle;
\fill[teal!86!white] (0.8571429,1.071429) -- (0.8571429,1.285714) -- (0.5714286,1.285714) -- cycle;
\draw[teal,thick] (0,0) -- (4,0) -- (0,3) -- cycle;
\draw[magenta!70,thin] (1.085411,1.002963) -- (3.6,2.35) -- (4.1,2.35);\filldraw[fill=magenta,draw=white,line width=.45pt] (1.085411,1.002963) circle(1.5pt) node[above right,font=\scriptsize,text=magenta,fill=white,inner sep=.6pt] {};
\node[anchor=west,text=magenta,font=\small] at(4.12,2.35){H0.01};\draw[gold!70,thin] (1,1) -- (3.6,1.8) -- (4.1,1.8);\filldraw[fill=gold,draw=white,line width=.45pt] (1,1) circle(1.5pt) node[above right,font=\scriptsize,text=gold,fill=white,inner sep=.6pt] {};
\node[anchor=west,text=gold,font=\small] at(4.12,1.8){$I$};\draw[navy!70,thin] (1.333333,1) -- (3.6,1.25) -- (4.1,1.25);\filldraw[fill=navy,draw=white,line width=.45pt] (1.333333,1) circle(1.5pt) node[above right,font=\scriptsize,text=navy,fill=white,inner sep=.6pt] {};
\node[anchor=west,text=navy,font=\small] at(4.12,1.25){$G$};\end{tikzpicture}
\caption[Early archived heat field]{Early archived heat field. Each field is normalized separately; this removes the cooling amplitude.}\label{fig:heat-H0.01}
\end{figure}
\begin{figure}[htbp]\centering
\begin{tikzpicture}[x=1.7cm,y=1.7cm]
\fill[teal!9!white] (3.142857,0.6428571) -- (2.857143,0.8571429) -- (2.857143,0.6428571) -- cycle;
\fill[teal!16!white] (1.714286,1.714286) -- (1.428571,1.928571) -- (1.428571,1.714286) -- cycle;
\fill[teal!16!white] (1.714286,1.714286) -- (1.714286,1.5) -- (2,1.5) -- cycle;
\fill[teal!9!white] (0.5714286,2.571429) -- (0.2857143,2.785714) -- (0.2857143,2.571429) -- cycle;
\fill[teal!13!white] (2.571429,1.071429) -- (2.285714,1.285714) -- (2.285714,1.071429) -- cycle;
\fill[teal!8!white] (3.714286,0.2142857) -- (3.428571,0.4285714) -- (3.428571,0.2142857) -- cycle;
\fill[teal!8!white] (3.714286,0.2142857) -- (3.714286,0) -- (4,0) -- cycle;
\fill[teal!13!white] (0.8571429,2.142857) -- (1.142857,2.142857) -- (0.8571429,2.357143) -- cycle;
\fill[teal!15!white] (1.142857,2.142857) -- (1.142857,1.928571) -- (1.428571,1.928571) -- cycle;
\fill[teal!8!white] (0,2.785714) -- (0.2857143,2.785714) -- (0,3) -- cycle;
\fill[teal!11!white] (0.5714286,2.571429) -- (0.5714286,2.357143) -- (0.8571429,2.357143) -- cycle;
\fill[teal!15!white] (2,1.285714) -- (2.285714,1.285714) -- (2,1.5) -- cycle;
\fill[teal!11!white] (2.571429,1.071429) -- (2.571429,0.8571429) -- (2.857143,0.8571429) -- cycle;
\fill[teal!8!white] (3.142857,0.4285714) -- (3.428571,0.4285714) -- (3.142857,0.6428571) -- cycle;
\fill[teal!9!white] (0.2857143,0) -- (0,0.2142857) -- (0,0) -- cycle;
\fill[teal!13!white] (0,0.2142857) -- (0.2857143,0) -- (0.2857143,0.2142857) -- cycle;
\fill[teal!14!white] (0.5714286,0.2142857) -- (0.2857143,0) -- (0.5714286,0) -- cycle;
\fill[teal!18!white] (0.2857143,0) -- (0.5714286,0.2142857) -- (0.2857143,0.2142857) -- cycle;
\fill[teal!25!white] (1.714286,1.5) -- (1.714286,1.714286) -- (1.428571,1.714286) -- cycle;
\fill[teal!40!white] (1.428571,1.5) -- (1.714286,1.5) -- (1.428571,1.714286) -- cycle;
\fill[teal!15!white] (0,1.928571) -- (0,1.714286) -- (0.2857143,1.928571) -- cycle;
\fill[teal!23!white] (0,1.714286) -- (0.2857143,1.714286) -- (0.2857143,1.928571) -- cycle;
\fill[teal!10!white] (0.2857143,2.357143) -- (0,2.571429) -- (0,2.357143) -- cycle;
\fill[teal!10!white] (0,2.571429) -- (0.2857143,2.357143) -- (0.2857143,2.571429) -- cycle;
\fill[teal!11!white] (0,2.142857) -- (0.2857143,2.357143) -- (0,2.357143) -- cycle;
\fill[teal!15!white] (0.2857143,2.357143) -- (0,2.142857) -- (0.2857143,2.142857) -- cycle;
\fill[teal!14!white] (0,2.142857) -- (0,1.928571) -- (0.2857143,1.928571) -- cycle;
\fill[teal!17!white] (0.2857143,2.142857) -- (0,2.142857) -- (0.2857143,1.928571) -- cycle;
\fill[teal!34!white] (0.2857143,1.714286) -- (0.5714286,1.928571) -- (0.2857143,1.928571) -- cycle;
\fill[teal!42!white] (0.5714286,1.714286) -- (0.5714286,1.928571) -- (0.2857143,1.714286) -- cycle;
\fill[teal!18!white] (0.2857143,1.5) -- (0,1.714286) -- (0,1.5) -- cycle;
\fill[teal!26!white] (0,1.714286) -- (0.2857143,1.5) -- (0.2857143,1.714286) -- cycle;
\fill[teal!42!white] (0.2857143,1.5) -- (0.5714286,1.714286) -- (0.2857143,1.714286) -- cycle;
\fill[teal!52!white] (0.5714286,1.714286) -- (0.2857143,1.5) -- (0.5714286,1.5) -- cycle;
\fill[teal!19!white] (0,1.285714) -- (0.2857143,1.5) -- (0,1.5) -- cycle;
\fill[teal!30!white] (0.2857143,1.5) -- (0,1.285714) -- (0.2857143,1.285714) -- cycle;
\fill[teal!20!white] (0,1.285714) -- (0,1.071429) -- (0.2857143,1.071429) -- cycle;
\fill[teal!32!white] (0.2857143,1.285714) -- (0,1.285714) -- (0.2857143,1.071429) -- cycle;
\fill[teal!8!white] (3.428571,0) -- (3.714286,0) -- (3.428571,0.2142857) -- cycle;
\fill[teal!8!white] (3.714286,0) -- (3.714286,0.2142857) -- (3.428571,0.2142857) -- cycle;
\fill[teal!17!white] (1.428571,0) -- (1.714286,0) -- (1.428571,0.2142857) -- cycle;
\fill[teal!25!white] (1.714286,0) -- (1.714286,0.2142857) -- (1.428571,0.2142857) -- cycle;
\fill[teal!16!white] (0.8571429,0) -- (0.5714286,0.2142857) -- (0.5714286,0) -- cycle;
\fill[teal!25!white] (0.5714286,0.2142857) -- (0.8571429,0) -- (0.8571429,0.2142857) -- cycle;
\fill[teal!18!white] (1.142857,0.2142857) -- (1.142857,0) -- (1.428571,0) -- cycle;
\fill[teal!27!white] (1.142857,0.2142857) -- (1.428571,0) -- (1.428571,0.2142857) -- cycle;
\fill[teal!18!white] (1.142857,0.2142857) -- (0.8571429,0) -- (1.142857,0) -- cycle;
\fill[teal!27!white] (0.8571429,0) -- (1.142857,0.2142857) -- (0.8571429,0.2142857) -- cycle;
\fill[teal!37!white] (1.142857,1.928571) -- (1.142857,1.714286) -- (1.428571,1.714286) -- cycle;
\fill[teal!23!white] (1.428571,1.928571) -- (1.142857,1.928571) -- (1.428571,1.714286) -- cycle;
\fill[teal!20!white] (0.8571429,2.142857) -- (1.142857,1.928571) -- (1.142857,2.142857) -- cycle;
\fill[teal!30!white] (1.142857,1.928571) -- (0.8571429,2.142857) -- (0.8571429,1.928571) -- cycle;
\fill[teal!9!white] (0,2.785714) -- (0,2.571429) -- (0.2857143,2.571429) -- cycle;
\fill[teal!9!white] (0.2857143,2.785714) -- (0,2.785714) -- (0.2857143,2.571429) -- cycle;
\fill[teal!47!white] (1.142857,1.928571) -- (0.8571429,1.714286) -- (1.142857,1.714286) -- cycle;
\fill[teal!43!white] (0.8571429,1.714286) -- (1.142857,1.928571) -- (0.8571429,1.928571) -- cycle;
\fill[teal!49!white] (0.8571429,1.714286) -- (0.5714286,1.928571) -- (0.5714286,1.714286) -- cycle;
\fill[teal!45!white] (0.5714286,1.928571) -- (0.8571429,1.714286) -- (0.8571429,1.928571) -- cycle;
\fill[teal!13!white] (0.2857143,2.357143) -- (0.5714286,2.357143) -- (0.2857143,2.571429) -- cycle;
\fill[teal!11!white] (0.5714286,2.357143) -- (0.5714286,2.571429) -- (0.2857143,2.571429) -- cycle;
\fill[teal!17!white] (0,0.6428571) -- (0,0.4285714) -- (0.2857143,0.6428571) -- cycle;
\fill[teal!24!white] (0,0.4285714) -- (0.2857143,0.4285714) -- (0.2857143,0.6428571) -- cycle;
\fill[teal!13!white] (0,0.4285714) -- (0,0.2142857) -- (0.2857143,0.2142857) -- cycle;
\fill[teal!20!white] (0.2857143,0.4285714) -- (0,0.4285714) -- (0.2857143,0.2142857) -- cycle;
\fill[teal!20!white] (2.285714,1.285714) -- (2,1.285714) -- (2.285714,1.071429) -- cycle;
\fill[teal!32!white] (2,1.285714) -- (2,1.071429) -- (2.285714,1.071429) -- cycle;
\fill[teal!40!white] (2,1.071429) -- (2,0.8571429) -- (2.285714,1.071429) -- cycle;
\fill[teal!36!white] (2,0.8571429) -- (2.285714,0.8571429) -- (2.285714,1.071429) -- cycle;
\fill[teal!24!white] (2.285714,0.8571429) -- (2.571429,0.8571429) -- (2.285714,1.071429) -- cycle;
\fill[teal!16!white] (2.571429,0.8571429) -- (2.571429,1.071429) -- (2.285714,1.071429) -- cycle;
\fill[teal!49!white] (1.428571,0.4285714) -- (1.714286,0.2142857) -- (1.714286,0.4285714) -- cycle;
\fill[teal!43!white] (1.714286,0.2142857) -- (1.428571,0.4285714) -- (1.428571,0.2142857) -- cycle;
\fill[teal!8!white] (3.142857,0.2142857) -- (3.428571,0) -- (3.428571,0.2142857) -- cycle;
\fill[teal!8!white] (3.142857,0.2142857) -- (3.142857,0) -- (3.428571,0) -- cycle;
\fill[teal!9!white] (3.142857,0) -- (3.142857,0.2142857) -- (2.857143,0) -- cycle;
\fill[teal!10!white] (3.142857,0.2142857) -- (2.857143,0.2142857) -- (2.857143,0) -- cycle;
\fill[teal!23!white] (1.714286,0) -- (2,0.2142857) -- (1.714286,0.2142857) -- cycle;
\fill[teal!15!white] (2,0.2142857) -- (1.714286,0) -- (2,0) -- cycle;
\fill[teal!31!white] (0.5714286,0.2142857) -- (0.5714286,0.4285714) -- (0.2857143,0.2142857) -- cycle;
\fill[teal!31!white] (0.5714286,0.4285714) -- (0.2857143,0.4285714) -- (0.2857143,0.2142857) -- cycle;
\fill[teal!41!white] (0.8571429,0.4285714) -- (0.5714286,0.2142857) -- (0.8571429,0.2142857) -- cycle;
\fill[teal!45!white] (0.8571429,0.4285714) -- (0.5714286,0.4285714) -- (0.5714286,0.2142857) -- cycle;
\fill[teal!54!white] (1.142857,1.5) -- (1.428571,1.5) -- (1.428571,1.714286) -- cycle;
\fill[teal!52!white] (1.142857,1.714286) -- (1.142857,1.5) -- (1.428571,1.714286) -- cycle;
\fill[teal!25!white] (0.5714286,2.142857) -- (0.2857143,2.142857) -- (0.2857143,1.928571) -- cycle;
\fill[teal!31!white] (0.5714286,1.928571) -- (0.5714286,2.142857) -- (0.2857143,1.928571) -- cycle;
\fill[teal!20!white] (0.5714286,2.142857) -- (0.2857143,2.357143) -- (0.2857143,2.142857) -- cycle;
\fill[teal!18!white] (0.5714286,2.142857) -- (0.5714286,2.357143) -- (0.2857143,2.357143) -- cycle;
\fill[teal!29!white] (0.8571429,2.142857) -- (0.5714286,2.142857) -- (0.8571429,1.928571) -- cycle;
\fill[teal!34!white] (0.5714286,2.142857) -- (0.5714286,1.928571) -- (0.8571429,1.928571) -- cycle;
\fill[teal!19!white] (0.5714286,2.142857) -- (0.8571429,2.142857) -- (0.8571429,2.357143) -- cycle;
\fill[teal!17!white] (0.5714286,2.357143) -- (0.5714286,2.142857) -- (0.8571429,2.357143) -- cycle;
\fill[teal!30!white] (1.714286,1.285714) -- (2,1.285714) -- (2,1.5) -- cycle;
\fill[teal!31!white] (1.714286,1.5) -- (1.714286,1.285714) -- (2,1.5) -- cycle;
\fill[teal!90!white] (1.142857,1.071429) -- (1.428571,0.8571429) -- (1.428571,1.071429) -- cycle;
\fill[teal!92!white] (1.428571,0.8571429) -- (1.142857,1.071429) -- (1.142857,0.8571429) -- cycle;
\fill[teal!9!white] (3.142857,0.4285714) -- (3.142857,0.2142857) -- (3.428571,0.2142857) -- cycle;
\fill[teal!8!white] (3.428571,0.4285714) -- (3.142857,0.4285714) -- (3.428571,0.2142857) -- cycle;
\fill[teal!11!white] (2.857143,0.4285714) -- (3.142857,0.4285714) -- (2.857143,0.6428571) -- cycle;
\fill[teal!10!white] (3.142857,0.4285714) -- (3.142857,0.6428571) -- (2.857143,0.6428571) -- cycle;
\fill[teal!10!white] (3.142857,0.2142857) -- (3.142857,0.4285714) -- (2.857143,0.2142857) -- cycle;
\fill[teal!11!white] (3.142857,0.4285714) -- (2.857143,0.4285714) -- (2.857143,0.2142857) -- cycle;
\fill[teal!15!white] (2.571429,0.8571429) -- (2.571429,0.6428571) -- (2.857143,0.8571429) -- cycle;
\fill[teal!14!white] (2.857143,0.8571429) -- (2.571429,0.6428571) -- (2.857143,0.6428571) -- cycle;
\fill[teal!13!white] (2.285714,0) -- (2.285714,0.2142857) -- (2,0) -- cycle;
\fill[teal!19!white] (2.285714,0.2142857) -- (2,0.2142857) -- (2,0) -- cycle;
\fill[teal!29!white] (2.285714,0.2142857) -- (2,0.4285714) -- (2,0.2142857) -- cycle;
\fill[teal!30!white] (2,0.4285714) -- (2.285714,0.2142857) -- (2.285714,0.4285714) -- cycle;
\fill[teal!54!white] (2,0.4285714) -- (1.714286,0.6428571) -- (1.714286,0.4285714) -- cycle;
\fill[teal!53!white] (1.714286,0.6428571) -- (2,0.4285714) -- (2,0.6428571) -- cycle;
\fill[teal!43!white] (1.714286,0.2142857) -- (2,0.4285714) -- (1.714286,0.4285714) -- cycle;
\fill[teal!34!white] (2,0.2142857) -- (2,0.4285714) -- (1.714286,0.2142857) -- cycle;
\fill[teal!19!white] (0.2857143,0.8571429) -- (0,0.8571429) -- (0,0.6428571) -- cycle;
\fill[teal!29!white] (0.2857143,0.8571429) -- (0,0.6428571) -- (0.2857143,0.6428571) -- cycle;
\fill[teal!32!white] (0,1.071429) -- (0.2857143,0.8571429) -- (0.2857143,1.071429) -- cycle;
\fill[teal!20!white] (0,0.8571429) -- (0.2857143,0.8571429) -- (0,1.071429) -- cycle;
\fill[teal!57!white] (0.2857143,1.5) -- (0.5714286,1.285714) -- (0.5714286,1.5) -- cycle;
\fill[teal!51!white] (0.5714286,1.285714) -- (0.2857143,1.5) -- (0.2857143,1.285714) -- cycle;
\fill[teal!48!white] (0.5714286,0.6428571) -- (0.2857143,0.8571429) -- (0.2857143,0.6428571) -- cycle;
\fill[teal!59!white] (0.2857143,0.8571429) -- (0.5714286,0.6428571) -- (0.5714286,0.8571429) -- cycle;
\fill[teal!42!white] (0.2857143,0.4285714) -- (0.5714286,0.6428571) -- (0.2857143,0.6428571) -- cycle;
\fill[teal!46!white] (0.5714286,0.4285714) -- (0.5714286,0.6428571) -- (0.2857143,0.4285714) -- cycle;
\fill[teal!57!white] (0.8571429,0.4285714) -- (0.5714286,0.6428571) -- (0.5714286,0.4285714) -- cycle;
\fill[teal!67!white] (0.5714286,0.6428571) -- (0.8571429,0.4285714) -- (0.8571429,0.6428571) -- cycle;
\fill[teal!46!white] (1.142857,0.4285714) -- (1.142857,0.2142857) -- (1.428571,0.2142857) -- cycle;
\fill[teal!54!white] (1.428571,0.4285714) -- (1.142857,0.4285714) -- (1.428571,0.2142857) -- cycle;
\fill[teal!46!white] (1.142857,0.2142857) -- (1.142857,0.4285714) -- (0.8571429,0.2142857) -- cycle;
\fill[teal!53!white] (1.142857,0.4285714) -- (0.8571429,0.4285714) -- (0.8571429,0.2142857) -- cycle;
\fill[teal!67!white] (0.8571429,0.4285714) -- (1.142857,0.4285714) -- (0.8571429,0.6428571) -- cycle;
\fill[teal!75!white] (1.142857,0.4285714) -- (1.142857,0.6428571) -- (0.8571429,0.6428571) -- cycle;
\fill[teal!55!white] (1.428571,1.285714) -- (1.714286,1.5) -- (1.428571,1.5) -- cycle;
\fill[teal!54!white] (1.428571,1.285714) -- (1.714286,1.285714) -- (1.714286,1.5) -- cycle;
\fill[teal!68!white] (1.142857,1.5) -- (1.428571,1.285714) -- (1.428571,1.5) -- cycle;
\fill[teal!78!white] (1.428571,1.285714) -- (1.142857,1.5) -- (1.142857,1.285714) -- cycle;
\fill[teal!86!white] (1.428571,1.285714) -- (1.142857,1.071429) -- (1.428571,1.071429) -- cycle;
\fill[teal!87!white] (1.142857,1.071429) -- (1.428571,1.285714) -- (1.142857,1.285714) -- cycle;
\fill[teal!64!white] (1.428571,0.6428571) -- (1.428571,0.4285714) -- (1.714286,0.4285714) -- cycle;
\fill[teal!66!white] (1.714286,0.6428571) -- (1.428571,0.6428571) -- (1.714286,0.4285714) -- cycle;
\fill[teal!69!white] (1.428571,0.6428571) -- (1.142857,0.4285714) -- (1.428571,0.4285714) -- cycle;
\fill[teal!76!white] (1.142857,0.4285714) -- (1.428571,0.6428571) -- (1.142857,0.6428571) -- cycle;
\fill[teal!87!white] (1.428571,0.6428571) -- (1.428571,0.8571429) -- (1.142857,0.8571429) -- cycle;
\fill[teal!86!white] (1.142857,0.6428571) -- (1.428571,0.6428571) -- (1.142857,0.8571429) -- cycle;
\fill[teal!47!white] (2,1.285714) -- (1.714286,1.071429) -- (2,1.071429) -- cycle;
\fill[teal!50!white] (1.714286,1.285714) -- (1.714286,1.071429) -- (2,1.285714) -- cycle;
\fill[teal!65!white] (1.428571,1.285714) -- (1.714286,1.071429) -- (1.714286,1.285714) -- cycle;
\fill[teal!76!white] (1.714286,1.071429) -- (1.428571,1.285714) -- (1.428571,1.071429) -- cycle;
\fill[teal!28!white] (2.285714,0.6428571) -- (2.571429,0.8571429) -- (2.285714,0.8571429) -- cycle;
\fill[teal!25!white] (2.285714,0.6428571) -- (2.571429,0.6428571) -- (2.571429,0.8571429) -- cycle;
\fill[teal!40!white] (2,0.8571429) -- (2.285714,0.6428571) -- (2.285714,0.8571429) -- cycle;
\fill[teal!46!white] (2.285714,0.6428571) -- (2,0.8571429) -- (2,0.6428571) -- cycle;
\fill[teal!36!white] (2.285714,0.6428571) -- (2,0.4285714) -- (2.285714,0.4285714) -- cycle;
\fill[teal!43!white] (2,0.4285714) -- (2.285714,0.6428571) -- (2,0.6428571) -- cycle;
\fill[teal!15!white] (2.571429,0.4285714) -- (2.857143,0.4285714) -- (2.857143,0.6428571) -- cycle;
\fill[teal!18!white] (2.571429,0.6428571) -- (2.571429,0.4285714) -- (2.857143,0.6428571) -- cycle;
\fill[teal!25!white] (2.285714,0.6428571) -- (2.571429,0.4285714) -- (2.571429,0.6428571) -- cycle;
\fill[teal!28!white] (2.571429,0.4285714) -- (2.285714,0.6428571) -- (2.285714,0.4285714) -- cycle;
\fill[teal!54!white] (0.2857143,0.8571429) -- (0.5714286,1.071429) -- (0.2857143,1.071429) -- cycle;
\fill[teal!62!white] (0.5714286,1.071429) -- (0.2857143,0.8571429) -- (0.5714286,0.8571429) -- cycle;
\fill[teal!54!white] (0.5714286,1.071429) -- (0.2857143,1.285714) -- (0.2857143,1.071429) -- cycle;
\fill[teal!63!white] (0.5714286,1.071429) -- (0.5714286,1.285714) -- (0.2857143,1.285714) -- cycle;
\fill[teal!63!white] (0.8571429,1.5) -- (0.5714286,1.714286) -- (0.5714286,1.5) -- cycle;
\fill[teal!61!white] (0.8571429,1.5) -- (0.8571429,1.714286) -- (0.5714286,1.714286) -- cycle;
\fill[teal!62!white] (0.8571429,1.714286) -- (0.8571429,1.5) -- (1.142857,1.714286) -- cycle;
\fill[teal!67!white] (0.8571429,1.5) -- (1.142857,1.5) -- (1.142857,1.714286) -- cycle;
\fill[teal!63!white] (1.714286,0.8571429) -- (1.714286,0.6428571) -- (2,0.6428571) -- cycle;
\fill[teal!58!white] (2,0.8571429) -- (1.714286,0.8571429) -- (2,0.6428571) -- cycle;
\fill[teal!57!white] (1.714286,0.8571429) -- (2,0.8571429) -- (2,1.071429) -- cycle;
\fill[teal!62!white] (1.714286,1.071429) -- (1.714286,0.8571429) -- (2,1.071429) -- cycle;
\fill[teal!74!white] (1.714286,0.8571429) -- (1.428571,0.6428571) -- (1.714286,0.6428571) -- cycle;
\fill[teal!80!white] (1.428571,0.6428571) -- (1.714286,0.8571429) -- (1.428571,0.8571429) -- cycle;
\fill[teal!82!white] (1.428571,0.8571429) -- (1.714286,0.8571429) -- (1.428571,1.071429) -- cycle;
\fill[teal!75!white] (1.714286,0.8571429) -- (1.714286,1.071429) -- (1.428571,1.071429) -- cycle;
\fill[teal!11!white] (2.571429,0) -- (2.571429,0.2142857) -- (2.285714,0) -- cycle;
\fill[teal!15!white] (2.571429,0.2142857) -- (2.285714,0.2142857) -- (2.285714,0) -- cycle;
\fill[teal!22!white] (2.285714,0.2142857) -- (2.571429,0.2142857) -- (2.285714,0.4285714) -- cycle;
\fill[teal!21!white] (2.571429,0.2142857) -- (2.571429,0.4285714) -- (2.285714,0.4285714) -- cycle;
\fill[teal!11!white] (2.857143,0.2142857) -- (2.571429,0.2142857) -- (2.857143,0) -- cycle;
\fill[teal!10!white] (2.571429,0.2142857) -- (2.571429,0) -- (2.857143,0) -- cycle;
\fill[teal!14!white] (2.857143,0.4285714) -- (2.571429,0.2142857) -- (2.857143,0.2142857) -- cycle;
\fill[teal!17!white] (2.571429,0.4285714) -- (2.571429,0.2142857) -- (2.857143,0.4285714) -- cycle;
\fill[teal!84!white] (1.142857,0.6428571) -- (0.8571429,0.8571429) -- (0.8571429,0.6428571) -- cycle;
\fill[teal!89!white] (0.8571429,0.8571429) -- (1.142857,0.6428571) -- (1.142857,0.8571429) -- cycle;
\fill[teal!74!white] (0.5714286,0.6428571) -- (0.8571429,0.8571429) -- (0.5714286,0.8571429) -- cycle;
\fill[teal!76!white] (0.8571429,0.8571429) -- (0.5714286,0.6428571) -- (0.8571429,0.6428571) -- cycle;
\fill[teal!83!white] (1.142857,1.5) -- (0.8571429,1.285714) -- (1.142857,1.285714) -- cycle;
\fill[teal!78!white] (0.8571429,1.5) -- (0.8571429,1.285714) -- (1.142857,1.5) -- cycle;
\fill[teal!73!white] (0.5714286,1.285714) -- (0.8571429,1.285714) -- (0.5714286,1.5) -- cycle;
\fill[teal!75!white] (0.8571429,1.285714) -- (0.8571429,1.5) -- (0.5714286,1.5) -- cycle;
\fill[teal!94!white] (1.142857,1.071429) -- (0.8571429,1.071429) -- (1.142857,0.8571429) -- cycle;
\fill[teal!92!white] (0.8571429,1.071429) -- (0.8571429,0.8571429) -- (1.142857,0.8571429) -- cycle;
\fill[teal!92!white] (0.8571429,1.071429) -- (1.142857,1.071429) -- (1.142857,1.285714) -- cycle;
\fill[teal!89!white] (0.8571429,1.285714) -- (0.8571429,1.071429) -- (1.142857,1.285714) -- cycle;
\fill[teal!79!white] (0.8571429,1.071429) -- (0.5714286,1.071429) -- (0.5714286,0.8571429) -- cycle;
\fill[teal!84!white] (0.8571429,0.8571429) -- (0.8571429,1.071429) -- (0.5714286,0.8571429) -- cycle;
\fill[teal!79!white] (0.5714286,1.071429) -- (0.8571429,1.071429) -- (0.5714286,1.285714) -- cycle;
\fill[teal!83!white] (0.8571429,1.071429) -- (0.8571429,1.285714) -- (0.5714286,1.285714) -- cycle;
\draw[teal,thick] (0,0) -- (4,0) -- (0,3) -- cycle;
\draw[magenta!70,thin] (1.095451,0.9975085) -- (3.6,2.35) -- (4.1,2.35);\filldraw[fill=magenta,draw=white,line width=.45pt] (1.095451,0.9975085) circle(1.5pt) node[above right,font=\scriptsize,text=magenta,fill=white,inner sep=.6pt] {};
\node[anchor=west,text=magenta,font=\small] at(4.12,2.35){H1.0};\draw[gold!70,thin] (1,1) -- (3.6,1.8) -- (4.1,1.8);\filldraw[fill=gold,draw=white,line width=.45pt] (1,1) circle(1.5pt) node[above right,font=\scriptsize,text=gold,fill=white,inner sep=.6pt] {};
\node[anchor=west,text=gold,font=\small] at(4.12,1.8){$I$};\draw[navy!70,thin] (1.333333,1) -- (3.6,1.25) -- (4.1,1.25);\filldraw[fill=navy,draw=white,line width=.45pt] (1.333333,1) circle(1.5pt) node[above right,font=\scriptsize,text=navy,fill=white,inner sep=.6pt] {};
\node[anchor=west,text=navy,font=\small] at(4.12,1.25){$G$};\end{tikzpicture}
\caption[Late archived heat field, dominated by the ground mode]{Late archived heat field, dominated by the ground mode. Each field is normalized separately; this removes the cooling amplitude.}\label{fig:heat-H1.0}
\end{figure}

\paragraph{Demonstration and investigation.} Pause the animation near \ensuremath{\Delta}=.01 and ask where the boundary layer is thickest. Resume to the late field and ask students to recognize the ground-state shape. Explain that normalized comparison slides may remove amplitude decay, and that this animation normalizes each frame to its own peak.


\section{Heat fingerprints}\label{sec:slide66}

A heat center depends on both dimensionless time and triangle shape. A fingerprint fixes the physical time parameter and varies a normalized triangle height. A center trajectory instead fixes the triangle and varies time. These are different curves. The diagram fixes \ensuremath{\Delta}=.03 and compares the isosceles fingerprint with its retained ETC neighbors and I, G and W. The derivative along the family measures shape response only in that special direction. It cannot replace the full three-vertex response test.


For a fixed \ensuremath{\Delta}, increasing a triangle height also changes $L^2$ in the physical time $\kappa t$=\ensuremath{\Delta}$L^2$. Holding the dimensional time fixed instead would define a different comparison.


\begin{figure}[htbp]\centering
\begin{tikzpicture}\begin{axis}[width=.95\linewidth,height=.56\linewidth,grid=major,grid style={black!10},xlabel={family parameter $x$},ylabel={$c_y$},tick label style={font=\scriptsize},label style={font=\small},legend style={font=\scriptsize,at={(0.5,-0.24)},anchor=north,legend columns=3,draw=none},scaled ticks=false]
\addplot[navy,thick,no marks] coordinates {(0.2,0.09530699) (0.228706,0.1082435) (0.2615321,0.1228303) (0.2990698,0.1392271) (0.3419952,0.1575953) (0.3910817,0.1782479) (0.4472136,0.2008468) (0.5114021,0.2258402) (0.5848035,0.2534594) (0.6687403,0.2836914) (0.7647245,0.316471) (0.8744853,0.3516631) (1,0.3897689) (1.14353,0.4308941) (1.30766,0.4749474) (1.495349,0.5220689) (1.709976,0.5723923) (1.732051,0.5773503) (1.955409,0.6256606) (2.236068,0.6817952) (2.55701,0.7405712) (2.924018,0.8016695) (3.343702,0.8652768) (3.823622,0.9306455) (4.372426,0.9967819) (5,1.065437)};
\addlegendentry{H0.03}
\addplot[teal,thick,no marks] coordinates {(0.2,0.09901951) (0.228706,0.1128955) (0.2615321,0.1286033) (0.2990698,0.1463328) (0.3419952,0.1662702) (0.3910817,0.1885865) (0.4472136,0.2134218) (0.5114021,0.2408662) (0.5848035,0.2709374) (0.6687403,0.3035587) (0.7647245,0.33854) (0.8744853,0.3755687) (1,0.4142136) (1.14353,0.4539441) (1.30766,0.4941653) (1.495349,0.5342616) (1.709976,0.5736416) (1.732051,0.5773503) (1.955409,0.6117774) (2.236068,0.6482315) (2.55701,0.682671) (2.924018,0.7148684) (3.343702,0.744694) (3.823622,0.7721018) (4.372426,0.7971139) (5,0.8198039)};
\addlegendentry{X1}
\addplot[magenta,thick,no marks] coordinates {(0.2,0.06666667) (0.228706,0.07623532) (0.2615321,0.08717737) (0.2990698,0.09968992) (0.3419952,0.1139984) (0.3910817,0.1303606) (0.4472136,0.1490712) (0.5114021,0.1704674) (0.5848035,0.1949345) (0.6687403,0.2229134) (0.7647245,0.2549082) (0.8744853,0.2914951) (1,0.3333333) (1.14353,0.3811766) (1.30766,0.4358868) (1.495349,0.4984496) (1.709976,0.569992) (1.732051,0.5773503) (1.955409,0.6518028) (2.236068,0.745356) (2.55701,0.8523368) (2.924018,0.9746726) (3.343702,1.114567) (3.823622,1.274541) (4.372426,1.457475) (5,1.666667)};
\addlegendentry{X2}
\addplot[gold,thick,no marks] coordinates {(0.2,0.05049024) (0.228706,0.05790523) (0.2615321,0.06646438) (0.2990698,0.07636846) (0.3419952,0.08786248) (0.3910817,0.1012476) (0.4472136,0.1168959) (0.5114021,0.135268) (0.5848035,0.1569331) (0.6687403,0.1825908) (0.7647245,0.2130923) (0.8744853,0.2494583) (1,0.2928932) (1.14353,0.3447929) (1.30766,0.4067476) (1.495349,0.4805436) (1.709976,0.5681672) (1.732051,0.5773503) (1.955409,0.6718156) (2.236068,0.7939182) (2.55701,0.9371697) (2.924018,1.104575) (3.343702,1.299504) (3.823622,1.52576) (4.372426,1.787656) (5,2.090098)};
\addlegendentry{X10}
\addplot[blue,thick,no marks] coordinates {(0.2,0.09108889) (0.228706,0.1037595) (0.2615321,0.1180676) (0.2990698,0.1341745) (0.3419952,0.1522395) (0.3910817,0.1724164) (0.4472136,0.1948491) (0.5114021,0.2196711) (0.5848035,0.2470077) (0.6687403,0.2769824) (0.7647245,0.3097249) (0.8744853,0.3453767) (1,0.3840905) (1.14353,0.4260193) (1.30766,0.471297) (1.495349,0.520012) (1.709976,0.5721818) (1.732051,0.5773503) (1.955409,0.6277305) (2.236068,0.6864745) (2.55701,0.7481168) (2.924018,0.8122493) (3.343702,0.8783634) (3.823622,0.9458677) (4.372426,1.014111) (5,1.082412)};
\addlegendentry{X29678}
\addplot[purple,thick,no marks] coordinates {(0.2,0.0946909) (0.228706,0.1073888) (0.2615321,0.121639) (0.2990698,0.1375885) (0.3419952,0.1553885) (0.3910817,0.1751927) (0.4472136,0.1971567) (0.5114021,0.2214392) (0.5848035,0.2482027) (0.6687403,0.2776161) (0.7647245,0.3098563) (0.8744853,0.3451068) (1,0.3835523) (1.14353,0.4253657) (1.30766,0.4706895) (1.495349,0.5196114) (1.709976,0.5721401) (1.732051,0.5773503) (1.955409,0.6281826) (2.236068,0.6875291) (2.55701,0.7498471) (2.924018,0.8146848) (3.343702,0.8814859) (3.823622,0.949612) (4.372426,1.018371) (5,1.087051)};
\addlegendentry{X2534}
\end{axis}\end{tikzpicture}
\caption[Isosceles fingerprint: the physical selector, two retained ETC location neighbors and the three classical controls]{Isosceles fingerprint: the physical selector, two retained ETC location neighbors and the three classical controls. The dimensionless heat time is fixed at .03. Similarity of sampled curves is not identity.}\label{fig:fp-H0.03}
\end{figure}

\section{Heat content and a temperature mean}\label{sec:slide67}

The diagram shows temperature means at \ensuremath{\Delta}=.03 and 1, computed by exact integration of the archived piecewise-linear field on its fine triangulation; these are numerical illustrations, not continuum certificates. A thermometer maximum locates H. A spatially integrated thermal image can also give a temperature-weighted mean. The total heat is proportional to \ensuremath{\int}u, and the first moment gives its balance-like center. At late time this mean tends to the $\psi_0$-weighted mean from the absorbing Brownian endpoint distribution, rather than the $\psi_0^2$ quantum expectation. If a uniform source balances losses at steady state, the field instead solves a Poisson or screened-Poisson equation. This makes a direct link to concentration and flow means.


\[C_{\rm temp}(t)=\frac{\int_Tpu\dd A}{\int_Tu\dd A}\ \longrightarrow\ \frac{\int_Tp\psi_0\dd A}{\int_T\psi_0\dd A}.\]


Insert $u\sim a_0e^{-\kappa\lambda_0t}\psi_0$ into the moment
quotient. The coefficient and exponential cancel, leaving
$\int_Tp\psi_0/\int_T\psi_0$.



\section{Production and removal: Abs(\ensuremath{\mu})}\label{sec:slide68}

Imagine a concentration generated uniformly in the triangle, diffusing with a fixed coefficient, removed at a uniform rate, and lost at the boundary. Dividing the steady balance by the diffusivity gives a screened-Poisson problem. Its positive concentration mean is Abs(\ensuremath{\mu}). At \ensuremath{\mu}=0 it becomes the torsion-volume and flow-flux mean. At very large \ensuremath{\mu} on a fixed triangle the bulk concentration is nearly uniform, with boundary layers, so the mean tends to G. The dimensionless parameter is removal rate times $L^2$ divided by diffusivity.


\[-\Delta u+\frac{\mu}{L^2}u=1,\quad u|_{\partial T}=0,\quad {\rm Abs}(\mu)=\frac{\int_Tpu\dd A}{\int_Tu\dd A}.\]



Uniform production $s_0$, diffusion coefficient $\kappa$ and removal
rate $k$ give $-\kappa\Delta c+kc=s_0$ with absorbing boundaries.
Put $\mu=kL^2/\kappa$ and $c=(s_0/\kappa)u$.
The source scale cancels from the centroid.
The energy is
\[
\mathcal J[u]=\frac12\int_T|\nabla u|^2
+\frac{\mu}{2L^2}\int_Tu^2-\int_Tu.
\]
Strict convexity gives a unique positive field.
Finite elements solve $(K+\mu M/L^2)v=f$.
At zero removal the mean is torsion weighted.
On a fixed triangle large $\mu$ produces a nearly uniform interior
with boundary thickness $L/\sqrt\mu$, so its mean tends to $G$.
This does not prove uniform attractivity over arbitrarily thin shapes.



\begin{figure}[htbp]\centering
\begin{tikzpicture}[x=1.7cm,y=1.7cm]
\fill[teal!18!white] (3.142857,0.6428571) -- (2.857143,0.8571429) -- (2.857143,0.6428571) -- cycle;
\fill[teal!21!white] (1.714286,1.714286) -- (1.428571,1.928571) -- (1.428571,1.714286) -- cycle;
\fill[teal!21!white] (1.714286,1.714286) -- (1.714286,1.5) -- (2,1.5) -- cycle;
\fill[teal!15!white] (0.5714286,2.571429) -- (0.2857143,2.785714) -- (0.2857143,2.571429) -- cycle;
\fill[teal!20!white] (2.571429,1.071429) -- (2.285714,1.285714) -- (2.285714,1.071429) -- cycle;
\fill[teal!13!white] (3.714286,0.2142857) -- (3.428571,0.4285714) -- (3.428571,0.2142857) -- cycle;
\fill[teal!8!white] (3.714286,0.2142857) -- (3.714286,0) -- (4,0) -- cycle;
\fill[teal!20!white] (0.8571429,2.142857) -- (1.142857,2.142857) -- (0.8571429,2.357143) -- cycle;
\fill[teal!21!white] (1.142857,2.142857) -- (1.142857,1.928571) -- (1.428571,1.928571) -- cycle;
\fill[teal!9!white] (0,2.785714) -- (0.2857143,2.785714) -- (0,3) -- cycle;
\fill[teal!18!white] (0.5714286,2.571429) -- (0.5714286,2.357143) -- (0.8571429,2.357143) -- cycle;
\fill[teal!21!white] (2,1.285714) -- (2.285714,1.285714) -- (2,1.5) -- cycle;
\fill[teal!19!white] (2.571429,1.071429) -- (2.571429,0.8571429) -- (2.857143,0.8571429) -- cycle;
\fill[teal!16!white] (3.142857,0.4285714) -- (3.428571,0.4285714) -- (3.142857,0.6428571) -- cycle;
\fill[teal!13!white] (0.2857143,0) -- (0,0.2142857) -- (0,0) -- cycle;
\fill[teal!24!white] (0,0.2142857) -- (0.2857143,0) -- (0.2857143,0.2142857) -- cycle;
\fill[teal!21!white] (0.5714286,0.2142857) -- (0.2857143,0) -- (0.5714286,0) -- cycle;
\fill[teal!31!white] (0.2857143,0) -- (0.5714286,0.2142857) -- (0.2857143,0.2142857) -- cycle;
\fill[teal!34!white] (1.714286,1.5) -- (1.714286,1.714286) -- (1.428571,1.714286) -- cycle;
\fill[teal!54!white] (1.428571,1.5) -- (1.714286,1.5) -- (1.428571,1.714286) -- cycle;
\fill[teal!27!white] (0,1.928571) -- (0,1.714286) -- (0.2857143,1.928571) -- cycle;
\fill[teal!44!white] (0,1.714286) -- (0.2857143,1.714286) -- (0.2857143,1.928571) -- cycle;
\fill[teal!21!white] (0.2857143,2.357143) -- (0,2.571429) -- (0,2.357143) -- cycle;
\fill[teal!27!white] (0,2.571429) -- (0.2857143,2.357143) -- (0.2857143,2.571429) -- cycle;
\fill[teal!23!white] (0,2.142857) -- (0.2857143,2.357143) -- (0,2.357143) -- cycle;
\fill[teal!37!white] (0.2857143,2.357143) -- (0,2.142857) -- (0.2857143,2.142857) -- cycle;
\fill[teal!26!white] (0,2.142857) -- (0,1.928571) -- (0.2857143,1.928571) -- cycle;
\fill[teal!40!white] (0.2857143,2.142857) -- (0,2.142857) -- (0.2857143,1.928571) -- cycle;
\fill[teal!63!white] (0.2857143,1.714286) -- (0.5714286,1.928571) -- (0.2857143,1.928571) -- cycle;
\fill[teal!71!white] (0.5714286,1.714286) -- (0.5714286,1.928571) -- (0.2857143,1.714286) -- cycle;
\fill[teal!28!white] (0.2857143,1.5) -- (0,1.714286) -- (0,1.5) -- cycle;
\fill[teal!45!white] (0,1.714286) -- (0.2857143,1.5) -- (0.2857143,1.714286) -- cycle;
\fill[teal!68!white] (0.2857143,1.5) -- (0.5714286,1.714286) -- (0.2857143,1.714286) -- cycle;
\fill[teal!76!white] (0.5714286,1.714286) -- (0.2857143,1.5) -- (0.5714286,1.5) -- cycle;
\fill[teal!29!white] (0,1.285714) -- (0.2857143,1.5) -- (0,1.5) -- cycle;
\fill[teal!47!white] (0.2857143,1.5) -- (0,1.285714) -- (0.2857143,1.285714) -- cycle;
\fill[teal!29!white] (0,1.285714) -- (0,1.071429) -- (0.2857143,1.071429) -- cycle;
\fill[teal!47!white] (0.2857143,1.285714) -- (0,1.285714) -- (0.2857143,1.071429) -- cycle;
\fill[teal!13!white] (3.428571,0) -- (3.714286,0) -- (3.428571,0.2142857) -- cycle;
\fill[teal!13!white] (3.714286,0) -- (3.714286,0.2142857) -- (3.428571,0.2142857) -- cycle;
\fill[teal!24!white] (1.428571,0) -- (1.714286,0) -- (1.428571,0.2142857) -- cycle;
\fill[teal!39!white] (1.714286,0) -- (1.714286,0.2142857) -- (1.428571,0.2142857) -- cycle;
\fill[teal!23!white] (0.8571429,0) -- (0.5714286,0.2142857) -- (0.5714286,0) -- cycle;
\fill[teal!37!white] (0.5714286,0.2142857) -- (0.8571429,0) -- (0.8571429,0.2142857) -- cycle;
\fill[teal!24!white] (1.142857,0.2142857) -- (1.142857,0) -- (1.428571,0) -- cycle;
\fill[teal!39!white] (1.142857,0.2142857) -- (1.428571,0) -- (1.428571,0.2142857) -- cycle;
\fill[teal!24!white] (1.142857,0.2142857) -- (0.8571429,0) -- (1.142857,0) -- cycle;
\fill[teal!38!white] (0.8571429,0) -- (1.142857,0.2142857) -- (0.8571429,0.2142857) -- cycle;
\fill[teal!53!white] (1.142857,1.928571) -- (1.142857,1.714286) -- (1.428571,1.714286) -- cycle;
\fill[teal!33!white] (1.428571,1.928571) -- (1.142857,1.928571) -- (1.428571,1.714286) -- cycle;
\fill[teal!32!white] (0.8571429,2.142857) -- (1.142857,1.928571) -- (1.142857,2.142857) -- cycle;
\fill[teal!51!white] (1.142857,1.928571) -- (0.8571429,2.142857) -- (0.8571429,1.928571) -- cycle;
\fill[teal!16!white] (0,2.785714) -- (0,2.571429) -- (0.2857143,2.571429) -- cycle;
\fill[teal!16!white] (0.2857143,2.785714) -- (0,2.785714) -- (0.2857143,2.571429) -- cycle;
\fill[teal!67!white] (1.142857,1.928571) -- (0.8571429,1.714286) -- (1.142857,1.714286) -- cycle;
\fill[teal!66!white] (0.8571429,1.714286) -- (1.142857,1.928571) -- (0.8571429,1.928571) -- cycle;
\fill[teal!76!white] (0.8571429,1.714286) -- (0.5714286,1.928571) -- (0.5714286,1.714286) -- cycle;
\fill[teal!72!white] (0.5714286,1.928571) -- (0.8571429,1.714286) -- (0.8571429,1.928571) -- cycle;
\fill[teal!36!white] (0.2857143,2.357143) -- (0.5714286,2.357143) -- (0.2857143,2.571429) -- cycle;
\fill[teal!25!white] (0.5714286,2.357143) -- (0.5714286,2.571429) -- (0.2857143,2.571429) -- cycle;
\fill[teal!27!white] (0,0.6428571) -- (0,0.4285714) -- (0.2857143,0.6428571) -- cycle;
\fill[teal!41!white] (0,0.4285714) -- (0.2857143,0.4285714) -- (0.2857143,0.6428571) -- cycle;
\fill[teal!22!white] (0,0.4285714) -- (0,0.2142857) -- (0.2857143,0.2142857) -- cycle;
\fill[teal!37!white] (0.2857143,0.4285714) -- (0,0.4285714) -- (0.2857143,0.2142857) -- cycle;
\fill[teal!33!white] (2.285714,1.285714) -- (2,1.285714) -- (2.285714,1.071429) -- cycle;
\fill[teal!52!white] (2,1.285714) -- (2,1.071429) -- (2.285714,1.071429) -- cycle;
\fill[teal!66!white] (2,1.071429) -- (2,0.8571429) -- (2.285714,1.071429) -- cycle;
\fill[teal!64!white] (2,0.8571429) -- (2.285714,0.8571429) -- (2.285714,1.071429) -- cycle;
\fill[teal!50!white] (2.285714,0.8571429) -- (2.571429,0.8571429) -- (2.285714,1.071429) -- cycle;
\fill[teal!31!white] (2.571429,0.8571429) -- (2.571429,1.071429) -- (2.285714,1.071429) -- cycle;
\fill[teal!69!white] (1.428571,0.4285714) -- (1.714286,0.2142857) -- (1.714286,0.4285714) -- cycle;
\fill[teal!61!white] (1.714286,0.2142857) -- (1.428571,0.4285714) -- (1.428571,0.2142857) -- cycle;
\fill[teal!21!white] (3.142857,0.2142857) -- (3.428571,0) -- (3.428571,0.2142857) -- cycle;
\fill[teal!16!white] (3.142857,0.2142857) -- (3.142857,0) -- (3.428571,0) -- cycle;
\fill[teal!18!white] (3.142857,0) -- (3.142857,0.2142857) -- (2.857143,0) -- cycle;
\fill[teal!28!white] (3.142857,0.2142857) -- (2.857143,0.2142857) -- (2.857143,0) -- cycle;
\fill[teal!38!white] (1.714286,0) -- (2,0.2142857) -- (1.714286,0.2142857) -- cycle;
\fill[teal!24!white] (2,0.2142857) -- (1.714286,0) -- (2,0) -- cycle;
\fill[teal!52!white] (0.5714286,0.2142857) -- (0.5714286,0.4285714) -- (0.2857143,0.2142857) -- cycle;
\fill[teal!53!white] (0.5714286,0.4285714) -- (0.2857143,0.4285714) -- (0.2857143,0.2142857) -- cycle;
\fill[teal!59!white] (0.8571429,0.4285714) -- (0.5714286,0.2142857) -- (0.8571429,0.2142857) -- cycle;
\fill[teal!64!white] (0.8571429,0.4285714) -- (0.5714286,0.4285714) -- (0.5714286,0.2142857) -- cycle;
\fill[teal!69!white] (1.142857,1.5) -- (1.428571,1.5) -- (1.428571,1.714286) -- cycle;
\fill[teal!68!white] (1.142857,1.714286) -- (1.142857,1.5) -- (1.428571,1.714286) -- cycle;
\fill[teal!56!white] (0.5714286,2.142857) -- (0.2857143,2.142857) -- (0.2857143,1.928571) -- cycle;
\fill[teal!63!white] (0.5714286,1.928571) -- (0.5714286,2.142857) -- (0.2857143,1.928571) -- cycle;
\fill[teal!50!white] (0.5714286,2.142857) -- (0.2857143,2.357143) -- (0.2857143,2.142857) -- cycle;
\fill[teal!46!white] (0.5714286,2.142857) -- (0.5714286,2.357143) -- (0.2857143,2.357143) -- cycle;
\fill[teal!55!white] (0.8571429,2.142857) -- (0.5714286,2.142857) -- (0.8571429,1.928571) -- cycle;
\fill[teal!64!white] (0.5714286,2.142857) -- (0.5714286,1.928571) -- (0.8571429,1.928571) -- cycle;
\fill[teal!40!white] (0.5714286,2.142857) -- (0.8571429,2.142857) -- (0.8571429,2.357143) -- cycle;
\fill[teal!38!white] (0.5714286,2.357143) -- (0.5714286,2.142857) -- (0.8571429,2.357143) -- cycle;
\fill[teal!44!white] (1.714286,1.285714) -- (2,1.285714) -- (2,1.5) -- cycle;
\fill[teal!44!white] (1.714286,1.5) -- (1.714286,1.285714) -- (2,1.5) -- cycle;
\fill[teal!94!white] (1.142857,1.071429) -- (1.428571,0.8571429) -- (1.428571,1.071429) -- cycle;
\fill[teal!94!white] (1.428571,0.8571429) -- (1.142857,1.071429) -- (1.142857,0.8571429) -- cycle;
\fill[teal!28!white] (3.142857,0.4285714) -- (3.142857,0.2142857) -- (3.428571,0.2142857) -- cycle;
\fill[teal!20!white] (3.428571,0.4285714) -- (3.142857,0.4285714) -- (3.428571,0.2142857) -- cycle;
\fill[teal!39!white] (2.857143,0.4285714) -- (3.142857,0.4285714) -- (2.857143,0.6428571) -- cycle;
\fill[teal!26!white] (3.142857,0.4285714) -- (3.142857,0.6428571) -- (2.857143,0.6428571) -- cycle;
\fill[teal!36!white] (3.142857,0.2142857) -- (3.142857,0.4285714) -- (2.857143,0.2142857) -- cycle;
\fill[teal!41!white] (3.142857,0.4285714) -- (2.857143,0.4285714) -- (2.857143,0.2142857) -- cycle;
\fill[teal!39!white] (2.571429,0.8571429) -- (2.571429,0.6428571) -- (2.857143,0.8571429) -- cycle;
\fill[teal!38!white] (2.857143,0.8571429) -- (2.571429,0.6428571) -- (2.857143,0.6428571) -- cycle;
\fill[teal!23!white] (2.285714,0) -- (2.285714,0.2142857) -- (2,0) -- cycle;
\fill[teal!37!white] (2.285714,0.2142857) -- (2,0.2142857) -- (2,0) -- cycle;
\fill[teal!57!white] (2.285714,0.2142857) -- (2,0.4285714) -- (2,0.2142857) -- cycle;
\fill[teal!63!white] (2,0.4285714) -- (2.285714,0.2142857) -- (2.285714,0.4285714) -- cycle;
\fill[teal!79!white] (2,0.4285714) -- (1.714286,0.6428571) -- (1.714286,0.4285714) -- cycle;
\fill[teal!80!white] (1.714286,0.6428571) -- (2,0.4285714) -- (2,0.6428571) -- cycle;
\fill[teal!68!white] (1.714286,0.2142857) -- (2,0.4285714) -- (1.714286,0.4285714) -- cycle;
\fill[teal!59!white] (2,0.2142857) -- (2,0.4285714) -- (1.714286,0.2142857) -- cycle;
\fill[teal!28!white] (0.2857143,0.8571429) -- (0,0.8571429) -- (0,0.6428571) -- cycle;
\fill[teal!45!white] (0.2857143,0.8571429) -- (0,0.6428571) -- (0.2857143,0.6428571) -- cycle;
\fill[teal!47!white] (0,1.071429) -- (0.2857143,0.8571429) -- (0.2857143,1.071429) -- cycle;
\fill[teal!29!white] (0,0.8571429) -- (0.2857143,0.8571429) -- (0,1.071429) -- cycle;
\fill[teal!78!white] (0.2857143,1.5) -- (0.5714286,1.285714) -- (0.5714286,1.5) -- cycle;
\fill[teal!71!white] (0.5714286,1.285714) -- (0.2857143,1.5) -- (0.2857143,1.285714) -- cycle;
\fill[teal!68!white] (0.5714286,0.6428571) -- (0.2857143,0.8571429) -- (0.2857143,0.6428571) -- cycle;
\fill[teal!77!white] (0.2857143,0.8571429) -- (0.5714286,0.6428571) -- (0.5714286,0.8571429) -- cycle;
\fill[teal!64!white] (0.2857143,0.4285714) -- (0.5714286,0.6428571) -- (0.2857143,0.6428571) -- cycle;
\fill[teal!67!white] (0.5714286,0.4285714) -- (0.5714286,0.6428571) -- (0.2857143,0.4285714) -- cycle;
\fill[teal!75!white] (0.8571429,0.4285714) -- (0.5714286,0.6428571) -- (0.5714286,0.4285714) -- cycle;
\fill[teal!82!white] (0.5714286,0.6428571) -- (0.8571429,0.4285714) -- (0.8571429,0.6428571) -- cycle;
\fill[teal!62!white] (1.142857,0.4285714) -- (1.142857,0.2142857) -- (1.428571,0.2142857) -- cycle;
\fill[teal!70!white] (1.428571,0.4285714) -- (1.142857,0.4285714) -- (1.428571,0.2142857) -- cycle;
\fill[teal!61!white] (1.142857,0.2142857) -- (1.142857,0.4285714) -- (0.8571429,0.2142857) -- cycle;
\fill[teal!69!white] (1.142857,0.4285714) -- (0.8571429,0.4285714) -- (0.8571429,0.2142857) -- cycle;
\fill[teal!81!white] (0.8571429,0.4285714) -- (1.142857,0.4285714) -- (0.8571429,0.6428571) -- cycle;
\fill[teal!86!white] (1.142857,0.4285714) -- (1.142857,0.6428571) -- (0.8571429,0.6428571) -- cycle;
\fill[teal!69!white] (1.428571,1.285714) -- (1.714286,1.5) -- (1.428571,1.5) -- cycle;
\fill[teal!69!white] (1.428571,1.285714) -- (1.714286,1.285714) -- (1.714286,1.5) -- cycle;
\fill[teal!80!white] (1.142857,1.5) -- (1.428571,1.285714) -- (1.428571,1.5) -- cycle;
\fill[teal!87!white] (1.428571,1.285714) -- (1.142857,1.5) -- (1.142857,1.285714) -- cycle;
\fill[teal!92!white] (1.428571,1.285714) -- (1.142857,1.071429) -- (1.428571,1.071429) -- cycle;
\fill[teal!93!white] (1.142857,1.071429) -- (1.428571,1.285714) -- (1.142857,1.285714) -- cycle;
\fill[teal!81!white] (1.428571,0.6428571) -- (1.428571,0.4285714) -- (1.714286,0.4285714) -- cycle;
\fill[teal!85!white] (1.714286,0.6428571) -- (1.428571,0.6428571) -- (1.714286,0.4285714) -- cycle;
\fill[teal!82!white] (1.428571,0.6428571) -- (1.142857,0.4285714) -- (1.428571,0.4285714) -- cycle;
\fill[teal!87!white] (1.142857,0.4285714) -- (1.428571,0.6428571) -- (1.142857,0.6428571) -- cycle;
\fill[teal!94!white] (1.428571,0.6428571) -- (1.428571,0.8571429) -- (1.142857,0.8571429) -- cycle;
\fill[teal!93!white] (1.142857,0.6428571) -- (1.428571,0.6428571) -- (1.142857,0.8571429) -- cycle;
\fill[teal!68!white] (2,1.285714) -- (1.714286,1.071429) -- (2,1.071429) -- cycle;
\fill[teal!68!white] (1.714286,1.285714) -- (1.714286,1.071429) -- (2,1.285714) -- cycle;
\fill[teal!79!white] (1.428571,1.285714) -- (1.714286,1.071429) -- (1.714286,1.285714) -- cycle;
\fill[teal!87!white] (1.714286,1.071429) -- (1.428571,1.285714) -- (1.428571,1.071429) -- cycle;
\fill[teal!61!white] (2.285714,0.6428571) -- (2.571429,0.8571429) -- (2.285714,0.8571429) -- cycle;
\fill[teal!59!white] (2.285714,0.6428571) -- (2.571429,0.6428571) -- (2.571429,0.8571429) -- cycle;
\fill[teal!72!white] (2,0.8571429) -- (2.285714,0.6428571) -- (2.285714,0.8571429) -- cycle;
\fill[teal!77!white] (2.285714,0.6428571) -- (2,0.8571429) -- (2,0.6428571) -- cycle;
\fill[teal!71!white] (2.285714,0.6428571) -- (2,0.4285714) -- (2.285714,0.4285714) -- cycle;
\fill[teal!76!white] (2,0.4285714) -- (2.285714,0.6428571) -- (2,0.6428571) -- cycle;
\fill[teal!50!white] (2.571429,0.4285714) -- (2.857143,0.4285714) -- (2.857143,0.6428571) -- cycle;
\fill[teal!54!white] (2.571429,0.6428571) -- (2.571429,0.4285714) -- (2.857143,0.6428571) -- cycle;
\fill[teal!63!white] (2.285714,0.6428571) -- (2.571429,0.4285714) -- (2.571429,0.6428571) -- cycle;
\fill[teal!66!white] (2.571429,0.4285714) -- (2.285714,0.6428571) -- (2.285714,0.4285714) -- cycle;
\fill[teal!71!white] (0.2857143,0.8571429) -- (0.5714286,1.071429) -- (0.2857143,1.071429) -- cycle;
\fill[teal!79!white] (0.5714286,1.071429) -- (0.2857143,0.8571429) -- (0.5714286,0.8571429) -- cycle;
\fill[teal!71!white] (0.5714286,1.071429) -- (0.2857143,1.285714) -- (0.2857143,1.071429) -- cycle;
\fill[teal!79!white] (0.5714286,1.071429) -- (0.5714286,1.285714) -- (0.2857143,1.285714) -- cycle;
\fill[teal!84!white] (0.8571429,1.5) -- (0.5714286,1.714286) -- (0.5714286,1.5) -- cycle;
\fill[teal!83!white] (0.8571429,1.5) -- (0.8571429,1.714286) -- (0.5714286,1.714286) -- cycle;
\fill[teal!81!white] (0.8571429,1.714286) -- (0.8571429,1.5) -- (1.142857,1.714286) -- cycle;
\fill[teal!82!white] (0.8571429,1.5) -- (1.142857,1.5) -- (1.142857,1.714286) -- cycle;
\fill[teal!86!white] (1.714286,0.8571429) -- (1.714286,0.6428571) -- (2,0.6428571) -- cycle;
\fill[teal!83!white] (2,0.8571429) -- (1.714286,0.8571429) -- (2,0.6428571) -- cycle;
\fill[teal!80!white] (1.714286,0.8571429) -- (2,0.8571429) -- (2,1.071429) -- cycle;
\fill[teal!81!white] (1.714286,1.071429) -- (1.714286,0.8571429) -- (2,1.071429) -- cycle;
\fill[teal!89!white] (1.714286,0.8571429) -- (1.428571,0.6428571) -- (1.714286,0.6428571) -- cycle;
\fill[teal!92!white] (1.428571,0.6428571) -- (1.714286,0.8571429) -- (1.428571,0.8571429) -- cycle;
\fill[teal!92!white] (1.428571,0.8571429) -- (1.714286,0.8571429) -- (1.428571,1.071429) -- cycle;
\fill[teal!88!white] (1.714286,0.8571429) -- (1.714286,1.071429) -- (1.428571,1.071429) -- cycle;
\fill[teal!22!white] (2.571429,0) -- (2.571429,0.2142857) -- (2.285714,0) -- cycle;
\fill[teal!35!white] (2.571429,0.2142857) -- (2.285714,0.2142857) -- (2.285714,0) -- cycle;
\fill[teal!54!white] (2.285714,0.2142857) -- (2.571429,0.2142857) -- (2.285714,0.4285714) -- cycle;
\fill[teal!57!white] (2.571429,0.2142857) -- (2.571429,0.4285714) -- (2.285714,0.4285714) -- cycle;
\fill[teal!31!white] (2.857143,0.2142857) -- (2.571429,0.2142857) -- (2.857143,0) -- cycle;
\fill[teal!21!white] (2.571429,0.2142857) -- (2.571429,0) -- (2.857143,0) -- cycle;
\fill[teal!46!white] (2.857143,0.4285714) -- (2.571429,0.2142857) -- (2.857143,0.2142857) -- cycle;
\fill[teal!52!white] (2.571429,0.4285714) -- (2.571429,0.2142857) -- (2.857143,0.4285714) -- cycle;
\fill[teal!92!white] (1.142857,0.6428571) -- (0.8571429,0.8571429) -- (0.8571429,0.6428571) -- cycle;
\fill[teal!94!white] (0.8571429,0.8571429) -- (1.142857,0.6428571) -- (1.142857,0.8571429) -- cycle;
\fill[teal!87!white] (0.5714286,0.6428571) -- (0.8571429,0.8571429) -- (0.5714286,0.8571429) -- cycle;
\fill[teal!88!white] (0.8571429,0.8571429) -- (0.5714286,0.6428571) -- (0.8571429,0.6428571) -- cycle;
\fill[teal!91!white] (1.142857,1.5) -- (0.8571429,1.285714) -- (1.142857,1.285714) -- cycle;
\fill[teal!90!white] (0.8571429,1.5) -- (0.8571429,1.285714) -- (1.142857,1.5) -- cycle;
\fill[teal!88!white] (0.5714286,1.285714) -- (0.8571429,1.285714) -- (0.5714286,1.5) -- cycle;
\fill[teal!89!white] (0.8571429,1.285714) -- (0.8571429,1.5) -- (0.5714286,1.5) -- cycle;
\fill[teal!94!white] (1.142857,1.071429) -- (0.8571429,1.071429) -- (1.142857,0.8571429) -- cycle;
\fill[teal!94!white] (0.8571429,1.071429) -- (0.8571429,0.8571429) -- (1.142857,0.8571429) -- cycle;
\fill[teal!94!white] (0.8571429,1.071429) -- (1.142857,1.071429) -- (1.142857,1.285714) -- cycle;
\fill[teal!94!white] (0.8571429,1.285714) -- (0.8571429,1.071429) -- (1.142857,1.285714) -- cycle;
\fill[teal!90!white] (0.8571429,1.071429) -- (0.5714286,1.071429) -- (0.5714286,0.8571429) -- cycle;
\fill[teal!92!white] (0.8571429,0.8571429) -- (0.8571429,1.071429) -- (0.5714286,0.8571429) -- cycle;
\fill[teal!90!white] (0.5714286,1.071429) -- (0.8571429,1.071429) -- (0.5714286,1.285714) -- cycle;
\fill[teal!92!white] (0.8571429,1.071429) -- (0.8571429,1.285714) -- (0.5714286,1.285714) -- cycle;
\draw[teal,thick] (0,0) -- (4,0) -- (0,3) -- cycle;
\draw[magenta!70,thin] (1.248144,1.000199) -- (3.6,2.35) -- (4.1,2.35);\filldraw[fill=magenta,draw=white,line width=.45pt] (1.248144,1.000199) circle(1.5pt) node[above right,font=\scriptsize,text=magenta,fill=white,inner sep=.6pt] {};
\node[anchor=west,text=magenta,font=\small] at(4.12,2.35){Abs100.0};\draw[gold!70,thin] (1,1) -- (3.6,1.8) -- (4.1,1.8);\filldraw[fill=gold,draw=white,line width=.45pt] (1,1) circle(1.5pt) node[above right,font=\scriptsize,text=gold,fill=white,inner sep=.6pt] {};
\node[anchor=west,text=gold,font=\small] at(4.12,1.8){$I$};\draw[navy!70,thin] (1.333333,1) -- (3.6,1.25) -- (4.1,1.25);\filldraw[fill=navy,draw=white,line width=.45pt] (1.333333,1) circle(1.5pt) node[above right,font=\scriptsize,text=navy,fill=white,inner sep=.6pt] {};
\node[anchor=west,text=navy,font=\small] at(4.12,1.25){$G$};\end{tikzpicture}
\caption[The screened-Poisson concentration]{The screened-Poisson concentration. Abs(100) is its weighted mean. Fields are archived numerical illustrations; teal saturation increases with the displayed value.}\label{fig:field-Abs100.0}
\end{figure}

\section{Thin domains change the concentration profile}\label{sec:slide69}

A very thin triangle is nearly a collection of vertical strips. Within one strip of height H, neglect the small longitudinal derivative and solve \ensuremath{-}u\ensuremath{^{\prime\prime}}+k\ensuremath{^2}u=1 with zero values at its two ends. Integrating the solution gives F(H), the concentration mass per unit longitudinal distance. For small kH, the field is parabolic and F is H cubed over twelve. For large kH, the nearly uniform bulk gives F\ensuremath{\approx}H/k\ensuremath{^2}. These two powers weight the triangular thickness very differently, so the center can shift horizontally while a vertex moves mostly vertically.


\[F(H)=H/k^2-2\tanh(kH/2)/k^3.\]


The strip solution is
$k^{-2}[1-\cosh(k(t-H/2))/\cosh(kH/2)]$.
Integrating from $0$ to $H$ gives the column mass $F(H)$.
This profile is also a useful thin-domain numerical basis.



\paragraph{Demonstration and investigation.} Use the vertical-strip schematic to sketch a parabolic profile and a flat-core profile on the board. Ask which one gives more relative weight to narrow parts of the triangle.


\section{A failure of the attractive-family conjecture}\label{sec:slide70}

The new follow-up studied 512 extreme shapes plus the old 100-shape grid. Direct Ritz sensitivities stayed positive for the earlier small and moderate parameters, but $\mu=10^4$ produced robust negative values. A different quadratic finite-element solver reproduced one witness through 192 subdivisions and three motion step sizes. This is strong numerical evidence at that exact parameter. A separate thin-domain energy argument proves that every sufficiently large \ensuremath{\mu} has some nonattractive triangle, although it does not provide an explicit threshold.


\[h\cdot\Delta C/|h|^2\simeq-0.366\quad(\mu=10^4),\quad\hbox{numerical witness}.\]


Set $\varepsilon=\mu^{-1/2}$ and compare heights
$\varepsilon\tau_1,\varepsilon\tau_2$ while moving the apex
horizontally by $\varepsilon$. A negative horizontal center change
of order one gives a dot product of order $-\varepsilon$;
the vertical contribution is only $O(\varepsilon^2)$.
Appendix A supplies the limiting field and energy estimate.



\section{The shape space of triangles}\label{sec:slide71}

Represent a triangle shape by normalized angles \ensuremath{\theta}\_i/\ensuremath{\pi}. These sum to one and therefore serve as barycentric coordinates in an equilateral angle diagram. Distances to its sides are proportional to the corresponding angles. Permuting the vertices identifies six equivalent sectors. Choose the fundamental region $\theta_1\le\theta_2\le\theta_3$. Its vertices correspond to the flat limits (0,0,\ensuremath{\pi}) and (0,\ensuremath{\pi}/2,\ensuremath{\pi}/2), and the equilateral shape. Sampling uniformly in its area is different from sampling side lengths or Cartesian apex positions uniformly.


\[\theta_1+\theta_2+\theta_3=\pi,\quad\theta_1\le\theta_2\le\theta_3.\]


In the chosen labeled sector the vertices have angle triples (0,0,1), (0,1/2,1/2) and (1/3,1/3,1/3). Boundary points are degenerate or symmetric limits. Generic numerical samples should remain strictly inside.


\begin{figure}[htbp]\centering
\begin{tikzpicture}[x=1.45cm,y=1.45cm]
\draw[teal,thick](0,0)--(4,0)--(2,3.4641)--cycle;\fill[teal!15](2,3.4641)--(3,1.73205)--(2,1.1547)--cycle;\fill[magenta] (2.033333,3.329387)circle(.65pt);\fill[magenta] (2.066667,3.194671)circle(.65pt);\fill[magenta] (2.033333,3.098446)circle(.65pt);\fill[magenta] (2.066667,2.963731)circle(.65pt);\fill[magenta] (2.033333,2.867506)circle(.65pt);\fill[magenta] (2.066667,2.732791)circle(.65pt);\fill[magenta] (2.033333,2.636566)circle(.65pt);\fill[magenta] (2.066667,2.501851)circle(.65pt);\fill[magenta] (2.033333,2.405626)circle(.65pt);\fill[magenta] (2.066667,2.270911)circle(.65pt);\fill[magenta] (2.033333,2.174686)circle(.65pt);\fill[magenta] (2.066667,2.039971)circle(.65pt);\fill[magenta] (2.033333,1.943746)circle(.65pt);\fill[magenta] (2.066667,1.809031)circle(.65pt);\fill[magenta] (2.033333,1.712806)circle(.65pt);\fill[magenta] (2.066667,1.578091)circle(.65pt);\fill[magenta] (2.033333,1.481866)circle(.65pt);\fill[magenta] (2.066667,1.347151)circle(.65pt);\fill[magenta] (2.033333,1.250926)circle(.65pt);\fill[magenta] (2.133333,3.156181)circle(.65pt);\fill[magenta] (2.166667,3.021466)circle(.65pt);\fill[magenta] (2.133333,2.925241)circle(.65pt);\fill[magenta] (2.166667,2.790526)circle(.65pt);\fill[magenta] (2.133333,2.694301)circle(.65pt);\fill[magenta] (2.166667,2.559586)circle(.65pt);\fill[magenta] (2.133333,2.463361)circle(.65pt);\fill[magenta] (2.166667,2.328646)circle(.65pt);\fill[magenta] (2.133333,2.232421)circle(.65pt);\fill[magenta] (2.166667,2.097706)circle(.65pt);\fill[magenta] (2.133333,2.001481)circle(.65pt);\fill[magenta] (2.166667,1.866766)circle(.65pt);\fill[magenta] (2.133333,1.770541)circle(.65pt);\fill[magenta] (2.166667,1.635826)circle(.65pt);\fill[magenta] (2.133333,1.539601)circle(.65pt);\fill[magenta] (2.166667,1.404886)circle(.65pt);\fill[magenta] (2.133333,1.308661)circle(.65pt);\fill[magenta] (2.233333,2.982976)circle(.65pt);\fill[magenta] (2.266667,2.848261)circle(.65pt);\fill[magenta] (2.233333,2.752036)circle(.65pt);\fill[magenta] (2.266667,2.617321)circle(.65pt);\fill[magenta] (2.233333,2.521096)circle(.65pt);\fill[magenta] (2.266667,2.386381)circle(.65pt);\fill[magenta] (2.233333,2.290156)circle(.65pt);\fill[magenta] (2.266667,2.155441)circle(.65pt);\fill[magenta] (2.233333,2.059216)circle(.65pt);\fill[magenta] (2.266667,1.924501)circle(.65pt);\fill[magenta] (2.233333,1.828276)circle(.65pt);\fill[magenta] (2.266667,1.693561)circle(.65pt);\fill[magenta] (2.233333,1.597336)circle(.65pt);\fill[magenta] (2.266667,1.462621)circle(.65pt);\fill[magenta] (2.233333,1.366396)circle(.65pt);\fill[magenta] (2.333333,2.809771)circle(.65pt);\fill[magenta] (2.366667,2.675056)circle(.65pt);\fill[magenta] (2.333333,2.578831)circle(.65pt);\fill[magenta] (2.366667,2.444116)circle(.65pt);\fill[magenta] (2.333333,2.347891)circle(.65pt);\fill[magenta] (2.366667,2.213176)circle(.65pt);\fill[magenta] (2.333333,2.116951)circle(.65pt);\fill[magenta] (2.366667,1.982236)circle(.65pt);\fill[magenta] (2.333333,1.886011)circle(.65pt);\fill[magenta] (2.366667,1.751296)circle(.65pt);\fill[magenta] (2.333333,1.655071)circle(.65pt);\fill[magenta] (2.366667,1.520356)circle(.65pt);\fill[magenta] (2.333333,1.424131)circle(.65pt);\fill[magenta] (2.433333,2.636566)circle(.65pt);\fill[magenta] (2.466667,2.501851)circle(.65pt);\fill[magenta] (2.433333,2.405626)circle(.65pt);\fill[magenta] (2.466667,2.270911)circle(.65pt);\fill[magenta] (2.433333,2.174686)circle(.65pt);\fill[magenta] (2.466667,2.039971)circle(.65pt);\fill[magenta] (2.433333,1.943746)circle(.65pt);\fill[magenta] (2.466667,1.809031)circle(.65pt);\fill[magenta] (2.433333,1.712806)circle(.65pt);\fill[magenta] (2.466667,1.578091)circle(.65pt);\fill[magenta] (2.433333,1.481866)circle(.65pt);\fill[magenta] (2.533333,2.463361)circle(.65pt);\fill[magenta] (2.566667,2.328646)circle(.65pt);\fill[magenta] (2.533333,2.232421)circle(.65pt);\fill[magenta] (2.566667,2.097706)circle(.65pt);\fill[magenta] (2.533333,2.001481)circle(.65pt);\fill[magenta] (2.566667,1.866766)circle(.65pt);\fill[magenta] (2.533333,1.770541)circle(.65pt);\fill[magenta] (2.566667,1.635826)circle(.65pt);\fill[magenta] (2.533333,1.539601)circle(.65pt);\fill[magenta] (2.633333,2.290156)circle(.65pt);\fill[magenta] (2.666667,2.155441)circle(.65pt);\fill[magenta] (2.633333,2.059216)circle(.65pt);\fill[magenta] (2.666667,1.924501)circle(.65pt);\fill[magenta] (2.633333,1.828276)circle(.65pt);\fill[magenta] (2.666667,1.693561)circle(.65pt);\fill[magenta] (2.633333,1.597336)circle(.65pt);\fill[magenta] (2.733333,2.116951)circle(.65pt);\fill[magenta] (2.766667,1.982236)circle(.65pt);\fill[magenta] (2.733333,1.886011)circle(.65pt);\fill[magenta] (2.766667,1.751296)circle(.65pt);\fill[magenta] (2.733333,1.655071)circle(.65pt);\fill[magenta] (2.833333,1.943746)circle(.65pt);\fill[magenta] (2.866667,1.809031)circle(.65pt);\fill[magenta] (2.833333,1.712806)circle(.65pt);\fill[magenta] (2.933333,1.770541)circle(.65pt);\node[above]at(2,3.4641){$(0,0,\pi)$};\node[below]at(0,0){$(\pi,0,0)$};\node[below]at(4,0){$(0,\pi,0)$};\end{tikzpicture}
\caption[The angle simplex and the actual 100 archived training shapes in the ordered sector]{The angle simplex and the actual 100 archived training shapes in the ordered sector. Distances to the sides are proportional to angles.}\label{fig:angle-space}
\end{figure}

\paragraph{Demonstration and investigation.} Point out the six sectors and the actual 100 archived sample points. Ask students to locate right triangles, where one angle equals \ensuremath{\pi}/2.


\section{Location ranks and behavioral fingerprints}\label{sec:slide72}

The supplied snapshot includes 72807 ETC entries. The restricted evaluator retained 70853 candidates finite on all 100 generic training shapes and excluded 1954 unsupported or singular expressions. The location score is RMS Euclidean difference divided by triangle diameter. Forty held-out shapes test transfer beyond the ranking sample. Isosceles and right fingerprints compare curves and derivatives. Symmetric vertex-Jacobian margins compare response behavior. These quantities remain separate, so a good location rank does not conceal a poor derivative match.


\[E_{\rm loc}(n)=\left(\frac1N\sum_k\frac{|\TC(T_k)-X_n(T_k)|^2}{D_{T_k}^2}\right)^{1/2}.\]



The search evaluated 72,807 ETC entries on 100 training shapes.
It retained 70,853 finite supported entries and set aside 1,954
unsupported or singular entries. Forty holdout shapes supplied
independent location scores. This is a broad screen, not an
exclusion proof for unevaluated catalogue points.

Location, derivative and response errors measure different things.
Declare separate feature normalizations and comparison weights.
Include $X_1,X_2,X_{10}$ as controls even when absent from top ranks.
For symbolic approximations fit permutation-equivariant weights or
invariant functions on angles. Enforce homogeneity and reflection
symmetry, check positivity if claimed, and report holdout and
thin-domain error. Pointwise closeness alone gives no derivative
or attractivity bound, and an approximate fit is not an identification.



\begin{center}\small\begin{tabular}{p{0.26\linewidth}p{0.28\linewidth}p{0.36\linewidth}}\toprule
Selector & Location neighbor & Holdout RMS / $D$\\\midrule
$\mathrm{Q0}$ & $X_{24936}$ & \(0.006263\)\\
F1 & $X_{5393}$ & \(0.00405\)\\
M6 & $X_{25086}$ & \(0.002682\)\\
E2 & $X_{59474}$ & \(0.001305\)\\
Dark1.0b1.0 & $X_{17210}$ & \(0.002795\)\\\bottomrule\end{tabular}\end{center}

\section{ETC membership is not an algebraic test}\label{sec:slide73}

Most studied PITCs have no known exact ETC identity, but we should not explain that by a categorical ban on transcendental formulas. X360 uses angles and already belongs to the ETC. The equal-exit center has an incomplete-beta representation. Other centers require solving a boundary-value problem, an eigenproblem or an integral equation. Their dimensionless parameters create continuous families, and some proposed maxima need uniqueness conventions. A numerical shortlist cannot establish nonmembership or transcendence. The useful claim is that the current atlas supplies definitions and computations beyond the currently identified exact ETC aliases.


Once normalized barycentric weights are obtained as functions of side lengths or angles, they can in principle define a center function, even if evaluation requires a numerical solver. Simplicity, historical cataloguing and proof of a clean identity are separate questions.



\section{Certificates and numerical evidence}\label{sec:slide74}

An analytical identity can prove that a physical center equals G or W. A positive-density argument can prove interiority without computing coordinates. An interval enclosure can prove the sign of one finite motion and therefore disprove a global attraction statement. Stable convergence of independent solvers is persuasive numerical evidence but does not provide the same guarantee. A no-failure screen supports a conjecture. A sign that changes under refinement remains unresolved. All these outcomes are useful if labeled correctly.


\[S_i=\tfrac12(J_i+J_i^T),\quad J_i=D_{V_i}\TC.\]



A negative floating-point value becomes a certificate only after
all relevant errors are enclosed. For PDE centers these include
discretization, quadrature, maximum location, branch choice and
parameter differentiation. Finer repeated calculations diagnose
error but do not bound the uncomputed remainder.

For a finite motion keep dimensionless parameters fixed. For a
derivative compare centered differences at several steps with a
differentiated solve. Use cluster projectors near multiple
eigenvalues and Hessian checks near flat maxima.
Screen the least eigenvalue of each $S_i$ over all vertices, then
optimize negative witnesses.




\paragraph{Demonstration and investigation.} Ask students to rank four claims: exact G response, E0 certified negative motion, Abs10000 refined negative motion, and M6 positive screen.


\section{The revised attractivity picture}\label{sec:slide75}

The proved attractive points are I, G, W and M3=X360, with their physical aliases. E0, E1, E2 and M1 have certified failures. The full Abs family fails for sufficiently large \ensuremath{\mu}. The solid-angle family fails for sufficiently small heights, by its M1 limit. The particular previously screened Abs(0), Abs(1), Abs(10), Abs(100), Solid(.1), Solid(1) and Solid(10) remain open. We will return to the optical family next. The central lesson is that attraction is a shape-sensitive property, and limiting to an attractive point on each fixed triangle does not settle a whole family.


\begin{center}\small\begin{tabular}{p{0.39\linewidth}p{0.51\linewidth}}\toprule
Evidence & Selectors\\\midrule
Proved attractive & $X_1,X_2,X_{10},X_{360}$\\
Certified failures & E0, E1, E2, M1\\
Strong positive conjecture & M6\\
Refinement remains unresolved & F1, Q0m\\
Numerical negative evidence & Q0--Q4, tested H, E3, Illum\\\bottomrule\end{tabular}\end{center}

\section*{Questions and exercises}\addcontentsline{toc}{section}{Questions and exercises}

\begin{enumerate}

\item\label{ex:62} Why is the incenter a plausible short-time candidate?

\item\label{ex:63} Why do many modes matter at very early time?

\item\label{ex:64} Does H(\ensuremath{\Delta}) equal I at every sufficiently small positive \ensuremath{\Delta}?

\item\label{ex:65} What disappears when every frame is normalized to its own maximum?

\item\label{ex:66} What must be written beside a heat fingerprint?

\item\label{ex:67} Does a late heat-image mean recover $\mathrm{Q0}$?

\item\label{ex:68} Does a fixed-triangle limit toward an attractive point prove all family members attractive?

\item\label{ex:69} Which weight is more concentrated near the thickest part?

\item\label{ex:70} Why is there no contradiction with Abs(\ensuremath{\mu})\ensuremath{\to}G on every fixed triangle?

\item\label{ex:71} Does one point in this space encode the triangle's size?

\item\label{ex:72} Why are excluded ETC entries not evidence of dissimilarity?

\item\label{ex:73} What exact ETC identity in this course uses a nonalgebraic function?

\item\label{ex:74} Which is enough to disprove a universal attractivity conjecture?

\item\label{ex:75} Which center would you prioritize for a positive theorem?

\end{enumerate}

\chapter{Light, imaging and research conclusions}\label{ch:6}

Optics adds two useful styles of mathematical problem. Ray-based illumination selects an extremum of a geometric integral. Fresnel diffraction selects a complex wave field, so interference and normalization matter. Some intuitively appealing dark-point definitions fail because there is no exact interior null or because an extremum is not unique. The clean alternatives use positive weighted means or a carefully defined finite part.


\section{Solid-angle illumination}\label{sec:slide77}

The solid angle combines inverse-square spreading with a projection cosine. Place an observation point a height h above the aperture plane. A small area dA at separation r subtends h dA/r cubed. Integrating gives the displayed field. Select a point in the projected triangle where the field is largest, and hold eta=h/L fixed under similarities. The maximum lies inside, but the current project does not claim global uniqueness for every shape and height. If several global maxima exist, the set is the intrinsic object until a justified selection convention is supplied.


\[\Omega_h(p)=\int_T\frac{h}{(h^2+|p-q|^2)^{3/2}}\dd A_q,\quad h=\eta L.\]



An observer at $(p,h)$ sees a patch under solid angle
$\cos\vartheta\,\dd A/r^2=h\dd A/(h^2+|p-q|^2)^{3/2}$.
Integrating gives angular source size. With $h=\eta L$, its maximum
is a candidate rule; general finite-height uniqueness remains open.
For large height,
\[
\Omega_h(p)=\frac{\mathcal A}{h^2}
-\frac3{2h^4}\int_T|p-q|^2\dd A_q+O(h^{-6}),
\]
whose leading position term selects $G$.
For small height polar integration gives
$\Omega_h=2\pi-h\int(R_p^2+h^2)^{-1/2}\dd\theta$.
The coefficient tends to the loop-field kernel.
All maximizing points approach M1, without assuming finite-height
uniqueness. Its certified negative motion implies eventual failure
for sufficiently small $\eta$.




\section{Solid-angle limits}\label{sec:slide78}

At large height, expand the kernel in inverse powers of h. The first position-dependent term minimizes mean squared distance to the aperture, giving G. At small height, almost the entire upward half-sphere is covered. The leading missing solid angle is h times the angular boundary integral \ensuremath{\int}1/R\_p(\ensuremath{\theta})d\ensuremath{\theta}, which is the loop-field function defining M1. Thus the two endpoints connect magnetostatics and mass geometry. Since M1 has a certified negative finite motion, all sufficiently small-height selections also fail attractivity on those two fixed triangles.


\[{\rm Solid}(\eta)\to{\rm M1}\quad(\eta\downarrow0),\qquad {\rm Solid}(\eta)\to G\quad(\eta\to\infty).\]


Polar integration yields
$\Omega_h=2\pi-h\int_0^{2\pi}(R^2+h^2)^{-1/2}\dd\theta$.
Compact-interior convergence and boundary estimates transfer the
strict M1 finite witness without requiring derivative convergence.



\section{Inverse-square illumination in the plane}\label{sec:slide79}

A planar integral of inverse-square illumination diverges near the observation point. Excluding an equal tiny disk at every trial point gives a common logarithmic divergence. The finite part is the angular integral of log ray length to the boundary. This gives Illum. For a triangle, each ray length is the minimum positive ratio of side distance to direction cosine. The logarithm is concave in the observation point, and the angular integral is strictly concave. It diverges downward at the boundary, so Illum has one unique interior maximum.


\[\mathcal I(p)=\int_0^{2\pi}\log R_p(\theta)\dd\theta,\quad {\rm Illum}=\argmax_T\mathcal I.\]



Inverse-square sources and observation in the same plane give a
divergent local integral. Remove a radius-$\varepsilon$ disk:
\[
\int_{T\setminus B_\varepsilon(p)}|p-q|^{-2}\dd A_q
=-2\pi\log\varepsilon+\int_0^{2\pi}\log R_p(\theta)\dd\theta.
\]
The finite part is $\mathcal I$, with no height parameter.
A ray hitting side $i$ has $R=d_i/(n_i\cdot e_\theta)$.
Differentiation gives
$\nabla\mathcal I=-\sum_i(\omega_i/d_i)n_i$.
Moving sector-endpoint terms cancel because ray length is continuous.

Within a sector the directional second derivative is
$-(v\cdot n_i)^2/d_i^2$. Sector switches supply nonpositive
contributions, because radial length is the minimum permitted
intersection. Normals span the plane, giving strict concavity.
Boundary approach sends $\mathcal I$ to $-\infty$.
Illum is therefore a unique interior maximum, despite its refined
numerical negative vertex response.




\section{Illum: unique but nonattractive numerically}\label{sec:slide80}

The finite-part objective has a useful derivative. Each side contributes its angular visibility theta\_i divided by distance d\_i, times the outward normal. Setting the vector sum to zero gives a stationary equation that can be solved directly instead of optimizing a noisy integral. Strict concavity establishes the unique global maximum. The new response screen found a finite-motion witness whose negative ratios repeat at 45 and 75 decimal digits as the step is halved. Those are high-precision numerical values, not outward-rounded interval enclosures.


\[\nabla\mathcal I(p)=-\sum_i\frac{\omega_i(p)}{d_i(p)}n_i,\quad n_i\ \hbox{outward}.\]


In the archived witness the three finite-motion ratios are approximately \ensuremath{-}.103971, \ensuremath{-}.101927 and \ensuremath{-}.100920. The response matrix's smallest symmetric eigenvalue is about \ensuremath{-}.099924. The limiting derivative sign is consistent with the finite motions.


\begin{figure}[htbp]\centering
\begin{tikzpicture}[x=1.7cm,y=1.7cm]
\fill[teal!46!white] (3.142857,0.6428571) -- (2.857143,0.8571429) -- (2.857143,0.6428571) -- cycle;
\fill[teal!49!white] (1.714286,1.714286) -- (1.428571,1.928571) -- (1.428571,1.714286) -- cycle;
\fill[teal!53!white] (1.714286,1.714286) -- (1.714286,1.5) -- (2,1.5) -- cycle;
\fill[teal!39!white] (0.5714286,2.571429) -- (0.2857143,2.785714) -- (0.2857143,2.571429) -- cycle;
\fill[teal!48!white] (2.571429,1.071429) -- (2.285714,1.285714) -- (2.285714,1.071429) -- cycle;
\fill[teal!32!white] (3.714286,0.2142857) -- (3.428571,0.4285714) -- (3.428571,0.2142857) -- cycle;
\fill[teal!8!white] (3.714286,0.2142857) -- (3.714286,0) -- (4,0) -- cycle;
\fill[teal!49!white] (0.8571429,2.142857) -- (1.142857,2.142857) -- (0.8571429,2.357143) -- cycle;
\fill[teal!49!white] (1.142857,2.142857) -- (1.142857,1.928571) -- (1.428571,1.928571) -- cycle;
\fill[teal!17!white] (0,2.785714) -- (0.2857143,2.785714) -- (0,3) -- cycle;
\fill[teal!46!white] (0.5714286,2.571429) -- (0.5714286,2.357143) -- (0.8571429,2.357143) -- cycle;
\fill[teal!42!white] (2,1.285714) -- (2.285714,1.285714) -- (2,1.5) -- cycle;
\fill[teal!49!white] (2.571429,1.071429) -- (2.571429,0.8571429) -- (2.857143,0.8571429) -- cycle;
\fill[teal!41!white] (3.142857,0.4285714) -- (3.428571,0.4285714) -- (3.142857,0.6428571) -- cycle;
\fill[teal!32!white] (0.2857143,0) -- (0,0.2142857) -- (0,0) -- cycle;
\fill[teal!54!white] (0,0.2142857) -- (0.2857143,0) -- (0.2857143,0.2142857) -- cycle;
\fill[teal!51!white] (0.5714286,0.2142857) -- (0.2857143,0) -- (0.5714286,0) -- cycle;
\fill[teal!62!white] (0.2857143,0) -- (0.5714286,0.2142857) -- (0.2857143,0.2142857) -- cycle;
\fill[teal!67!white] (1.714286,1.5) -- (1.714286,1.714286) -- (1.428571,1.714286) -- cycle;
\fill[teal!79!white] (1.428571,1.5) -- (1.714286,1.5) -- (1.428571,1.714286) -- cycle;
\fill[teal!58!white] (0,1.928571) -- (0,1.714286) -- (0.2857143,1.928571) -- cycle;
\fill[teal!72!white] (0,1.714286) -- (0.2857143,1.714286) -- (0.2857143,1.928571) -- cycle;
\fill[teal!49!white] (0.2857143,2.357143) -- (0,2.571429) -- (0,2.357143) -- cycle;
\fill[teal!58!white] (0,2.571429) -- (0.2857143,2.357143) -- (0.2857143,2.571429) -- cycle;
\fill[teal!52!white] (0,2.142857) -- (0.2857143,2.357143) -- (0,2.357143) -- cycle;
\fill[teal!66!white] (0.2857143,2.357143) -- (0,2.142857) -- (0.2857143,2.142857) -- cycle;
\fill[teal!56!white] (0,2.142857) -- (0,1.928571) -- (0.2857143,1.928571) -- cycle;
\fill[teal!69!white] (0.2857143,2.142857) -- (0,2.142857) -- (0.2857143,1.928571) -- cycle;
\fill[teal!81!white] (0.2857143,1.714286) -- (0.5714286,1.928571) -- (0.2857143,1.928571) -- cycle;
\fill[teal!85!white] (0.5714286,1.714286) -- (0.5714286,1.928571) -- (0.2857143,1.714286) -- cycle;
\fill[teal!60!white] (0.2857143,1.5) -- (0,1.714286) -- (0,1.5) -- cycle;
\fill[teal!73!white] (0,1.714286) -- (0.2857143,1.5) -- (0.2857143,1.714286) -- cycle;
\fill[teal!83!white] (0.2857143,1.5) -- (0.5714286,1.714286) -- (0.2857143,1.714286) -- cycle;
\fill[teal!87!white] (0.5714286,1.714286) -- (0.2857143,1.5) -- (0.5714286,1.5) -- cycle;
\fill[teal!60!white] (0,1.285714) -- (0.2857143,1.5) -- (0,1.5) -- cycle;
\fill[teal!74!white] (0.2857143,1.5) -- (0,1.285714) -- (0.2857143,1.285714) -- cycle;
\fill[teal!60!white] (0,1.285714) -- (0,1.071429) -- (0.2857143,1.071429) -- cycle;
\fill[teal!74!white] (0.2857143,1.285714) -- (0,1.285714) -- (0.2857143,1.071429) -- cycle;
\fill[teal!33!white] (3.428571,0) -- (3.714286,0) -- (3.428571,0.2142857) -- cycle;
\fill[teal!35!white] (3.714286,0) -- (3.714286,0.2142857) -- (3.428571,0.2142857) -- cycle;
\fill[teal!56!white] (1.428571,0) -- (1.714286,0) -- (1.428571,0.2142857) -- cycle;
\fill[teal!70!white] (1.714286,0) -- (1.714286,0.2142857) -- (1.428571,0.2142857) -- cycle;
\fill[teal!53!white] (0.8571429,0) -- (0.5714286,0.2142857) -- (0.5714286,0) -- cycle;
\fill[teal!68!white] (0.5714286,0.2142857) -- (0.8571429,0) -- (0.8571429,0.2142857) -- cycle;
\fill[teal!56!white] (1.142857,0.2142857) -- (1.142857,0) -- (1.428571,0) -- cycle;
\fill[teal!70!white] (1.142857,0.2142857) -- (1.428571,0) -- (1.428571,0.2142857) -- cycle;
\fill[teal!56!white] (1.142857,0.2142857) -- (0.8571429,0) -- (1.142857,0) -- cycle;
\fill[teal!69!white] (0.8571429,0) -- (1.142857,0.2142857) -- (0.8571429,0.2142857) -- cycle;
\fill[teal!78!white] (1.142857,1.928571) -- (1.142857,1.714286) -- (1.428571,1.714286) -- cycle;
\fill[teal!66!white] (1.428571,1.928571) -- (1.142857,1.928571) -- (1.428571,1.714286) -- cycle;
\fill[teal!64!white] (0.8571429,2.142857) -- (1.142857,1.928571) -- (1.142857,2.142857) -- cycle;
\fill[teal!76!white] (1.142857,1.928571) -- (0.8571429,2.142857) -- (0.8571429,1.928571) -- cycle;
\fill[teal!41!white] (0,2.785714) -- (0,2.571429) -- (0.2857143,2.571429) -- cycle;
\fill[teal!43!white] (0.2857143,2.785714) -- (0,2.785714) -- (0.2857143,2.571429) -- cycle;
\fill[teal!84!white] (1.142857,1.928571) -- (0.8571429,1.714286) -- (1.142857,1.714286) -- cycle;
\fill[teal!83!white] (0.8571429,1.714286) -- (1.142857,1.928571) -- (0.8571429,1.928571) -- cycle;
\fill[teal!87!white] (0.8571429,1.714286) -- (0.5714286,1.928571) -- (0.5714286,1.714286) -- cycle;
\fill[teal!85!white] (0.5714286,1.928571) -- (0.8571429,1.714286) -- (0.8571429,1.928571) -- cycle;
\fill[teal!65!white] (0.2857143,2.357143) -- (0.5714286,2.357143) -- (0.2857143,2.571429) -- cycle;
\fill[teal!56!white] (0.5714286,2.357143) -- (0.5714286,2.571429) -- (0.2857143,2.571429) -- cycle;
\fill[teal!57!white] (0,0.6428571) -- (0,0.4285714) -- (0.2857143,0.6428571) -- cycle;
\fill[teal!70!white] (0,0.4285714) -- (0.2857143,0.4285714) -- (0.2857143,0.6428571) -- cycle;
\fill[teal!51!white] (0,0.4285714) -- (0,0.2142857) -- (0.2857143,0.2142857) -- cycle;
\fill[teal!66!white] (0.2857143,0.4285714) -- (0,0.4285714) -- (0.2857143,0.2142857) -- cycle;
\fill[teal!66!white] (2.285714,1.285714) -- (2,1.285714) -- (2.285714,1.071429) -- cycle;
\fill[teal!78!white] (2,1.285714) -- (2,1.071429) -- (2.285714,1.071429) -- cycle;
\fill[teal!84!white] (2,1.071429) -- (2,0.8571429) -- (2.285714,1.071429) -- cycle;
\fill[teal!83!white] (2,0.8571429) -- (2.285714,0.8571429) -- (2.285714,1.071429) -- cycle;
\fill[teal!76!white] (2.285714,0.8571429) -- (2.571429,0.8571429) -- (2.285714,1.071429) -- cycle;
\fill[teal!64!white] (2.571429,0.8571429) -- (2.571429,1.071429) -- (2.285714,1.071429) -- cycle;
\fill[teal!85!white] (1.428571,0.4285714) -- (1.714286,0.2142857) -- (1.714286,0.4285714) -- cycle;
\fill[teal!82!white] (1.714286,0.2142857) -- (1.428571,0.4285714) -- (1.428571,0.2142857) -- cycle;
\fill[teal!51!white] (3.142857,0.2142857) -- (3.428571,0) -- (3.428571,0.2142857) -- cycle;
\fill[teal!42!white] (3.142857,0.2142857) -- (3.142857,0) -- (3.428571,0) -- cycle;
\fill[teal!45!white] (3.142857,0) -- (3.142857,0.2142857) -- (2.857143,0) -- cycle;
\fill[teal!60!white] (3.142857,0.2142857) -- (2.857143,0.2142857) -- (2.857143,0) -- cycle;
\fill[teal!70!white] (1.714286,0) -- (2,0.2142857) -- (1.714286,0.2142857) -- cycle;
\fill[teal!55!white] (2,0.2142857) -- (1.714286,0) -- (2,0) -- cycle;
\fill[teal!76!white] (0.5714286,0.2142857) -- (0.5714286,0.4285714) -- (0.2857143,0.2142857) -- cycle;
\fill[teal!76!white] (0.5714286,0.4285714) -- (0.2857143,0.4285714) -- (0.2857143,0.2142857) -- cycle;
\fill[teal!80!white] (0.8571429,0.4285714) -- (0.5714286,0.2142857) -- (0.8571429,0.2142857) -- cycle;
\fill[teal!82!white] (0.8571429,0.4285714) -- (0.5714286,0.4285714) -- (0.5714286,0.2142857) -- cycle;
\fill[teal!85!white] (1.142857,1.5) -- (1.428571,1.5) -- (1.428571,1.714286) -- cycle;
\fill[teal!85!white] (1.142857,1.714286) -- (1.142857,1.5) -- (1.428571,1.714286) -- cycle;
\fill[teal!78!white] (0.5714286,2.142857) -- (0.2857143,2.142857) -- (0.2857143,1.928571) -- cycle;
\fill[teal!81!white] (0.5714286,1.928571) -- (0.5714286,2.142857) -- (0.2857143,1.928571) -- cycle;
\fill[teal!74!white] (0.5714286,2.142857) -- (0.2857143,2.357143) -- (0.2857143,2.142857) -- cycle;
\fill[teal!72!white] (0.5714286,2.142857) -- (0.5714286,2.357143) -- (0.2857143,2.357143) -- cycle;
\fill[teal!78!white] (0.8571429,2.142857) -- (0.5714286,2.142857) -- (0.8571429,1.928571) -- cycle;
\fill[teal!82!white] (0.5714286,2.142857) -- (0.5714286,1.928571) -- (0.8571429,1.928571) -- cycle;
\fill[teal!68!white] (0.5714286,2.142857) -- (0.8571429,2.142857) -- (0.8571429,2.357143) -- cycle;
\fill[teal!67!white] (0.5714286,2.357143) -- (0.5714286,2.142857) -- (0.8571429,2.357143) -- cycle;
\fill[teal!73!white] (1.714286,1.285714) -- (2,1.285714) -- (2,1.5) -- cycle;
\fill[teal!74!white] (1.714286,1.5) -- (1.714286,1.285714) -- (2,1.5) -- cycle;
\fill[teal!94!white] (1.142857,1.071429) -- (1.428571,0.8571429) -- (1.428571,1.071429) -- cycle;
\fill[teal!94!white] (1.428571,0.8571429) -- (1.142857,1.071429) -- (1.142857,0.8571429) -- cycle;
\fill[teal!59!white] (3.142857,0.4285714) -- (3.142857,0.2142857) -- (3.428571,0.2142857) -- cycle;
\fill[teal!50!white] (3.428571,0.4285714) -- (3.142857,0.4285714) -- (3.428571,0.2142857) -- cycle;
\fill[teal!68!white] (2.857143,0.4285714) -- (3.142857,0.4285714) -- (2.857143,0.6428571) -- cycle;
\fill[teal!55!white] (3.142857,0.4285714) -- (3.142857,0.6428571) -- (2.857143,0.6428571) -- cycle;
\fill[teal!66!white] (3.142857,0.2142857) -- (3.142857,0.4285714) -- (2.857143,0.2142857) -- cycle;
\fill[teal!70!white] (3.142857,0.4285714) -- (2.857143,0.4285714) -- (2.857143,0.2142857) -- cycle;
\fill[teal!69!white] (2.571429,0.8571429) -- (2.571429,0.6428571) -- (2.857143,0.8571429) -- cycle;
\fill[teal!67!white] (2.857143,0.8571429) -- (2.571429,0.6428571) -- (2.857143,0.6428571) -- cycle;
\fill[teal!54!white] (2.285714,0) -- (2.285714,0.2142857) -- (2,0) -- cycle;
\fill[teal!69!white] (2.285714,0.2142857) -- (2,0.2142857) -- (2,0) -- cycle;
\fill[teal!80!white] (2.285714,0.2142857) -- (2,0.4285714) -- (2,0.2142857) -- cycle;
\fill[teal!82!white] (2,0.4285714) -- (2.285714,0.2142857) -- (2.285714,0.4285714) -- cycle;
\fill[teal!89!white] (2,0.4285714) -- (1.714286,0.6428571) -- (1.714286,0.4285714) -- cycle;
\fill[teal!89!white] (1.714286,0.6428571) -- (2,0.4285714) -- (2,0.6428571) -- cycle;
\fill[teal!85!white] (1.714286,0.2142857) -- (2,0.4285714) -- (1.714286,0.4285714) -- cycle;
\fill[teal!81!white] (2,0.2142857) -- (2,0.4285714) -- (1.714286,0.2142857) -- cycle;
\fill[teal!59!white] (0.2857143,0.8571429) -- (0,0.8571429) -- (0,0.6428571) -- cycle;
\fill[teal!72!white] (0.2857143,0.8571429) -- (0,0.6428571) -- (0.2857143,0.6428571) -- cycle;
\fill[teal!74!white] (0,1.071429) -- (0.2857143,0.8571429) -- (0.2857143,1.071429) -- cycle;
\fill[teal!60!white] (0,0.8571429) -- (0.2857143,0.8571429) -- (0,1.071429) -- cycle;
\fill[teal!88!white] (0.2857143,1.5) -- (0.5714286,1.285714) -- (0.5714286,1.5) -- cycle;
\fill[teal!85!white] (0.5714286,1.285714) -- (0.2857143,1.5) -- (0.2857143,1.285714) -- cycle;
\fill[teal!83!white] (0.5714286,0.6428571) -- (0.2857143,0.8571429) -- (0.2857143,0.6428571) -- cycle;
\fill[teal!87!white] (0.2857143,0.8571429) -- (0.5714286,0.6428571) -- (0.5714286,0.8571429) -- cycle;
\fill[teal!81!white] (0.2857143,0.4285714) -- (0.5714286,0.6428571) -- (0.2857143,0.6428571) -- cycle;
\fill[teal!83!white] (0.5714286,0.4285714) -- (0.5714286,0.6428571) -- (0.2857143,0.4285714) -- cycle;
\fill[teal!86!white] (0.8571429,0.4285714) -- (0.5714286,0.6428571) -- (0.5714286,0.4285714) -- cycle;
\fill[teal!89!white] (0.5714286,0.6428571) -- (0.8571429,0.4285714) -- (0.8571429,0.6428571) -- cycle;
\fill[teal!82!white] (1.142857,0.4285714) -- (1.142857,0.2142857) -- (1.428571,0.2142857) -- cycle;
\fill[teal!86!white] (1.428571,0.4285714) -- (1.142857,0.4285714) -- (1.428571,0.2142857) -- cycle;
\fill[teal!82!white] (1.142857,0.2142857) -- (1.142857,0.4285714) -- (0.8571429,0.2142857) -- cycle;
\fill[teal!85!white] (1.142857,0.4285714) -- (0.8571429,0.4285714) -- (0.8571429,0.2142857) -- cycle;
\fill[teal!89!white] (0.8571429,0.4285714) -- (1.142857,0.4285714) -- (0.8571429,0.6428571) -- cycle;
\fill[teal!91!white] (1.142857,0.4285714) -- (1.142857,0.6428571) -- (0.8571429,0.6428571) -- cycle;
\fill[teal!86!white] (1.428571,1.285714) -- (1.714286,1.5) -- (1.428571,1.5) -- cycle;
\fill[teal!86!white] (1.428571,1.285714) -- (1.714286,1.285714) -- (1.714286,1.5) -- cycle;
\fill[teal!90!white] (1.142857,1.5) -- (1.428571,1.285714) -- (1.428571,1.5) -- cycle;
\fill[teal!92!white] (1.428571,1.285714) -- (1.142857,1.5) -- (1.142857,1.285714) -- cycle;
\fill[teal!94!white] (1.428571,1.285714) -- (1.142857,1.071429) -- (1.428571,1.071429) -- cycle;
\fill[teal!94!white] (1.142857,1.071429) -- (1.428571,1.285714) -- (1.142857,1.285714) -- cycle;
\fill[teal!90!white] (1.428571,0.6428571) -- (1.428571,0.4285714) -- (1.714286,0.4285714) -- cycle;
\fill[teal!91!white] (1.714286,0.6428571) -- (1.428571,0.6428571) -- (1.714286,0.4285714) -- cycle;
\fill[teal!90!white] (1.428571,0.6428571) -- (1.142857,0.4285714) -- (1.428571,0.4285714) -- cycle;
\fill[teal!92!white] (1.142857,0.4285714) -- (1.428571,0.6428571) -- (1.142857,0.6428571) -- cycle;
\fill[teal!94!white] (1.428571,0.6428571) -- (1.428571,0.8571429) -- (1.142857,0.8571429) -- cycle;
\fill[teal!94!white] (1.142857,0.6428571) -- (1.428571,0.6428571) -- (1.142857,0.8571429) -- cycle;
\fill[teal!85!white] (2,1.285714) -- (1.714286,1.071429) -- (2,1.071429) -- cycle;
\fill[teal!85!white] (1.714286,1.285714) -- (1.714286,1.071429) -- (2,1.285714) -- cycle;
\fill[teal!89!white] (1.428571,1.285714) -- (1.714286,1.071429) -- (1.714286,1.285714) -- cycle;
\fill[teal!92!white] (1.714286,1.071429) -- (1.428571,1.285714) -- (1.428571,1.071429) -- cycle;
\fill[teal!81!white] (2.285714,0.6428571) -- (2.571429,0.8571429) -- (2.285714,0.8571429) -- cycle;
\fill[teal!80!white] (2.285714,0.6428571) -- (2.571429,0.6428571) -- (2.571429,0.8571429) -- cycle;
\fill[teal!86!white] (2,0.8571429) -- (2.285714,0.6428571) -- (2.285714,0.8571429) -- cycle;
\fill[teal!88!white] (2.285714,0.6428571) -- (2,0.8571429) -- (2,0.6428571) -- cycle;
\fill[teal!85!white] (2.285714,0.6428571) -- (2,0.4285714) -- (2.285714,0.4285714) -- cycle;
\fill[teal!87!white] (2,0.4285714) -- (2.285714,0.6428571) -- (2,0.6428571) -- cycle;
\fill[teal!75!white] (2.571429,0.4285714) -- (2.857143,0.4285714) -- (2.857143,0.6428571) -- cycle;
\fill[teal!77!white] (2.571429,0.6428571) -- (2.571429,0.4285714) -- (2.857143,0.6428571) -- cycle;
\fill[teal!82!white] (2.285714,0.6428571) -- (2.571429,0.4285714) -- (2.571429,0.6428571) -- cycle;
\fill[teal!83!white] (2.571429,0.4285714) -- (2.285714,0.6428571) -- (2.285714,0.4285714) -- cycle;
\fill[teal!85!white] (0.2857143,0.8571429) -- (0.5714286,1.071429) -- (0.2857143,1.071429) -- cycle;
\fill[teal!88!white] (0.5714286,1.071429) -- (0.2857143,0.8571429) -- (0.5714286,0.8571429) -- cycle;
\fill[teal!85!white] (0.5714286,1.071429) -- (0.2857143,1.285714) -- (0.2857143,1.071429) -- cycle;
\fill[teal!88!white] (0.5714286,1.071429) -- (0.5714286,1.285714) -- (0.2857143,1.285714) -- cycle;
\fill[teal!90!white] (0.8571429,1.5) -- (0.5714286,1.714286) -- (0.5714286,1.5) -- cycle;
\fill[teal!90!white] (0.8571429,1.5) -- (0.8571429,1.714286) -- (0.5714286,1.714286) -- cycle;
\fill[teal!89!white] (0.8571429,1.714286) -- (0.8571429,1.5) -- (1.142857,1.714286) -- cycle;
\fill[teal!90!white] (0.8571429,1.5) -- (1.142857,1.5) -- (1.142857,1.714286) -- cycle;
\fill[teal!91!white] (1.714286,0.8571429) -- (1.714286,0.6428571) -- (2,0.6428571) -- cycle;
\fill[teal!90!white] (2,0.8571429) -- (1.714286,0.8571429) -- (2,0.6428571) -- cycle;
\fill[teal!89!white] (1.714286,0.8571429) -- (2,0.8571429) -- (2,1.071429) -- cycle;
\fill[teal!90!white] (1.714286,1.071429) -- (1.714286,0.8571429) -- (2,1.071429) -- cycle;
\fill[teal!93!white] (1.714286,0.8571429) -- (1.428571,0.6428571) -- (1.714286,0.6428571) -- cycle;
\fill[teal!93!white] (1.428571,0.6428571) -- (1.714286,0.8571429) -- (1.428571,0.8571429) -- cycle;
\fill[teal!94!white] (1.428571,0.8571429) -- (1.714286,0.8571429) -- (1.428571,1.071429) -- cycle;
\fill[teal!92!white] (1.714286,0.8571429) -- (1.714286,1.071429) -- (1.428571,1.071429) -- cycle;
\fill[teal!52!white] (2.571429,0) -- (2.571429,0.2142857) -- (2.285714,0) -- cycle;
\fill[teal!67!white] (2.571429,0.2142857) -- (2.285714,0.2142857) -- (2.285714,0) -- cycle;
\fill[teal!77!white] (2.285714,0.2142857) -- (2.571429,0.2142857) -- (2.285714,0.4285714) -- cycle;
\fill[teal!79!white] (2.571429,0.2142857) -- (2.571429,0.4285714) -- (2.285714,0.4285714) -- cycle;
\fill[teal!63!white] (2.857143,0.2142857) -- (2.571429,0.2142857) -- (2.857143,0) -- cycle;
\fill[teal!50!white] (2.571429,0.2142857) -- (2.571429,0) -- (2.857143,0) -- cycle;
\fill[teal!72!white] (2.857143,0.4285714) -- (2.571429,0.2142857) -- (2.857143,0.2142857) -- cycle;
\fill[teal!76!white] (2.571429,0.4285714) -- (2.571429,0.2142857) -- (2.857143,0.4285714) -- cycle;
\fill[teal!93!white] (1.142857,0.6428571) -- (0.8571429,0.8571429) -- (0.8571429,0.6428571) -- cycle;
\fill[teal!94!white] (0.8571429,0.8571429) -- (1.142857,0.6428571) -- (1.142857,0.8571429) -- cycle;
\fill[teal!91!white] (0.5714286,0.6428571) -- (0.8571429,0.8571429) -- (0.5714286,0.8571429) -- cycle;
\fill[teal!91!white] (0.8571429,0.8571429) -- (0.5714286,0.6428571) -- (0.8571429,0.6428571) -- cycle;
\fill[teal!93!white] (1.142857,1.5) -- (0.8571429,1.285714) -- (1.142857,1.285714) -- cycle;
\fill[teal!92!white] (0.8571429,1.5) -- (0.8571429,1.285714) -- (1.142857,1.5) -- cycle;
\fill[teal!92!white] (0.5714286,1.285714) -- (0.8571429,1.285714) -- (0.5714286,1.5) -- cycle;
\fill[teal!92!white] (0.8571429,1.285714) -- (0.8571429,1.5) -- (0.5714286,1.5) -- cycle;
\fill[teal!94!white] (1.142857,1.071429) -- (0.8571429,1.071429) -- (1.142857,0.8571429) -- cycle;
\fill[teal!94!white] (0.8571429,1.071429) -- (0.8571429,0.8571429) -- (1.142857,0.8571429) -- cycle;
\fill[teal!94!white] (0.8571429,1.071429) -- (1.142857,1.071429) -- (1.142857,1.285714) -- cycle;
\fill[teal!94!white] (0.8571429,1.285714) -- (0.8571429,1.071429) -- (1.142857,1.285714) -- cycle;
\fill[teal!92!white] (0.8571429,1.071429) -- (0.5714286,1.071429) -- (0.5714286,0.8571429) -- cycle;
\fill[teal!93!white] (0.8571429,0.8571429) -- (0.8571429,1.071429) -- (0.5714286,0.8571429) -- cycle;
\fill[teal!92!white] (0.5714286,1.071429) -- (0.8571429,1.071429) -- (0.5714286,1.285714) -- cycle;
\fill[teal!93!white] (0.8571429,1.071429) -- (0.8571429,1.285714) -- (0.5714286,1.285714) -- cycle;
\draw[teal,thick] (0,0) -- (4,0) -- (0,3) -- cycle;
\draw[magenta!70,thin] (1.132007,0.9959148) -- (3.6,2.35) -- (4.1,2.35);\filldraw[fill=magenta,draw=white,line width=.45pt] (1.132007,0.9959148) circle(1.5pt) node[above right,font=\scriptsize,text=magenta,fill=white,inner sep=.6pt] {};
\node[anchor=west,text=magenta,font=\small] at(4.12,2.35){Illum};\draw[gold!70,thin] (1,1) -- (3.6,1.8) -- (4.1,1.8);\filldraw[fill=gold,draw=white,line width=.45pt] (1,1) circle(1.5pt) node[above right,font=\scriptsize,text=gold,fill=white,inner sep=.6pt] {};
\node[anchor=west,text=gold,font=\small] at(4.12,1.8){$I$};\draw[navy!70,thin] (1.333333,1) -- (3.6,1.25) -- (4.1,1.25);\filldraw[fill=navy,draw=white,line width=.45pt] (1.333333,1) circle(1.5pt) node[above right,font=\scriptsize,text=navy,fill=white,inner sep=.6pt] {};
\node[anchor=west,text=navy,font=\small] at(4.12,1.25){$G$};\end{tikzpicture}
\caption[Finite-part planar illumination, characterized by a strictly concave logarithmic ray-length integral]{Finite-part planar illumination, characterized by a strictly concave logarithmic ray-length integral. Fields are archived numerical illustrations; teal saturation increases with the displayed value.}\label{fig:field-Illum}
\end{figure}

\section{Fresnel diffraction through a triangle}\label{sec:slide81}

Within scalar paraxial diffraction, the field at a screen point is a common prefactor times the quadratic-phase integral over the aperture. The Fresnel number F=$L^2$/(lambda z) controls phase variation, with propagation distance z and wavelength lambda. Intensity is proportional to |U|\ensuremath{^2}. Overall phase and positive amplitude normalization do not move intensity extrema. The triangle is a two-dimensional aperture, while the screen coordinate is a projected point in the same plane. The approximation assumes a coherent uniformly phased incident wave and the usual paraxial regime.


\[A_F(p)=\mathcal A^{-1}\int_Te^{i\pi F|p-q|^2/L^2}\dd A_q,\quad F=L^2/(\lambda z).\]



The paraxial approximation is
$\sqrt{z^2+|p-q|^2}=z+|p-q|^2/(2z)+\cdots$.
With wavenumber $2\pi/\lambda$, the variable phase is
$\pi|p-q|^2/(\lambda z)$. Uniform coherent illumination gives the
integral after removing a common phase and amplitude.
Assumptions are scalar waves, uniform aperture phase and paraxial
observation. This is not automatically the Fraunhofer limit
\cite{fresnel}.

Use analytic edge reductions or adaptive oscillatory quadrature.
A mesh resolving torsion can fail at large $F$.
Double phase resolution, compare complex amplitudes, check symmetry
and verify $A_F\to1$ as $F\downarrow0$.




\paragraph{Demonstration and investigation.} Reveal a few complex phase contributions before the intensity field. Remind the class that intensity must be computed after summing amplitudes.


\section{There is no exact interior Fresnel null}\label{sec:slide82}

For a point inside the triangle, rays fill a star-shaped radial region. Integrate exp(ikr\ensuremath{^2}) r dr exactly along each ray. The imaginary part of the resulting angular expression is the integral of [1\ensuremath{-}cos(kR\ensuremath{^2})]/(2k). It is nonnegative for every angle. It cannot vanish identically for a triangle at finite positive k, because its ray lengths vary continuously and are not all locked to phase multiples of two pi. Thus the uniformly phased triangular Fresnel aperture has no exact zero at an interior point. A "dark point" must mean a positive local minimum or a positive darkness-weighted mean.


\[\Im\int_Te^{ik|p-q|^2}\dd A_q=\int_0^{2\pi}\frac{1-\cos(kR_p(\theta)^2)}{2k}\dd\theta>0.\]



Radial integration gives
$\int_0^R e^{ikr^2}r\dd r=(e^{ikR^2}-1)/(2ik)$.
For $k>0$ its imaginary part is nonnegative.
It cannot vanish almost everywhere in direction: on any side-view
interval $R=d/\cos(\theta-\theta_n)$ varies continuously, whereas
$kR^2\in2\pi\mathbb Z$ is discrete. The angular integral is strictly
positive, so no exact interior null exists under these assumptions.
Low-intensity minima can still switch with parameters.
A positive darkness-weighted mean avoids the ambiguity of selecting
one minimum without claiming an exact optical zero.




\section{A clean darkness-weighted center}\label{sec:slide83}

The atlas uses an explicitly calibrated exponential detector response, not a reciprocal power of intensity. Divide intensity by its area mean and set w=exp(\ensuremath{-}beta J). The weight is continuous, bounded and strictly positive on the closed triangle for every finite beta, even if a different optical model has exact zeros. Its first moment is unique and strictly interior because all three barycentric coordinates have positive weighted integrals. Beta controls preference for darkness. Weak weighting often passed retained response screens, while strong weighting at several Fresnel numbers has stable negative numerical witnesses.


\[J_F=|A_F|^2/\mean_T|A_F|^2,\quad w=e^{-\beta J_F},\quad {\rm Dark}(F,\beta)=\frac{\int_Tpw\dd A}{\int_Tw\dd A}.\]



Normalization gives $\mean_TJ_F=1$.
For finite parameters $e^{-\beta J_F}$ is bounded and positive.
Its mean is strictly interior: barycentric weights are integrals
of positive barycentric functions against positive density.
At $\beta=0$ it is $G$, and
\[
\mathrm{Dark}(F,\beta)=G-\beta
\mathbb E_{\rm area}[(p-G)J_F]+O(\beta^2).
\]
Large $\beta$ concentrates on deepest minima, with multiple minima
weighted by their local geometry. These are not attractivity proofs.
A persistent-darkness variant averages normalized intensity over
a specified $\log F$ interval, then exponentiates and takes a mean.
Averaging amplitudes, intensities or selected centers are different
experiments. Specify the averaging order.



\begin{figure}[htbp]\centering
\begin{tikzpicture}[x=1.7cm,y=1.7cm]
\fill[teal!54!white] (3.142857,0.6428571) -- (2.857143,0.8571429) -- (2.857143,0.6428571) -- cycle;
\fill[teal!12!white] (1.714286,1.714286) -- (1.428571,1.928571) -- (1.428571,1.714286) -- cycle;
\fill[teal!9!white] (1.714286,1.714286) -- (1.714286,1.5) -- (2,1.5) -- cycle;
\fill[teal!77!white] (0.5714286,2.571429) -- (0.2857143,2.785714) -- (0.2857143,2.571429) -- cycle;
\fill[teal!20!white] (2.571429,1.071429) -- (2.285714,1.285714) -- (2.285714,1.071429) -- cycle;
\fill[teal!94!white] (3.714286,0.2142857) -- (3.428571,0.4285714) -- (3.428571,0.2142857) -- cycle;
\fill[teal!94!white] (3.714286,0.2142857) -- (3.714286,0) -- (4,0) -- cycle;
\fill[teal!34!white] (0.8571429,2.142857) -- (1.142857,2.142857) -- (0.8571429,2.357143) -- cycle;
\fill[teal!20!white] (1.142857,2.142857) -- (1.142857,1.928571) -- (1.428571,1.928571) -- cycle;
\fill[teal!94!white] (0,2.785714) -- (0.2857143,2.785714) -- (0,3) -- cycle;
\fill[teal!53!white] (0.5714286,2.571429) -- (0.5714286,2.357143) -- (0.8571429,2.357143) -- cycle;
\fill[teal!12!white] (2,1.285714) -- (2.285714,1.285714) -- (2,1.5) -- cycle;
\fill[teal!34!white] (2.571429,1.071429) -- (2.571429,0.8571429) -- (2.857143,0.8571429) -- cycle;
\fill[teal!80!white] (3.142857,0.4285714) -- (3.428571,0.4285714) -- (3.142857,0.6428571) -- cycle;
\fill[teal!38!white] (0.2857143,0) -- (0,0.2142857) -- (0,0) -- cycle;
\fill[teal!34!white] (0,0.2142857) -- (0.2857143,0) -- (0.2857143,0.2142857) -- cycle;
\fill[teal!26!white] (0.5714286,0.2142857) -- (0.2857143,0) -- (0.5714286,0) -- cycle;
\fill[teal!28!white] (0.2857143,0) -- (0.5714286,0.2142857) -- (0.2857143,0.2142857) -- cycle;
\fill[teal!10!white] (1.714286,1.5) -- (1.714286,1.714286) -- (1.428571,1.714286) -- cycle;
\fill[teal!10!white] (1.428571,1.5) -- (1.714286,1.5) -- (1.428571,1.714286) -- cycle;
\fill[teal!62!white] (0,1.928571) -- (0,1.714286) -- (0.2857143,1.928571) -- cycle;
\fill[teal!55!white] (0,1.714286) -- (0.2857143,1.714286) -- (0.2857143,1.928571) -- cycle;
\fill[teal!84!white] (0.2857143,2.357143) -- (0,2.571429) -- (0,2.357143) -- cycle;
\fill[teal!82!white] (0,2.571429) -- (0.2857143,2.357143) -- (0.2857143,2.571429) -- cycle;
\fill[teal!78!white] (0,2.142857) -- (0.2857143,2.357143) -- (0,2.357143) -- cycle;
\fill[teal!70!white] (0.2857143,2.357143) -- (0,2.142857) -- (0.2857143,2.142857) -- cycle;
\fill[teal!67!white] (0,2.142857) -- (0,1.928571) -- (0.2857143,1.928571) -- cycle;
\fill[teal!64!white] (0.2857143,2.142857) -- (0,2.142857) -- (0.2857143,1.928571) -- cycle;
\fill[teal!47!white] (0.2857143,1.714286) -- (0.5714286,1.928571) -- (0.2857143,1.928571) -- cycle;
\fill[teal!41!white] (0.5714286,1.714286) -- (0.5714286,1.928571) -- (0.2857143,1.714286) -- cycle;
\fill[teal!54!white] (0.2857143,1.5) -- (0,1.714286) -- (0,1.5) -- cycle;
\fill[teal!51!white] (0,1.714286) -- (0.2857143,1.5) -- (0.2857143,1.714286) -- cycle;
\fill[teal!42!white] (0.2857143,1.5) -- (0.5714286,1.714286) -- (0.2857143,1.714286) -- cycle;
\fill[teal!36!white] (0.5714286,1.714286) -- (0.2857143,1.5) -- (0.5714286,1.5) -- cycle;
\fill[teal!50!white] (0,1.285714) -- (0.2857143,1.5) -- (0,1.5) -- cycle;
\fill[teal!44!white] (0.2857143,1.5) -- (0,1.285714) -- (0.2857143,1.285714) -- cycle;
\fill[teal!45!white] (0,1.285714) -- (0,1.071429) -- (0.2857143,1.071429) -- cycle;
\fill[teal!41!white] (0.2857143,1.285714) -- (0,1.285714) -- (0.2857143,1.071429) -- cycle;
\fill[teal!94!white] (3.428571,0) -- (3.714286,0) -- (3.428571,0.2142857) -- cycle;
\fill[teal!94!white] (3.714286,0) -- (3.714286,0.2142857) -- (3.428571,0.2142857) -- cycle;
\fill[teal!20!white] (1.428571,0) -- (1.714286,0) -- (1.428571,0.2142857) -- cycle;
\fill[teal!19!white] (1.714286,0) -- (1.714286,0.2142857) -- (1.428571,0.2142857) -- cycle;
\fill[teal!22!white] (0.8571429,0) -- (0.5714286,0.2142857) -- (0.5714286,0) -- cycle;
\fill[teal!20!white] (0.5714286,0.2142857) -- (0.8571429,0) -- (0.8571429,0.2142857) -- cycle;
\fill[teal!18!white] (1.142857,0.2142857) -- (1.142857,0) -- (1.428571,0) -- cycle;
\fill[teal!17!white] (1.142857,0.2142857) -- (1.428571,0) -- (1.428571,0.2142857) -- cycle;
\fill[teal!18!white] (1.142857,0.2142857) -- (0.8571429,0) -- (1.142857,0) -- cycle;
\fill[teal!18!white] (0.8571429,0) -- (1.142857,0.2142857) -- (0.8571429,0.2142857) -- cycle;
\fill[teal!16!white] (1.142857,1.928571) -- (1.142857,1.714286) -- (1.428571,1.714286) -- cycle;
\fill[teal!16!white] (1.428571,1.928571) -- (1.142857,1.928571) -- (1.428571,1.714286) -- cycle;
\fill[teal!27!white] (0.8571429,2.142857) -- (1.142857,1.928571) -- (1.142857,2.142857) -- cycle;
\fill[teal!28!white] (1.142857,1.928571) -- (0.8571429,2.142857) -- (0.8571429,1.928571) -- cycle;
\fill[teal!94!white] (0,2.785714) -- (0,2.571429) -- (0.2857143,2.571429) -- cycle;
\fill[teal!92!white] (0.2857143,2.785714) -- (0,2.785714) -- (0.2857143,2.571429) -- cycle;
\fill[teal!21!white] (1.142857,1.928571) -- (0.8571429,1.714286) -- (1.142857,1.714286) -- cycle;
\fill[teal!25!white] (0.8571429,1.714286) -- (1.142857,1.928571) -- (0.8571429,1.928571) -- cycle;
\fill[teal!33!white] (0.8571429,1.714286) -- (0.5714286,1.928571) -- (0.5714286,1.714286) -- cycle;
\fill[teal!31!white] (0.5714286,1.928571) -- (0.8571429,1.714286) -- (0.8571429,1.928571) -- cycle;
\fill[teal!68!white] (0.2857143,2.357143) -- (0.5714286,2.357143) -- (0.2857143,2.571429) -- cycle;
\fill[teal!66!white] (0.5714286,2.357143) -- (0.5714286,2.571429) -- (0.2857143,2.571429) -- cycle;
\fill[teal!38!white] (0,0.6428571) -- (0,0.4285714) -- (0.2857143,0.6428571) -- cycle;
\fill[teal!34!white] (0,0.4285714) -- (0.2857143,0.4285714) -- (0.2857143,0.6428571) -- cycle;
\fill[teal!38!white] (0,0.4285714) -- (0,0.2142857) -- (0.2857143,0.2142857) -- cycle;
\fill[teal!34!white] (0.2857143,0.4285714) -- (0,0.4285714) -- (0.2857143,0.2142857) -- cycle;
\fill[teal!14!white] (2.285714,1.285714) -- (2,1.285714) -- (2.285714,1.071429) -- cycle;
\fill[teal!13!white] (2,1.285714) -- (2,1.071429) -- (2.285714,1.071429) -- cycle;
\fill[teal!14!white] (2,1.071429) -- (2,0.8571429) -- (2.285714,1.071429) -- cycle;
\fill[teal!17!white] (2,0.8571429) -- (2.285714,0.8571429) -- (2.285714,1.071429) -- cycle;
\fill[teal!23!white] (2.285714,0.8571429) -- (2.571429,0.8571429) -- (2.285714,1.071429) -- cycle;
\fill[teal!25!white] (2.571429,0.8571429) -- (2.571429,1.071429) -- (2.285714,1.071429) -- cycle;
\fill[teal!15!white] (1.428571,0.4285714) -- (1.714286,0.2142857) -- (1.714286,0.4285714) -- cycle;
\fill[teal!15!white] (1.714286,0.2142857) -- (1.428571,0.4285714) -- (1.428571,0.2142857) -- cycle;
\fill[teal!94!white] (3.142857,0.2142857) -- (3.428571,0) -- (3.428571,0.2142857) -- cycle;
\fill[teal!94!white] (3.142857,0.2142857) -- (3.142857,0) -- (3.428571,0) -- cycle;
\fill[teal!83!white] (3.142857,0) -- (3.142857,0.2142857) -- (2.857143,0) -- cycle;
\fill[teal!73!white] (3.142857,0.2142857) -- (2.857143,0.2142857) -- (2.857143,0) -- cycle;
\fill[teal!23!white] (1.714286,0) -- (2,0.2142857) -- (1.714286,0.2142857) -- cycle;
\fill[teal!27!white] (2,0.2142857) -- (1.714286,0) -- (2,0) -- cycle;
\fill[teal!25!white] (0.5714286,0.2142857) -- (0.5714286,0.4285714) -- (0.2857143,0.2142857) -- cycle;
\fill[teal!27!white] (0.5714286,0.4285714) -- (0.2857143,0.4285714) -- (0.2857143,0.2142857) -- cycle;
\fill[teal!19!white] (0.8571429,0.4285714) -- (0.5714286,0.2142857) -- (0.8571429,0.2142857) -- cycle;
\fill[teal!20!white] (0.8571429,0.4285714) -- (0.5714286,0.4285714) -- (0.5714286,0.2142857) -- cycle;
\fill[teal!12!white] (1.142857,1.5) -- (1.428571,1.5) -- (1.428571,1.714286) -- cycle;
\fill[teal!14!white] (1.142857,1.714286) -- (1.142857,1.5) -- (1.428571,1.714286) -- cycle;
\fill[teal!54!white] (0.5714286,2.142857) -- (0.2857143,2.142857) -- (0.2857143,1.928571) -- cycle;
\fill[teal!47!white] (0.5714286,1.928571) -- (0.5714286,2.142857) -- (0.2857143,1.928571) -- cycle;
\fill[teal!59!white] (0.5714286,2.142857) -- (0.2857143,2.357143) -- (0.2857143,2.142857) -- cycle;
\fill[teal!57!white] (0.5714286,2.142857) -- (0.5714286,2.357143) -- (0.2857143,2.357143) -- cycle;
\fill[teal!37!white] (0.8571429,2.142857) -- (0.5714286,2.142857) -- (0.8571429,1.928571) -- cycle;
\fill[teal!39!white] (0.5714286,2.142857) -- (0.5714286,1.928571) -- (0.8571429,1.928571) -- cycle;
\fill[teal!41!white] (0.5714286,2.142857) -- (0.8571429,2.142857) -- (0.8571429,2.357143) -- cycle;
\fill[teal!48!white] (0.5714286,2.357143) -- (0.5714286,2.142857) -- (0.8571429,2.357143) -- cycle;
\fill[teal!9!white] (1.714286,1.285714) -- (2,1.285714) -- (2,1.5) -- cycle;
\fill[teal!9!white] (1.714286,1.5) -- (1.714286,1.285714) -- (2,1.5) -- cycle;
\fill[teal!9!white] (1.142857,1.071429) -- (1.428571,0.8571429) -- (1.428571,1.071429) -- cycle;
\fill[teal!10!white] (1.428571,0.8571429) -- (1.142857,1.071429) -- (1.142857,0.8571429) -- cycle;
\fill[teal!88!white] (3.142857,0.4285714) -- (3.142857,0.2142857) -- (3.428571,0.2142857) -- cycle;
\fill[teal!93!white] (3.428571,0.4285714) -- (3.142857,0.4285714) -- (3.428571,0.2142857) -- cycle;
\fill[teal!60!white] (2.857143,0.4285714) -- (3.142857,0.4285714) -- (2.857143,0.6428571) -- cycle;
\fill[teal!64!white] (3.142857,0.4285714) -- (3.142857,0.6428571) -- (2.857143,0.6428571) -- cycle;
\fill[teal!74!white] (3.142857,0.2142857) -- (3.142857,0.4285714) -- (2.857143,0.2142857) -- cycle;
\fill[teal!65!white] (3.142857,0.4285714) -- (2.857143,0.4285714) -- (2.857143,0.2142857) -- cycle;
\fill[teal!37!white] (2.571429,0.8571429) -- (2.571429,0.6428571) -- (2.857143,0.8571429) -- cycle;
\fill[teal!43!white] (2.857143,0.8571429) -- (2.571429,0.6428571) -- (2.857143,0.6428571) -- cycle;
\fill[teal!36!white] (2.285714,0) -- (2.285714,0.2142857) -- (2,0) -- cycle;
\fill[teal!30!white] (2.285714,0.2142857) -- (2,0.2142857) -- (2,0) -- cycle;
\fill[teal!27!white] (2.285714,0.2142857) -- (2,0.4285714) -- (2,0.2142857) -- cycle;
\fill[teal!28!white] (2,0.4285714) -- (2.285714,0.2142857) -- (2.285714,0.4285714) -- cycle;
\fill[teal!16!white] (2,0.4285714) -- (1.714286,0.6428571) -- (1.714286,0.4285714) -- cycle;
\fill[teal!16!white] (1.714286,0.6428571) -- (2,0.4285714) -- (2,0.6428571) -- cycle;
\fill[teal!18!white] (1.714286,0.2142857) -- (2,0.4285714) -- (1.714286,0.4285714) -- cycle;
\fill[teal!22!white] (2,0.2142857) -- (2,0.4285714) -- (1.714286,0.2142857) -- cycle;
\fill[teal!40!white] (0.2857143,0.8571429) -- (0,0.8571429) -- (0,0.6428571) -- cycle;
\fill[teal!35!white] (0.2857143,0.8571429) -- (0,0.6428571) -- (0.2857143,0.6428571) -- cycle;
\fill[teal!38!white] (0,1.071429) -- (0.2857143,0.8571429) -- (0.2857143,1.071429) -- cycle;
\fill[teal!41!white] (0,0.8571429) -- (0.2857143,0.8571429) -- (0,1.071429) -- cycle;
\fill[teal!33!white] (0.2857143,1.5) -- (0.5714286,1.285714) -- (0.5714286,1.5) -- cycle;
\fill[teal!36!white] (0.5714286,1.285714) -- (0.2857143,1.5) -- (0.2857143,1.285714) -- cycle;
\fill[teal!28!white] (0.5714286,0.6428571) -- (0.2857143,0.8571429) -- (0.2857143,0.6428571) -- cycle;
\fill[teal!25!white] (0.2857143,0.8571429) -- (0.5714286,0.6428571) -- (0.5714286,0.8571429) -- cycle;
\fill[teal!27!white] (0.2857143,0.4285714) -- (0.5714286,0.6428571) -- (0.2857143,0.6428571) -- cycle;
\fill[teal!24!white] (0.5714286,0.4285714) -- (0.5714286,0.6428571) -- (0.2857143,0.4285714) -- cycle;
\fill[teal!19!white] (0.8571429,0.4285714) -- (0.5714286,0.6428571) -- (0.5714286,0.4285714) -- cycle;
\fill[teal!17!white] (0.5714286,0.6428571) -- (0.8571429,0.4285714) -- (0.8571429,0.6428571) -- cycle;
\fill[teal!14!white] (1.142857,0.4285714) -- (1.142857,0.2142857) -- (1.428571,0.2142857) -- cycle;
\fill[teal!13!white] (1.428571,0.4285714) -- (1.142857,0.4285714) -- (1.428571,0.2142857) -- cycle;
\fill[teal!15!white] (1.142857,0.2142857) -- (1.142857,0.4285714) -- (0.8571429,0.2142857) -- cycle;
\fill[teal!15!white] (1.142857,0.4285714) -- (0.8571429,0.4285714) -- (0.8571429,0.2142857) -- cycle;
\fill[teal!14!white] (0.8571429,0.4285714) -- (1.142857,0.4285714) -- (0.8571429,0.6428571) -- cycle;
\fill[teal!13!white] (1.142857,0.4285714) -- (1.142857,0.6428571) -- (0.8571429,0.6428571) -- cycle;
\fill[teal!9!white] (1.428571,1.285714) -- (1.714286,1.5) -- (1.428571,1.5) -- cycle;
\fill[teal!8!white] (1.428571,1.285714) -- (1.714286,1.285714) -- (1.714286,1.5) -- cycle;
\fill[teal!10!white] (1.142857,1.5) -- (1.428571,1.285714) -- (1.428571,1.5) -- cycle;
\fill[teal!11!white] (1.428571,1.285714) -- (1.142857,1.5) -- (1.142857,1.285714) -- cycle;
\fill[teal!9!white] (1.428571,1.285714) -- (1.142857,1.071429) -- (1.428571,1.071429) -- cycle;
\fill[teal!10!white] (1.142857,1.071429) -- (1.428571,1.285714) -- (1.142857,1.285714) -- cycle;
\fill[teal!12!white] (1.428571,0.6428571) -- (1.428571,0.4285714) -- (1.714286,0.4285714) -- cycle;
\fill[teal!12!white] (1.714286,0.6428571) -- (1.428571,0.6428571) -- (1.714286,0.4285714) -- cycle;
\fill[teal!12!white] (1.428571,0.6428571) -- (1.142857,0.4285714) -- (1.428571,0.4285714) -- cycle;
\fill[teal!11!white] (1.142857,0.4285714) -- (1.428571,0.6428571) -- (1.142857,0.6428571) -- cycle;
\fill[teal!9!white] (1.428571,0.6428571) -- (1.428571,0.8571429) -- (1.142857,0.8571429) -- cycle;
\fill[teal!10!white] (1.142857,0.6428571) -- (1.428571,0.6428571) -- (1.142857,0.8571429) -- cycle;
\fill[teal!10!white] (2,1.285714) -- (1.714286,1.071429) -- (2,1.071429) -- cycle;
\fill[teal!9!white] (1.714286,1.285714) -- (1.714286,1.071429) -- (2,1.285714) -- cycle;
\fill[teal!8!white] (1.428571,1.285714) -- (1.714286,1.071429) -- (1.714286,1.285714) -- cycle;
\fill[teal!8!white] (1.714286,1.071429) -- (1.428571,1.285714) -- (1.428571,1.071429) -- cycle;
\fill[teal!25!white] (2.285714,0.6428571) -- (2.571429,0.8571429) -- (2.285714,0.8571429) -- cycle;
\fill[teal!30!white] (2.285714,0.6428571) -- (2.571429,0.6428571) -- (2.571429,0.8571429) -- cycle;
\fill[teal!19!white] (2,0.8571429) -- (2.285714,0.6428571) -- (2.285714,0.8571429) -- cycle;
\fill[teal!18!white] (2.285714,0.6428571) -- (2,0.8571429) -- (2,0.6428571) -- cycle;
\fill[teal!25!white] (2.285714,0.6428571) -- (2,0.4285714) -- (2.285714,0.4285714) -- cycle;
\fill[teal!21!white] (2,0.4285714) -- (2.285714,0.6428571) -- (2,0.6428571) -- cycle;
\fill[teal!49!white] (2.571429,0.4285714) -- (2.857143,0.4285714) -- (2.857143,0.6428571) -- cycle;
\fill[teal!42!white] (2.571429,0.6428571) -- (2.571429,0.4285714) -- (2.857143,0.6428571) -- cycle;
\fill[teal!34!white] (2.285714,0.6428571) -- (2.571429,0.4285714) -- (2.571429,0.6428571) -- cycle;
\fill[teal!31!white] (2.571429,0.4285714) -- (2.285714,0.6428571) -- (2.285714,0.4285714) -- cycle;
\fill[teal!30!white] (0.2857143,0.8571429) -- (0.5714286,1.071429) -- (0.2857143,1.071429) -- cycle;
\fill[teal!26!white] (0.5714286,1.071429) -- (0.2857143,0.8571429) -- (0.5714286,0.8571429) -- cycle;
\fill[teal!32!white] (0.5714286,1.071429) -- (0.2857143,1.285714) -- (0.2857143,1.071429) -- cycle;
\fill[teal!29!white] (0.5714286,1.071429) -- (0.5714286,1.285714) -- (0.2857143,1.285714) -- cycle;
\fill[teal!28!white] (0.8571429,1.5) -- (0.5714286,1.714286) -- (0.5714286,1.5) -- cycle;
\fill[teal!26!white] (0.8571429,1.5) -- (0.8571429,1.714286) -- (0.5714286,1.714286) -- cycle;
\fill[teal!21!white] (0.8571429,1.714286) -- (0.8571429,1.5) -- (1.142857,1.714286) -- cycle;
\fill[teal!17!white] (0.8571429,1.5) -- (1.142857,1.5) -- (1.142857,1.714286) -- cycle;
\fill[teal!12!white] (1.714286,0.8571429) -- (1.714286,0.6428571) -- (2,0.6428571) -- cycle;
\fill[teal!13!white] (2,0.8571429) -- (1.714286,0.8571429) -- (2,0.6428571) -- cycle;
\fill[teal!12!white] (1.714286,0.8571429) -- (2,0.8571429) -- (2,1.071429) -- cycle;
\fill[teal!10!white] (1.714286,1.071429) -- (1.714286,0.8571429) -- (2,1.071429) -- cycle;
\fill[teal!10!white] (1.714286,0.8571429) -- (1.428571,0.6428571) -- (1.714286,0.6428571) -- cycle;
\fill[teal!9!white] (1.428571,0.6428571) -- (1.714286,0.8571429) -- (1.428571,0.8571429) -- cycle;
\fill[teal!9!white] (1.428571,0.8571429) -- (1.714286,0.8571429) -- (1.428571,1.071429) -- cycle;
\fill[teal!8!white] (1.714286,0.8571429) -- (1.714286,1.071429) -- (1.428571,1.071429) -- cycle;
\fill[teal!48!white] (2.571429,0) -- (2.571429,0.2142857) -- (2.285714,0) -- cycle;
\fill[teal!41!white] (2.571429,0.2142857) -- (2.285714,0.2142857) -- (2.285714,0) -- cycle;
\fill[teal!37!white] (2.285714,0.2142857) -- (2.571429,0.2142857) -- (2.285714,0.4285714) -- cycle;
\fill[teal!39!white] (2.571429,0.2142857) -- (2.571429,0.4285714) -- (2.285714,0.4285714) -- cycle;
\fill[teal!61!white] (2.857143,0.2142857) -- (2.571429,0.2142857) -- (2.857143,0) -- cycle;
\fill[teal!58!white] (2.571429,0.2142857) -- (2.571429,0) -- (2.857143,0) -- cycle;
\fill[teal!56!white] (2.857143,0.4285714) -- (2.571429,0.2142857) -- (2.857143,0.2142857) -- cycle;
\fill[teal!49!white] (2.571429,0.4285714) -- (2.571429,0.2142857) -- (2.857143,0.4285714) -- cycle;
\fill[teal!13!white] (1.142857,0.6428571) -- (0.8571429,0.8571429) -- (0.8571429,0.6428571) -- cycle;
\fill[teal!12!white] (0.8571429,0.8571429) -- (1.142857,0.6428571) -- (1.142857,0.8571429) -- cycle;
\fill[teal!20!white] (0.5714286,0.6428571) -- (0.8571429,0.8571429) -- (0.5714286,0.8571429) -- cycle;
\fill[teal!17!white] (0.8571429,0.8571429) -- (0.5714286,0.6428571) -- (0.8571429,0.6428571) -- cycle;
\fill[teal!14!white] (1.142857,1.5) -- (0.8571429,1.285714) -- (1.142857,1.285714) -- cycle;
\fill[teal!17!white] (0.8571429,1.5) -- (0.8571429,1.285714) -- (1.142857,1.5) -- cycle;
\fill[teal!25!white] (0.5714286,1.285714) -- (0.8571429,1.285714) -- (0.5714286,1.5) -- cycle;
\fill[teal!23!white] (0.8571429,1.285714) -- (0.8571429,1.5) -- (0.5714286,1.5) -- cycle;
\fill[teal!12!white] (1.142857,1.071429) -- (0.8571429,1.071429) -- (1.142857,0.8571429) -- cycle;
\fill[teal!14!white] (0.8571429,1.071429) -- (0.8571429,0.8571429) -- (1.142857,0.8571429) -- cycle;
\fill[teal!13!white] (0.8571429,1.071429) -- (1.142857,1.071429) -- (1.142857,1.285714) -- cycle;
\fill[teal!15!white] (0.8571429,1.285714) -- (0.8571429,1.071429) -- (1.142857,1.285714) -- cycle;
\fill[teal!21!white] (0.8571429,1.071429) -- (0.5714286,1.071429) -- (0.5714286,0.8571429) -- cycle;
\fill[teal!18!white] (0.8571429,0.8571429) -- (0.8571429,1.071429) -- (0.5714286,0.8571429) -- cycle;
\fill[teal!22!white] (0.5714286,1.071429) -- (0.8571429,1.071429) -- (0.5714286,1.285714) -- cycle;
\fill[teal!20!white] (0.8571429,1.071429) -- (0.8571429,1.285714) -- (0.5714286,1.285714) -- cycle;
\draw[teal,thick] (0,0) -- (4,0) -- (0,3) -- cycle;
\draw[magenta!70,thin] (1.3664,1.013379) -- (3.6,2.35) -- (4.1,2.35);\filldraw[fill=magenta,draw=white,line width=.45pt] (1.3664,1.013379) circle(1.5pt) node[above right,font=\scriptsize,text=magenta,fill=white,inner sep=.6pt] {};
\node[anchor=west,text=magenta,font=\small] at(4.12,2.35){Dark1.0b1.0};\draw[gold!70,thin] (1,1) -- (3.6,1.8) -- (4.1,1.8);\filldraw[fill=gold,draw=white,line width=.45pt] (1,1) circle(1.5pt) node[above right,font=\scriptsize,text=gold,fill=white,inner sep=.6pt] {};
\node[anchor=west,text=gold,font=\small] at(4.12,1.8){$I$};\draw[navy!70,thin] (1.333333,1) -- (3.6,1.25) -- (4.1,1.25);\filldraw[fill=navy,draw=white,line width=.45pt] (1.333333,1) circle(1.5pt) node[above right,font=\scriptsize,text=navy,fill=white,inner sep=.6pt] {};
\node[anchor=west,text=navy,font=\small] at(4.12,1.25){$G$};\end{tikzpicture}
\caption[The weight $\exp(-J)$ at $F=1$, $\beta=1$]{The weight $\exp(-J)$ at $F=1$, $\beta=1$. Greater saturation denotes greater darkness weight, rather than greater optical intensity. Fields are archived numerical illustrations; teal saturation increases with the displayed value.}\label{fig:field-Dark1.0b1.0}
\end{figure}

\section{Wavelength independence needs a convention}\label{sec:slide84}

A triangle center cannot depend on an unscaled external wavelength if it is supposed to be intrinsic to the triangle alone. Holding the Fresnel number fixed removes that length ambiguity, but physically this requires adjusting the screen distance as the wavelength changes. It is not an achromatic point on one fixed screen. The atlas persistence mean first averages the normalized intensity J over log F in [1/2,2], then applies exp(\ensuremath{-}beta Jbar) and takes the first moment. The order of operations, range and measure are part of its definition.


\[F=L^2/(\lambda z)\ \hbox{fixed}\quad\Longleftrightarrow\quad\lambda z\ \hbox{fixed for a fixed aperture}.\]


For the same triangle, z=$L^2$/(lambda F). Thus doubling lambda halves z at fixed F. A log-F average uses dF/F, not dF. Averaging center positions or exponential response images would give different constructions. Persistent dark means with beta one and beta eight have different response evidence.



\section{Gaussian images and nonlinear fluorescence}\label{sec:slide85}

Convolve a uniform triangular image with a centered isotropic Gaussian point-spread function. On the entire plane the first moment remains G, because adding an independent zero-mean blur does not change expectation. Covariances add, so the image covariance is the triangle covariance plus s squared times the identity. The same principal axes remain, while the axis ratio moves toward one. This connects optical imaging to the Steiner covariance structure. A detector response proportional to the square of the blurred intensity gives a different nonlinear mean, called Opt(sigma) in the atlas.


\[\mean(\mathbf1_T*g_s)=G,\qquad\Sigma_{\rm blurred}=\Sigma_{\rm area}+s^2\mathbf1.\]



Normalize the source indicator by area and the Gaussian by mass.
The blurred image is the density of $X+Z$, with $X$ uniform on $T$
and $Z$ independent, centered and of covariance $s^2\mathbf1$.
Expectation and covariance addition prove the displayed formulas.
Principal directions remain unchanged while the ellipse becomes
rounder. This uses the whole plane; cropping changes the moments.

A nonlinear detector responding to squared intensity instead sees
$(\mathbf1_T*g_s)^2$. This is not a convolution probability,
and its mean generally differs from $G$.
Specify width relative to $L$, response power and observation window
to obtain a reproducible new family.




\section{Optical fingerprints and ETC neighbors}\label{sec:slide86}

A complex wave integral leads to a positive weight and then an ordinary barycentric mean. Its shape fingerprints can be compared with the same familiar centers and ETC candidates. Keep Fresnel number and beta fixed across the family. The diagrams show the actual archived darkness center and its location neighbors. A different beta may move the center substantially and change response signs. Numerical field convergence must be checked independently of the ETC fit.


For the diffraction integral use boundary-adapted or radial quadrature and compare independent formulations. Highly concentrated darkness weights can magnify small intensity errors. The project's polar and area integrals agree on selected checks, but this does not automatically certify every shape or parameter.


\begin{figure}[htbp]\centering
\begin{tikzpicture}\begin{axis}[width=.95\linewidth,height=.56\linewidth,grid=major,grid style={black!10},xlabel={family parameter $x$},ylabel={$c_y$},tick label style={font=\scriptsize},label style={font=\small},legend style={font=\scriptsize,at={(0.5,-0.24)},anchor=north,legend columns=3,draw=none},scaled ticks=false]
\addplot[navy,thick,no marks] coordinates {(0.2,0.06067713) (0.228706,0.06945129) (0.2615321,0.07951623) (0.2990698,0.09107199) (0.3419952,0.1043538) (0.3910817,0.1196398) (0.4472136,0.1372603) (0.5114021,0.1576084) (0.5848035,0.1811529) (0.6687403,0.2084494) (0.7647245,0.2401498) (0.8744853,0.2770036) (1,0.3198471) (1.14353,0.3695797) (1.30766,0.4271338) (1.495349,0.4934566) (1.709976,0.5695265) (1.732051,0.5773503) (1.955409,0.6564131) (2.236068,0.7553669) (2.55701,0.8679087) (2.924018,0.9958947) (3.343702,1.14155) (3.823622,1.307488) (4.372426,1.496726) (5,1.71272)};
\addlegendentry{Dark1.0b1.0}
\addplot[teal,thick,no marks] coordinates {(0.2,0.09901951) (0.228706,0.1128955) (0.2615321,0.1286033) (0.2990698,0.1463328) (0.3419952,0.1662702) (0.3910817,0.1885865) (0.4472136,0.2134218) (0.5114021,0.2408662) (0.5848035,0.2709374) (0.6687403,0.3035587) (0.7647245,0.33854) (0.8744853,0.3755687) (1,0.4142136) (1.14353,0.4539441) (1.30766,0.4941653) (1.495349,0.5342616) (1.709976,0.5736416) (1.732051,0.5773503) (1.955409,0.6117774) (2.236068,0.6482315) (2.55701,0.682671) (2.924018,0.7148684) (3.343702,0.744694) (3.823622,0.7721018) (4.372426,0.7971139) (5,0.8198039)};
\addlegendentry{X1}
\addplot[magenta,thick,no marks] coordinates {(0.2,0.06666667) (0.228706,0.07623532) (0.2615321,0.08717737) (0.2990698,0.09968992) (0.3419952,0.1139984) (0.3910817,0.1303606) (0.4472136,0.1490712) (0.5114021,0.1704674) (0.5848035,0.1949345) (0.6687403,0.2229134) (0.7647245,0.2549082) (0.8744853,0.2914951) (1,0.3333333) (1.14353,0.3811766) (1.30766,0.4358868) (1.495349,0.4984496) (1.709976,0.569992) (1.732051,0.5773503) (1.955409,0.6518028) (2.236068,0.745356) (2.55701,0.8523368) (2.924018,0.9746726) (3.343702,1.114567) (3.823622,1.274541) (4.372426,1.457475) (5,1.666667)};
\addlegendentry{X2}
\addplot[gold,thick,no marks] coordinates {(0.2,0.05049024) (0.228706,0.05790523) (0.2615321,0.06646438) (0.2990698,0.07636846) (0.3419952,0.08786248) (0.3910817,0.1012476) (0.4472136,0.1168959) (0.5114021,0.135268) (0.5848035,0.1569331) (0.6687403,0.1825908) (0.7647245,0.2130923) (0.8744853,0.2494583) (1,0.2928932) (1.14353,0.3447929) (1.30766,0.4067476) (1.495349,0.4805436) (1.709976,0.5681672) (1.732051,0.5773503) (1.955409,0.6718156) (2.236068,0.7939182) (2.55701,0.9371697) (2.924018,1.104575) (3.343702,1.299504) (3.823622,1.52576) (4.372426,1.787656) (5,2.090098)};
\addlegendentry{X10}
\addplot[blue,thick,no marks] coordinates {(0.2,0.06141333) (0.228706,0.07029244) (0.2615321,0.08047638) (0.2990698,0.09216611) (0.3419952,0.1055971) (0.3910817,0.1210458) (0.4472136,0.1388382) (0.5114021,0.159358) (0.5848035,0.1830571) (0.6687403,0.2104645) (0.7647245,0.2421954) (0.8744853,0.2789567) (1,0.32155) (1.14353,0.3708713) (1.30766,0.4279081) (1.495349,0.4937384) (1.709976,0.5695336) (1.732051,0.5773503) (1.955409,0.6565703) (2.236068,0.7562522) (2.55701,0.8701433) (2.924018,1.000011) (3.343702,1.147878) (3.823622,1.316076) (4.372426,1.507306) (5,1.724704)};
\addlegendentry{X17210}
\addplot[purple,thick,no marks] coordinates {(0.2,0.05974366) (0.228706,0.06841511) (0.2615321,0.07837621) (0.2990698,0.08983196) (0.3419952,0.1030246) (0.3910817,0.118241) (0.4472136,0.1358217) (0.5114021,0.15617) (0.5848035,0.1797613) (0.6687403,0.2071511) (0.7647245,0.2389798) (0.8744853,0.2759744) (1,0.3189452) (1.14353,0.3687806) (1.30766,0.4264422) (1.495349,0.492965) (1.709976,0.5694673) (1.732051,0.5773503) (1.955409,0.6571704) (2.236068,0.7574303) (2.55701,0.8717775) (2.924018,1.001964) (3.343702,1.150011) (3.823622,1.318265) (4.372426,1.509451) (5,1.726729)};
\addlegendentry{X17248}
\end{axis}\end{tikzpicture}
\caption[Isosceles fingerprint: the physical selector, two retained ETC location neighbors and the three classical controls]{Isosceles fingerprint: the physical selector, two retained ETC location neighbors and the three classical controls. The optical parameters are fixed at $F=1$, $\beta=1$. Similarity of sampled curves is not identity.}\label{fig:fp-Dark1.0b1.0}
\end{figure}
\begin{figure}[htbp]\centering
\begin{tikzpicture}\begin{axis}[width=.95\linewidth,height=.56\linewidth,grid=major,grid style={black!10},xlabel={apex height $x$},ylabel={$dc_y/dx$},tick label style={font=\scriptsize},label style={font=\small},legend style={font=\scriptsize,at={(0.5,-0.24)},anchor=north,legend columns=3,draw=none},scaled ticks=false]
\addplot[navy,thick,no marks] coordinates {(0.2,0.3052455) (0.228706,0.3060842) (0.2615321,0.3071646) (0.2990698,0.3085497) (0.3419952,0.3103138) (0.3910817,0.3125408) (0.4472136,0.3153188) (0.5114021,0.3187283) (0.5848035,0.3228207) (0.6687403,0.3275854) (0.7647245,0.3329056) (0.8744853,0.3385149) (1,0.3439792) (1.14353,0.3487416) (1.30766,0.3522528) (1.495349,0.3541576) (1.709976,0.3544413) (1.732051,0.3543934) (1.955409,0.3534364) (2.236068,0.351676) (2.55701,0.3496895) (2.924018,0.3478553) (3.343702,0.3463613) (3.823622,0.3452453) (4.372426,0.3444613) (5,0.3439338)};
\addlegendentry{Dark1.0b1.0}
\addplot[teal,thick,no marks] coordinates {(0.2,0.4833837) (0.228706,0.4807868) (0.2615321,0.4752065) (0.2990698,0.4681264) (0.3419952,0.4592189) (0.3910817,0.4481301) (0.4472136,0.434504) (0.5114021,0.4180222) (0.5848035,0.3984569) (0.6687403,0.3757348) (0.7647245,0.3499968) (0.8744853,0.3216382) (1,0.2913098) (1.14353,0.2598701) (1.30766,0.2282923) (1.495349,0.1975472) (1.709976,0.1820376) (1.732051,0.1553818) (1.955409,0.1406326) (2.236068,0.1178415) (2.55701,0.09686305) (2.924018,0.0788403) (3.343702,0.06362058) (3.823622,0.05095623) (4.372426,0.04054999) (5,0.03615511)};
\addlegendentry{X1}
\addplot[magenta,thick,no marks] coordinates {(0.2,0.3333333) (0.228706,0.3333333) (0.2615321,0.3333333) (0.2990698,0.3333333) (0.3419952,0.3333333) (0.3910817,0.3333333) (0.4472136,0.3333333) (0.5114021,0.3333333) (0.5848035,0.3333333) (0.6687403,0.3333333) (0.7647245,0.3333333) (0.8744853,0.3333333) (1,0.3333333) (1.14353,0.3333333) (1.30766,0.3333333) (1.495349,0.3333333) (1.709976,0.3333333) (1.732051,0.3333333) (1.955409,0.3333333) (2.236068,0.3333333) (2.55701,0.3333333) (2.924018,0.3333333) (3.343702,0.3333333) (3.823622,0.3333333) (4.372426,0.3333333) (5,0.3333333)};
\addlegendentry{X2}
\addplot[gold,thick,no marks] coordinates {(0.2,0.2583082) (0.228706,0.2596066) (0.2615321,0.2623967) (0.2990698,0.2659368) (0.3419952,0.2703905) (0.3910817,0.275935) (0.4472136,0.282748) (0.5114021,0.2909889) (0.5848035,0.3007716) (0.6687403,0.3121326) (0.7647245,0.3250016) (0.8744853,0.3391809) (1,0.3543451) (1.14353,0.370065) (1.30766,0.3858539) (1.495349,0.4012264) (1.709976,0.4089812) (1.732051,0.4223091) (1.955409,0.4296837) (2.236068,0.4410793) (2.55701,0.4515685) (2.924018,0.4605799) (3.343702,0.4681897) (3.823622,0.4745219) (4.372426,0.479725) (5,0.4819224)};
\addlegendentry{X10}
\addplot[blue,thick,no marks] coordinates {(0.2,0.3093125) (0.228706,0.3098067) (0.2615321,0.3108655) (0.2990698,0.3122012) (0.3419952,0.313869) (0.3910817,0.3159248) (0.4472136,0.318418) (0.5114021,0.3213817) (0.5848035,0.3248195) (0.6687403,0.3286907) (0.7647245,0.3328983) (0.8744853,0.3372838) (1,0.3416334) (1.14353,0.3456995) (1.30766,0.3492339) (1.495349,0.352026) (1.709976,0.3532369) (1.732051,0.3546257) (1.955409,0.354952) (2.236068,0.3550072) (2.55701,0.354327) (2.924018,0.3530416) (3.343702,0.3513367) (3.823622,0.3493917) (4.372426,0.3473616) (5,0.3464103)};
\addlegendentry{X17210}
\addplot[purple,thick,no marks] coordinates {(0.2,0.3020782) (0.228706,0.3028103) (0.2615321,0.3043732) (0.2990698,0.306331) (0.3419952,0.3087536) (0.3910817,0.3117055) (0.4472136,0.3152338) (0.5114021,0.319352) (0.5848035,0.3240219) (0.6687403,0.3291358) (0.7647245,0.3345078) (0.8744853,0.3398801) (1,0.3449476) (1.14353,0.3494013) (1.30766,0.3529781) (1.495349,0.3555048) (1.709976,0.3565042) (1.732051,0.3573409) (1.955409,0.3572894) (2.236068,0.3567261) (2.55701,0.3554523) (2.924018,0.3536754) (3.343702,0.3516007) (3.823622,0.3494036) (4.372426,0.3472214) (5,0.3462186)};
\addlegendentry{X17248}
\end{axis}\end{tikzpicture}
\caption[Isosceles derivative comparison]{Isosceles derivative comparison. Neighbor derivatives use sampled differences, and do not replace the full vertex-response matrices.}\label{fig:deriv-Dark1.0b1.0}
\end{figure}

\section{The most useful physical definitions}\label{sec:slide87}

Review the most intuitive examples without trying to memorize labels. G and W are means of prescribed measures. E2 is a mean of an equilibrium measure. M1 is a scalar-field minimum. F1 is a Poisson maximum that appears in several disciplines. $\mathrm{Q0}$ is a quantum probability mean. H is a time-dependent field maximum. Dark is a positive mean of a coherent-wave intensity transform. The mathematical challenge follows the definition: integration, singular kernels, nonlinear optimization, sparse PDE systems, spectral gaps or oscillatory quadrature.


\begin{center}\small\begin{tabular}{p{0.4\linewidth}p{0.5\linewidth}}\toprule
Physical problem & Mathematical problem\\\midrule
Prescribed mass or charge & First and second moments\\
Mobile conductor charge & Singular energy minimization\\
Membrane, duct, exit time & Poisson equation\\
Quantum confinement & Dirichlet eigenproblem\\
Cooling; removal & Evolution; screened Poisson\\
Diffraction & Oscillatory integral and moments\\\bottomrule\end{tabular}\end{center}

\paragraph{Demonstration and investigation.} Ask students to choose one row and explain its first-principles equation in a sentence. A correct physical assumption matters more than remembering an atlas index.


\section{Demonstrations students can reproduce}\label{sec:slide88}

Offer a practical sequence scaled to an undergraduate project. First locate G and W mechanically with suspension or balancing. Next map a triangular current loop above its plane, documenting wire radius, leads and probe height. A numerical thermal image can be produced by the supplied heat solver and displayed with a time slider. The quantum and Brownian examples need only a Dirichlet mesh and a random-number generator. For diffraction, a triangle aperture and a camera can compare patterns at matched Fresnel number, while the center extraction is done in software. Each experiment should record its physical regularization and its statistical or numerical uncertainty.



\paragraph{Demonstration and investigation.} Open the offline demonstrations from the course package. They need no internet. Drag a vertex in the geometry app, cool the triangle, trace a billiard path, then adjust optical scale. Use each demo only after stating the governing assumptions.


\section{Attractivity after the new investigation}\label{sec:slide89}

The latest record matters more than an attractive-looking contour. The four distinct points I, G, W and X360 remain proved attractive. M6 remains the strongest new positive conjecture after targeted thin-shape testing. E0 has a new interval-certified negative finite motion, correcting the earlier conjecture. The Abs family has a proof of failure for sufficiently large removal, with independent numerical evidence at mu ten thousand. The solid-angle family fails at sufficiently small height through its M1 limit. The moderate parameter members retain separate open statuses.



Four distinct centers have general attractivity proofs here:
$X_1,X_2,X_{10},X_{360}$, including their physical aliases.
E0, E1, E2 and M1 have certified failures.
M6 has strong positive evidence, not a theorem.
F1 and Q0m remain refinement sensitive.
The state means Q0--Q4, E3, tested heat centers and Illum have
numerical negative witnesses.

Abs$(0,1,10,100)$ denotes four positively screened parameter choices,
not a theorem for a parameter interval. All sufficiently large
removal values fail on some thin triangles, with no explicit
rigorous threshold. The $\mu=10^4$ result is independent Ritz/P2
finite-element evidence, not an interval certificate.
Solid$(\eta)$ similarly fails eventually as $\eta\downarrow0$.




\section{A physical atlas of triangle geometry}\label{sec:slide90}

The course began with the temptation to speak of "the center" of a triangle. We now have a more precise vocabulary. A distribution gives a mean, an energy gives a variational solution, a field gives an extremum, and symmetries can give lines, conics or a focal pair. Their mathematical definitions can be clear even when their barycentric functions require numerical evaluation. The atlas connects them by fingerprints and ETC neighbors, while keeping equality proofs and shape-response evidence distinct. This opens research problems small enough for undergraduates and deep enough to connect geometry with physics.



\paragraph{Demonstration and investigation.} Finish with a two-minute retrieval discussion. Ask students to explain the membrane--flow--Brownian equality, the magnetic Steiner foci, one mean-versus-maximum distinction, and one failed attractivity conjecture.


\section*{Questions and exercises}\addcontentsline{toc}{section}{Questions and exercises}

\begin{enumerate}

\item\label{ex:77} What happens to scale invariance if h stays fixed while the triangle grows?

\item\label{ex:78} Does the small-height theorem refute Solid(.1), Solid(1), or Solid(10)?

\item\label{ex:79} Why must the local exclusion disk be uniform across trial positions?

\item\label{ex:80} What extra work would promote this numerical failure to a certificate?

\item\label{ex:81} Why cannot we add the intensities of the aperture elements?

\item\label{ex:82} Would changing the incident phase mask invalidate the proof?

\item\label{ex:83} Does increasing beta always improve the center's geometric behavior?

\item\label{ex:84} What needs to be written beside a "lambda-independent" optical center?

\item\label{ex:85} Would restricting the blurred image to the original aperture preserve G?

\item\label{ex:86} Why should an optical center's ETC search use the same shape normalization as a membrane center's search?

\item\label{ex:87} Which two entries share exactly the same scalar PDE?

\item\label{ex:88} What should a project report include besides a plotted point?

\item\label{ex:89} Which result changed the interpretation of E0?

\item\label{ex:90} What is the next step for M6?

\end{enumerate}

\appendix

\chapter{Proofs, conventions and a reproducible numerical workflow}
\label{app:proofs}

\section{The maximum-area inellipse}
It is enough to prove the result in an equilateral triangle, because
an invertible affine map multiplies every area by the same determinant.
Put the centroid at zero, the inradius at $r$, and the three outward
unit normals at $n_1,n_2,n_3$, separated by $120^\circ$.
An ellipse is $c+B\mathbb D$ with $B$ symmetric positive definite.
Containment is equivalent to
\[
n_i\cdot c+|Bn_i|\le r,\qquad i=1,2,3.
\]
These inequalities are convex in $(c,B)$. Rotate a feasible pair
through $0,120^\circ,240^\circ$ and average. The averaged center is
zero and averaged matrix is $\overline B=(\tr B/2)\mathbf1$.
It is still feasible, so $\tr B/2\le r$.
By the arithmetic--geometric mean inequality,
$\det B\le(\tr B/2)^2\le r^2$. The ellipse area is $\pi\det B$.
Equality requires $B=r\mathbf1$. The three original containment
inequalities then give $n_i\cdot c\le0$; their sum is zero, forcing
$c=0$. The incircle is the unique maximum-area ellipse.
Its affine image is the unique Steiner inellipse.
Thus midpoint tangency, maximum area and the area-covariance formula
characterize the same conic by different routes.

\section{An elementary proof of the focal relation}
Translate the centroid to zero and represent the vertices by complex
numbers $z_1,z_2,z_3$, so $\sum z_i=0$.
The cubic derivative has roots $\pm f$, with
\[
f^2=-\tfrac13\sum_{i<j}z_i z_j
=\tfrac16\sum_i z_i^2.
\]
The Steiner shape matrix is
$Q=2\Sigma_{\rm area}=\frac16\sum_i V_iV_i^T$.
Consequently its anisotropy is
\[
Q_{xx}-Q_{yy}+2iQ_{xy}=\tfrac16\sum_i z_i^2=f^2.
\]
If its eigenvalues are $q_+\ge q_-$ and major direction angle is
$\vartheta$, the left side is
$(q_+-q_-)e^{2i\vartheta}$.
Ellipse foci are therefore
$\pm\sqrt{q_+-q_-}\,e^{i\vartheta}=\pm f$.
This proves the relation between the cubic derivative, covariance
and the two magnetic zeros without assuming that an affine map
preserves foci.

\section{Four exact attractivity results}
All derivatives below keep the other two vertices fixed.
The triangle inequalities are strict. Integrating the positive
quadratic form along any nondegenerate straight vertex path gives
the finite-motion result.

\subsection{The centroid}
$D_AG=\mathbf1/3$, so
$h\cdot[G(A+h)-G(A)]=|h|^2/3>0$ for $h\ne0$.
Its physical aliases include uniform lamina mean, equal vertex-mass
mean, prescribed uniform area-charge mean, and the reflecting
Brownian stationary mean.

\subsection{The incenter}
Set $A=0$, $p=(B-A)/c$, $q=(C-A)/b$.
Use an orthonormal frame along the internal angle bisector, so
$p=(\cos t,\sin t)$ and $q=(\cos t,-\sin t)$,
where $t=\theta_A/2$.
The incenter is $(2bc\cos t/P,0)$.
Direct differentiation of $I=(aA+bB+cC)/P$ gives
\[
D_AI=\frac aP\mathbf1-
\frac{(B-I)q^T+(C-I)p^T}{P}.
\]
Put $m=b+c$ and $d=b-c$. The cosine law gives
$\sin^2t=(a^2-d^2)/(m^2-d^2)$.
The diagonal entries of the symmetric part are
\[
(S_A)_{11}=\frac{m(a^2-d^2)}{(a+m)(m^2-d^2)},\qquad
(S_A)_{22}=\frac{a+m\sin^2t}{P}.
\]
Its off-diagonal entries vanish; the asymmetry of $b,c$ only produces
an antisymmetric part of the full derivative.
Since $|d|<a<m$, both displayed eigenvalues are positive.
The incenter is strictly attractive. The same theorem applies to
the restricted three-face deepest-apex construction.

\subsection{The uniform-wire mean}
Set $u=p\cdot q$. Differentiating the length-weighted midpoint mean
gives
\[
\begin{split}
S_A(W)={}&\frac{b+c}{2P}\mathbf1
-\frac{bc}{2P^2}(pp^T+qq^T)\\
&+\frac{c(a+c)+b(a+b)}{4P^2}(pq^T+qp^T).
\end{split}
\]
Directions $p+q$ and $p-q$ are orthogonal eigenvectors. Their
eigenvalues are
\[
\begin{split}
\lambda_+&=
\frac{2P(b+c)+(1+u)[(b-c)^2+a(b+c)]}{4P^2},\\
\lambda_-&=\frac{(b+c)(1+u)}{4P}.
\end{split}
\]
Both are positive because $-1<u<1$.
Thus $X_{10}$ is strictly attractive, including its specified
thin-wire charge and loop-field cutoff aliases.

\subsection{The angle-weighted vertex mean}
For counterclockwise vertices put
$u=(A-B)/|A-B|$, $v=(A-C)/|A-C|$ and let $\mathsf R$ rotate $90^\circ$.
Differentiating angle functions and their weighted mean gives
\[
D_A X_{360}
=\frac1\pi[\theta_A\mathbf1-u(\mathsf Ru)^T+v(\mathsf Rv)^T].
\]
The symmetric part has eigenvalues
$[\theta_A\pm\sin\theta_A]/\pi$.
For $0<\theta_A<\pi$, both are strictly positive.
This proves attractivity of M3 for its equal-current equal-core
interior-energy convention. Very thin angles make the smaller
eigenvalue tiny; numerical differentiation would be a weaker argument.

\section{Why the loop minimum is unique}
Let $R_p(\theta)$ be the distance from an interior point to the boundary
along direction $e_\theta$. Up to a positive constant, the loop
field is
\[
b(p)=\int_0^{2\pi}\frac1{R_p(\theta)}\dd\theta.
\]
One can obtain this identity from the boundary Biot--Savart integral,
or as the small-height solid-angle coefficient.
If $d_i(p)$ is side distance,
\[
\frac1{R_p(\theta)}
=\max_i\frac{\max(n_i\cdot e_\theta,0)}{d_i(p)}.
\]
Each positive reciprocal affine-distance function is convex in $p$;
the maximum and integral preserve convexity.
For any nonzero direction of motion, at least one side distance
changes. That side is first hit on an open angular interval, where
its reciprocal has strictly positive second derivative.
Thus $b$ is strictly convex. It diverges at the boundary and
has exactly one interior minimum. This proves that M1 is a point
rule; it does not prove attractive response to moving vertices.

\section{A finite-motion certificate for E0}
The explicit witness uses
\[
A=(0,0),\quad B=(1,0),\quad
C=(0.9963640131614389,0.0015991738862318396),
\]
and
$h=10^{-5}(-0.9997966839583786,-0.02016409546273338)$.
Both triangles are nondegenerate.
The project certificate encloses
\[
\frac{h\cdot[\mathrm{E0}(A+h,B,C)-\mathrm{E0}(A,B,C)]}{|h|^2}
\in[-3.13100322614,-3.13100322613]\,10^{-7}.
\]
Outward-rounded incomplete-beta series evaluations, a degree-250
local series, a degree-160 continuation and explicit remainder
bounds produce the enclosure. The finite-motion margin is small,
but its entire interval is negative. This single witness disproves
global attractivity. It does not say that E0 responds negatively
on most triangles. Rounded printed center coordinates alone would
not reproduce the certificate's sign; use the bundled interval record.

\section{Why large removal eventually fails}
Here is the mechanism behind the analytical family result.
Put $\varepsilon=\mu^{-1/2}$ and consider
$A=(0,0),B=(1,0),C=(x,\varepsilon\tau)$ with $1/2<x<1$.
Rescale $y=\varepsilon t$ and write $u=\varepsilon^2w$.
Apart from an overall factor, the energy becomes
\[
\mathcal J_\varepsilon[w]=
\int_{D_\tau}
\left[\frac{\varepsilon^2}{2}|w_s|^2+
\frac12|w_t|^2+\frac{\kappa_\varepsilon^2}{2}w^2-w\right]\dd s\dd t,
\quad \kappa_\varepsilon^2=L_\varepsilon^{-2}\to k^2.
\]
The horizontal derivative vanishes in the limit.
On a vertical interval of height $H$, solve
$-w_{tt}+k^2w=1$, $w(0)=w(H)=0$:
\[
w(t)=k^{-2}\left[1-
\frac{\cosh(k(t-H/2))}{\cosh(kH/2)}\right].
\]
Its column mass is
$F(H)=H/k^2-2\tanh(kH/2)/k^3$.
The triangular height is $\tau q(s)$, where $q(s)=s/x$ to the left
of the apex and $q(s)=(1-s)/(1-x)$ to its right.
Changing variables on each side shows
\[
c_x(\tau)=1-x+(2x-1)R_k(\tau),\qquad
R_k(\tau)=
\frac{\int_0^1qF(\tau q)\dd q}{\int_0^1F(\tau q)\dd q}.
\]
Since $F(H)\sim H^3/12$ for small $H$ and $F(H)\sim H/k^2$
for large $H$,
$R_k(\tau)\to4/5$ as $\tau\downarrow0$ and $R_k(\tau)\to2/3$
as $\tau\to\infty$.
For $x>1/2$ choose fixed $\tau_1,\tau_2$ giving
$c_x(\tau_2)-c_x(\tau_1)=-\delta<0$.

Move the apex from $(x,\varepsilon\tau_1)$ to
$(x+\varepsilon,\varepsilon\tau_2)$.
The horizontal motion is positive, but the horizontal center change
tends to $-\delta$. The vertical contribution is $O(\varepsilon^2)$,
so the dot product is $-\varepsilon\delta+o(\varepsilon)<0$.
Energy comparison controls convergence; for the vertical limiting
solution with $\kappa_\varepsilon$,
\[
\|(w_\varepsilon-w_0)_t\|_2^2+
\kappa_\varepsilon^2\|w_\varepsilon-w_0\|_2^2
\le\varepsilon^2\|(w_0)_s\|_2^2.
\]
This is a family theorem for every sufficiently large $\mu$,
with triangles depending on $\mu$. It supplies no explicit threshold.
The numerical $\mu=10^4$ witness is supporting evidence, not a
certified threshold value.

\section{A practical computation and evidence protocol}
\begin{enumerate}
\item Declare the model, boundary conditions, dimensionless parameters,
normalization and any finite-thickness or core convention.
\item Prove symmetry equivariance. Separate a unique field from a
unique extremum; use a positive mean or invariant set if needed.
\item Solve the reference triangle with refinement. Compare exact
controls and analytic limits. Integrate probability and moments
with quadrature appropriate to the representation.
\item Sample ordered-angle shape space and the four special families.
Keep training and holdout sets separate.
\item Search ETC locations, then compare derivatives and vertex
responses independently. Record unsupported catalogue entries.
\item Optimize potential negative witnesses. Refine both shapes
and compare finite differences with differentiated equations.
\item Certify a sign only with an enclosure for every error source.
Maintain separate labels for theorem, interval witness, numerical
witness, positive screen and open conjecture.
\end{enumerate}

\clearpage\section{Glossary of the main point rules}
\begin{center}\small\begin{tabular}{p{.18\linewidth}p{.70\linewidth}}
\toprule Symbol & Defining rule\\\midrule
$X_1=I$ & Equal side distances; largest contained circle center.\\
$X_2=G$ & Uniform area mean and equal vertex-mass mean.\\
$X_{10}=W$ & Uniform boundary-arclength mean.\\
$X_{360}=\mathrm{M3}$ & Interior-angle weighted vertex mean.\\
E0 & Equal harmonic measures of the three sides.\\
E1 & Maximum conformal radius; least small-electrode capacity.\\
E2 & Logarithmic equilibrium charge mean; exterior-map constant term.\\
E3 & Three-dimensional Newtonian equilibrium charge mean of a lamina.\\
M1 & Minimum interior magnitude of a perimeter-current normal field.\\
M6 & Interior mean weighted by field magnitude of a volume-current prism.\\
F1 & Unique maximum of the Dirichlet torsion function.\\
Pgrav & Candidate maximum of the clamped gravitational plate solution.\\
Q0m & Maximum positive ground-state amplitude.\\
Q$j$ & Mean weighted by normalized $|\psi_j|^2$; use clusters at degeneracy.\\
H$(\Delta_h)$ & Maximum of uniformly initialized Dirichlet temperature.\\
Abs$(\mu)$ & Mean of a uniformly produced, diffusively removed field.\\
Solid$(\eta)$ & Candidate maximum solid angle at dimensionless height $\eta$.\\
Illum & Maximum logarithmic ray-length finite-part illumination.\\
Dark$(F,\beta)$ & Mean weighted by $e^{-\beta J_F}$.\\
\bottomrule
\end{tabular}\end{center}
Physical constants that only multiply a positive field disappear from
a normalized mean or an extremum location. Parameters that change
the field shape must remain in the definition.

\clearpage\section{Numerical landmarks in the reference triangle}

Coordinates use $A=(0,0),B=(4,0),C=(0,3)$. Decimal values are archived numerical estimates, not interval enclosures. The classical controls have the exact coordinates given in Chapter 1.

\begin{center}\small\begin{tabular}{p{0.3\linewidth}p{0.3\linewidth}p{0.3\linewidth}}\toprule
Rule & $c_x$ & $c_y$\\\midrule
X1 & 1.000000 & 1.000000\\
X2 & 1.333333 & 1.000000\\
X10 & 1.500000 & 1.000000\\
E0 & 0.930688 & 0.960667\\
E2 & 1.502468 & 1.015772\\
E3 & 1.417142 & 1.002754\\
M1 & 1.103977 & 0.998576\\
M6 & 1.310681 & 1.006307\\
F1 & 1.094988 & 0.998228\\
Pgrav & 1.075314 & 0.999031\\
$\mathrm{Q0m}$ & 1.095451 & 0.997508\\
$\mathrm{Q0}$ & 1.126561 & 1.001531\\
Q1 & 1.413362 & 0.987726\\
Q2 & 0.958195 & 1.050956\\
Q3 & 1.526410 & 0.905603\\
Q4 & 1.191916 & 1.127577\\\bottomrule\end{tabular}\end{center}

\clearpage\section{Selected ETC location comparisons}

Two top training-location neighbors are retained for each rule below. The last column gives holdout RMS for the first neighbor, divided by diameter. These are sample ranks, not identities or attractivity comparisons. Unsupported entries are excluded as described in Chapter 5.

\begin{center}\small\begin{tabular}{p{0.27\linewidth}p{0.39\linewidth}p{0.24\linewidth}}\toprule
Rule & ETC neighbors & Holdout RMS / $D$\\\midrule
E0 & $X_{7707}$, $X_{30421}$ & 0.00041995\\
E1 & $X_{53588}$, $X_{66575}$ & 0.0024644\\
E2 & $X_{59474}$, $X_{5235}$ & 0.0013054\\
E3 & $X_{27760}$, $X_{31351}$ & 0.0043591\\
M1 & $X_{2534}$, $X_{5393}$ & 0.0065758\\
M6 & $X_{25086}$, $X_{26104}$ & 0.0026819\\
F1 & $X_{5393}$, $X_{71701}$ & 0.0040502\\
$\mathrm{Q0m}$ & $X_{5393}$, $X_{5308}$ & 0.0068985\\
$\mathrm{Q0}$ & $X_{24936}$, $X_{39622}$ & 0.0062635\\
Q1 & $X_{7324}$, $X_{21939}$ & 0.014209\\
Q2 & $X_{43441}$, $X_{43469}$ & 0.03281\\
Q3 & $X_{24762}$, $X_{43469}$ & 0.05539\\
Q4 & $X_{4179}$, $X_{50094}$ & 0.048552\\
H0.03 & $X_{29678}$, $X_{2534}$ & 0.0054628\\
H1.0 & $X_{5393}$, $X_{5308}$ & 0.0069927\\
Abs0.0 & $X_{30105}$, $X_{37870}$ & 0.00072444\\
Abs100.0 & $X_{32009}$, $X_{70563}$ & 0.0031274\\
Illum & $X_{37884}$, $X_{39619}$ & 0.0078736\\
Solid1.0 & $X_{55753}$, $X_{60182}$ & 0.00011148\\
Dark1.0b1.0 & $X_{17210}$, $X_{17248}$ & 0.0027955\\
Persist1.0 & $X_{61001}$, $X_{19870}$ & 0.0017181\\\bottomrule\end{tabular}\end{center}

\chapter{Supplementary lecture material}\label{app:supplement}

\section{Coordinates and normalization}\label{sec:slide91}

Keep the coordinate convention beside every numerical comparison. The base normalization A=(-1,0), B=(1,0) used for shape fingerprints differs from the 3--4--5 laboratory coordinates. Shape means side ratios or angles, not absolute length. A normalized RMS divided by diameter is dimensionless. External scales such as wire radius, screen distance and diffusion time must be declared.


\[p=\alpha A+\beta B+\gamma C,\quad\alpha+\beta+\gamma=1.\]


\section{Fingerprints need more than one family}\label{sec:slide92}

In the isosceles family only the vertical coordinate varies by reflection symmetry. The right-triangle family supplies two coordinates and breaks this symmetry. Differentiate numerical curves only after a convergence check. Near an eigenvalue crossing, replace a single eigenstate by the entire cluster projector. Thin limits reveal behavior missed by moderate shapes. This appendix shows the $\mathrm{Q0}$ right-family position traces; the bundled JSON also contains derivative traces.


\begin{figure}[htbp]\centering
\begin{tikzpicture}\begin{axis}[width=.95\linewidth,height=.56\linewidth,grid=major,grid style={black!10},xlabel={family parameter $x$},ylabel={$c_x$},tick label style={font=\scriptsize},label style={font=\small},legend style={font=\scriptsize,at={(0.5,-0.24)},anchor=north,legend columns=3,draw=none},scaled ticks=false]
\addplot[navy,thick,no marks] coordinates {(0.2,0.08366923) (0.228706,0.09421154) (0.2615321,0.1059427) (0.2990698,0.1189638) (0.3419952,0.1333763) (0.3910817,0.149279) (0.4472136,0.1667661) (0.5114021,0.1859223) (0.5848035,0.2068206) (0.6687403,0.2295182) (0.7647245,0.2540545) (0.8744853,0.28045) (1,0.3087068) (1.14353,0.3388122) (1.30766,0.3707423) (1.495349,0.4044681) (1.709976,0.4399606) (1.955409,0.4771961) (2.236068,0.5161592) (2.55701,0.5568458) (2.924018,0.5992598) (3.343702,0.643426) (3.823622,0.689373) (4.372426,0.7371403) (5,0.7867845)};
\addlegendentry{$\mathrm{Q0}$}
\addplot[teal,thick,no marks] coordinates {(0.2,0.09009805) (0.228706,0.101443) (0.2615321,0.1139491) (0.2990698,0.127653) (0.3419952,0.1425658) (0.3910817,0.1586645) (0.4472136,0.1758842) (0.5114021,0.1941113) (0.5848035,0.2131792) (0.6687403,0.2328692) (0.7647245,0.2529173) (0.8744853,0.273028) (1,0.2928932) (1.14353,0.3122156) (1.30766,0.33073) (1.495349,0.3482207) (1.709976,0.3645313) (1.955409,0.3795669) (2.236068,0.3932891) (2.55701,0.4057068) (2.924018,0.4168649) (3.343702,0.4268336) (3.823622,0.4356983) (4.372426,0.4435522) (5,0.4504902)};
\addlegendentry{X1}
\addplot[magenta,thick,no marks] coordinates {(0.2,0.06666667) (0.228706,0.07623532) (0.2615321,0.08717737) (0.2990698,0.09968992) (0.3419952,0.1139984) (0.3910817,0.1303606) (0.4472136,0.1490712) (0.5114021,0.1704674) (0.5848035,0.1949345) (0.6687403,0.2229134) (0.7647245,0.2549082) (0.8744853,0.2914951) (1,0.3333333) (1.14353,0.3811766) (1.30766,0.4358868) (1.495349,0.4984496) (1.709976,0.569992) (1.955409,0.6518028) (2.236068,0.745356) (2.55701,0.8523368) (2.924018,0.9746726) (3.343702,1.114567) (3.823622,1.274541) (4.372426,1.457475) (5,1.666667)};
\addlegendentry{X2}
\addplot[gold,thick,no marks] coordinates {(0.2,0.05495098) (0.228706,0.06363146) (0.2615321,0.0737915) (0.2990698,0.08570837) (0.3419952,0.0997147) (0.3910817,0.1162086) (0.4472136,0.1356647) (0.5114021,0.1586454) (0.5848035,0.1858122) (0.6687403,0.2179356) (0.7647245,0.2559036) (0.8744853,0.3007287) (1,0.3535534) (1.14353,0.4156571) (1.30766,0.4884652) (1.495349,0.5735641) (1.709976,0.6727223) (1.955409,0.7879208) (2.236068,0.9213894) (2.55701,1.075652) (2.924018,1.253576) (3.343702,1.458434) (3.823622,1.693962) (4.372426,1.964437) (5,2.274755)};
\addlegendentry{X10}
\addplot[blue,thick,no marks] coordinates {(0.2,0.08406177) (0.228706,0.09446653) (0.2615321,0.1060318) (0.2990698,0.1188656) (0.3419952,0.1330816) (0.3910817,0.1487961) (0.4472136,0.1661265) (0.5114021,0.1851866) (0.5848035,0.2060819) (0.6687403,0.228904) (0.7647245,0.2537238) (0.8744853,0.2805852) (1,0.3094988) (1.14353,0.3404363) (1.30766,0.3733269) (1.495349,0.4080544) (1.709976,0.4444566) (1.955409,0.4823257) (2.236068,0.5214118) (2.55701,0.5614274) (2.924018,0.6020542) (3.343702,0.642953) (3.823622,0.6837735) (4.372426,0.7241665) (5,0.7637955)};
\addlegendentry{X24936}
\addplot[purple,thick,no marks] coordinates {(0.2,0.08415365) (0.228706,0.09457791) (0.2615321,0.1061741) (0.2990698,0.1190581) (0.3419952,0.1333545) (0.3910817,0.1491952) (0.4472136,0.1667168) (0.5114021,0.1860563) (0.5848035,0.207344) (0.6687403,0.2306951) (0.7647245,0.256198) (0.8744853,0.2839032) (1,0.3138115) (1.14353,0.345866) (1.30766,0.3799474) (1.495349,0.4158749) (1.709976,0.4534111) (1.955409,0.4922709) (2.236068,0.5321319) (2.55701,0.5726468) (2.924018,0.6134544) (3.343702,0.6541914) (3.823622,0.694503) (4.372426,0.7340531) (5,0.7725328)};
\addlegendentry{X39622}
\end{axis}\end{tikzpicture}
\par\vspace{5mm}
\begin{tikzpicture}\begin{axis}[width=.95\linewidth,height=.56\linewidth,grid=major,grid style={black!10},xlabel={family parameter $x$},ylabel={$c_y$},tick label style={font=\scriptsize},label style={font=\small},legend style={font=\scriptsize,at={(0.5,-0.24)},anchor=north,legend columns=3,draw=none},scaled ticks=false]
\addplot[navy,thick,no marks] coordinates {(0.2,0.1573569) (0.228706,0.1685884) (0.2615321,0.1802932) (0.2990698,0.1924293) (0.3419952,0.204944) (0.3910817,0.2177722) (0.4472136,0.2308334) (0.5114021,0.2440391) (0.5848035,0.2572905) (0.6687403,0.2704841) (0.7647245,0.2835157) (0.8744853,0.2962863) (1,0.3087068) (1.14353,0.3207029) (1.30766,0.3322171) (1.495349,0.3432098) (1.709976,0.3536583) (1.955409,0.3635541) (2.236068,0.3729003) (2.55701,0.3817081) (2.924018,0.3899947) (3.343702,0.3977795) (3.823622,0.4050848) (4.372426,0.411933) (5,0.4183462)};
\addlegendentry{$\mathrm{Q0}$}
\addplot[teal,thick,no marks] coordinates {(0.2,0.09009805) (0.228706,0.101443) (0.2615321,0.1139491) (0.2990698,0.127653) (0.3419952,0.1425658) (0.3910817,0.1586645) (0.4472136,0.1758842) (0.5114021,0.1941113) (0.5848035,0.2131792) (0.6687403,0.2328692) (0.7647245,0.2529173) (0.8744853,0.273028) (1,0.2928932) (1.14353,0.3122156) (1.30766,0.33073) (1.495349,0.3482207) (1.709976,0.3645313) (1.955409,0.3795669) (2.236068,0.3932891) (2.55701,0.4057068) (2.924018,0.4168649) (3.343702,0.4268336) (3.823622,0.4356983) (4.372426,0.4435522) (5,0.4504902)};
\addlegendentry{X1}
\addplot[magenta,thick,no marks] coordinates {(0.2,0.3333333) (0.228706,0.3333333) (0.2615321,0.3333333) (0.2990698,0.3333333) (0.3419952,0.3333333) (0.3910817,0.3333333) (0.4472136,0.3333333) (0.5114021,0.3333333) (0.5848035,0.3333333) (0.6687403,0.3333333) (0.7647245,0.3333333) (0.8744853,0.3333333) (1,0.3333333) (1.14353,0.3333333) (1.30766,0.3333333) (1.495349,0.3333333) (1.709976,0.3333333) (1.955409,0.3333333) (2.236068,0.3333333) (2.55701,0.3333333) (2.924018,0.3333333) (3.343702,0.3333333) (3.823622,0.3333333) (4.372426,0.3333333) (5,0.3333333)};
\addlegendentry{X2}
\addplot[gold,thick,no marks] coordinates {(0.2,0.454951) (0.228706,0.4492785) (0.2615321,0.4430255) (0.2990698,0.4361735) (0.3419952,0.4287171) (0.3910817,0.4206678) (0.4472136,0.4120579) (0.5114021,0.4029443) (0.5848035,0.3934104) (0.6687403,0.3835654) (0.7647245,0.3735413) (0.8744853,0.363486) (1,0.3535534) (1.14353,0.3438922) (1.30766,0.334635) (1.495349,0.3258897) (1.709976,0.3177344) (1.955409,0.3102165) (2.236068,0.3033554) (2.55701,0.2971466) (2.924018,0.2915676) (3.343702,0.2865832) (3.823622,0.2821508) (4.372426,0.2782239) (5,0.2747549)};
\addlegendentry{X10}
\addplot[blue,thick,no marks] coordinates {(0.2,0.1527591) (0.228706,0.1656212) (0.2615321,0.1788287) (0.2990698,0.1922878) (0.3419952,0.2058996) (0.3910817,0.219564) (0.4472136,0.2331825) (0.5114021,0.2466624) (0.5848035,0.2599198) (0.6687403,0.2728824) (0.7647245,0.2854922) (0.8744853,0.2977065) (1,0.3094988) (1.14353,0.3208575) (1.30766,0.3317846) (1.495349,0.3422913) (1.709976,0.3523951) (1.955409,0.3621155) (2.236068,0.3714703) (2.55701,0.3804733) (2.924018,0.3891328) (3.343702,0.3974512) (3.823622,0.4054254) (4.372426,0.413048) (5,0.4203089)};
\addlegendentry{X24936}
\addplot[purple,thick,no marks] coordinates {(0.2,0.1545066) (0.228706,0.1678823) (0.2615321,0.1816348) (0.2990698,0.1956489) (0.3419952,0.2097984) (0.3910817,0.2239517) (0.4472136,0.2379766) (0.5114021,0.2517483) (0.5848035,0.2651564) (0.6687403,0.2781123) (0.7647245,0.2905551) (0.8744853,0.3024547) (1,0.3138115) (1.14353,0.3246518) (1.30766,0.3350201) (1.495349,0.3449696) (1.709976,0.3545533) (1.955409,0.3638161) (2.236068,0.3727902) (2.55701,0.3814937) (2.924018,0.389931) (3.343702,0.3980948) (3.823622,0.4059696) (4.372426,0.4135349) (5,0.4207682)};
\addlegendentry{X39622}
\end{axis}\end{tikzpicture}
\caption[The two position coordinates of Q0 on the archived right-triangle family]{The two position coordinates of Q0 on the archived right-triangle family.}\label{fig:right-fp}
\end{figure}

\section{Why the Steiner inellipse has maximum area}\label{sec:slide93}

Every inscribed ellipse is sent to an inscribed ellipse and every area is multiplied by the same determinant. In the equilateral reference triangle the incircle is the unique maximum-area inscribed ellipse. Its affine image touches all three side midpoints, is centered at G and is the Steiner inellipse. This argument explains affine invariance and uniqueness; a full proof of the reference maximum is available in the mass chapter. Ellipse Q describes (p-G)\ensuremath{^T}Q\ensuremath{^{-}}\ensuremath{^1}(p-G)=1. It is not an inertia tensor.


\[Q_{\rm Steiner}=2\Sigma_{\rm area}.\]


\section{An axis ratio is not a moment ratio}\label{sec:slide94}

For a distribution in the plane, I=M(tr(\ensuremath{\Sigma})1\ensuremath{-}\ensuremath{\Sigma}). If \ensuremath{\Sigma} has eigenvalues \ensuremath{\lambda}1\ensuremath{\ge}\ensuremath{\lambda}2, then the moments about the corresponding principal lines are M\ensuremath{\lambda}2 and M\ensuremath{\lambda}1. Thus the ellipse major axis and the smaller inertia moment share an eigenvector. This reversal matters when comparing geometric axis ratios with rotational measurements. Isotropic covariance removes the distinction between principal directions.


\[a_{\rm ellipse}/b_{\rm ellipse}=\sqrt{\lambda_{\max}(\Sigma)/\lambda_{\min}(\Sigma)}.\]


\section{The skeleton has a covariance ellipse}\label{sec:slide95}

A wire is a mixture of uniform segments. Each segment has midpoint m\_i, length mass w\_i and covariance e\_i e\_i\ensuremath{^T}/12, where e\_i is its full vector. The total covariance is the weighted mean of segment covariances plus (m\_i-W)(m\_i-W)\ensuremath{^T}. A covariance ellipse always exists. What fails is the coincidence with the midpoint-tangent, centroid-centered Steiner ellipse. The 3--4--5 wire center is W=(3/2,1), already different from G.


\[\Sigma_{\rm wire}=\begin{pmatrix}7/4&-2/3\\-2/3&1\end{pmatrix}.\]


\section{M3: angle-weighted vertex mean}\label{sec:slide96}

For three equal axial vertex currents, magnetic energy diverges logarithmically inside each interior angular sector. Equal circular core radii give weights equal to those angles. The symmetric derivative at A has eigenvalues (\ensuremath{\alpha}+sin \ensuremath{\alpha})/\ensuremath{\pi} and (\ensuremath{\alpha}\ensuremath{-}sin \ensuremath{\alpha})/\ensuremath{\pi}, both positive for 0<\ensuremath{\alpha}<\ensuremath{\pi}. Integrating along a vertex displacement that remains in the same nondegenerate half-plane proves strict attraction. The whole-plane source neighborhoods instead have equal full angular weights and select G. The full derivation is in the magnetostatics chapter.


\[\lambda_\pm(D_AX_{360})=(\theta_A\pm\sin\theta_A)/\pi>0.\]


\section{Least capacity and the Robin function}\label{sec:slide97}

For a conformal map from T to the unit disk taking p to zero, the inverse derivative gives conformal radius. Maximizing it minimizes the corresponding logarithmic self-energy after subtraction of the pole singularity. Constants and signs depend on the Green-function convention. In n\ensuremath{\ge}3 use a Newtonian Green function and its regular part; planar conformal equivalence does not survive unchanged in higher dimension. Distinguish a moving point charge self-energy from the equilibrium distribution on a whole conductor.


\[r_T(p)=1/|f_p^\prime(p)|,\quad f_p(p)=0.\]


\section{Finite elements in one slide}\label{sec:slide98}

Choose continuous piecewise-polynomial basis functions. Dirichlet boundary degrees of freedom are removed. A Poisson solve uses Ku=f. A reaction solve uses (K+\ensuremath{\mu}M)u=f. A quantum solve uses Kv=\ensuremath{\lambda}Mv, not an ordinary Euclidean eigenproblem. The eigenvector normalization is v\ensuremath{^T}Mv=1. Curved maxima require local optimization or interpolation; a largest mesh node alone is not a certified maximum. Compare meshes and polynomial orders before reporting response signs.


\[K_{ij}=\int_T\nabla\phi_i\cdot\nabla\phi_j\dd A,\quad M_{ij}=\int_T\phi_i\phi_j\dd A.\]


\section{Degenerate quantum states}\label{sec:slide99}

Inside a degenerate eigenspace, orthogonal basis rotations change an individual squared eigenfunction and its expectation. The sum of squared amplitudes is invariant because it is the diagonal of the orthogonal projector. Divide by dimension m to normalize the density. At an exact crossing this cluster convention restores geometric meaning. A single numbered excited center can have discontinuous limits if the definition ignores degeneracy.


\[\rho_{\rm cluster}=m^{-1}\sum_j|\psi_j|^2.\]


\section{What a counterexample certificate proves}\label{sec:slide100}

The archived E0 interval certificate encloses both centers
before and after the finite motion, and their normalized dot product
is strictly negative. Appendix A gives the exact input decimals,
the enclosing interval and the series-error convention.
High-precision agreement is useful corroboration, but the interval
enclosure supplies the proof. A small negative floating-point number
alone would not suffice.


\[h\cdot\Delta\mathrm{E0}<0.\]


\section{Why large removal can fail attractivity}\label{sec:slide101}

In a thin triangle of height \ensuremath{\varepsilon}\ensuremath{\tau} the leading reaction field is a vertical strip profile. Its normalized horizontal first moment contains R(\ensuremath{\tau}), which decreases from 4/5 toward 2/3. For x>1/2, increasing \ensuremath{\tau} while moving the apex slightly right can move the center left. The horizontal dot-product term dominates the smaller vertical term. An energy estimate controls the strip limit. This proves failure for every sufficiently large \ensuremath{\mu}, without giving an explicit threshold. Independent numerical methods support a witness at $\mu=10^4$ but do not certify that exact parameter.


\[\varepsilon=\mu^{-1/2},\quad c_x(\tau)\to1-x+(2x-1)R_k(\tau).\]


\section{Continue the investigation}\label{sec:slide102}

Use the atlas and six chapters as a starting point. Good undergraduate projects include estimating $\mathrm{Q0}$ with a refined mesh, checking the membrane--flow--Brownian identity, reproducing the wire covariance, comparing matched Fresnel numbers and investigating M6 response. Keep source conventions, solver tolerances and evidence status beside each result. The bundle manifest gives the exact data sources and distinguishes archival field illustrations from new certificates.


\chapter{Answers and routes to solutions}\label{app:answers}

\section{Lecture 1}

\begin{enumerate}

\item G, because their average position is (A+B+C)/3.

\item No. The dimensionless ratio of height to triangle size must be held fixed.

\item The midpoint of AB.

\item Normalization cancels the common factor.

\item No. Different physical experiments can produce the same X-number.

\item No. The dot product condition permits a sideways component.

\item W at x=3/2.

\item Equal mass per edge, which differs from equal mass per unit length.

\item No. In-plane inertia about e uses the other principal covariance eigenvalue.

\item Rotational symmetry can fix the point and make every direction equally valid.

\item They cannot transform an equilateral triangle into a general scalene triangle. An affine map can.

\item Generally no. Different edge directions stretch by different factors.

\item No. Several ellipses exist, but the mass covariance does not reproduce the Steiner coincidence.

\item The foci of the Steiner inellipse.

\end{enumerate}

\section{Lecture 2}

\begin{enumerate}

\item Interiority when the distribution is not confined to one side. It does not guarantee vertex attractivity.

\item Their interaction kernels differ, so the energy-minimizing distributions differ.

\item No. It supplies convergence evidence. A validated error estimate is needed for certification.

\item The incenter. It generally differs from E0.

\item Changing its size changes the divergent self contribution and can dominate the location effect.

\item No. Uniqueness concerns a fixed triangle. Attractivity compares different triangles.

\item The two integrals diverge. The physical regularization is part of the definition.

\item The inellipse is a circle, and the two foci merge at the common center.

\item Each vertex then sees a full angular neighborhood, rather than its unequal interior sector.

\item It proves interiority of the mean, but gives no universal sign for its response to changing the domain.

\item One strict negative finite motion with rigorously enclosed center coordinates.

\item Only that no failure was found in those cases. The strict certified counterexample refutes a global conjecture.

\item A global M6 response inequality, or a validated counterexample if one exists.

\end{enumerate}

\section{Lecture 3}

\begin{enumerate}

\item The imposed tension is uniform and changes little under the small deflection.

\item No. It multiplies the solution by two and preserves its maximum location.

\item It restricts the test functions and removes the boundary term.

\item I. That maximin distance problem differs from the Poisson problem.

\item No. Here the u and |\ensuremath{\nabla}u|\ensuremath{^2} distributions give an explicit counterexample.

\item G. Sampling particles by flux gives the velocity-weighted mean instead.

\item No. Equal exit probabilities select E0.

\item It must specify whether the slope is clamped or whether bending moments are prescribed.

\item The operator and boundary conditions differ, and fourth-order positivity is more delicate.

\item Material law, support constraints, thickness, temperature, and which density measure is intended.

\item No. Without additional mechanical or geometric restrictions the fixed-area gravity problem can be unbounded.

\item A mean, because it integrates the field. A maximum depends on local derivatives and localization.

\item Yes. A one-parameter symmetry family is only a slice of shape space.

\item Replacing slope energy by curvature energy and imposing clamped plate boundary conditions.

\end{enumerate}

\section{Lecture 4}

\begin{enumerate}

\item No. It changes the energy scale but preserves the normalized spatial eigenfunction.

\item Otherwise multiplying the wavefunction by zero would make the unnormalized energy vanish.

\item Yes, for an excited state. The first moment averages probability on both sides of nodes.

\item No. Symmetry may force the same expectation for several states.

\item The full two-dimensional eigenspace, unless additional physical state preparation is specified.

\item It is the average over the distribution, while each individual measurement is random.

\item The basis is not orthonormal in physical area. Probability normalization uses M.

\item An independent identity theorem, not merely a close numerical rank.

\item The nearly collinear family, especially when a vertex also approaches another vertex.

\item Total probability decays, and a uniform stationary density is no longer available.

\item At least G, the $\psi_0$ mean, $\mathrm{Q0}$, F1 and E0, because the experimental questions differ.

\item No. It constrains time along a flight but does not prove ergodicity of the direction.

\item Errors in opposite directions can cancel in the vector average.

\item The full eigenspace projector or a properly defined thermal mixture.

\end{enumerate}

\section{Lecture 5}

\begin{enumerate}

\item It lies farthest from the boundary and initially loses heat least strongly.

\item The initial constant interior field with a zero boundary requires a sharp boundary layer that a few smooth modes cannot represent.

\item The statement here is a limiting one. Exact equality at finite time is not asserted.

\item The actual cooling amplitude. Shape comparisons and energy decay need different visual conventions.

\item The triangle normalization and the fixed dimensionless time parameter.

\item Generally no. $\mathrm{Q0}$ uses $\psi_0^2$; the heat image uses $\psi_0$ itself.

\item No. Attraction requires a statement uniform over changes of triangle shape.

\item H cubed, compared with a weight proportional to H.

\item The counterexample triangles become thinner as \ensuremath{\mu} grows, so the fixed-shape limit is not uniform over shape space.

\item No. It encodes shape up to similarity.

\item The evaluator could not test them reliably. Exclusion is a computational limitation.

\item M3=X360, with angle barycentrics.

\item A strict rigorously enclosed negative finite motion.

\item M6, while keeping open the possibility of a difficult counterexample.

\end{enumerate}

\section{Lecture 6}

\begin{enumerate}

\item Eta changes. It is a different member of the family.

\item No explicit threshold is provided. Those particular parameters retain their open no-failure status.

\item A position-dependent cutoff would add a position-dependent finite term and define a different center.

\item Rigorous enclosures of the stationary coordinates and propagated finite-motion dot product.

\item Coherent contributions interfere, so their complex amplitudes must be summed first.

\item Yes. The radial integral would no longer have this simple nonnegative imaginary part.

\item No. It can amplify narrow dark regions and adverse shape sensitivity.

\item The dimensionless parameter convention and the physical change of screen distance or averaging procedure.

\item Generally no. Truncation removes an asymmetric part of the blurred tails.

\item A common normalization makes the location and derivative comparisons meaningful.

\item The linear membrane and fully developed Newtonian duct, both producing the torsion field.

\item Definition, parameter conventions, numerical convergence, relevant classical comparisons, and a clearly labeled evidence level.

\item A rigorously enclosed negative finite dot product, rather than another numerical location comparison.

\item Control the field-redistribution term in its shape derivative globally, or find a validated counterexample.

\end{enumerate}

\backmatter

\renewcommand{\bibname}{References and research record}
\begin{thebibliography}{99}
\addcontentsline{toc}{chapter}{References and research record}
\item[] The mathematical definitions in these notes are independent of the
project files. Decimal illustrations and catalogue rankings reproduce
the archived teaching dataset. The evidence snapshot is dated
30 September 2026. The companion lecture package includes the atlas,
selected certificates and the analytical addendum; the large vendor
ETC database is referenced rather than redistributed.
\bibitem{etc} C. Kimberling, \emph{Encyclopedia of Triangle Centers}.
\url{https://faculty.evansville.edu/ck6/encyclopedia/ETC.html}.
Local project snapshot: X1--X72807, Wolfram WXF; numerical search
coverage stated in Chapter 5.
\bibitem{attractive} Project notes, \emph{Attractive Centers},
definition and Parts I--IV.
\url{https://ismaa3iil.fyi/attractive-centers/}.
The derivative proofs needed here are reproduced in Appendix A.
\bibitem{iannaccone} A. Iannaccone,
\emph{The Conformal Center of a Triangle or a Quadrilateral},
Harvey Mudd College thesis (2003).
\url{https://scholarship.claremont.edu/hmc_theses/149/}.
\bibitem{finch} S. R. Finch, \emph{Appell $F_1$ and conformal mapping}
(2014). \url{https://arxiv.org/abs/1408.1074}.
\bibitem{wire} M. H. Partovi and J. D. Griffiths,
\emph{Equilibrium charge density on a thin curved wire},
American Journal of Physics \textbf{77}, 1173 (2009).
\url{https://doi.org/10.1119/1.3231081};
\url{https://arxiv.org/abs/0904.4279}.
\bibitem{marden} E. Badertscher, \emph{A simple direct proof of Marden's
theorem}, American Mathematical Monthly \textbf{121}, 547--548 (2014).
\url{https://doi.org/10.4169/amer.math.monthly.121.06.547}.
\bibitem{torsion} W. Borrelli, S. Mosconi and M. Squassina,
\emph{Uniqueness of the critical point for solutions to some
$p$-Laplace equations in the plane} (2022 preprint).
\url{https://arxiv.org/abs/2201.12788}.
\bibitem{endo} R. Endo and X. Liu, \emph{The second Dirichlet eigenvalue
is simple on every non-equilateral triangle, Part II: Nearly
Equilateral Triangles}, arXiv v6 (2025).
\url{https://arxiv.org/abs/2305.14063}.
\bibitem{brascamp} H. J. Brascamp and E. H. Lieb,
\emph{On extensions of the Brunn--Minkowski and Prekopa--Leindler
theorems, including inequalities for log concave functions, and with
an application to the diffusion equation},
Journal of Functional Analysis \textbf{22}, 366--389 (1976).
\url{https://doi.org/10.1016/0022-1236(76)90004-5}.
\bibitem{heatconc} K. Ishige, P. Salani and A. Takatsu,
\emph{Characterization of $F$-concavity preserved by the Dirichlet
heat flow}, arXiv v2 (2023).
\url{https://arxiv.org/abs/2207.13449}.
\bibitem{kms} S. Kerckhoff, H. Masur and J. Smillie,
\emph{Ergodicity of billiard flows and quadratic differentials},
Annals of Mathematics \textbf{124}, 293--311 (1986).
\url{https://doi.org/10.2307/1971280}.
\bibitem{marklof} J. Marklof and Z. Rudnick,
\emph{Almost all eigenfunctions of a rational polygon are uniformly
distributed}, Journal of Spectral Theory \textbf{2}, 107--113 (2012).
\url{https://ems.press/content/serial-article-files/33468}.
\bibitem{fresnel} University of Tennessee, Knoxville,
\emph{Fresnel diffraction}, Physics 421 teaching module.
\url{https://labs.phys.utk.edu/mbreinig/phys421/modules/m6/Fresnel.html}.
\bibitem{project} \emph{Physically Inspired Triangle Centers},
six-lecture teaching dataset and attractivity addendum
(30 September 2026). Companion materials:
course content, scientific visualization data, atlas catalogue,
E0/E1/E2/M1 interval records and refined PDE witnesses.
File names and checksums are supplied in the lecture package.
\end{thebibliography}

\end{document}